% Papiers de Sophie Germain — Français 9115, lecture modernisée du volume entier.
%
% Pass: Opus 5.5 (claude-opus-5-5), 2026-09-23 — staged modernisation, unchecked by a human.
%
% This is an interpretation, not a record. Where this file and the
% transcriptions differ in substance, the transcriptions (batch-01.fr.tex …
% batch-38.fr.tex) are the record.
%
% Why it was made in stages. The transcription of this volume — 38 batches,
% views 1 to 750, about 1.9 MB of LaTeX — exceeds what one pass can hold. With
% the user's approval the reading was therefore made in stages: first a plan,
% drawn from the header comments of the 38 batch files; then seven section
% writers, each of whom read its own batches in full, twice, and wrote one
% group of sections; then an assembler, who read the plan, the seven fragments
% and the writers' reports, joined them, and wrote the résumé, the map of the
% volume and the closing list. No single context read the whole volume. The
% argument therefore runs continuously within each group, but not necessarily
% across groups; the cross-references between groups were made at assembly, by
% section title and views.
%
% The groups, in the order in which they appear below:
%   G1  views 3-114, 136-142, 465, 483, 650: the « Exposition des principes »,
%       the « Mémoire sur la courbure des surfaces » and its drafts, and the
%       scattered leaves on curvature;
%   G2  views 144-334: Euler's 1779 memoir on laminae (Latin copy) and her
%       critique of his § 47, paragraphs 1 to 7;
%   G3  views 246-396: the critique continued, paragraphs 8 to 25, and the
%       figure leaves;
%   G4  views 398-490 and 593-626: the long memoir on elastic surfaces, her
%       pages (1) to (80), and the other states of its opening;
%   G5  views 491-592 and 600-613: the long memoir, her pages (81) to (165),
%       and a second redaction of its comparison with Chladni;
%   G7  views 549 and 630-716: heat in the sphere, the cone and the cylinder;
%   G6  views 115-134 and 704-749: Manuscript C, the notes on quadratic
%       residues, Legendre's printed memoir, and the extracts from the Turin
%       memoirs.
%
% Sources read: the 38 batch transcriptions of this volume and nothing else.
% No image was consulted, and no printed edition of any piece in it; where a
% transcription has a gap, the reading says so and does not fill it.
%
% What the volume turned out to be: an unsorted binding. Elastic plates and
% the curvature beneath them make the largest mass; then heat after Fourier;
% then number theory, including Manuscript C (ff. 348r-349r, views 708-710).
% Sophie Germain's theorem on auxiliary primes is not on these leaves.
%
% Notation translated, each with a footnote at its first use in its section:
% the straight d of partial derivatives into the partial sign; her fourth-order
% plate equation into the bilaplacian; her « somme des raisons inverses des deux
% rayons de courbure » into 2H; Euler's parenthesised differentials; her
% « résidu » / « non résidu » into the Legendre symbol; her Pi (a parameter)
% kept apart from pi; her stilet written stylet.
%
% Inferences made by this reading, each footnoted where it is made: that the
% long memoir (views 398-592) is her « mémoire pour 1813 »; that views 598 and
% 614-618 are one state of its opening and the oldest of three; that views
% 620-626 draft the crowned memoir « présenté en 1815 »; that the slip of view
% 416 corrects § 6 of the long memoir; that view 142 renders the opening of
% Gauss's 1827 memoir; that the heat leaves post-date Fourier's Théorie (1822);
% the reading order 639-636-637 and the page 3 of the sphere piece (view 660).
%
% Assembly: duplicated modern statements were kept once, where the plan put
% them (Kirchhoff's edge conditions in the long memoir's N° 9; the null
% Lagrangian in the Exposition; the beam on an intermediate support at the end
% of the critique's first part; the ring law with her ring; the wax
% equivalence with the long memoir's § 6; the node spacing of the free lamina
% in the critique), the others pointing there. Contradictions between
% fragments were settled by the mathematics or the transcription, as the
% closing section and the footnotes say.
\documentclass[11pt,a4paper]{article}
% The preamble shared by every transcription of the mathematicians' manuscripts in Gallica.
%
% It rests on one idea: **uncertainty compounds, it does not resolve.** A
% transcription that smooths over a doubtful word has destroyed the only thing
% it was made for — a reader must be able to tell, at every point, what was on
% the page and what was supplied. The apparatus macros below are therefore the
% critical apparatus, not decoration.
%
% The same commands are recognised by `scripts/render.mjs`, which produces the
% site's reading view, and by `scripts/tei.mjs`. A macro added here must be
% added there too, or rendering fails loudly — which is the wanted behaviour.
%
% Comments here are English, like the rest of the repository. The documents
% this preamble sets are French, and that is deliberate: see the skills.

\ProvidesPackage{germain}[2026/09/15 Transcriptions of mathematicians' manuscripts in Gallica]

\usepackage{amsmath,amssymb,amsthm}
\usepackage{mathrsfs}
% Kept from the parent preamble so the renderer's diagram support stays
% available; these manuscripts draw no commutative diagrams, and nothing here uses it.
\usepackage{tikz-cd}
\usepackage[english,french]{babel}
\usepackage{fontspec}

% --- Greek ------------------------------------------------------------------
% The text face stays fontspec's default, Latin Modern, which has no Greek:
% XeTeX drops every glyph it lacks with a line in the log and nothing on the
% page, and the Greek words the manuscripts quote vanished from the PDFs. So
% the Greek and Greek Extended blocks get a XeTeX character class of their own,
% and entering or leaving a run of it switches the family to CMU Serif — the
% same Computer Modern design, with polytonic Greek — and back. Only the
% family changes: a Greek word in \textit or \textbf keeps its shape and series.
% The transitions to and from every other class carry the tokens that class
% has towards an ordinary letter, so French spacing before « ; » or inside
% « » still holds after a Greek word. No grouping, so no brace or \endgroup
% can come out unbalanced; maths is left alone, since XeTeX inserts nothing
% there. Kept identical in `exercices.sty`.
\makeatletter
\newfontfamily\germainGreekfont{cmun}[
  Extension=.otf, NFSSFamily=germaingreek,
  UprightFont=*rm, ItalicFont=*ti, SlantedFont=*sl,
  BoldFont=*bx, BoldItalicFont=*bi, BoldSlantedFont=*bl]
\newXeTeXintercharclass\germain@greek
\def\germain@greekrange#1#2{%
  \@tempcnta=#1\relax
  \loop
    \XeTeXcharclass\@tempcnta=\germain@greek
    \ifnum\@tempcnta<#2\advance\@tempcnta\@ne\repeat}
\germain@greekrange{"0370}{"03FF}
\germain@greekrange{"1F00}{"1FFF}
\def\germain@greekprev{\rmdefault}
\def\germain@greekin{%
  \edef\germain@greekprev{\f@family}\fontfamily{germaingreek}\selectfont}
\def\germain@greekout{\fontfamily{\germain@greekprev}\selectfont}
\def\germain@greekjoin{%
  \XeTeXinterchartokenstate=\@ne
  \@tempcnta=\z@
  \loop
    \ifnum\@tempcnta=\germain@greek\else
      \XeTeXinterchartoks\germain@greek\@tempcnta=\expandafter{%
        \expandafter\germain@greekout\the\XeTeXinterchartoks\z@\@tempcnta}%
      \XeTeXinterchartoks\@tempcnta\germain@greek=\expandafter{%
        \the\XeTeXinterchartoks\@tempcnta\z@\germain@greekin}%
    \fi
    \ifnum\@tempcnta<4095\advance\@tempcnta\@ne\repeat}
% babel-french sets its own classes, and saves and restores the interchar
% state, each time a language is selected: join again after it does.
\addto\extrasfrench{\germain@greekjoin}
\addto\extrasenglish{\germain@greekjoin}
\AtBeginDocument{\germain@greekjoin}
\makeatother

\usepackage{geometry}
\geometry{a4paper,margin=25mm,marginparwidth=28mm}
\usepackage{marginnote}
\usepackage{xcolor}
\usepackage[normalem]{ulem}
\setlength{\emergencystretch}{3em}
\usepackage{hyperref}
\hypersetup{colorlinks=true,linkcolor=black,urlcolor=black,pdfborder={0 0 0}}

% --- Batch metadata ---------------------------------------------------------
% Taken as they stand by the rendering script, which shows them at the head of
% the reading view. They fix what the file claims to cover. The macro names are
% the parent project's, so that the scripts stay shared; \folder here means the
% volume's slug — `fr-9115` — and \pages counts Gallica views.

\newcommand{\folder}[1]{\gdef\grothFolder{#1}}
\newcommand{\batch}[1]{\gdef\grothBatch{#1}}
\newcommand{\pages}[2]{\gdef\grothFirst{#1}\gdef\grothLast{#2}}
\newcommand{\foldertitle}[1]{\gdef\grothFoldertitle{#1}}

% The holder's own shelfmark, printed as they write it: « Français 9115 ».
\newcommand{\shelfmark}[1]{\gdef\grothShelfmark{#1}}

% The dating, copied verbatim from the holder's catalogue — never inferred
% here. « XIXe siècle » is the BnF's; a pass that dates a leaf from its content
% says so in a \note{}, as its own inference.
\newcommand{\dating}[1]{\gdef\grothDating{#1}}

% The status notice, printed in the title block and repeated as a diagonal
% watermark on every page of the PDF. The manuscripts are in the public
% domain; what this line declares is not a right but a quality — an
% unverified first pass by a machine. The skills write the line; a file
% without it gets no watermark, so leaving it out is visible in the output.
\newcommand{\watermark}[1]{\gdef\grothWatermark{#1}}

\gdef\grothFolder{?}\gdef\grothBatch{}\gdef\grothFirst{?}\gdef\grothLast{?}%
\gdef\grothFoldertitle{}\gdef\grothShelfmark{}\gdef\grothDating{}\gdef\grothWatermark{}

% --- The résumé -------------------------------------------------------------
\newenvironment{resume}
  {\small\setlength{\parskip}{0.45em}}
  {\par\medskip\hrule\medskip}

% --- Keywords ---------------------------------------------------------------
% In English, closing the modernised reading's résumé: the modern vocabulary
% under which to search for what the leaves construct. The manifest extracts
% them to tag the volume — this line is the single source of the tags.
\newcommand{\keywords}[1]{\par\smallskip\noindent
  {\footnotesize\textsc{Keywords} --- \textit{#1}}\par}

% --- The critical apparatus -------------------------------------------------

% The Gallica view number, in the margin. This is the anchor that lets a
% disagreement over a formula be settled by looking at *one* image rather than
% a volume of 750. The site uses it to turn the facsimile as the transcription
% is scrolled. A view is one scanned image: usually one side of a leaf.
\newcommand{\page}[1]{%
  \marginnote{\footnotesize\color{gray!70}\textsf{v.\,#1}}%
  \label{page:#1}\ignorespaces}

% The BnF's pencilled foliation, where the leaf carries it and a pass can read
% it: « 198r ». This is what the literature cites — « ff. 198r–208v » — and
% the one line that turns a citation into a view. Recorded, never inferred:
% a folio a pass cannot read is not written.
\newcommand{\folio}[1]{%
  \marginnote{\footnotesize\color{gray!50}\textsf{f.\,#1}}\ignorespaces}

% The modernised reading's coarser counterpart to \page{}: which views a
% section is drawn from.
\newcommand{\pagerange}[2]{%
  \marginnote{\footnotesize\color{gray!70}\textsf{v.\,#1--#2}}%
  \ignorespaces}

% Illegible: never guessed.
\newcommand{\ill}{\textcolor{red!60!black}{[\,\ldots\,]}}

% A doubtful reading: the word is given, and flagged as uncertain.
\newcommand{\uncertain}[1]{\uline{#1}}

% An editorial addition — a missing letter, an implied word.
\newcommand{\add}[1]{\textcolor{blue!50!black}{[#1]}}

% Struck out by the writer. What was crossed out is part of the document.
\newcommand{\struck}[1]{\sout{#1}}

% The transcriber's note, on the state of the page or the mathematics.
\newcommand{\note}[1]{\par\smallskip{\small\color{gray!60!black}\textit{[#1]}}\par\smallskip}

% A marginal note of hers, distinct from the body of the text.
\newcommand{\marginal}[1]{\par\smallskip\noindent\hspace{1em}%
  {\small\color{gray!60!black}#1}\par\smallskip}

% --- Title ------------------------------------------------------------------
% The title block is French, like the documents themselves.

\AddToHook{shipout/background}{%
  \ifx\grothWatermark\empty\else
    \begin{tikzpicture}[remember picture, overlay]
      \node[rotate=52, scale=3.4, align=center, text=gray!18,
            font=\sffamily\bfseries]
        at (current page.center) {\grothWatermark};
    \end{tikzpicture}%
  \fi}

\AtBeginDocument{%
  \begin{center}
    {\large\bfseries Manuscrits de mathématiciens (Gallica) ---
      \ifx\grothShelfmark\empty\grothFolder\else\grothShelfmark\fi}\\[2pt]
    {\itshape\grothFoldertitle}\\[4pt]
    {\small vues \grothFirst--\grothLast
      \ifx\grothBatch\empty\else\ (lot \grothBatch)\fi}\\[2pt]
    \ifx\grothDating\empty\else
      {\small\color{gray!60!black}Datation du catalogue : \grothDating}\\[2pt]
    \fi
    \ifx\grothWatermark\empty\else
      {\small\bfseries\color{red!45!gray}\grothWatermark}\\[2pt]
    \fi
    {\footnotesize\color{gray!60!black}Transcription établie d'après les images IIIF de Gallica.
     Source gallica.bnf.fr / Bibliothèque nationale de France.\\
     Les lectures incertaines sont soulignées, les illisibles notées [\,\ldots\,],
     les ajouts entre crochets ; \textsf{v.} numérote les vues de Gallica,
     \textsf{f.} la foliotation au crayon de la BnF.}
  \end{center}
  \vspace{1em}}

\endinput


\folder{fr-9115}
\shelfmark{Français 9115}
\pages{1}{750}
\dating{1801-1900}
\watermark{Lecture automatique\\interprétation non vérifiée}
\foldertitle{Papiers de Sophie GERMAIN. — Recueil de dissertations et problèmes mathématiques et physiques — lecture modernisée du volume entier}

\begin{document}

\section*{Résumé}

\begin{resume}
Ces sept cent cinquante vues ne forment pas un livre. C'est un volume relié sans
tri, où des pièces de Sophie Germain sur des sujets différents se suivent sans
ordre de sujet ni de date, avec quelques pièces d'autres mains. La plus grande
masse traite des vibrations des plaques élastiques, la question que l'Institut
mit au concours et pour laquelle l'Académie couronna son mémoire de janvier 1816 : un
exposé de ses principes, un long mémoire de cent soixante-cinq pages et plusieurs
états de son ouverture, une copie latine du mémoire d'Euler de 1779 sur les lames
vibrantes, et une critique de ce mémoire en trois états. Autour, la géométrie qui
porte sa théorie : un mémoire sur la courbure des surfaces et ses brouillons. Plus
loin, une quarantaine de pages sur la propagation de la chaleur dans la sphère, le
cône et le cylindre, à la suite de Fourier. À la fin, de l'arithmétique : trois
pages sur le théorème de Fermat pour les exposants pairs, que la littérature
appelle le « manuscrit C », des notes sur les résidus quadratiques, un mémoire
imprimé de Legendre, et des extraits des Mémoires de Turin. Le théorème qui porte
son nom, sur les nombres premiers auxiliaires, n'est pas sur ces feuillets.

Le problème des plaques est posé par une expérience. Chladni répand du sable sur
une plaque de verre et la fait vibrer à l'archet : le sable quitte les parties qui
vibrent et se range le long des lignes immobiles, les lignes nodales, qui
dessinent des figures régulières, chacune avec son son. Pour une lame, plaque
étroite qu'on peut traiter comme une ligne, Euler avait l'équation du mouvement,
du quatrième ordre, et les conditions aux bouts ; pour une plaque, on n'avait
rien. La question de l'Institut était : « Donner la théorie mathématique des
vibrations des surfaces élastiques et la comparer à l'expérience ».

Son idée est une analogie. Dans la lame, la force élastique est proportionnelle à
la courbure ; dans une surface, elle la suppose proportionnelle à la somme des
deux courbures principales, c'est-à-dire, en termes actuels, au double de la
courbure moyenne. Il en sort, par une dérivation qui n'est pas une
démonstration, une équation où figure le bilaplacien,
$\partial_t^2 z + b\tau^2\,\Delta^2 z = 0$, qui est l'équation actuelle des plaques
minces. Elle la défend contre l'hypothèse d'un rival qu'elle ne nomme pas, « un
membre de la classe » : les deux hypothèses diffèrent par un terme de courbure de
Gauss, qui, pour une plaque naturellement plane, ne change pas l'équation à
l'intérieur et change les conditions au bord libre. Pour justifier le choix de la somme, elle écrit un
mémoire où elle veut « signaler et définir un genre de quantités dont l'existence
ne paroît pas avoir été soupçonné » : la courbure moyenne, moyenne des courbures
normales dans toutes les directions, dont la répartition autour d'un point est
« calquée » sur la loi de la chaleur dans l'anneau de Fourier. Pour les figures de
Chladni, elle invente les « limites analytiques » : les lignes où les conditions du
bord sont remplies découpent la plaque en plaques plus petites, et des sommes de
fonctions de la lame d'Euler, ou des produits de sinus, rendent compte, dit-elle,
de plus de la moitié des figures de la plaque carrée. Contre Euler, elle conteste
la seconde solution de la lame appuyée à ses bouts et en son milieu, en supposant
que l'appui du milieu n'exerce aucune force. Sur la chaleur, elle étend les solutions de Fourier au
solide échauffé « d'une manière entièrement arbitraire ». En arithmétique, elle
cherche à démontrer, par des décompositions en sommes de deux carrés, que
$x^{2n}+y^{2n}=2z^{2n}$ est impossible pour $n$ impair, et en tire le théorème de
Fermat pour les exposants pairs.

Elle dit elle-même ce qu'elle tente, et ce qu'il lui en coûte. L'Exposition
s'ouvre sur ce que l'autre hypothèse « a pour appui un nom justement célèbre »,
« une forte raison de se défier de celle qui m'appartient » ; son point de vue,
écrit-elle, « est entièrement neuf ; et c'est pour moi un grand sujet de
défiance », et elle finit : « Un pareil prestige n'est pas de mon côté ». Devant
Euler, un brouillon qu'elle barre ensuite dit : « Je n'ose admettre ce resultat
contraire a l'opinion d'Euler » ; la mise au net, qu'elle n'admet pas une opinion
si contraire « sans la soumettre à un examen plus approfondi ». Le long mémoire
s'adresse à « la classe », dont il espère avoir rempli le programme « au gré de mes
juges », et se ferme sur ces mots : « Je termine enfin ce long mémoire dans lequel
je n'ai pourtant examiné qu'une partie des différens cas que l'expérience
présente ». Un brouillon sur les surfaces courbes avoue : « Je me suis résolue
beaucoup trop tard à reprendre le sujet qui fait l'objet du présent mémoire ».

Où cela mène. Son équation est celle de la théorie des plaques de Kirchhoff
(1850) ; mais ses conditions au bord libre sont celles de Kirchhoff pour un
coefficient de Poisson $\nu = 1$, et celles de son rival pour $\nu = -1$ : la loi
retenue aujourd'hui est entre les deux. Ni ses produits de sinus ni ses sommes de
fonctions de la lame ne satisfont, pour une matière réelle, les conditions d'un
bord libre, et
c'est la cause qu'on donne aujourd'hui à l'écart d'environ un ton qu'elle
constate entre deux familles de figures ; les figures de la plaque libre n'ont été
calculées qu'avec la méthode de Ritz (1909), qui part des mêmes fonctions. Sur la
lame appuyée en son milieu, Euler avait raison contre elle ; sur l'anneau, Chladni avait
raison contre Euler et contre elle. Les séries qu'elle écrit pour la plaque
circulaire et pour le cylindre chauffé sont des fonctions de Bessel ; sa sphère
échauffée arbitrairement demande les harmoniques sphériques. Sa démonstration du
manuscrit C a des lacunes, et ses deux énoncés, corrigés de la solution triviale,
n'ont été démontrés que bien plus tard (Darmon et Merel, Wiles). Les noms à
chercher : équation des plaques de Sophie Germain (ou de Germain--Lagrange),
bilaplacien, conditions aux limites de Kirchhoff, lagrangien nul, figures de
Chladni, poutre d'Euler--Bernoulli, anneau vibrant, courbure moyenne et formule
d'Euler, loi de Laplace, fonctions de Bessel, équation de la chaleur en
coordonnées sphériques et cylindriques, théorème de Fermat pour les exposants
pairs, résidus quadratiques et loi de réciprocité, racines primitives des nombres
premiers de Fermat.

\keywords{Sophie Germain's plate equation, biharmonic operator, Kirchhoff plate boundary conditions, null Lagrangian, Chladni figures, Euler–Bernoulli beam, beam on an intermediate support, vibrating ring, Bessel functions, mean curvature, Euler's curvature formula, Young–Laplace equation, heat equation in spherical coordinates, heat equation in cylindrical coordinates, Fermat's Last Theorem for even exponents, quadratic residues, quadratic reciprocity, primitive roots of Fermat primes, cyclotomic polynomials, vibrating string controversy}
\end{resume}

\section*{Ce que le volume contient}

\pagerange{1}{750}
Sept cent cinquante vues : la reliure (vue 1) ; un feuillet de garde qui porte,
d'une autre main, plus grande, « Manuscrits de mémoires imprimés » (vue
2)\footnote{Ce libellé invite à chercher lesquelles des pièces du volume
répondent à des mémoires imprimés de l'auteur. Cette lecture ne l'a pas fait, et
ne s'en sert ni pour identifier une pièce ni pour combler une lacune.} ; puis des
feuillets de toutes tailles, montés, souvent sur onglets, pièce après pièce, sans
ordre de sujet ni de date. Le volume n'est pas trié, et cette lecture ne prétend
pas y trouver un fil.

Le catalogue de la Bibliothèque compte 372 feuillets. Une série de nombres tracés à
l'encre en haut à droite, un par feuillet, va de 1 (vue 2) à 372 (vue 748), et un
« 373 » barré est porté sur le feuillet blanc de la vue 750. Elle compte les petits
feuillets volants et la plupart des feuillets blancs ; elle porte un « 55 bis »
(vue 113) et laisse sans numéro le feuillet blanc des vues 732--733. Si l'un et
l'autre sont des feuillets, le volume en a 374 : la lecture le constate, et ne le
concilie pas avec le chiffre du catalogue. Cette série est tracée à la plume, et
les transcriptions ne la donnent pas pour une foliotation ; aussi la lecture
cite-t-elle des vues, et non des folios, avec une seule exception : aux vues 708 à
710, les nombres 348 et 349 de la série répondent exactement aux ff. 348r--349r où
Laubenbacher et Pengelley situent le « manuscrit C ».

Les pièces, dans l'ordre de la reliure :
\begin{itemize}
\item \emph{« Exposition des principes qui servent de base à la théorie des
surfaces élastiques »} (vues 3 à 41), mise au net continue, qui s'ouvre sur son
histoire du prix.
\item \emph{« Mémoire sur la courbure des surfaces »} (vues 43 à 89), paginé 1 à
44, avec une planche ; puis un feuillet de titre, « Moyennes courbures », et les
brouillons d'une introduction (vues 91 à 114).
\item \emph{Un mémoire imprimé de Legendre} sur l'équation
$4(x^n-1)=(x-1)(Y^2\pm nZ^2)$, lu à l'Académie le 11 octobre 1830, collé page à page
sur des feuillets de montage (vues 115 à 134).
\item \emph{« Idées à vérifier »}, sur la surface pesante et la chaînette (vues
136 à 140), et une page qui rend en français l'ouverture d'un texte de Gauss sur
la courbure (vue 142).
\item \emph{La copie latine du mémoire d'Euler} « Investigatio motuum quibus
laminæ et virgæ elasticæ contremiscunt » (Acta de Pétersbourg pour 1779), pages
(103) à (161) de l'imprimé (vues 144 à 206).
\item \emph{La critique du § 47 de ce mémoire}, en trois états : un premier
brouillon (vues 210 à 216), un second brouillon paginé 1 à 45 (vues 218 à 316),
une mise au net paginée (1) à (35) et inachevée (vues 320 à 389), sous deux
feuillets de titre (vues 208 et 318) ; puis deux feuillets de figures et deux
petits calculs (vues 391 à 396).
\item \emph{Le long mémoire sur les surfaces élastiques}, « sur la question
proposée par la première classe de l'Institut », paginé (1) à (165), mise au net
corrigée (vues 398 à 592). Sont reliés parmi ses pages sans en faire partie : un
feuillet de calculs (vues 410--411) ; le petit feuillet de la vue 416, qui corrige
le calcul de la cire du § 6 du même mémoire ; deux petits feuillets de calcul
(vues 428 et 483) ; un fragment sur la mesure de la courbure (vue 465) ; un
brouillon sur la chaleur (vue 549) ; des petits feuillets sur la plaque circulaire
(vues 553 à 564) ; et une note plus tardive sur Fourier, Poisson et Navier (vue
575).
\item \emph{Les autres états du long mémoire} (vues 593 à 626) : une seconde mise
au net de son ouverture et de ses pages (23)--(24) (vues 593 à 597) ; un
troisième état de l'ouverture (vue 598), que la vue 614 continue mot pour mot, de
sorte que les vues 598 et 614 à 618 forment un seul état, selon toute apparence le
plus ancien, et non l'ouverture d'un autre mémoire ; reliée entre les deux, une
seconde rédaction de la comparaison des sons, pages (61) à (72) (vues 600 à 613) ;
enfin quatre pages sur des surfaces cylindriques, qui renvoient au « N.o 11 de mon
memoire pour 1813 » (vues 620 à 626).
\item \emph{La chaleur} (vues 630 à 701 et 712 à 716, avec la vue 549) : des
calculs sur le cône, la sphère et le cylindre ; « Équation du mouvement varié de
la chaleur dans une sphère solide échauffée d'une manière entièrement
arbitraire » (vues 652 à 670) ; la même pour le cylindre, en brouillon et en copie
incomplète (vues 672 à 695) ; un dernier état pour la sphère, le cône et le
cylindre (vues 696 à 701). Au milieu, une feuille sur la somme des courbures de la
sphère (vue 650), qui n'est pas de la chaleur.
\item \emph{L'arithmétique} (vues 704 à 730) : des notes sur les résidus
quadratiques, et, aux vues 708 à 710, le manuscrit C. Les feuillets de la chaleur
des vues 712 à 716 sont reliés parmi elles.
\item \emph{Des extraits des Mémoires de Turin} : Lagrange, Euler, Condorcet (vues
734 à 749). La vue 750 est un feuillet blanc.
\end{itemize}

La lecture reprend ces pièces dans cet ordre, à trois exceptions près : elle lit
ensemble les feuillets d'un même sujet que la reliure a dispersés (les feuillets
épars sur la courbure, les pièces sur la chaleur) ; elle place les autres états
de l'ouverture du long mémoire entre les deux parties de celui-ci ; et elle lit
l'imprimé de Legendre avec l'arithmétique, après le manuscrit C.

\emph{Ce qui est relié hors d'ordre, et ce qui est photographié deux fois.} Dans la
copie d'Euler, la page (128) est reliée entre la page (125) et la page (126) (vue
167), ce que la copie signale elle-même. Dans le long mémoire, ses pages (9) à
(12) sont écrites deux fois (vues 406 à 409 et 412--413), et deux pages portent le
numéro (132) (vues 542 et 543). Les calculs sur la chaleur se lisent dans l'ordre
639, 636, 637 ; la « Note » de la vue 656 et un feuillet blanc sont reliés entre la
deuxième et la troisième page de la pièce sur la sphère ; la copie de la pièce sur
le cylindre n'a pas ses pages 3 à 6. Plusieurs billets montés à deux par feuille
sont numérotés de bas en haut (vues 553, 726 et 728). Sont des secondes prises de
vue d'une même page : 53 (de 51), 122 (de 121), 191 (de 190), 196--197 (de
194--195), 204 (de 203), 281 (de 280, béquet relevé), 302 (de 301), 359 (de 358),
431 (de 429), 559 (de 557), 594 (de 593), 611--612 (de 609--610).

\emph{Les mains.} La plupart des pièces sont de sa main, en deux écritures : une
mise au net régulière (l'Exposition, le mémoire sur la courbure, le long mémoire)
et une écriture rapide de travail (les brouillons, les calculs, les notes
d'arithmétique). La copie d'Euler et la mise au net de la critique sont d'une main
droite et régulière, que les transcriptions attribuent diversement ; cette lecture
ne tranche pas. La seconde mise au net des vues 593 à 597 est d'une écriture plus
grande et très régulière ; le manuscrit C, d'une main posée ; les extraits de
Turin, d'une main penchée que les feuillets ne désignent pas, et dont les remarques
entre crochets sont dites ici « du copiste ». Les transcriptions ne se prononcent
pas sur la main des « Idées à vérifier », écrites à la première personne.

\emph{Ce que les feuillets datent.} Aucune pièce n'est datée de sa main. Les
feuillets citent : le mémoire d'Euler de 1779, et sa \emph{Methodus} « imprimée en
1744 » ; les \emph{Leçons sur le calcul des fonctions} de Lagrange, « an 1806 », et
la nouvelle édition de sa \emph{Mécanique} ; « mon memoire pour 1813 » (vue 626) ;
les concours de janvier 1814 et de janvier 1816 (vues 3 à 9, 31) ; « le mémoire
couronné qui fut présenté en 1815 » et « mes recherches … publiées en 1821 » (vue
103) ; « cette hypothese que j'ai donné en 1811 » et un « N.o d'aout » d'un
périodique (vue 575) ; la théorie de la chaleur de Fourier, par articles et par
pages ; l'imprimé de Legendre porte « Lu à l'Académie, le 11 octobre 1830 ». Toute
date qu'une pièce ne porte pas est, dans ce qui suit, une inférence dite comme
telle. Le \emph{1801--1900} de l'en-tête est la datation du catalogue.

\emph{Ce qui n'est pas dans le volume.} Le théorème de Sophie Germain sur les
nombres premiers auxiliaires, et son grand plan sur le théorème de Fermat ; le
mémoire pour 1816, dont l'Exposition cite le N\textsuperscript{o} 15 ; les
planches d'Euler (Tab. II et III) ; la planche de Chladni « jointe à ce mémoire »
et la « fig. A » du long mémoire (vue 464) ; les pages 3 à 6 de la copie sur le
cylindre ; la suite des pages (72) de la seconde rédaction et (35) de la mise au net
de la critique.

\section*{« Exposition des principes qui servent de base à la théorie des surfaces élastiques »}

\pagerange{3}{41}
Le volume s'ouvre sur une pièce entière et continue : une mise au net, d'une
seule main régulière, de trente-neuf pages sans lacune, qui compare deux
manières de fonder la théorie des plaques et des surfaces élastiques — la
sienne, qui fait dépendre les forces d'élasticité de la somme des courbures
principales, et celle d'un rival qu'elle ne nomme jamais autrement que « un
membre de la classe » ou « le savant auteur du mémoire sur les surfaces
élastiques ».\footnote{Vues 3 à 41. Mise au net, écriture régulière, sans
rature notable. La série à l'encre numérote les feuillets 2 à 21 (vues 3 à
41) ; sur les rectos, un second nombre, biffé, accompagne le premier. Le
feuillet de garde de la vue 2, numéroté 1, porte d'une main plus grande et plus
grasse « Manuscrits de mémoires imprimés ». Ce libellé invite à chercher si
cette pièce et le mémoire sur la courbure qui la suit correspondent à des
mémoires imprimés de l'auteur ; la présente lecture ne l'a pas vérifié et ne
s'en sert pour rien. La vue 42 est le verso blanc du dernier feuillet.} Le
nom de ce rival n'est pas sur ces feuillets, et la lecture ne l'y met
pas.\footnote{Le feuillet de la vue 138, plus loin dans le volume, nomme « le
mémoire de M poisson sur la surface élastique » ; le titre est voisin de celui
que l'Exposition donne au mémoire de son rival (vues 6, 19, 25), mais
l'Exposition elle-même ne nomme personne, et l'identification n'est pas faite
ici.} Tout au long, les dérivées partielles sont écrites avec un $d$
droit.\footnote{Le feuillet écrit $\frac{d^2z}{dx^2}$, $\frac{d^4z}{dx^2dy^2}$ ;
on écrit $\frac{\partial^2 z}{\partial x^2}$, $\frac{\partial^4 z}{\partial x^2
\partial y^2}$. Elle pose aussi, « comme à l'ordinaire », $r = \frac{d^2
z}{dx^2}$, $s = \frac{d^2 z}{dx\,dy}$, $t = \frac{d^2 z}{dy^2}$ ; ces lettres
sont gardées dans le commentaire, avec $p = \partial z/\partial x$,
$q = \partial z/\partial y$.}

Les notations de la lecture sont celles de tout ce groupe : $\kappa_1 =
1/\rho$, $\kappa_2 = 1/\rho'$ les courbures principales de la figure forcée,
$1/R$, $1/R'$ celles de la figure naturelle, $H = \frac12(\kappa_1+\kappa_2)$
la courbure moyenne, $K = \kappa_1\kappa_2$ la courbure de Gauss. Sa « somme
des raisons inverses des deux rayons de courbures principales » est donc $2H$.

\subsection*{L'Avertissement : l'histoire de la question, racontée par elle}

\pagerange{3}{9}
Elle dit d'emblée ce qu'elle cherche : deux hypothèses « essentiellement
différentes » sont proposées, l'une « a pour appui un nom justement célèbre »,
ce qui est « une forte raison de se défier de celle qui m'appartient » ; elle
expose ses raisons pour que « le public éclairé les pèse », et sollicite
l'examen critique de la sienne.\footnote{Vues 3 et 4, première page de la pièce,
avec le cachet « Bibliothèque impériale — Mss » au bas de la vue 3.}

Puis l'histoire, dans ses mots. Les premières expériences de Chladni lui
firent croire que l'analyse pouvait en trouver les lois ; « un grand géomètre
dont les premiers travaux avoient été consacrés à la théorie du son » lui
apprit que la question contenait des difficultés insoupçonnées, et elle cessa
d'y penser.\footnote{Le feuillet ne nomme pas ce géomètre. La description
conviendrait à Lagrange, dont les premiers mémoires de Turin traitent du son,
et dont le volume contient justement des extraits sur ce sujet (vues 734 à
739) ; c'est une inférence de cette lecture, non une donnée du feuillet.} Le
séjour de Chladni à Paris la ramena au sujet ; elle étudia le mémoire d'Euler
sur le cas linéaire, « non pas certainement dans l'intention de concourir au
prix extraordinaire » que proposa l'Institut.\footnote{« extraordinaire » est
ajouté en marge (vue 4).} Dans le cas linéaire les forces d'élasticité sont
supposées proportionnelles à la courbure ; l'hypothèse qui met à sa place la
somme des deux courbures principales d'une surface « me frappa par son
analogie et par sa simplicité ».

Au concours de janvier 1814 elle présenta un mémoire contenant l'hypothèse,
l'équation de la plaque vibrante et, par quelques intégrales particulières, la
comparaison avec l'expérience ; l'Institut jugea la comparaison satisfaisante,
mais elle reconnaît qu'une « ignorance la moins excusable » l'avait empêchée de
déduire régulièrement l'équation de l'hypothèse, faute qui ne permettait pas à
la classe une approbation entière.\footnote{Vues 4 à 6. Une note au bas des
vues 5 et 6 cite deux passages du programme du prix (« p. 4 » et « p. 5 et
6 ») : on n'avait pas même « les équations différentielles du mouvement pour
cette espèce de vibration », et les simplifications d'Euler n'étaient pas
applicables aux surfaces métalliques. Elle en retient qu'après son mémoire on
était « en possession d'une hypothèse qui introduisoit deux dimensions dans le
calcul ». Que la pièce de janvier 1814 soit le long mémoire des vues 398 à 592
est une inférence tirée d'autres feuillets (vue 626, « mon memoire pour
1813 ») ; voir la section « Le mémoire sur la question proposée par la première classe de l'Institut » (vues 398 à 490) et la fin de la section « Les autres états de l'ouverture » (vues 620 à 626).} En août 1814, « un membre de
la classe » lut un mémoire sur les surfaces élastiques, fondé sur une hypothèse
nouvelle, borné aux surfaces naturellement planes, et parvenant à la même
équation ; il paraissait reconnaître qu'on y arrivait aussi par la somme des
inverses des rayons principaux. Pour le concours de janvier 1816 elle tenta de
justifier son hypothèse \emph{a priori} et de l'étendre à une surface de figure
quelconque ; l'Académie couronna ce mémoire, mais « a pourtant mis une
certaine restriction à son approbation », restriction dont elle dit n'avoir jamais bien
compris le sens. Le mémoire de son rival fut publié ensuite sans mention de
son dernier travail ; elle en conclut qu'il n'adoptait pas son sentiment, et
dit avoir « longtems attendu » la suite promise de ses recherches.

Elle annonce enfin la division de la pièce (examen \emph{a priori} ;
conclusions analytiques pour les plaques ; comparaison avec l'expérience) et
une correction : le signe du terme qui représente la courbure naturelle est
contraire à celui qu'elle avait adopté, « en suivant l'indication d'Euler
(\emph{de sono campanorum}) », dans le mémoire pour 1816, de sorte que les
différences de son observées sur des plaques avant et après leur courbure,
qu'elle avait dites contraires à sa théorie au N\textsuperscript{o} 15 de ce
mémoire, « y sont entièrement conformes ».\footnote{Vues 8 et 9. Le
N\textsuperscript{o} 15 cité est celui du mémoire pour 1816, que le volume ne
contient pas ; la lecture ne peut donc pas contrôler ce changement de signe.}

\subsection*{Deux genres d'hypothèses}

\pagerange{9}{11}
Les unes ne sont que l'énoncé de la question et transmettent leur évidence à
toutes leurs conséquences ; les autres sont un choix de point de vue fait pour
rendre l'analyse applicable, et leurs conséquences demandent une confirmation
venue d'ailleurs.\footnote{Vues 9 à 11 : « De la nature des hypothèses
mathématiques ». À la vue 10, le mot « observation » est écrit sur un grattage
et lu avec doute ; à la vue 11, « §§ » est lu avec doute.} Les mathématiques
pures n'admettent que les premières ; les mathématiques mixtes doivent souvent
employer les secondes, et les géomètres s'efforcent de s'en affranchir — ainsi
Biot a montré que la théorie des surfaces tendues ne dépend pas de la
supposition à l'aide de laquelle Euler l'avait d'abord établie. Elle place sa
thèse ici : son hypothèse serait du premier genre, « la définition complette de
l'état de la question », et celle de son rival ne s'appliquerait que là où elle
équivaut à la sienne. La suite de la lecture montre que ni l'une ni l'autre
n'est la loi que la théorie actuelle retient (voir plus bas).

\subsection*{Les deux hypothèses examinées a priori (N\textsuperscript{os} 1 à 6)}

\pagerange{12}{24}
\textbf{Les deux termes.} Elle prend la définition de l'élasticité de
d'Alembert (« propriété … au moyen de laquelle [les corps] se rétablissent
dans la figure et l'étendue que quelque cause extérieure leur avoit fait
perdre ») et en tire que les forces ont pour mesure la différence $E - N$ entre
la figure forcée et la figure naturelle ; leur terme dans l'équation des
vitesses virtuelles est $-(E-N)\,\delta(E-N)$.\footnote{Vues 12 et 13,
N\textsuperscript{o} 1 ; la citation vient de l'Encyclopédie méthodique,
article « élasticité ». Au N\textsuperscript{o} 3 le même terme est écrit
$-(N-E)\,\delta(N-E)$, lettres inversées ; la valeur est la même.} Pour
l'hypothèse de son rival, une surface d'épaisseur égale dont les points se
repoussent dans une sphère d'activité de rayon insensible, l'action se réduit à
détruire « l'irrégularité » $I$ de la surface, d'où un terme
$-I\,\delta I$.\footnote{Vue 13, N\textsuperscript{o} 2 ; « expansibilité » est
lu avec doute.}

\textbf{La courbure d'une surface (N\textsuperscript{o} 3).} Elle part du
passage d'Euler (Recherches sur la courbure des surfaces, Mémoires de Berlin
pour 1760, p. 119) où la question de la courbure d'une surface est dite « fort
équivoque », puisqu'il y a en chaque point une infinité de courbures, et se
propose de sommer les courbures de toutes les sections planes par le point. Elle
écarte d'abord les sections obliques par l'argument que deux sections inclinées
de $\theta$ et de $\pi - \theta$ sur la même tangente auraient des courbures
$\cos\theta/\rho_p$ et $\cos(\pi-\theta)/\rho_p$, de somme nulle. Cet argument
ne tient pas : par le théorème de Meusnier, la section oblique dont le plan
fait l'angle $\theta$ avec la normale a pour rayon $\rho_p\cos\theta$, et donc
pour courbure $1/(\rho_p\cos\theta)$, la même pour $\theta$ et pour
$\pi-\theta$ ; rien ne se compense.\footnote{Vue 15. Le feuillet écrit
$\frac{1}{R} = \frac{\cos\theta}{\rho_p}$, $\frac{1}{R'} =
\frac{\cos(\pi-\theta)}{\rho_p}$ et $\frac1R + \frac1{R'} = 0$ : la relation de
Meusnier y est renversée (le cosinus est au numérateur au lieu du
dénominateur), et la somme nulle vient de ce renversement joint à une
convention de signe. Meusnier (mémoire de 1776, publié en 1785) n'est pas cité
par le feuillet ; ce sont les « relations connues » qu'elle invoque.} Ce qui
est vrai, et qui justifie la conclusion qu'elle en tire, c'est que la courbure
d'une section oblique est déterminée par celle de la section normale de même
tangente : seules les sections normales portent de l'information.

Pour celles-ci elle écrit la formule d'Euler : si $\varphi$ est l'angle d'une
section normale avec la direction principale de courbure $\kappa_1$,
\[
\kappa(\varphi) = \kappa_1\cos^2\varphi + \kappa_2\sin^2\varphi ,
\qquad
\kappa(\varphi) + \kappa\!\left(\varphi + \tfrac{\pi}{2}\right) = \kappa_1 +
\kappa_2 = 2H ,
\]
la somme des courbures de deux sections normales perpendiculaires ne dépendant
pas de leur position.\footnote{Vue 16. Le feuillet mesure l'angle $A$ à partir
du plan de la courbure $1/\rho$ et écrit $\frac{1}{\rho_p} =
\frac{1}{\rho}\sin^2 A + \frac{1}{\rho'}\cos^2 A$ : sinus et cosinus y sont
échangés (pour $A = 0$ on doit retrouver $1/\rho$). La somme, seule utilisée, ne
change pas.} Elle conclut que la somme sur toutes les sections est « la somme
des raisons inverses des deux rayons de principales courbures prise une
infinité de fois », l'infini n'étant qu'une répétition de la même mesure. Une
somme infinie de cette sorte n'a pas de sens ; son énoncé correct est une
moyenne : la valeur moyenne de $\kappa(\varphi)$ sur toutes les directions est
$H$. C'est exactement ce que le mémoire sur la courbure établira plus loin par
ses « surfaces des distances ». Pour la figure naturelle et la figure forcée,
elle pose $N = \frac1R + \frac1{R'}$ et $E = \frac1\rho + \frac1{\rho'}$, d'où
son terme
\[
-\bigl(E - N\bigr)\,\delta\bigl(E - N\bigr),
\qquad E - N = \Bigl(\frac1\rho + \frac1{\rho'}\Bigr) - \Bigl(\frac1R +
\frac1{R'}\Bigr).
\]

\textbf{L'irrégularité (N\textsuperscript{o} 4).} Elle mesure l'irrégularité
d'une surface par l'écart de courbure entre sections normales perpendiculaires.
En coordonnées où les dérivées premières sont nulles au point, elle écrit
\[
\frac{1}{\rho_p} - \frac{1}{\rho'_p} = (r-t)\cos 2A + 2s\sin 2A ,
\]
et pour le couple complémentaire (angle $\frac\pi2 - A$) $-(r-t)\cos 2A + 2s\sin
2A$ ; la somme des deux écarts vaut $4s\sin 2A$ et s'annule quand les axes
sont les directions principales ($s = 0$).\footnote{Vues 17 à 19. Les formules
sont justes (linéarisées). À la vue 19, « $\rho_p - \rho'_p = (r-t)\cos 2A = 0$ »
pour $A = 45^\circ$ porte sur les rayons là où le calcul porte sur leurs
inverses ; et le terme final est écrit $-(\rho-\rho')\,\delta(\rho-\rho')$, sans
les raisons inverses, que la vue 20 rétablit.} Elle en conclut que la seule
différence qui ne disparaît pas est celle des courbures principales, et prend
pour terme de son rival $-(\frac1\rho - \frac1{\rho'})\,\delta(\frac1\rho -
\frac1{\rho'})$. Le raisonnement par couples est une heuristique ; ce qui est
vrai et le remplace est la décomposition
\[
\kappa(\varphi) = H + \frac{\kappa_1 - \kappa_2}{2}\cos 2\varphi ,
\qquad
\Bigl(\frac{\kappa_1-\kappa_2}{2}\Bigr)^2 = H^2 - K ,
\]
qui sépare en chaque point une partie isotrope, $H$, et une partie
anisotrope, d'amplitude $\frac12|\kappa_1-\kappa_2|$, nulle aux seuls points
ombilics. Prendre $\kappa_1 - \kappa_2$ pour mesure de l'« irrégularité » est
donc légitime.

En note, elle tire de la section à $45^\circ$ un théorème qu'elle juge
remarquable parce qu'il contredit la lettre d'Euler : la courbure d'une surface
« aura pour mesure la courbure de la sphère dont le rayon est égal au rayon de
la courbe formée par la section normale qui fait un angle de $45^\circ$ avec
les plans de plus grande et de moindre courbures ». C'est exact : $\kappa(\pi/4)
= H$, et la sphère en question, de rayon $1/H$, a la même somme de courbures
$2H$ que la surface.\footnote{Vue 19, note au bas de la page, appelée par un
astérisque ; « quelque » corrigé en marge en « quelle ».}

\textbf{Le texte du rival (N\textsuperscript{o} 5) et la comparaison
(N\textsuperscript{o} 6).} Elle identifie l'hypothèse de ses
N\textsuperscript{os} 2 et 4 à celle du « savant auteur du mémoire sur les
surfaces élastiques », dont elle cite le passage sur une force répulsive
sensible seulement à des distances imperceptibles, proportionnelle au produit
des épaisseurs, et bornée aux surfaces également épaisses.\footnote{Vues 19 et
20. Référence du feuillet : « M. de l'Institut 1812. seconde partie p 192 », et
« §§ 11 N\textsuperscript{o} 8 », lu avec doute ; elle renvoie aussi aux N\textsuperscript{os} 12
et suivants de ce mémoire, où la différence $\frac1\rho - \frac1{\rho'}$ est
introduite en prenant pour plans coordonnés ceux des courbures principales.}
Puis elle compare. Son terme vaut pour toute figure naturelle ; celui du rival
ne s'épuise que quand $\kappa_1 = \kappa_2$ : si la figure naturelle n'est pas
ombilicale, ses forces répulsives agiraient encore quand la surface a repris sa
figure naturelle, et l'hypothèse ne s'applique qu'aux surfaces dont les rayons
naturels sont « égaux et dirigés du même côté ».\footnote{Vues 20 à 25. « de l'hypothèse » (vue 22) et « pour
chacun » (vue 25) sont ajoutés dans l'interligne et lus avec doute ; à la vue 23, la phrase
« Cependans aucun des points intérieurs … se trouve animé » manque de la
négation attendue.} Pour la lame droite le choix est indifférent, les deux
différences valant $1/\rho$ ; pour la plaque plane, les points intérieurs
reçoivent de leurs voisins une action « sphérique » équivalente à des forces
propres, et le choix redevient indifférent pour eux, mais non pour les points
du bord, dont la sphère d'activité est incomplète.

En termes actuels, les deux densités d'énergie de flexion sont, à un facteur
de rigidité près,
\[
W_{\mathrm{G}} = \bigl(\kappa_1 + \kappa_2 - N\bigr)^2 ,
\qquad
W_{\mathrm{r}} = \bigl(\kappa_1 - \kappa_2\bigr)^2 .
\]
Son objection est juste pour ces deux densités : $W_{\mathrm{r}}$ n'est nulle
que sur une surface ombilicale, et une coque naturellement courbe mais non
sphérique ne serait pas en équilibre dans sa figure naturelle. Pour une plaque
naturellement plane ($N = 0$),
\[
W_{\mathrm{G}} - W_{\mathrm{r}} = (\kappa_1+\kappa_2)^2 - (\kappa_1-\kappa_2)^2
= 4\kappa_1\kappa_2 = 4K ,
\]
et tout se joue sur la courbure de Gauss (paragraphe suivant). La théorie
actuelle des plaques minces (Kirchhoff, 1850) donne pour densité
$(\kappa_1+\kappa_2)^2 - 2(1-\nu)\kappa_1\kappa_2$, où $\nu$ est le coefficient
de Poisson du matériau.\footnote{Kirchhoff (1850) est postérieur à ces feuillets ; le
coefficient $\nu$, qu'on nomme aujourd'hui coefficient de Poisson, n'y paraît
pas. L'énoncé complet des conditions de Kirchhoff au bord libre est donné une fois, pour tout le volume, dans la sous-section « Les limites de la plaque : côtés, coins, plaque ronde » du long mémoire (vues 433 à 438).} Son hypothèse est
le cas $\nu = 1$, celle de son rival le cas $\nu = -1$. Ce sont les deux bornes
de l'intervalle $-1 < \nu < 1$ où cette densité est définie positive, et
chacune est dégénérée : $W_{\mathrm{G}}$ s'annule pour une torsion pure
$\kappa_1 = -\kappa_2$ (la plaque $z = xy$ se tordrait sans effort), et
$W_{\mathrm{r}}$ pour une flexion sphérique $\kappa_1 = \kappa_2$. Les
matériaux usuels ont $\nu$ voisin de $0{,}3$ : ni l'une ni l'autre hypothèse
n'est la loi retenue aujourd'hui, et la vraie est entre les deux.

\subsection*{Ce qui les sépare : une variation exacte (N\textsuperscript{os} 7 et 8)}

\pagerange{21}{28}
Elle reprend une remarque de son rival : pour les termes qui restent sous le
double signe après les intégrations par parties, on a identiquement
\[
\delta\iint \frac{k\,dx\,dy}{\rho\rho'} = \iint \delta\Bigl(\frac{k}{\rho\rho'}
\Bigr)\,dx\,dy = 0 ,
\]
et puisque les deux hypothèses diffèrent par $\iint \delta(4k/\rho\rho')\,dx\,dy$,
elles donnent la même équation pour les points intérieurs ; elles diffèrent
pour les points du bord.\footnote{Vue 24. Le coefficient est lu $k$ dans la
première intégrale et $4k$ dans la seconde ; les deux lectures sont celles du
feuillet, et la seconde est la bonne, puisque la différence des deux densités
est $4\kappa_1\kappa_2$.} C'est exact, et c'est le point central de la comparaison. En
langage actuel, pour une surface graphe $z(x,y)$, la courbure de Gauss
linéarisée est $rt - s^2$, et l'intégrale $\iint (rt - s^2)\,dx\,dy$ est un
\emph{lagrangien nul} : son équation d'Euler--Lagrange est identiquement
nulle, et sa variation se réduit à des termes de bord. La propriété tient même
sans linéariser : $(rt - s^2)\,f(p,q)$ est un lagrangien nul pour toute
fonction $f$ (c'est le jacobien de l'application $(x,y) \mapsto (p,q)$ multiplié
par une fonction de $(p,q)$), ce qui couvre $K\,dA$, où $f = (1+p^2+q^2)^{-3/2}$,
aussi bien que $K\,dx\,dy$, où $f = (1+p^2+q^2)^{-2}$.\footnote{Vérifié par
calcul symbolique pour ces deux choix de $f$ et pour un $f$ arbitraire.
L'expression « lagrangien nul » est moderne ; le feuillet dit « on a
identiquement ».} Le pendant global est le théorème de Gauss--Bonnet : l'intégrale
de $K\,dA$ sur une portion de surface ne dépend que de son bord.\footnote{Gauss
(1827) pour la courbure totale, Bonnet (1848) pour la forme avec bord ; le
feuillet ne cite ni l'un ni l'autre, et la forme avec bord lui est
postérieure.} Donc, pour une plaque naturellement plane, les deux hypothèses
donnent la même équation intérieure ; si les bords sont encastrés
($\delta z = 0$ et $\delta\,\partial_n z = 0$), les termes de bord de la
différence s'annulent aussi et les deux problèmes coïncident entièrement ; aux
bords libres ils diffèrent. Pour une surface naturellement courbe, au
contraire, $W_{\mathrm{G}} - W_{\mathrm{r}} = 4K - 2N(\kappa_1+\kappa_2) + N^2$
n'est plus un lagrangien nul, et les deux hypothèses donnent des équations
intérieures différentes.

Elle cite ensuite son rival lui-même sur les points du contour : « détermination
très délicate, sur laquelle je me propose de revenir par la suite mais dont il
ne sera pas question dans ce mémoire ».\footnote{Vues 25 et 26, citation de
« p. 202 N\textsuperscript{o} 14 » du mémoire du rival.} Réduite à ses « seules lumières », elle
croit que l'hypothèse des forces répulsives conduit à regarder la plaque comme
une portion de sphère de rayon infini, donc comme infinie, et elle avance
l'idée qu'elle juge fondamentale : les variations de courbure ne sont que le
résultat des déplacements des points matériels, qui tendent chacun à reprendre
leur position naturelle ; « je supplie les géomètres de lui accorder
quelqu'attention ».\footnote{Vue 27 : « Je ne me dissimule pas que le point de
vue sous lequel j'ai envisagé la question est entièrement neuf ; et c'est pour
moi un grand sujet de défiance. »}

Au N\textsuperscript{o} 8, elle rapporte encore, d'après son rival, que la plaque
en équilibre rend stationnaire, parmi les surfaces de même étendue,
$\iint(\frac1\rho + \frac1{\rho'})^2\,k\,dx\,dy$, et étend l'énoncé à une surface
de figure quelconque avec $\iint (E-N)^2\,k\,dx\,dy$.\footnote{Vue 28 ; la page
citée, « N\textsuperscript{o} 23 p. 221 », est lue avec doute. Le feuillet dit « un
\emph{maximum} ou un \emph{minimum} ».} Qu'un partisan de $(\kappa_1-\kappa_2)^2$
puisse énoncer ce principe avec $(\kappa_1+\kappa_2)^2$ n'a rien de
contradictoire : les deux intégrales diffèrent de $4\iint K$, qui ne change
pas l'équation intérieure. Pour une surface naturellement courbe, en revanche,
l'extension proposée ne décrit pas la flexion des coques telle qu'on la
comprend aujourd'hui : l'énergie y dépend de la variation de toute la seconde
forme fondamentale, et non de la seule variation de sa trace. Une densité en
$(E-N)^2$ ne coûterait rien, par exemple, pour changer un cylindre de rayon $a$
en un cylindre de même rayon dont l'axe a tourné d'un angle droit au point
considéré ($E = N = 1/a$), ce qui est une flexion véritable.

\subsection*{L'équation de la plaque vibrante (N\textsuperscript{os} 9 à 11)}

\pagerange{29}{31}
Pour une plaque naturellement plane ($\frac1R + \frac1{R'} = 0$), son terme se
réduit, aux quantités du second ordre près, à $-\Delta z\,\delta(\Delta z)$
avec $\Delta z = r + t$.\footnote{Vue 29, « Conclusions analytiques déduites de
l'une et l'autre hypothèses ». Le feuillet écrit à la seconde accolade
$\delta\{\frac1\rho + \frac1{\rho'} + (\frac1R + \frac1{R'})\}$, avec un $+$ là où
la vue 28 écrit $-$ ; c'est sans effet ici, puisque ce terme est nul.} Deux
intégrations par parties donnent
\begin{align*}
\iint \Delta z\,\delta\Delta z\,dx\,dy
&= \int\Bigl(\Delta z\,\delta\frac{\partial z}{\partial x} - \frac{\partial \Delta
z}{\partial x}\,\delta z\Bigr)dy
 + \int\Bigl(\Delta z\,\delta\frac{\partial z}{\partial y} - \frac{\partial \Delta
z}{\partial y}\,\delta z\Bigr)dx \\
&\quad + \iint \Delta^2 z\,\delta z\,dx\,dy ,
\end{align*}
où $\Delta^2 z = \frac{\partial^4 z}{\partial x^4} + 2\frac{\partial^4
z}{\partial x^2\partial y^2} + \frac{\partial^4 z}{\partial y^4}$ est le
bilaplacien. Pour son rival, le terme linéarisé est
$-\frac12\,\delta\bigl[(r+t)^2 - 4(rt-s^2)\bigr]$, soit
\[
-(r+t)\,\delta(r+t) + 2\bigl(r\,\delta t + t\,\delta r - 2s\,\delta s\bigr) ;
\]
le feuillet porte $-2$ devant l'accolade, là où la différentiation donne
$+2$.\footnote{Vues 29 et 30. Le signe $-2$ est lu deux fois (vues 29 et 30).
Les conditions de bord qu'elle obtient à la vue 35 sont néanmoins celles du
signe correct (voir plus bas), de sorte que l'erreur est d'écriture plus que de
calcul. À la vue 30, un terme de la dernière accolade est biffé puis récrit
sans qu'aucune couche se laisse lire, et un signe est sous une tache d'encre ;
deux groupes y sont rayés.} Les termes sous le double signe sont les mêmes pour
les deux hypothèses, comme le paragraphe précédent l'établit ; multipliés par
une constante $n^2$ et joints au terme d'inertie, ils donnent
\[
\frac{\partial^2 z}{\partial t^2} + n^2\,\Delta^2 z = 0 ,
\]
l'équation de la plaque vibrante, « calculée d'après mon hypothèse » et qui « se
trouve dans la pièce que j'ai envoyée pour le concours de janvier
1814 ».\footnote{Vue 31. Le feuillet écrit $\frac{d^2z}{dt^2} +
n^2\bigl(\frac{d^4z}{dx^4} + 2\frac{d^4z}{dx^2dy^2} + \frac{d^4z}{dy^4}\bigr)
= 0$ ; on écrit $\partial_t^2 z + n^2\Delta^2 z = 0$. La même équation (A) est
au début du long mémoire (vues 399 à 404).} C'est bien l'équation actuelle de la
plaque mince en flexion. Le feuillet dit $n^2$ « proportionnel à l'épaisseur et
à l'élasticité naturelle » ; la théorie actuelle le donne proportionnel au \emph{carré} de l'épaisseur,
comme la constante $b\tau^2$ du long mémoire, dont le § 1 (vues 399 à 404) donne
le coefficient de Kirchhoff.\footnote{Vue 30, fin du N\textsuperscript{o} 11.}

Pour les termes de bord, elle note qu'ils disparaîtraient pour une plaque
infinie ; dans tout autre cas « il me reste a examiner commens la différence
entre les mêmes termes s'accomode avec l'experience ».\footnote{Vue 31 : « on
auroit $x$ et $y$ pareillement infinis et par conséquent $dx = 0$, $dy = 0$ ».
Ce qui est vrai : pour une plaque sans bord, ou dont le déplacement décroît à
l'infini, les termes de bord s'annulent.}

\subsection*{Le bord, les lignes nodales, la plaque chargée (N\textsuperscript{os} 12 à 15)}

\pagerange{31}{41}
Elle part de l'approbation de l'Institut « comme d'un point convenu » et ne
retient que les phénomènes qui touchent au choix de l'hypothèse.\footnote{Vues
31 et 32, N\textsuperscript{o} 12, « Comparaison entre les résultats du calcul et
les phénomènes observés ».} Sur un bord droit $x = \mathrm{const}$, son
hypothèse demande
\[
\Delta z\,\delta\frac{\partial z}{\partial x} - \frac{\partial\Delta z}{\partial
x}\,\delta z = 0 ,
\]
qu'on satisfait soit par $\Delta z = 0$ et $\delta z = 0$, soit par $\partial_x
\Delta z = 0$ et $\delta\,\partial_x z = 0$.\footnote{Vues 32 à 34, premier
N\textsuperscript{o} 13 (le numéro 13 sert deux fois : vues 32 et 35).} Avec
l'intégrale $z = \sin\frac{m\pi y}{A}\sin\frac{n\pi x}{A}$ de son mémoire de
1814, les droites $x = 0$, $x = A$ vérifient $z = 0$ et $\Delta z = 0$ : ce sont
les lignes nodales ; la droite $x = A/2$ vérifie $\partial_x z = 0$ et
$\partial_x\Delta z = 0$, pourvu que $n$ soit impair, et se meut parallèlement à
elle-même : elle l'appelle « libre », par analogie avec les extrémités libres
d'Euler.\footnote{Vue 34. Le feuillet écrit « $\cos n\frac\pi2 = 0$ » sans dire
que $n$ doit être impair ; les points $\ldots$ qu'il met devant l'intégrale
tiennent la place d'un facteur omis. Ce qu'elle appelle ici ligne libre est, en termes
actuels, un bord guidé (pente nulle, effort tranchant nul), non un bord libre ;
voir le § 3 du long mémoire (vues 414 à 432), où elle distingue « libre au sens d'Euler » et « libre au sens nouveau ».}

Elle ajoute que si la plaque est coupée suivant ces lignes, chaque partie rend
le même son et porte la partie correspondante de la figure nodale. C'est vrai à
condition que les bords coupés soient maintenus dans les conditions que la
ligne satisfaisait : appuyés sur les lignes nodales ($z = 0$, $\Delta z = 0$ :
appui simple sur un bord droit), guidés sur les lignes « libres ». Laissé
réellement libre, un bord coupé obéit à d'autres conditions, et le mode
$\sin\cdot\sin$ ne les satisfait pas : sur $x = 0$ il donne $\partial_x\Delta z
\neq 0$.

\textbf{L'hypothèse du rival au bord (second N\textsuperscript{o} 13).} Pour
l'autre hypothèse elle trouve sur le même bord un terme de plus,
$4\int \frac{\partial^2 z}{\partial x\partial y}\,\delta\frac{\partial
z}{\partial y}\,dy$ ; sur la ligne nodale, les conditions $\delta z = 0$ et
$\Delta z = 0$ entraînent $\partial_x^2 z - \partial_y^2 z = 0$, « mais elles ne
feront pas disparoître le terme » ; sur la ligne « libre », elles entraînent
$\partial_x(\partial_x^2 z + 3\partial_y^2 z) = 0$ et $\partial_x\partial_y z =
0$, et les deux hypothèses sont satisfaites par les mêmes
conditions.\footnote{Vues 33 à 35. Aux vues 33 et 35 le signe qui joint les deux
dérivées secondes sous la dérivation est couvert d'un pâté d'encre ; le calcul
demande $-$ (inférence de cette lecture). Le passage final de la vue 35 est
remanié en deux couches ; la couche annulée disait notamment « cas dans lesquels
la ligne de repos ne satisfait plus aux conditions » et « le choix de
l'hypothèse est indifférent ».} Ses conditions sont exactement celles de
Kirchhoff pour $\nu = -1$ : moment $\partial_x^2 z + \nu\,\partial_y^2 z =
\partial_x^2 z - \partial_y^2 z$, effort tranchant réduit $\partial_x^3 z +
(2-\nu)\,\partial_x\partial_y^2 z = \partial_x(\partial_x^2 z + 3\partial_y^2
z)$, et le terme de torsion $2(1-\nu)\int \partial_x\partial_y z\;\delta\,
\partial_y z\,dy$ de coefficient $4$. Le terme supplémentaire qu'elle isole est
donc le terme de torsion des conditions de Kirchhoff.\footnote{Vérifié par
calcul symbolique de la variation de $\frac12[(r+t)^2 - 2(1-\nu)(rt-s^2)]$. Pour
$\nu = 1$ on retrouve ses propres conditions, $\Delta z = 0$ et
$\partial_x\Delta z = 0$ au bord libre.}

Mais sa conclusion sur la ligne nodale est inexacte. Sur un bord $x =
\mathrm{const}$ où $\delta z = 0$ pour tout $y$, la dérivée tangentielle
$\delta\,\partial_y z = \partial_y(\delta z)$ est nulle aussi, et le terme de
torsion disparaît. Pour l'intégrale $\sin\cdot\sin$, les deux hypothèses
donnent donc les mêmes conditions sur les deux sortes de lignes, et ce mode
n'est pas un phénomène capable de les départager. Elles ne diffèrent qu'aux
bords réellement libres, où la valeur de $\nu$ entre dans les conditions.

\textbf{L'intégrale « linéaire » (N\textsuperscript{o} 14).} Elle rappelle
l'intégrale
\[
z = \mathcal{A}\,f(x) + B\,f(y), \qquad f(u) = a\,e^{\omega u/A} + b\,e^{-\omega
u/A} + c\sin\frac{\omega u}{A} + d\cos\frac{\omega u}{A},
\]
formée des fonctions d'Euler pour la lame, qui lui a servi à expliquer comment,
le son restant le même, des figures nodales « les plus bizares » remplacent les simples
droites.\footnote{Vues 35 à 37. Le feuillet note $A$ à la fois la longueur du
côté et le coefficient de $f(x)$ ; on écrit ici $\mathcal{A}$ pour le
coefficient, comme dans le long mémoire. À la vue 37 le signe entre $z$ et $0$
est sous un pâté d'encre ; la même condition est écrite $z = 0$ quelques
lignes plus haut.} Cette intégrale vérifie bien $\Delta^2 z = (\omega/A)^4 z$,
donc l'équation de la plaque avec un facteur périodique en temps, puisque le
terme croisé $\partial_x^2\partial_y^2$ s'annule sur une somme de fonctions
d'une seule variable. Les points où $f(x) = 0$ et $f(y) = 0$ à la fois sont
nodaux pour tous les rapports $\mathcal{A}/B$ : ses « points invariables » sont
réels dans cette famille. Mais la famille ne satisfait pas les conditions du
bord libre de la plaque : ni les siennes, puisque $\Delta z = B f''(y)$ subsiste
sur $x = \mathrm{const}$, ni celles de Kirchhoff pour une matière réelle, où le
moment de flexion vaut $\nu B f''(y)$ ; seul le cas limite $\nu = 0$ fait
exception (voir le § 4 du long mémoire, vues 439 à 468). L'accord avec les figures
de Chladni ne peut être que qualitatif.\footnote{La méthode qui, en 1909, a donné
de bonnes approximations des modes de la plaque carrée libre à partir de ces
mêmes fonctions de poutre est celle de Ritz ; elle est postérieure et n'est
nommée ici que pour situer la question.} Elle demande enfin comment des forces
répulsives, qui agissent sur tous les points de leur sphère d'activité,
deviendraient « linéaires » au point de n'agir que selon les fils croisés
qu'Euler imaginait pour les surfaces tendues : c'est une objection de sa part,
que la lecture rapporte.

\textbf{La plaque chargée (N\textsuperscript{o} 15).} Un corps mou appliqué
également le long d'un bord donne, dit-elle, la même figure et le même son que
si la plaque était réellement prolongée d'une quantité de même poids ; plus le
poids augmente, plus la ligne nodale voisine se rapproche du bord chargé ; elle
a obtenu un son juste d'une plaque chargée au point que l'intervalle entre deux
lignes nodales intérieures fût dix fois la bande comprise entre le bord chargé
et la ligne voisine, et elle a observé le même fait sur des surfaces
courbes.\footnote{Vues 38 à 40 ; « (M. pour janvier 1814) ». À la vue 39, la
phrase « parceque le mouvement des parties voisines devant alors sensible » est
inachevée sur le feuillet. Le long mémoire traite du même effet à son § 6
(vues 538 à 547), avec des chiffres ; le feuillet volant de la vue 416, qui
corrige un calcul sur la cire, se rapporte à ce § 6, non au présent
N\textsuperscript{o} 15, qui n'en donne pas (voir « § 6. L'effet d'un corps mou :
la cire »).} Elle l'explique par des forces propres à chaque point : la matière inerte
remplace des points élastiques sans changer la tendance des autres ; des forces
répulsives ne permettraient pas de le comprendre.

L'équivalence est juste au premier ordre et fausse au-delà : la
« substitution » d'une masse à une longueur n'est exacte que pour de petites
charges, et, pour une charge forte, le bord chargé se comporte de plus en plus
comme un appui, la ligne nodale voisine s'en approchant, ce qu'elle observe. La
raison, et un calcul sur la lame, sont donnés une fois, avec l'expérience chiffrée
du long mémoire, à la sous-section « § 6. L'effet d'un corps mou : la cire »
(vues 538 à 547).

La pièce se termine sur le poids des noms : « Un pareil prestige n'est pas de
mon côté, et loin delà, il est contre moi », et sur le souhait que les géomètres
ne se détournent pas « d'une théorie, si simple, si générale, si naturelle,
… et à laquelle il n'auroit manqué, pour paroître susceptible des plus belles
applications, que d'être tombée en des mains plus habiles ».\footnote{Vues 40 et
41 ; la phrase finale n'a pas de point, et un trait horizontal ferme la
pièce.}

\section*{« Mémoire sur la courbure des surfaces »}

\pagerange{43}{89}
Le mémoire qui suit est une seconde mise au net, de la même main, paginée par
elle de 1 à 44, suivie d'une planche de quatre figures ; son objet n'est plus
l'élasticité mais la géométrie qui la sous-tend : une quantité qu'elle dit
n'avoir pas été remarquée, la courbure moyenne d'une surface en un point, et la
loi selon laquelle la courbure se répartit autour de ce point.\footnote{Vues 43 à
89. Série à l'encre 22 à 44 (vues 43 à 89) ; pagination de l'auteur 1 à 44, les
numéros impairs biffés à partir de la page 17. La vue 52 montre seule un béquet
de six lignes sur la disparition des sections obliques dans la somme générale,
qui pourrait remplacer le passage biffé de la vue 49 ; la vue 53 est une
seconde photographie de la page de la vue 51. Le mémoire cite la \emph{Théorie
de la chaleur} de Fourier par numéros d'articles et de pages ; ce livre a paru
en 1822 (contexte, non tiré des feuillets), ce qui situe le mémoire après
cette date. Un « Mémoire sur la courbure des surfaces » de Sophie Germain a
paru imprimé en 1831 (contexte, non tiré des feuillets) ; si cette mise au net
en est la copie, la lecture ne l'a pas vérifié.} Le mémoire ne contient aucune
dérivée partielle ; ses outils sont la formule d'Euler et une construction.

\subsection*{Observations préliminaires : une quantité qu'on n'a pas remarquée}

\pagerange{43}{47}
C'est ici qu'elle dit le plus nettement ce qu'elle veut faire. Elle craint moins
la sécheresse du sujet que « la trop grande nouveauté de l'aspect » : elle se
propose « de signaler et de définir un genre de quantités dont l'existence ne
paroît pas avoir été soupçonné ».\footnote{Vues 43 à 47, ses pages 1 à 5. À la
vue 44, plusieurs mots sont repassés d'une encre plus grasse, sans rature.} Tant
qu'on regarde une surface pour elle-même, figure et courbure sont synonymes ;
mais quand la courbure entre en comparaison avec des quantités dynamiques,
résistance ou force, rien n'empêche que la force née de la courbure dépende non
de sa répartition autour du point mais de sa valeur moyenne, de sorte que deux
surfaces d'aspect entièrement différent posséderaient « des quantités égales
de courbures ». Elle annonce que la loi de répartition est « calquée » sur la
loi de la chaleur dans l'anneau de Fourier, et que l'analogie n'est pas
fortuite : de part et d'autre, une quantité moyenne est unie à un espace
circulaire. Elle annonce aussi la construction : la courbure de la sphère est
proportionnelle à la surface convexe d'un petit cylindre compris entre le plan
tangent et un plan parallèle ; celle d'une surface dont les rayons principaux
sont de même sens, à la surface convexe d'un segment cylindrique appuyé sur le
plan tangent et terminé par une courbe gauche à quatre parties symétriques ; les
deux surfaces convexes sont égales quand la sphère est la « sphère de courbure
moyenne ».\footnote{Vue 46 : la base du segment y est dite « d'ellipse » ; aux
vues 69 et 72, une correction de sa main établit qu'elle peut être prise pour un
cercle parfait, ce qui est exact à l'ordre considéré. La page 5 (vue 47) s'arrête
sur une phrase inachevée, écrite très petit au ras du bord, suivie d'un double
trait.}

\subsection*{L'ambiguïté de la courbure, et sa loi de distribution comparée à l'anneau de Fourier (N\textsuperscript{os} 1 à 3)}

\pagerange{48}{63}
Elle repart du mémoire d'Euler de 1760, qui « contient tout ce que l'on sait
d'important » : la courbure d'une surface est « fort équivoque », mais on peut
se borner aux sections normales, dont « l'assemblage de tous ces rayons … nous
donnera la juste mesure de la courbure ».\footnote{Vues 48 et 49, citations des
pages 119 et 120 du mémoire d'Euler.} Puis son histoire : quand elle proposa de
représenter la courbure par la somme des inverses des rayons principaux,
l'objection d'Euler lui fut opposée, comme si le choix des courbures principales
était arbitraire ; dans le mémoire pour le concours de 1816, elle prit la somme
sur toutes les sections et la réduisit à la somme principale « répétée un
nombre infini de fois » ; cette démonstration « fut formellement désapprouvée »,
et la légitimité de l'hypothèse resta contestée.\footnote{Vues 49 et 50. Deux
rédactions biffées se superposent à la vue 49 (« Cette somme se réduit
d'abord… », puis une reprise lue « Or, soit » avec doute), et une note
marginale verticale dit que la formule des sections obliques « renferme le
cosinus de l'inclinaison ». À la vue 50 un mot est couvert par une tache
d'encre ; à la vue 51, le chiffre abrégé après « Enfin » (3 ou 5) et deux amas
de mots biffés n'ont pas été lus. Elle ajoute à la vue 51 que les principales
conséquences de ces observations « ont été déjà publiées », sans dire où.} Ce
jugement, la lecture le partage pour la forme de l'argument (voir la section
précédente, N\textsuperscript{o} 3) ; le mémoire le remplace par un énoncé
correct.

\textbf{La formule (N\textsuperscript{o} 2).} Avec $f$ et $g$ le plus grand et le
moindre rayon principal, elle écrit
\[
r = \frac{2fg}{f + g - (f-g)\cos 2\varphi},
\qquad
\frac1r = \frac{f\sin^2\varphi + g\cos^2\varphi}{fg},
\qquad
\frac1{r'} = \frac{f\cos^2\varphi + g\sin^2\varphi}{fg},
\]
pour deux sections normales perpendiculaires.\footnote{Vue 54. Les formules sont
justes, mais pour $\varphi = 0$ elles donnent $r = f$, le plus grand rayon :
l'angle $\varphi$ est donc compté à partir du plan de \emph{moindre} courbure,
alors que le feuillet dit « avec celui de plus grande courbure ». Un ou deux mots
biffés avant « plan normal » n'ont pas été lus.} C'est la formule d'Euler,
$\kappa(\varphi) = \kappa_1\cos^2\varphi + \kappa_2\sin^2\varphi$ avec $\kappa_1 =
1/f$. Elle se fait à elle-même une objection : il s'agit de la courbure des
surfaces, et les formules employées sont celles de la courbure linéaire ; elle y
répondra au N\textsuperscript{o} 7.

\textbf{La loi de répartition (N\textsuperscript{o} 3).} Elle l'énonce en trois
propositions, chacune suivie de la proposition « correspondante » de Fourier sur
l'anneau.\footnote{Vues 55 à 63. Elle cite la \emph{Théorie de la chaleur},
« N\textsuperscript{o} 247 p. 277-27… » (vue 56 ; l'abréviation avant « IV » et les derniers
chiffres, perdus au bord du feuillet, ne sont pas lus) et « p. 278-279 »
(vue 58). La parenthèse insérée dans la citation de Fourier à la vue 58 est
d'elle.} Réunies en une formule :
\[
\kappa(\varphi) = H + \frac{\kappa_1-\kappa_2}{2}\cos 2\varphi .
\]

\begin{enumerate}
\item \emph{Deux plans normaux perpendiculaires portent une somme de courbures
constante}, égale à $\kappa_1 + \kappa_2$. Pour l'anneau : dans l'état qui
s'établit au bout d'un certain temps, la demi-somme des températures de deux
points diamétralement opposés est la même pour tout diamètre.\footnote{Vue 57.
Elle écrit « la demi somme des courbures sera une quantité $\frac1f +
\frac1g$ », sans le facteur $\frac12$ que la suite emploie ; la demi-somme est
$\frac12(\frac1f + \frac1g) = H$.}
\item \emph{La section à $45^\circ$ des plans principaux porte la courbure
moyenne} $\frac12(\frac1f + \frac1g)$, ainsi que sa perpendiculaire ; ces deux
plans partagent le contour du point en quatre régions « dont l'état est pareil
mais alternativement de signe opposé », l'excès sur la moyenne dans l'une égalant
le défaut dans la voisine. Pour l'anneau : les deux points à mi-chemin des
températures extrêmes ont la température moyenne, et partagent l'anneau en deux
moitiés de signe opposé.\footnote{Vues 57 à 59. « régions » est lu deux fois
avec doute ; à la vue 59, deux ou trois mots biffés après « dans les lois que
nous » ne sont pas lus.}
\item \emph{Il y a en général quatre directions autour du point où les courbes
normales sont semblables, deux seulement pour les courbures principales} :
$\kappa(\varphi) = \kappa(-\varphi) = \kappa(\pi \pm \varphi)$. Avec $\varphi =
\frac\pi4 + \varphi'$ et $\varphi = \frac{3\pi}{4} - \varphi'$, elle trouve
\[
\frac1r = \frac{f + g + (f-g)\sin 2\varphi'}{2fg},
\]
et les deux plans équivalents font entre eux l'angle $\frac\pi2 - 2\varphi'$.
Pour l'anneau, deux points de même température sont séparés par $\pi -
2\varphi'$.\footnote{Vues 60 à 62. Le feuillet écrit $\sin(45^\circ + \varphi') =
\frac12(\cos\varphi' + \sin\varphi')$ et trois identités semblables ; le facteur
exact est $\frac{\sqrt2}{2}$, et c'est lui que le calcul suivant emploie. Pour
$\varphi' = 45^\circ$ elle écrit que le second membre « devient $\frac1f$ » ; il
devient $\frac{2f}{2fg} = \frac1g$ (la direction obtenue, à $90^\circ$ du point
de départ, est l'autre direction principale). Elle reconnaît à la vue 61 que
cette troisième proposition n'est pas textuellement dans Fourier mais « résulte
des formules qui y sont données ».}
\end{enumerate}

Où l'analogie tient et où elle n'est qu'analogie. Les deux lois sont de la même
forme : une valeur moyenne plus un seul harmonique du cercle. La répartition
de la courbure est exactement $H + \frac12(\kappa_1-\kappa_2)\cos 2\varphi$,
harmonique d'ordre $2$, d'où le quart de circonférence entre les extrêmes. La
température de l'anneau est, dans l'état final, la moyenne plus l'harmonique
d'ordre $1$, $\bar v + a\,e^{-\lambda t}\cos(x - x_0)$, d'où la demi-circonférence ;
mais cet état n'est atteint qu'asymptotiquement, les harmoniques supérieurs
s'éteignant plus vite. La première proposition est exacte à tout instant pour la
courbure ; pour l'anneau, la demi-somme de deux températures diamétralement
opposées ne vaut la moyenne que si les harmoniques pairs sont absents, ce qui
n'est vrai qu'à la limite. Et les deux lois n'ont pas la même origine : l'une
vient d'une forme quadratique (la seconde forme fondamentale), l'autre d'une
équation d'évolution. L'analogie est donc une analogie de forme, qu'elle
explique elle-même par « l'emploi des quantités moyennes … uni à la
considération d'un espace circulaire », et qui est juste en ce sens-là.

Une conséquence : la sphère de rayon $1/H$ tangente au point, centrée sur la
normale, a ses quatre quadrants alternativement au dedans et au dehors de la
surface ; les sphères de rayons compris entre les deux rayons principaux coupent
la surface, et celles qui ont un rayon principal ne la coupent pas.\footnote{Vues
62 et 63. Ces énoncés portent sur les termes du second ordre ; pour les sphères
de rayon principal, le contact est d'ordre supérieur le long d'une direction, et
c'est l'ordre suivant qui déciderait. À la vue 62, un amas biffé après « l'arc »
n'est pas lu.} Elle ferme ce numéro par une réflexion
sur l'ordre nécessaire des choses : l'existence des quantités moyennes « forme le lien
entre l'uniformité et la variété », et un phénomène qui se développe autour
d'un centre donne, « par la substitution de l'espace au tems, une construction
de la périodicité ».\footnote{Vues 63 et 64.}

\subsection*{Ligne de distance, surface des distances, surface des distances moyennes (N\textsuperscript{os} 4 à 8)}

\pagerange{64}{84}
\textbf{La ligne de distance (N\textsuperscript{o} 4).} On juge qu'un arc est
plus courbe quand son extrémité s'écarte davantage de la tangente menée à
l'autre extrémité ; elle nomme \emph{ligne de distance} la perpendiculaire
abaissée sur la tangente, égale au sinus verse. Pour un arc de longueur $s$ sur
un cercle de rayon $\rho$,
\[
d = \rho\Bigl(1 - \cos\frac{s}{\rho}\Bigr) = \frac{s^2}{2\rho} + O(s^4),
\]
de sorte que, à longueur d'arc égale et infiniment petite, la ligne de distance
est proportionnelle à la courbure.\footnote{Vues 64 à 66. Elle prend des rayons
$\rho m$, $\rho n$ et des arcs $\varphi n$, $\varphi m$ de même longueur. Le
feuillet porte « l'arc le plus courbe $m\rho$ » là où la définition demande
$m\varphi$ (vue 65), et l'argument $(\varphi m)$ dans la seconde paire
d'équations là où le développement qui suit, $d = \frac12 m n^2\rho\varphi^2$,
demande $(\varphi n)$ (vue 66). Le résultat, $D$ et $d$ en raison inverse des
rayons, est juste.} Cette manière de parler serait « bizarre et surabondante »
pour une courbe, dit-elle ; elle sert pour les surfaces.

\textbf{Les deux surfaces (N\textsuperscript{os} 5 et 6).} Autour du point, on
prend sur chaque section normale un arc de même longueur $s$ ; la courbe qui
joint leurs extrémités est la \emph{directrice} ; la ligne de distance, qui
change de grandeur en la parcourant, engendre la \emph{surface des distances},
segment cylindrique qui a pour base, dans le plan tangent, une petite courbe
fermée.\footnote{Vues 66 à 69. La question est posée à la vue 67 en ces termes :
« trouver l'expression de la distance moyenne entre le plan tangent et les points
de la surface qui environnent le point de tangence » ; un astérisque et un mot au
crayon, d'une autre main, sont en marge. Le mot sous « assujettis » est perdu
sous la correction ; un mot biffé après « de l'arc » (vue 68) n'est pas lu.
L'addition marginale de la vue 69 sur la base circulaire se termine par deux
lignes serrées qui n'ont pas été lues.} Pour la sphère de courbure moyenne, la
ligne de distance est constante et engendre la \emph{surface des distances
moyennes}, cylindre droit sur un cercle. Elle affirme que les deux ont la même
étendue, et que la courbure de la surface est proportionnelle à cette étendue.
C'est exact. La ligne de distance dans la direction $\varphi$ est
$\frac12\kappa(\varphi)s^2$, la base est un cercle de rayon $s$ à l'ordre
considéré, et
\[
\int_0^{2\pi} \tfrac12\kappa(\varphi)\,s^2\cdot s\,d\varphi
= \tfrac12 s^3\cdot 2\pi H
= 2\pi s\cdot\tfrac12 H s^2 ,
\]
ce qui est l'aire du cylindre de hauteur $\frac12 H s^2$. L'égalité des deux
étendues est donc l'énoncé que \emph{la moyenne de la courbure normale sur
toutes les directions est $H$}. C'est l'une des définitions actuelles de la
courbure moyenne : le nom moderne s'applique ici exactement.

Elle le montre de deux façons. D'abord parce que les lignes de distance ont
entre elles les mêmes relations que les courbures (somme constante sur deux
directions perpendiculaires, moyenne à $45^\circ$).\footnote{Vues 71 et 72 ; les
crochets sont d'elle. Une note encerclée de trois mots après
« N\textsuperscript{o} 5 » (vue 72) et un mot en exposant et une abréviation dans
l'addition interlinéaire de la vue 71 n'ont pas été lus.} Ensuite par le
développement des deux surfaces sur un plan : le cylindre moyen devient le
parallélogramme $ABDC$, la surface des distances une figure bornée par une ligne
brisée $Bqhq'h'D$ en dents, et l'égalité d'aire se ramène à celle des triangles
$Bqn$, $nhn'$, … que chaque dent retranche ou ajoute.\footnote{Vues 72 à 75 et
planche de la vue 89. Le parallélogramme est nommé $ABDC$ (vues 73, 80) et $ABCD$
(vue 74). Les points $s$, $t$, $b$ de la vue 75 ne sont introduits nulle part, et
aucune figure ne les porte ; « Il s'agit à démontrer » est d'elle.} Le vrai
développement de la directrice n'est pas une ligne brisée mais une sinusoïde,
$\frac12 s^2\bigl(H + \frac12(\kappa_1-\kappa_2)\cos 2\varphi\bigr)$ ; son
argument vaut pourtant pour elle, parce que l'écart à la moyenne est
antisymétrique autour de chaque direction de courbure moyenne,
$\kappa(\frac\pi4 + u) + \kappa(\frac\pi4 - u) = 2H$ : l'excès d'un quart de tour
compense exactement le défaut du quart voisin.

\textbf{Les propositions redémontrées (N\textsuperscript{o} 7).} Elle reprend
les trois propositions du N\textsuperscript{o} 3 sur le développement, dans
l'ordre 2, 3, 1 : $mn$ moyenne entre $gh$ et $pq$ ($pq = mn - sn$, $gh = mn +
sn$), quatre lignes de distance égales en général, deux pour les lignes
principales, et somme constante de deux lignes séparées d'un quart de
$AC$.\footnote{Vues 75 à 77. Le paragraphe sur l'objection du
N\textsuperscript{o} 2 est marqué en marge d'un long trait courbe.} Elle répond
alors à l'objection du N\textsuperscript{o} 2 : dire qu'une surface a une
courbure proportionnelle à sa surface des distances, c'est la comparer à son
plan tangent comme on compare une courbe à sa tangente. Et elle en tire la
distinction qui porte tout le mémoire : la ligne de distance fixe à la fois la
quantité de courbure et la figure d'une courbe, tandis que l'étendue de la
surface des distances peut rester la même quand sa forme change.

\textbf{Une famille de figures.} Avec $pq = mn - x$ et $gh = mn + x$, $x$ croissant
de $0$ à $mn$, soit en termes actuels
\[
\kappa_1 = H + x', \qquad \kappa_2 = H - x', \qquad 0 \le x' \le H ,
\]
on passe de la sphère ($x' = 0$) à la surface dont une courbure principale est
nulle ($x' = H$, le cylindre), par une infinité de figures qui ont toutes la
même courbure moyenne. « Si la figure de la surface donnée dépend de la
différence entre les courbures principales, la courbure moyenne ne dépend, au
contraire, que de la somme ».\footnote{Vues 77 à 79 ; « fig 1 et 2 ». Un mot
biffé sous « accompagneront » et un mot au crayon d'une autre main, en marge, ne
sont pas lus.}

\textbf{Les rayons de sens contraires (N\textsuperscript{o} 8).} La sphère d'une
part, la surface dont les rayons principaux sont égaux et de signes contraires de
l'autre, « comprennent entr'elles toutes les figures possibles de la surface
courbe ». C'est exact : à similitude et au choix de la normale près, la figure
locale est fixée par le rapport $\kappa_2/\kappa_1 \in [-1, 1]$ (en prenant
$|\kappa_1| \ge |\kappa_2|$), qui vaut $1$ pour la sphère, $0$ pour le cylindre et
$-1$ pour $\kappa_2 = -\kappa_1$ (hors le plan). Pour cette seconde moitié, elle
garde fixe la somme des lignes principales « abstraction faite du signe »,
c'est-à-dire l'amplitude $\frac12(\kappa_1 - \kappa_2)$, et fait varier $H$ de
$\frac12|\kappa_1|$ à $0$ ; la directrice garde alors sa forme, une partie passe
de l'autre côté du plan tangent, et la surface des distances moyennes diminue
jusqu'à s'anéantir quand les lignes de distance des deux côtés sont
égales.\footnote{Vues 80 à 84 ; figures 2, 3, 4 de la planche. À la vue 84, les
quatre points du développement qui tombent sur $AC$ sont nommés $v$, $t$, $v'$,
$t'$ dans un alinéa et $v$, $u$, $v'$, $u'$ dans le suivant ; la planche porte
$v$, $u$, $v'$, $u'$. « il en est » (vue 83) est lu avec doute.} Tout cela est
juste : l'aire signée de la surface des distances est toujours $\pi H s^3$, et
elle est nulle quand $H = 0$. Elle conclut que la sphère de courbure moyenne est
alors de rayon infini et se confond avec le plan tangent. Un point où
$\kappa_1 = -\kappa_2$ est un point de courbure moyenne nulle ; une surface qui
l'est en tout point est une \emph{surface minimale}, au sens de Lagrange et de
Meusnier.\footnote{Le feuillet n'emploie pas ce nom et ne parle que de la figure
en un point ; le nom ne s'applique qu'à une surface où la condition a lieu
partout.} La loi de répartition, montre-t-elle encore, subsiste entière dans ce
cas.

\subsection*{La courbure comme quantité dynamique ; voûtes et dômes (N\textsuperscript{o} 9)}

\pagerange{84}{89}
Voici l'usage qu'elle attendait de la courbure moyenne. Si des forces
proportionnelles à la courbure agissent sur un point de la surface, elles sont
réparties sur tout le contour, en nombre infini, chacune infiniment petite et
proportionnelle à la ligne de distance dans son plan ; leur action est donc
indépendante de leur répartition et égale à celle d'une répartition uniforme,
c'est-à-dire à ce qu'elle serait sur la sphère de courbure moyenne. Et comme
toute force se décompose selon deux plans perpendiculaires de position
arbitraire, les résultantes dans ces deux plans sont proportionnelles aux lignes
de distance qu'ils contiennent, dont la somme ne dépend pas du choix des
plans.\footnote{Vues 84 à 87. À la vue 85, la phrase se lit sur le feuillet « Le
nombre de ces forces est infini ; d'elles chacune infiniment petite, elle est
proportionnelle à la ligne de distance contenue dans le plan sa direction » ; un
mot biffé sous « chacune » n'est pas lu. À la vue 86, le bord droit du feuillet
passe sous le montage et cache la fin de plusieurs mots.}

Pour une membrane soumise à une tension $T$ égale dans toutes les directions,
c'est exact, et c'est la \emph{loi de Laplace} : la résultante normale par unité
d'aire est $T(\kappa_1 + \kappa_2) = 2TH$. On le voit exactement comme elle le
voit : la composante normale de la tension sur le bord d'un petit disque de
rayon $\varepsilon$ vaut $T\int_0^{2\pi}\kappa(\varphi)\,\varepsilon^2\,d\varphi =
2\pi\varepsilon^2\,TH$.\footnote{Laplace a donné cette loi dans sa théorie de la
capillarité (1805--1806) ; le feuillet ne le cite pas, et ne parle ni de
capillarité ni de tension superficielle. Le nom s'applique à l'énoncé pour une
tension isotrope ; il ne s'applique pas à une plaque élastique.} Il n'en va pas
de même pour une plaque élastique : la force de rappel n'y est pas
proportionnelle à la courbure moyenne, mais à son laplacien ($n^2\Delta^2 z =
n^2\Delta(\Delta z)$ en théorie linéaire), et l'énergie dépend aussi de la
courbure de Gauss dès que $\nu \neq 1$. Aussi, quand elle dit que, « pour
justifier l'hypothèse qui a servi à trouver, la première fois, l'équation de ces
surfaces, il suffiroit de répéter les diverses propositions qui ont été
démontrées », la lecture doit dire que ce n'est pas une justification : le
mémoire établit la loi de répartition de la courbure et le rôle de $H$ pour une
force isotrope, non que l'élasticité d'une plaque ne dépende que de
$H$.\footnote{Vues 87 et 88 ; « a » est lu avec doute dans « qui a servi ».}

Pour les voûtes, la simplicité et le petit nombre des figures employées ont
masqué aux géomètres la nécessité d'exprimer en général les résistances nées de
la courbure ; comparer les voûtes en dômes aux voûtes en berceaux demanderait de
comparer les deux genres de courbure. Elle termine par un exemple : le cylindre
circonscrit à une sphère de rayon $a$ et le quart de la sphère de rayon $2a$ ont
même étendue et même courbure moyenne en tout point. C'est exact : les deux
aires valent $4\pi a^2$, et les deux courbures moyennes $\frac1{2a}$ ; la sphère de
courbure moyenne du cylindre de rayon $a$ a bien le rayon $2a$.\footnote{Vue 88 ;
la fin de la ligne passe sous le montage (« employés » lu avec doute, la
ponctuation de l'interrogation non lue) ; « de moyenne courbure », à la dernière
ligne, est souligné deux fois. La vue 89 est la planche des quatre figures, à
l'encre, sans texte.}

\section*{« Moyennes courbures » : les brouillons de l'introduction}

\pagerange{91}{114}
Après la mise au net vient un feuillet de titre, « Moyennes courbures », puis
des brouillons d'une écriture rapide et très corrigée, un côté écrit par feuillet,
qui travaillent l'introduction d'un exposé où la courbure moyenne sert à la
théorie des surfaces élastiques.\footnote{Vues 91 à 114. Série à l'encre 45
(vue 91) à 55 bis (vue 113, au-dessus d'un « 66 » biffé) ; les versos 90, 92, 94,
96, 98, 99, 102, 104, 106, 108, 110 sont blancs, et les vues 111 et 113 ne
portent que leur numéro, le texte de ces feuillets étant aux vues 112 et 114. Des
lettres de sa main servent de repères : « E » (vue 95), « A » (vues 95 et 101),
« B » (vues 95 et 103), « D » (vue 107).} Leur ordre n'est pas établi par les
feuillets ; on les lit ici dans l'ordre des vues.

\subsection*{Le solide élastique réduit à sa surface}

\pagerange{93}{112}
\textbf{Les dômes.} Dans une voûte en dôme engendrée par la révolution d'une
demi-courbe, chaque zone comprise entre deux plans parallèles a une courbure
partout semblable ; elle lui associe une portion de surface conique ou
cylindrique de même courbure moyenne, et, en portant zone après zone ces rayons
comme ordonnées, une surface cylindrique dont les éléments ont la même courbure
moyenne que ceux du dôme.\footnote{Vue 93. Le feuillet est déchiré et tout le bas
manque ; « Aussi on formera », « prenoit » et « à » sont lus avec doute ; les
couples de lettres $A$——$B$, $A'$——$B'$ sont écrits sur deux segments tracés à la
plume, qu'une accolade relie.} Ce qui est exact : un cylindre de rayon
$\frac12\varrho_m$ a pour rayon de courbure moyenne $\varrho_m$ ; pour un cône
de révolution, il faut entendre par « rayon » le segment de normale compris entre
le point et l'axe, dont la moitié de la longueur est de même l'inverse de la
courbure moyenne.

\textbf{Du solide à sa surface.} Les feuillets suivants cherchent à dire ce que
désigne le mot « surface » dans la question des plaques. Le solide élastique est
« un système de points matériels mobiles » ; le déplacement d'un point comprend
le changement de ses distances aux voisins et le mouvement commun qu'il partage
avec eux, et elle écrit que les forces de rappel sont proportionnelles au
« déplacement absolu », mesuré à partir de points fixes.\footnote{Vue 101, en
tête « N\textsuperscript{o} 1 » en marge et un « A » isolé ; passage repris en tous sens, dont
l'état final n'est pas cohérent : « symétrie » est lu avec doute au-dessus de
« liaison » biffé, et « il sera facile d'imprimer » se lit là où le sens
attendrait « d'exprimer ». Deux additions interlinéaires sont biffées et
illisibles.} Pris à la lettre, c'est inexact : dans la théorie actuelle de
l'élasticité, un déplacement d'ensemble (translation ou rotation) ne produit
aucune force ; les forces dépendent des déplacements relatifs, c'est-à-dire de
la déformation. Elle en arrive d'ailleurs, sur les feuillets suivants, à ne
mesurer que la position relative des points voisins par rapport au plan
tangent.

Au feuillet marqué « B », elle refait l'histoire à la première personne : les
expériences de Chladni ne portaient que sur des plaques ; les programmes de
l'Académie employaient le mot « surface » sans le définir ; « le mémoire publié
en 1814 par un des membres de l'accademie » ne traite que des surfaces planes ;
« le mémoire couronné qui fut présenté en 1815 » contient les premiers essais
sur les surfaces courbes, par le cylindre ; « mes recherches sur la théorie des
surfaces publiées en 1821 » reprennent cet exemple et généralisent l'hypothèse
des surfaces planes ; elle reconnaît enfin « avec une véritable satisfaction »
ce que les expériences de Savart l'ont aidée à développer.\footnote{Vue 103.
Toutes ces dates sont sur le feuillet. Le mémoire couronné est dit « présenté en
1815 » ici et « envoyé pour le concours du mois de janvier 1816 » dans
l'Exposition (vue 7) ; les deux formules ne se contredisent pas. « reconnoîtrai »
et « opposition » sont lus avec doute ; un mot biffé devant « variété » et un
autre après « Savart » ne sont pas lus.} Puis elle définit : le mot « surface »
désigne un solide formé de couches infiniment minces ; tous les points d'une
même perpendiculaire au plan tangent parcourent des espaces égaux, sans changer
de distance entre eux ; la figure est donc la même pour toutes les couches et
indépendante de l'épaisseur, ce qu'atteste l'identité des figures nodales sur
les deux faces.\footnote{Vue 105. Un passage biffé disait que, si l'on y avait
réfléchi, l'objection qu'il ne s'agissait que d'une surface d'épaisseur
infiniment petite « n'auroit pas été proposée à l'auteur du mémoire publié en
1814 et ce savant géomètre ne l'auroit pas accepté ». Le haut du feuillet porte
deux états superposés ; un mot après « Je crois utile d' » et deux mots biffés
ne sont pas lus. La note marginale, en partie biffée, s'arrête sur « les figures
qu'affectent la surface ».} Cette condition — le déplacement transversal ne
dépend pas de la position dans l'épaisseur, et il se fait selon la normale — est
l'une des hypothèses cinématiques sur lesquelles Kirchhoff fondera la théorie des
plaques minces.\footnote{Kirchhoff, 1850, postérieur aux feuillets ; le nom ne
couvre qu'une partie de l'hypothèse de Kirchhoff, qui demande aussi que les
normales restent normales au feuillet moyen.}

Au feuillet « D », elle en tire l'argument qui lui importe : son hypothèse ne
portait que sur la courbure d'un élément de surface plane, et pourtant une
simple opération de calcul en a tiré une équation « dont personne ne conteste
l'exactitude » ; si le solide n'était pas formé de couches qui suivent la couche
superficielle, comment l'équation de la surface deviendrait-elle celle du solide
par la seule multiplication par un coefficient dépendant de
l'épaisseur ?\footnote{Vues 107 et 97. À la vue 97, le coefficient est dit
« composé des deux facteurs qui expriment l'influence de l'épaisseur et celle de
l'élasticité naturelle » ; en marge, « cette uniformité elle même résulte de
l'invariabilité de l'épaisseur ». Les quatre lignes ajoutées en tête de la vue
97, d'une plume très fine, ne se lisent qu'en partie.} Le raisonnement est
juste dans sa forme : en théorie des plaques, l'épaisseur n'entre que par le
coefficient (qui en contient le carré ; voir l'Exposition, N\textsuperscript{os} 9 à 11, et le § 1 du long mémoire). Les
vues 109 et 112 réservent expressément la direction du mouvement
(perpendiculaire aux plans tangents), renvoient l'épaisseur variable à plus tard,
et ramènent les espaces parcourus au changement des distances des points voisins
au plan tangent, où « la considération de la courbure moyenne … est ici de
grand secours ».\footnote{Vues 95, 100, 109, 112. À la vue 95, « Commencement de
la p. A » est une indication de l'auteur à elle-même ; les trois mots après
« N\textsuperscript{o} » sont perdus sous la rature, un mot devant « supposerons »
n'est pas lu, et « simplificatrices » et « substituer » sont lus avec doute. La dernière ligne de la vue 100 n'est pas lue.
À la vue 109, une phrase biffée mentionne un mémoire « présenté depuis longtems à
l'accademie », dont deux mots ne sont pas lus ; en marge, « les N\textsuperscript{os} suivans
seront consacrés à démontrer son existence et à développer ses propriétés »
(« les N\textsuperscript{os} suivans » et « consacrés » lus avec doute).} La vue 112 dit enfin
que « le cas dont je me suis spécialement occupée » est celui des surfaces
élastiques vibrantes, et que l'identité des figures nodales sur les faces
opposées, confirmée par « les ingénieuses expériences de Mr Savart », est
« le signe distinctif de ce genre de vibrations ».\footnote{Vue 112 ;
« confirmée » est lu avec doute, d'une encre pâle ; le mot devant « choix » et
trois mots biffés ne sont pas lus ; un signe de renvoi après « mêmes » n'a pas
d'appel sur le feuillet.}

\subsection*{Trois sections normales voisines}

\pagerange{114}{114}
D'une plume plus fine, un calcul sur trois sections normales faisant les angles
$A - a$, $A$, $A + a$ avec la section de plus grande courbure. En termes
actuels, avec $\kappa(\theta) = \kappa_1\cos^2\theta + \kappa_2\sin^2\theta$,
\begin{align*}
\kappa(A) - \kappa(A+a) &= (\kappa_1-\kappa_2)\sin a\,\sin(2A+a), \\
\kappa(A-a) - \kappa(A) &= (\kappa_1-\kappa_2)\sin a\,\sin(2A-a), \\
\kappa(A-a) + \kappa(A+a) - 2\kappa(A) &= -2(\kappa_1-\kappa_2)\sin^2 a\,\cos 2A .
\end{align*}
Ce sont ses équations (1), (2) et (3), aux lettres et aux signes
près.\footnote{Vue 114, feuillet « 55 bis ». Les trois équations sont repassées
à l'encre noire et empâtées : le second dénominateur de (1) se lit seulement
$R$, le signe entre $\frac1r$ et $\frac1{r'}$ dans (3) est sous une tache, et le
facteur de (2) se lit « $\sin(2A-a)$ » avec doute. Le premier membre de (2)
est écrit $\frac1R - \frac1{R'}$ là où, avec l'ordre qu'elle donne ensuite (les
sections $R'$, $R$, $R''$ aux angles $A - a$, $A$, $A + a$), la diminution de
courbure est $\frac1{R'} - \frac1R$. Vérifié par calcul symbolique.} Le second
membre de (3) s'annule, dit-elle, pour la sphère, pour $a = 0$ ou $180^\circ$, et
pour $A = 45^\circ$ : c'est exact. Elle en conclut que « la somme des rayon de
courbure des deux sections est égale à la somme des rayons de courbure
principale » ; ce qui est vrai porte sur les courbures, non sur les rayons :
$\kappa(45^\circ - a) + \kappa(45^\circ + a) = \kappa_1 + \kappa_2$ pour tout $a$.

Sur la vitesse de variation, le feuillet dit deux choses opposées. La
comparaison de (1) et (2) — la diminution entre $A - a$ et $A$ est plus petite
qu'entre $A$ et $A + a$ — est juste, et sa dernière phrase aussi : la courbure
varie d'autant plus vite, pour un même angle, qu'on approche de la section de
courbure moyenne ; en effet $\kappa'(\theta) = -(\kappa_1-\kappa_2)\sin 2\theta$
est maximal en valeur absolue à $45^\circ$ et nul sur les directions
principales. Mais la phrase intermédiaire, selon laquelle, en partant de
$45^\circ$, l'augmentation serait d'autant plus rapide que les deux plans font un
angle plus petit avec celui de plus grande courbure, dit le contraire, et elle
est fausse.\footnote{Vue 114. Deux lignes barrées d'un trait oblique précèdent ce
passage.} Enfin la courbe $\kappa(\theta) - H$ a un point d'inflexion à
$45^\circ$ et y est symétrique par rapport à ce point ; c'est ce qu'elle
décrit en disant ses deux parties « symétriques autour de l'ordonnée de moyenne
courbure », l'une convexe, l'autre concave vers le diamètre qui tient lieu du
plan tangent.

\section*{Feuillets épars sur la courbure}

\pagerange{136}{650}
Cinq morceaux, reliés à des endroits divers du volume, touchent au même sujet
que les pièces qui précèdent sans appartenir à aucune.\footnote{Les dérivées y
sont écrites avec un $d$ droit ; on écrit $\partial$. On note $p = \partial_x
z$, $q = \partial_y z$, $W = \sqrt{1+p^2+q^2}$, et $\mathcal{N} = (1+q^2)\,r -
2pq\,s + (1+p^2)\,t$, de sorte que $\kappa_1+\kappa_2 = \mathcal{N}/W^3$ au signe
près.}

\subsection*{« Idées à vérifier » : la chaînette et la surface pesante}

\pagerange{136}{140}
La question posée : dans une surface pesante d'épaisseur égale, chaque section
normale est-elle une chaînette ?\footnote{Vues 136 et 137, feuillet numéroté 67,
recto et verso ; vue 138, feuillet 68 ; vue 140, feuillet 69 ; versos 139 et
141 blancs. L'en-tête du lot 7 ne se prononce pas sur la main de ces trois
feuillets, alors que la note de la vue 136 les dit « de la main de l'auteur » ;
ils sont écrits à la première personne. La lettre $e$ qui suit $g$ dans les
équations est bien un $e$ ; que ce soit l'épaisseur est une inférence de cette
lecture.} Elle écrit l'équilibre de la surface $ge + T\cdot\frac2R = 0$, où $R$
est le rayon de la sphère de courbure moyenne ($\frac2R = \kappa_1+\kappa_2$) :
c'est l'équation d'une membrane pesante soumise à une tension isotrope $T$, la
loi de Laplace du N\textsuperscript{o} 9 du mémoire sur la courbure (vues 84 à 89), avec la pesanteur pour pression. Puis
elle demande si la section à $45^\circ$, isolée du reste, garderait sa forme :
il faudrait pour cela, dit-elle, que le retranchement diminue la tension dans le
rapport de un à deux. C'est l'inverse : la section à $45^\circ$ a la courbure
$\frac1R = H$, la moitié de $\kappa_1+\kappa_2$ ; isolée avec la même charge, elle
ne resterait en équilibre dans sa forme que si sa tension était doublée — ou, ce
qui revient au même, si elle ne portait que la moitié de la charge, l'autre
moitié étant portée dans la direction perpendiculaire. Elle-même trouve
insatisfaisant son premier argument, par les appuis $B'$, $B''$, $B'''$ d'une
figure, et écrit : « je n'ai pas trouvé la manière de détruire cette
objection ».\footnote{Vue 136. « conserve », « Cela » et « complémens » sont lus
avec doute ; la ligne « en $B$ $B''$ $B''$ » est lue avec les accents qu'elle
porte, là où la figure a $B'$, $B''$, $B'''$.}

Au verso, une seconde voie : sur un châssis circulaire, la tension se partage
également entre deux plans normaux, et, en posant $T = 2T'$ et un poids
$4g''$ pour la surface, $2g''$ pour une courbe, elle arrive à $g''e + T'\frac1R
= 0$ pour chaque section normale, et conclut qu'elle est une chaînette. La
conclusion n'est pas établie. L'équation de la surface ne contraint que la
somme $\kappa_1+\kappa_2$ ; en écrivant pour chaque section la même courbure
$\frac1R$, elle prête à toute section normale la courbure moyenne, ce qui n'est
vrai que des deux sections à $45^\circ$ (et de toutes au seul cas ombilical).
Pour une surface peu inclinée, tout se réduit à $T\,\Delta z = w$ (charge $w$
par unité d'aire) ; sur un châssis circulaire, la solution est un paraboloïde
de révolution, dont les sections normales au centre sont des paraboles, forme
qu'a toute chaînette peu tendue à cet ordre : à cet ordre la question ne peut
donc pas être tranchée par ce calcul.\footnote{Vue 137 ; « élévation » est lu
d'après le sens, le mot étant tracé plutôt « ébration » ; l'équation
« $ge + T.\frac2R$ » y est écrite sans « $= 0$ ». La page s'arrête sur le cas
d'appuis inégaux, « et l'abbaissement le plus grand ».}

Le feuillet suivant (vue 138) dit pourquoi : « dans la vue d'étendre les conséquences qu'on
peut tirer de l'existence de la sphère de moyenne courbure », elle a cherché de
tout côté ce qui concerne la chaînette, « à cause de son rapport avec la voûte
qui se soutient d'elle même », et ne l'a trouvé que dans l'histoire des
mathématiques, l'Encyclopédie et « le mémoire de M poisson sur la surface
élastique ». De l'article de d'Alembert elle retient qu'une courbe pressée en
chaque point par une force normale est en équilibre si la force est en raison
inverse du rayon de courbure — vrai pour un fil sous tension constante, $p =
T\kappa$ — et elle l'étend à la surface : « il faut que la moitié de cette
puissance soit en raison inverse du rayon de sa sphère de moyenne courbure ».
C'est exactement $p = T(\kappa_1+\kappa_2) = 2TH$, la loi de Laplace pour une
tension isotrope.\footnote{Vue 138. L'énoncé général est souligné ligne à ligne ;
« courbure » y remplace « surface » biffé, et « raison » est barré dans « soit en
raison inverse double raison inverse ». Le poids à donner au mot « moitié » est
celui de la seconde rédaction, « ou ce qui est plus clair ».}

Celui de la vue 140, « Voila ce qui doit être examiné », rapproche l'élément
$dx\,dy\,W$ de sa projection $dx\,dy$, puis prend la variation de $W$ et
l'intègre par parties : « et en intégrant par partie $\frac2R\,\delta z$ ». C'est
exact : la variation de l'aire est
\[
\delta\iint W\,dx\,dy = -\iint \frac{\mathcal{N}}{W^3}\,\delta z\,dx\,dy +
\text{termes de bord},
\]
et $\mathcal{N}/W^3 = \kappa_1+\kappa_2 = \frac2R$ au signe près.\footnote{Vue 140 ;
« je crois », biffé, est lu avec doute. La page s'arrête sur « ainsi la
variation ». Vérifié par calcul symbolique.} Réunis, les trois feuillets
contiennent la chose même : l'équilibre d'une membrane pesante sous tension
isotrope est le point stationnaire de $T\cdot\text{aire} + \text{énergie
potentielle}$, et son équation est la loi de Laplace ; les surfaces où $H = 0$
en sont le cas sans charge (surfaces minimales). Le feuillet s'arrête avant de
l'écrire.

\subsection*{Une page de Gauss}

\pagerange{142}{142}
Un petit feuillet intitulé « Gauss courbure des surfaces » rend en français,
entre guillemets, deux articles d'un texte : la situation d'un point est donnée
par ses distances à trois plans fixes perpendiculaires ; on décrit autour d'un
centre arbitraire une sphère de rayon 1, et la direction d'une droite est
représentée par le point de la sphère où aboutit le rayon qui lui est parallèle ;
les trois axes, par trois points séparés par des quarts de cercle ; la
situation d'un plan, par le grand cercle que découpe le plan parallèle mené par
le centre.\footnote{Vue 142, feuillet numéroté 70, verso blanc (vue 143).
« Gauss », « soit » et le « 2 » marginal sont lus avec doute ; « tous les
points » est biffé en tête ; le feuillet écrit « un cercle de rayon 1 » puis
« le point de la sphère ». Le texte s'arrête sur « parallèle ».} C'est, par
inférence de cette lecture, l'ouverture des \emph{Disquisitiones generales circa
superficies curvas} de Gauss (1827), où la représentation sphérique des
directions prépare l'application qu'on appelle aujourd'hui \emph{application de
Gauss} et la mesure de la courbure par l'aire de l'image sphérique.\footnote{Le
feuillet ne donne ni le titre ni l'année ; l'identification est de cette lecture,
faite sur le seul contenu des deux articles. Le feuillet ne va pas jusqu'à la
courbure elle-même.} La note ne dit rien de ce qu'elle en pensait.

\subsection*{Deux théorèmes proposés aux géomètres}

\pagerange{465}{465}
Un fragment déchiré, relié au milieu du long mémoire sur les plaques sans en
faire partie, porte : « je proposerai d'abord a l'examen des g… les deux
théorèmes suivants. La courbure d'une surface a pour mesure la somme des raisons
inverses des deux principaux rayons de courbure. L'irrégularité dela courbure
d'une surface a pour mesure la difference des raisons inverses de ses deux rayons
principaux de courbure ».\footnote{Vue 465, fragment numéroté 228 dans la série à
l'encre, d'une écriture plus grande, sans pagination ; la déchirure emporte une
ligne entière et le début de la suivante (« …ment », le reste illisible) ;
« suivants » est lu avec doute ; le dernier mot s'arrête sur « g ». Le dos (vue
466) est blanc. La pagination du mémoire passe de (58) à (59) par-dessus.} Ce sont
les deux mesures des N\textsuperscript{os} 3 et 4 de l'Exposition : $2H$, et
$\kappa_1 - \kappa_2$, dont le carré vaut $4(H^2-K)$. Comme définitions, les deux
sont légitimes, et la décomposition $\kappa(\varphi) = H +
\frac12(\kappa_1-\kappa_2)\cos 2\varphi$ montre qu'à elles deux elles déterminent la
courbure normale dans toutes les directions. Comme « mesures » des forces
élastiques, elles sont les deux hypothèses dont la théorie actuelle retient une
combinaison (section de l'Exposition).

\subsection*{Un calcul de variation}

\pagerange{483}{483}
Un feuillet de calcul sans un mot, relié lui aussi dans le long mémoire, calcule
la variation d'une intégrale double dont le premier membre se lit $\frac{ds}{(1/R
+ 1/R')^2}$.\footnote{Vue 483, feuillet numéroté 237, sans pagination ; dos (vue
484) blanc. La transcription ne décide pas si le trait sous $ds$ est une barre de
fraction ou un trait de réunion. Deux lignes lourdement biffées, le premier signe
de la grande fraction et un signe intérieur (des taches d'encre) ne sont pas lus ;
des essais de plume sont en marge.} Elle écrit d'abord la somme des courbures,
$\frac1R + \frac1{R'} = -\mathcal{N}/W^3$, qui est juste. Mais la quantité dont elle
prend ensuite la variation, à en juger par sa structure, est
\[
\iint \Bigl(\frac{\mathcal{N}}{W^2}\Bigr)^2 dx\,dy ,
\qquad
\delta\Bigl(\frac{\mathcal{N}}{W^2}\Bigr)^2 = 2\,\frac{\mathcal{N}}{W^2}\cdot
\frac{W^2\,\delta\mathcal{N} - \mathcal{N}\,\delta(W^2)}{W^4},
\]
et elle en range les coefficients selon $\delta z$, $\delta p$, $\delta q$,
$\delta r$, $\delta s$, $\delta t$ sous les lettres $N = 0$ (la quantité ne
dépend pas de $z$), $P$, $P'$, $Q$, $Q'$, $Q''$.\footnote{À la troisième ligne le
feuillet écrit l'intégrande $-\mathcal{N}^2\,dx\,dy/W^2$, et la variation qui suit
est celle de $\mathcal{N}^2/W^4$ : les puissances de $1+p^2+q^2$ ne s'accordent pas
d'une ligne à l'autre. Dans le coefficient de $\delta p$, le calcul donne
$2\{q^2p\,t - (1+q^2)p\,r - q(1+q^2-p^2)\,s\}$ ; le feuillet omet le facteur $q$
du dernier terme, et de même le facteur $p$ dans $P'$ (vérifié par calcul
symbolique).} Cette intégrande n'est ni $ds/(\kappa_1+\kappa_2)^2 = W^7\,dx\,dy
/\mathcal{N}^2$, ni $(\kappa_1+\kappa_2)^2\,ds = \mathcal{N}^2\,dx\,dy/W^5$ ; aux
petites pentes, toutes se réduisent à $(\Delta z)^2\,dx\,dy$, l'énergie de son
hypothèse pour la plaque. Si le calcul visait $\iint(\kappa_1+\kappa_2)^2\,ds$,
forme exacte de son énergie pour une surface naturellement plane, c'est
l'intégrale qu'on nomme aujourd'hui énergie de Willmore
($4\iint H^2\,dA$).\footnote{Willmore, 1965 ; nom postérieur, donné ici à titre
d'hypothèse sur l'intention du calcul, que le feuillet ne permet pas
d'établir.} Le feuillet ne dit pas à quelle pièce il servait, et la lecture ne
l'attache à aucune.

\subsection*{La somme des courbures sur la sphère}

\pagerange{650}{650}
Parmi les calculs sur la chaleur, un feuillet intitulé « Somme des raisons
inverses des deux rayons de principale courbure » écrit juste la formule
\[
\frac1R + \frac1{R'} = \frac{(1+q^2)\,r - 2pq\,s + (1+p^2)\,t}{(1+p^2+q^2)^{3/2}},
\]
puis pose $x = r\sin\theta\cos\varphi$, $y = r\sin\theta\sin\varphi$ avec $r$
constant, « pour avoir seulement le cas relatif à la surface de la sphère », et
calcule $\partial_x z$, $\partial_y z$, le facteur
$K^3 = (1+p^2+q^2)^{3/2}$ et $\partial_x\partial_y z$ en $\theta$ et
$\varphi$.\footnote{Vue 650, feuillet numéroté 319, relié entre les calculs sur la
chaleur de la sphère ; dos (vue 651) blanc. Le coefficient de
$\cos^2\theta\,(\partial_\theta z)^2$ dans $K^3$ est lu « 2 » avec doute ; avec ses
propres $p$ et $q$ le calcul donne $1$. Un premier départ de $\partial_x\partial_y
z$ est barré ; les accolades de droite ne se ferment pas, le dernier facteur
touchant le bord.} Ses dérivées premières sont
\[
\frac{\partial z}{\partial x} = \frac{\cos\theta\cos\varphi}{r}\,\frac{\partial
z}{\partial\theta} - \frac{\sin\varphi}{r\sin\theta}\,\frac{\partial
z}{\partial\varphi}
\]
et sa symétrique ; le coefficient $\cos\theta\cos\varphi/r$ est celui du
changement de variables à trois dimensions, où $r$ varie aussi. Si $r$ est tenu
constant et que $(\theta,\varphi)$ servent seuls de coordonnées dans le plan des
$xy$, la règle de dérivation donne $\partial\theta/\partial x =
\cos\varphi/(r\cos\theta)$, et le calcul, tel qu'il est écrit, ne donne pas la
somme des courbures d'une surface décrite au-dessus de ce plan.\footnote{Vérifié
par calcul symbolique des deux jacobiens. Le feuillet ne dit pas quelle surface
il vise, et la lecture ne tranche pas.} Le feuillet, relié parmi les calculs de
la chaleur en coordonnées sphériques, emprunte sans doute leurs formules ; il ne
va pas au-delà de $\partial_x\partial_y z$.

\section*{Le mémoire d'Euler sur les vibrations des lames et des verges élastiques (copie latine)}

\pagerange{144}{206}
De la vue 144 à la vue 206, le volume contient une copie latine complète d'un mémoire d'Euler, dont le premier feuillet donne lui-même la source : « Acta academiæ scientiarum imperialis Petropolitanæ pro anno 1779 Pars prior », « Investigatio motuum quibus laminæ et virgæ elasticæ contremiscunt auctore L. Eulero ».\footnote{Vues 144 à 206. Copie au net, écrite au recto et au verso, d'une main droite et régulière que la transcription rapproche de celle des feuillets français ; la main de ce volume n'est pas attribuée ici. Chaque page porte en tête un nombre entre parenthèses, de (103) à (161), que la transcription tient, par inférence, pour la page de l'imprimé copié ; ses propres critiques citent le mémoire « P.I p. 103 et seq. » (vues 210 et 320), ce qui s'accorde avec le premier de ces nombres. La série à l'encre du volume va de 71 (vue 144) à 100 (vue 206). Les pages (126) à (128) sont reliées hors d'ordre : la vue 167 porte (128), entre (125) à la vue 166 et (126)--(127) aux vues 168--169, et la copie le signale elle-même au bas de la vue 166, « (voyez p. 126 après 128) ». Sont photographiées deux fois : la vue 191 (= 190), les vues 196--197 (= 194--195), la vue 204 (= 203). La planche que la copie cite en marge (« Tab. II fig. 3 » à « Tab. III fig. 10 ») n'est pas dans le volume. La copie finit à la vue 206 sur le mot français « fin. » et un paraphe.} Euler y reprend, dit-il, un sujet déjà traité par Daniel Bernoulli et par lui-même, parce que les principes de ces mouvements et « la partie de l'analyse qui concerne les fonctions de deux variables » sont désormais mieux établis ; il se propose de déduire toute la question des premiers principes et d'exposer tous les genres de mouvement qui peuvent se produire.\footnote{Vue 144 : « neque ea analyseos pars, quæ circa functiones binarum variabilium versatur, satis explorata ». Le nom lu « Illustriss. D. Bernoulli » est une abréviation dont l'initiale ressemble à un M (lu avec doute).}

\subsection*{L'équation du mouvement de la verge élastique}

\pagerange{144}{152}
Deux lemmes posent le cadre.\footnote{Vues 144 à 147, pages (103) à (106) de la copie. Euler écrit les différentielles partielles entre parenthèses, $\left(\frac{ddy}{dt^2}\right)$, « quibus uncinulis ( ) indicatur » que seul le temps varie (vue 147) ; on écrit ici $\partial_t^2 y$, et $y'$, $y''$, $y'''$, $y''''$ pour les dérivées par rapport à l'abscisse. Dans toute la copie, $g$ est la hauteur de chute des graves en une seconde, de sorte que l'accélération de la pesanteur vaut $2g$.} Le premier donne l'équilibre d'une verge élastique sous des forces quelconques : le moment d'élasticité vaut $V/r$, $r$ étant le rayon de courbure et $V$ l'« élasticité absolue », et l'équilibre s'écrit
\[
\int dy\int P\,ds-\int dx\int Q\,ds=V\,\frac{dy\,d^2x-dx\,d^2y}{ds^3},
\]
où $P\,ds$ et $Q\,ds$ sont les composantes de la force appliquée à l'élément ; si la verge est naturellement courbe, de rayon de courbure naturel $\varsigma$, on retranche $V/\varsigma$ au second membre, puisque l'élasticité n'agit qu'autant que la courbure naturelle est augmentée ou diminuée (§ 2). Le second lemme passe au mouvement en remplaçant $P$ et $Q$ par $P-\frac{S}{2g}\,\partial_t^2x$ et $Q-\frac{S}{2g}\,\partial_t^2y$, $S$ étant la masse par unité de longueur : c'est, en substance, le principe de d'Alembert, que le feuillet ne nomme pas.

Le problème 1 (§ 4) applique ces formules à une verge droite, d'épaisseur et d'élasticité uniformes, de longueur $a$ et de section $c^2$, dans l'hypothèse des vibrations très petites.\footnote{Vues 147 à 150, pages (106) à (109). Euler mesure forces et masses par des volumes de la matière de la verge ; l'élasticité absolue est prise « quadrato crassitiei proportionalis », la « crassities » étant la section $c^2$, d'où $V=bc^4$, où $b$ est une longueur qui dépend de la matière. À la vue 149, la phrase sur les forces appliquées au bout $E$ porte trois « alteram » pour deux forces, avec des directions lues $EE$, $EF$, $Ef$ ; la transcription laisse ces lettres telles qu'elles sont lues, et la suite du calcul n'en dépend pas.} En dérivant deux fois l'équation du mouvement, il obtient
\[
E\,\partial_s^2 y+\frac{c^2}{2g}\,\partial_t^2 y=-\,bc^4\,\partial_s^4 y,
\]
où $E$ est la force longitudinale appliquée au bout, comptée positivement quand elle comprime. C'est l'équation de la poutre d'Euler--Bernoulli avec effort axial,\footnote{Le nom « Euler--Bernoulli » est postérieur à ce mémoire, qui ne l'emploie évidemment pas ; il désigne aujourd'hui exactement cette équation de flexion des poutres minces, sans inertie de rotation ni cisaillement.}
\[
\rho A\,\partial_t^2 y+B\,\partial_s^4 y-T\,\partial_s^2 y=0,
\]
avec le dictionnaire $\rho A=c^2/2g$ (masse par unité de longueur, les poids étant comptés en volumes), $B=bc^4$ (rigidité de flexion, le produit du module d'Young par le moment d'inertie de la section) et $T=-E$ (tension). Le § 5 note que, sans force extérieure, tout se réduit à
\[
\partial_t^2 y=-\,2g\,bc^2\,\partial_s^4 y,
\]
et Euler avoue que l'intégrale générale de cette équation du quatrième ordre « n'a encore pu être trouvée d'aucune manière » : il cherchera des solutions particulières.\footnote{Vue 150 : « ejus integrale nullo adhuc modo inveniri potuisse ». La seconde forme de l'équation porte $bc^2$ là où la première porte $bc^4$ ; c'est juste, la seconde étant la première divisée par $c^2$.} Le § 6 traite de la force qui comprime (la question des colonnes, qu'Euler dit avoir résolue autrefois) ; le § 7, d'une tension $c^2h$ donnée par un poids passé sur une poulie : l'équation devient $\frac{1}{2g}\,\partial_t^2 y=h\,\partial_s^2 y-bc^2\,\partial_s^4 y$, qui redonne la corde vibrante pour $b=0$ et la verge élastique pour $h=0$, le son étant « mixte » entre les deux.

Le scholie du § 8 écarte les forces longitudinales et garde, à côté de l'équation finale, les deux intégrales dont elle provient :\footnote{Vues 151 et 152, pages (110) et (111).}
\[
\text{I.}\quad -Fx+\frac{c^2}{2g}\int dx\int \partial_t^2 y\,ds=-\,bc^4\,y'',\qquad
\text{II.}\quad -F+\frac{c^2}{2g}\int \partial_t^2 y\,dx=-\,bc^4\,y''',
\]
car une force normale $F$ y entre « quand la verge n'est pas entièrement libre, mais fixée en un ou plusieurs points comme au moyen d'un stylet, autour duquel elle reste pourtant mobile ».\footnote{Vue 152 : « quando virga non omnino est libera, sed in uno pluribusve punctis quasi ope styli est fixa, circa quem tamen sit mobilis ».} En termes actuels, II est le bilan de l'effort tranchant : là où un stylet appuie, $y'''$ saute d'une quantité proportionnelle à la force qu'il exerce. Ce point commande toute la discussion qui suit, et Euler l'a posé lui-même dès le § 8.

\subsection*{Les vibrations régulières et les six cas}

\pagerange{153}{192}
Le § 9 cherche les vibrations « régulières », c'est-à-dire isochrones à un pendule simple de longueur $k$ : ce sont, en termes actuels, les modes propres, obtenus par séparation des variables.\footnote{Vues 153 à 156, pages (112) à (115). La copie écrit le facteur du temps $\sin\bigl(\zeta+t\frac{\sqrt{2g}}{k}\bigr)$, le radical ne couvrant que $2g$ ; la durée d'une vibration qu'elle en tire, $\pi\sqrt{k/2g}$ (vue 155), et la forme $\sqrt{2g/k}$ des vues 161 et 167 demandent $\sin\bigl(\zeta+t\sqrt{2g/k}\bigr)$, qu'on écrit. Le coefficient arbitraire, une majuscule calligraphiée, est rendu $\mathfrak{M}$ par la transcription ; « clausulis » (vue 153) est lu avec doute.} En posant $f^4=bc^2k$, $\omega=a/f$ et $u=s/a$, on a
\[
y=C\sin\bigl(\zeta+t\sqrt{2g/k}\bigr)\,Y(u),\qquad
Y=\alpha e^{\omega u}+\beta e^{-\omega u}+\gamma\sin\omega u+\delta\cos\omega u,
\]
et le nombre de vibrations par seconde, qu'Euler prend pour mesure du son, vaut $\omega^2c\sqrt{2gb}/\pi a^2$. Une « vibration » est ici une oscillation simple, une demi-période : le nombre d'Euler est le double de la fréquence actuelle $\frac{\omega^2}{2\pi a^2}\sqrt{B/\rho A}$, ce qui ne change aucun rapport. Les sons d'une même verge sont donc entre eux comme les carrés $\omega^2$ des racines d'une équation que fixent les extrémités ; les §§ 10 à 12 donnent les séries, les dérivées et la vitesse.

Le scholie I (§ 13) distingue trois états d'une extrémité, chacun donnant deux conditions tirées des équations I et II du § 8 :\footnote{Vues 156 à 159, pages (115) à (118). « portabit », « evolvi » et « extensione » sont lus avec doute.} libre ($y''=0$, $y'''=0$) ; « simpliciter fixus », appuyée ou articulée ($y=0$, $y''=0$, la force du stylet valant $F=bc^4y'''$, « qu'il nous importe peu de connaître, puisqu'elle n'influe pas sur le mouvement de vibration ») ; « infixus », encastrée, comme murée ($y=0$, $y'=0$).\footnote{Vue 158 : « quam autem nosse parum nobis refert ; quia in motum vibratorium non influit ». Pour un appui placé au bout, c'est exact : la force n'entre dans aucune des quatre conditions. Pour un appui intermédiaire, elle entre dans les conditions de raccordement ; voir plus bas le § 47.} Ce sont les conditions actuelles. Le scholie 2 (§ 14) en tire six cas, que l'on garde sous les chiffres romains d'Euler. Leurs équations aux fréquences, vérifiées sur la copie, et leurs deux premières racines, calculées, sont :\footnote{Racines calculées pour cette lecture par dichotomie. Les modes que la copie écrit pour chaque cas (vues 161, 167, 176, 185, 188) ont été vérifiés : ils satisfont les conditions aux deux bouts pour les racines exactes.}
\[
\begin{array}{llll}
\text{I} & \text{libre, libre} & \cos\omega\cosh\omega=1 & 4{,}7300 ;\ 7{,}8532\\
\text{II} & \text{libre, appuyée} & \tan\omega=\tanh\omega & 3{,}9266 ;\ 7{,}0686\\
\text{III} & \text{libre, encastrée} & \cos\omega\cosh\omega=-1 & 1{,}8751 ;\ 4{,}6941\\
\text{IV} & \text{appuyée, appuyée} & \sin\omega=0 & \pi ;\ 2\pi\\
\text{V} & \text{appuyée, encastrée} & \tan\omega=\tanh\omega & 3{,}9266 ;\ 7{,}0686\\
\text{VI} & \text{encastrée, encastrée} & \cos\omega\cosh\omega=1 & 4{,}7300 ;\ 7{,}8532
\end{array}
\]
Euler donne ces équations sous les formes $\cos\omega=2/(e^\omega+e^{-\omega})$, $\tan\omega=(e^\omega-e^{-\omega})/(e^\omega+e^{-\omega})$ et $\cos\omega=-2/(e^\omega+e^{-\omega})$ : ce sont les mêmes.

\emph{Cas I, les deux bouts libres (§§ 15 à 22).} L'élimination des quatre coefficients mène à $\cos\omega\cosh\omega=1$.\footnote{Vue 160. La copie écrit $2+pp+qq+(p-q)(e^\omega+e^{-\omega})=0$ puis $p-q=-2\cos\omega$, avec $p=\sin\omega+\cos\omega$ et $q=\sin\omega-\cos\omega$ ; le calcul donne $(q-p)$ dans la première et $p-q=+2\cos\omega$. Les deux fautes de signe se compensent, et le résultat $2-\cos\omega\,(e^\omega+e^{-\omega})=0$ est juste. La transcription laisse ces signes tels qu'elle les lit.} Aucune racine non nulle dans le premier quadrant (§ 16) — la série tronquée $\cos\omega\,(e^\omega+e^{-\omega})=2(1-\omega^4/6)$ n'en est pas une preuve, mais la conclusion est vraie : $\cos\omega\cosh\omega<1$ sur $]0\,;\,4{,}73[$. La première racine est $3\pi/2+\varphi$ avec $\varphi\approx 2/e^{3\pi/2}\approx 1/56$, « à peine un degré » (§ 17) : la racine exacte est $3\pi/2+1^\circ 01'$. La suivante est « un peu moins que $5\pi/2$ », de quelques minutes (§ 18) : il s'en faut de $2{,}7'$. Les sons sont donc sensiblement comme $9,25,49,81\dots$ (§ 19) ; le rapport exact des deux premiers est $2{,}757$ au lieu de $25/9=2{,}778$. Euler nomme l'intervalle « une octave avec un triton » : $25/9$ vaut une octave et $25/18$, quarte augmentée de l'intonation juste ($569$ cents).\footnote{Vue 163 : « Unde si sonus fundamentalis fuerit b, sequens futurus erit cis ». Un triton au-dessus d'une note nommée b ne donne pas cis, quelle que soit la valeur donnée à b dans la notation allemande ; le triton sous cis est g, la note qu'Euler donne au cas I au § 27. La lecture « b » serait à revoir sur l'image ; elle n'est pas corrigée ici.} Le mouvement général est une somme de tels modes (§ 20). Le son fondamental a au moins deux nœuds (§ 21) : sans nœud, le centre de gravité ne pourrait rester immobile ; avec un seul, la verge tournerait sans fin autour de lui. L'argument est juste : en termes actuels, un mode de la verge libre est orthogonal aux deux mouvements rigides, $\int y\,ds=0$ et $\int s\,y\,ds=0$, ce qui exige au moins deux changements de signe ; les nœuds sont en $u=0{,}2242$ et $u=0{,}7758$. Le § 22 calcule les ordonnées aux bouts, égales, et au milieu, $2e^{\omega/2}-\sqrt2\,e^\omega<0$ ; la pente au milieu, trouvée $-\sqrt2$, « serait nulle si le calcul était plus exact » — elle l'est, le mode étant symétrique.\footnote{Vues 161 à 166 et 168, pages (120) à (126). Avec la racine exacte, l'ordonnée du milieu vaut $-139$ pour des ordonnées $2(1+e^\omega)=229$ aux bouts ; Euler trouve $21-111\sqrt2\approx-136$. La suite du calcul de la vue 166 est à la vue 168, page (126).}

\emph{Cas II, un bout libre, l'autre appuyé (§§ 23 à 29).} L'équation est $\tan\omega=\tanh\omega$.\footnote{Vues 168, 169, 167, 170 à 174, dans cet ordre de lecture : pages (126), (127), (128), (129) à (133). À la vue 169, un signe du second dénominateur de la proportion est lu « + » avec doute là où la ligne précédente demande « $-$ » ; les valeurs de $\sin\omega$ et $\cos\omega$ y sont écrites avec $\sqrt2(e^{2\omega}+e^{-2\omega})$ au dénominateur, et le calcul demande que le radical couvre tout : $\sqrt{2(e^{2\omega}+e^{-2\omega})}$. À la vue 167, deux exposants sont écrits sans le facteur 2 que portent les lignes voisines.} La première racine est $\pi+\pi/4-\varphi$ avec $\tan\varphi=e^{-2\omega}\approx 1/2576$, soit « à peine une minute » (§ 24) : il s'en faut exactement de $1{,}34'$, et $\omega=3{,}9266$. Les suivantes sont voisines de $9\pi/4$, $13\pi/4$, …, d'où des sons comme $25,81,169,289$ (§§ 25--26).\footnote{Vue 171. La copie écrit « quartum » pour le quadrant de $2\pi+45^\circ$, qui est le cinquième, et « $\zeta=\frac{\pi}{4}$ » alors que $\zeta$ désigne depuis le § 16 l'angle droit $\pi/2$ ; ni l'un ni l'autre n'affecte les valeurs $5\pi/4$, $9\pi/4$, … qui suivent.} Le rapport $81/25=3{,}24$ (exact : $3{,}2406$) vaut une octave et une quinte augmentée $25/16$, plus un reste de $1{,}0368$, soit $63$ cents : Euler dit « un comma », mot large pour trois commas syntoniques. Le fondamental du cas II est à celui du cas I comme $25:36$ (exact : $0{,}689$ contre $0{,}694$), une « fausse quinte » (§ 27). La figure du fondamental (§ 29), calculée pour $\omega=5\pi/4$, n'a qu'un nœud intérieur, entre $u=\frac15$ et $u=\frac25$ : il est en $u=0{,}264$.\footnote{Vues 173 et 174, pages (132) et (133). Les ordonnées justes, pour $\omega=5\pi/4$ et $e^{2\omega}=2576$, sont $-5150$, $-1171$, $2046$, $3409$, $2487$ et $0$ en $u=0,\frac15,\dots,1$. La copie a les deux premières et la cinquième ; elle écrit « $0{,}2079$ » pour la troisième, que le produit $2576\times0{,}7921$ qu'elle vient d'écrire ne donne pas, et « $3387$ » pour la quatrième. Elle écrit aussi $s=\frac45\omega$ pour $\frac45a$, « $\frac35\omega=\pi$ » pour $\frac45\omega$, et « $\omega=\frac{9\pi}{2}$ » pour la seconde racine, qui est $\frac{9\pi}{4}$.}

\emph{Cas III, un bout libre, l'autre encastré (§§ 30 à 34).} L'équation est $\cos\omega\cosh\omega=-1$, et la racine cherchée est dans le deuxième quadrant, $\omega=\pi/2+\varphi$ avec $\sin\varphi=2/(e^{\pi/2+\varphi}+e^{-\pi/2-\varphi})$.\footnote{Vues 175 à 180, pages (134) à (139). À la vue 176, $k$ est écrit $a^4/bc\omega^4$ et le numérateur de $\sqrt{2g/k}$ sans le $c$ que portent les mêmes formules aux vues 161 et 167 ; la vue 177 le rétablit dans la durée et le son.} Euler trouve d'abord $\varphi=23^\circ 29'$ par la série au premier ordre, puis $16^\circ 58'$ en gardant $\varphi^2$ ; il essaie ensuite $\varphi=12^\circ$ et $\varphi=15^\circ$ par les logarithmes, qui sont justes ($9{,}31788$ contre $9{,}51572$, puis $9{,}41300$ contre $9{,}49590$), et conclut, par une moyenne entre les deux sinus, à $\varphi=16^\circ 33'$, qu'il porte à $16^\circ 45'$, d'où « $\omega=106^\circ 45'$ » (§ 32). La racine est en réalité $\varphi=17^\circ 26'$, $\omega=107^\circ 26'=1{,}8751$ : son interpolation reste en deçà de la racine.\footnote{Vues 178 et 179. Au second essai, la copie tient $1/3{,}4262$ pour $\sin\varphi$, d'où $16^\circ 58'$ ; pris pour $\varphi$ lui-même, ce nombre donnerait $16^\circ 43'$. Le second logarithme est écrit « $495905$ », sans caractéristique ; « semisse » est lu avec doute. L'arrondi qu'Euler adopte ensuite, $\omega\approx\frac35\pi=108^\circ$, est plus près de la racine vraie que sa valeur calculée.} Il en déduit que ce fondamental est au fondamental du cas I comme $4:25$, soit $1:6\frac14$ (exact : $1:6{,}363$), un peu plus de deux octaves ; les racines suivantes, $3\pi/2-\varphi$ avec $\varphi$ d'environ un degré (exact : $1^\circ 03'$), puis $5\pi/2$, $7\pi/2$, …, donnent les mêmes sons que le cas I (§ 33).

\emph{Cas IV, les deux bouts appuyés (§§ 35 à 38).} Ici $\alpha=\beta=\delta=0$ et $\sin\omega=0$ : $\omega=\pi,2\pi,3\pi,\dots$ exactement, les modes sont $\sin(n\pi u)$, la figure d'une corde, et les sons croissent comme $1,4,9,16,\dots$\footnote{Vues 181 à 184, pages (140) à (143). À la vue 181, le signe entre $e^\omega$ et $e^{-\omega}$ est lu « + » avec doute ; la substitution $\beta=-\alpha$ demande « $-$ », et la somme qui suit l'a. Les noms de notes de la vue 182 ($C$, $c$, $d$, $c$, $gis$, $d$ pour $1,4,9,16,25,36$) s'accordent avec ces rapports ; les barres d'octave comptées par la transcription sur les deux derniers en donnent une de plus que le calcul.} C'est, remarque Euler, le seul cas où toutes les racines sont exactes, les exponentielles disparaissant du calcul (§ 38) ; il en conclut qu'une telle verge rendrait des sons plus purs qu'une corde, n'ajoutant guère au fondamental que la double octave. Le § 37 compare les quatre fondamentaux, notés gis, d, C, fis, comme $\frac94$, $\frac{25}{16}$, $\frac{9}{25}$, $1$ ; les valeurs exactes sont $2{,}267$, $1{,}562$, $0{,}356$ et $1$.

\emph{Cas V et VI (§§ 39 à 44).} Le cas V (appuyé, encastré) a l'équation du cas II, le cas VI (deux bouts encastrés) celle du cas I : mêmes sons, figures différentes. Euler calcule la figure du fondamental du cas V, $y=e^{u\omega}-e^{-u\omega}+\Omega\sin u\omega$ avec $\Omega=\sqrt{2(e^{2\omega}+e^{-2\omega})}\approx72$, et celle du cas VI aux tiers de la longueur.\footnote{Vues 184 à 190, pages (143) à (149). Cas V : la copie écrit $\surd2(e^{2\omega}+e^{-2\omega})$, sans barre, et la valeur $72=\sqrt{2\times2576}$ montre que le radical porte sur tout le produit ($71{,}75$ exactement). Ses ordonnées en $u=\frac15,\frac25,\frac35,\frac45$ sont « 56 », 76, 61, 23 ; la première vaut $52{,}7$ avec ses propres nombres. Dans la formule générale de la vue 187, les exposants n'ont pas le facteur $u$ que porte le sinus et le temps $t$ manque aux deux dernières lignes. Cas VI : les ordonnées $138$ et $140$ en $u=\frac13$ et $\frac23$ se vérifient pour $\omega=3\pi/2$ ; la copie écrit « ipsi $\omega$ » pour les valeurs données à $u$, et le chiffre de page de la vue 189 est sous une tache.} Pour ce dernier, il trouve $138$ et $140$ et note qu'un calcul plus exact les donnerait égales : c'est vrai, le fondamental du cas VI étant symétrique.

\emph{Verges semblables et sections quelconques (§§ 45 et 46).} Sous des circonstances semblables, les sons sont comme $c/a^2$, directement comme l'épaisseur et inversement comme le carré de la longueur : c'est la loi actuelle, la fréquence étant proportionnelle au rayon de giration de la section divisé par le carré de la longueur. Pour une section rectangulaire $ABCD$, $c$ est le côté situé dans le plan de flexion et l'autre côté n'entre pas dans le calcul : exact aussi, la largeur multipliant à la fois la rigidité et la masse. Pour une section circulaire de diamètre $AB$, la copie donne « $c=\frac34AB$ ». L'équivalence juste s'obtient en égalant les rayons de giration, $c/\sqrt{12}$ pour le rectangle et $AB/4$ pour le cercle : $c=\frac{\sqrt3}{2}AB\approx0{,}866\,AB$, c'est-à-dire $c^2=\frac34AB^2$.\footnote{Vues 190 et 192, pages (149) et (150). La valeur $\frac34$ de la copie est celle du rapport des carrés, $c^2/AB^2$ ; la lecture n'essaie pas de dire si l'écart est d'Euler, de l'imprimé ou de la copie. Juste avant, après « fuerit », une tache couvre une ou deux lettres, et le mot lu avec doute « parest » ne fait pas sens.}

\subsection*{La verge appuyée en un point intermédiaire (§ 47 à 51)}

\pagerange{192}{202}
Euler ajoute un cas « pour faire comprendre la méthode » : la verge appuyée à ses deux bouts $E$ et $F$ l'est encore en un point $L$, avec $EL=\lambda a$, $0<\lambda<1$.\footnote{Vues 192 à 200 et 201--202, pages (150) à (158). La lettre $\lambda$ est une petite lettre crochue, lue grecque par la transcription ; la marge cite « fig. 7 » et « Tab. III fig. 8 ».} Si la verge était encastrée en $L$, dit-il, les deux portions seraient sans communication ; appuyée sur un stylet autour duquel elle peut tourner, elle transmet le mouvement d'une portion à l'autre, mais la continuité de la courbe est interrompue et chaque portion a sa propre expression, $\alpha,\beta,\gamma,\delta$ pour $EL$ et $\alpha',\beta',\gamma',\delta'$ pour $LF$, avec le même facteur du temps. En $L$ il impose $y=0$ des deux côtés, et l'égalité des pentes et des rayons osculateurs, $y'$ et $y''$ continus ; aux bouts, deux conditions chacun. Huit équations homogènes pour huit coefficients : la condition de compatibilité est l'équation finale
\[
0=2-2e^{2\omega}-\bigl(e^{\lambda\omega}-e^{(2-\lambda)\omega}\bigr)\bigl(e^{\lambda\omega}-e^{-\lambda\omega}\bigr)\frac{\sin\omega}{\sin\lambda\omega\,\sin(1-\lambda)\omega}.
\]
Le décompte d'Euler est exact : $y'''$ n'est pas supposé continu en $L$, et son saut est la force du stylet, l'effort $F$ de l'équation II du § 8. Avec $\mu=1-\lambda$, l'équation finale s'écrit
\[
\coth\lambda\omega+\coth\mu\omega=\cot\lambda\omega+\cot\mu\omega,
\qquad\text{ou}\qquad
\sinh\omega\,\sin\lambda\omega\,\sin\mu\omega=\sin\omega\,\sinh\lambda\omega\,\sinh\mu\omega,
\]
car le second membre d'Euler vaut exactement
\[
-\,\frac{4e^\omega\bigl(\sinh\omega\,\sin\lambda\omega\,\sin\mu\omega-\sin\omega\,\sinh\lambda\omega\,\sinh\mu\omega\bigr)}{\sin\lambda\omega\,\sin\mu\omega} ;
\]
c'est l'équation aux fréquences qu'on écrit aujourd'hui pour la poutre continue sur trois appuis.\footnote{Vérifié pour cette lecture : les racines de cette équation coïncident avec celles du déterminant $8\times8$ des conditions du § 47, calculé directement, pour $\lambda=\frac12,\frac13,0{,}3,1/\sqrt2,0{,}05,0{,}01$. À la vue 194, deux équations portent le numéro III ; à la vue 195, la parenthèse fermante d'un membre est écrite dans le numérateur. Ce sont des accidents de copie sans effet.} Euler déclare ne pas oser en chercher les racines, « à cause de la complication des formules ».\footnote{Vue 198 : « quem quidem laborem, ob formulas tantopere complicatas, suscipere non ausim » (« ausim » lu avec doute).}

Il remarque ensuite que $\lambda=0$ ou $\lambda=1$ ramène l'équation à $0=2(1-e^{2\omega})$, donc à $\omega=0$ et au repos. La conclusion vient de ce qu'il fait $\lambda=0$ dans une expression où $\sin\lambda\omega$ est au dénominateur ; la limite de l'équation quand $\lambda\to0$ est $\coth\omega=\cot\omega$, c'est-à-dire $\tan\omega=\tanh\omega$, le cas V : un appui collé contre un bout appuyé encastre ce bout, ce qu'Euler dit lui-même pour $\lambda=1$, « tum enim terminus $F$ erit quasi muro infixus ».\footnote{Vue 198. Pour $\lambda=0{,}01$, la plus petite racine est $3{,}953$, déjà voisine de la racine $3{,}9266$ du cas V. Ce qui tend vers le repos, c'est la petite portion $EL$, non la verge.}

Pour $\lambda=\frac12$ (§ 48), l'équation se ramène à
\[
\sin\tfrac12\omega\,\Bigl[(1+e^\omega)\sin\tfrac12\omega-(e^\omega-1)\cos\tfrac12\omega\Bigr]=0,
\]
d'où deux familles.\footnote{Vue 199, page (155). Dans la troisième équation, l'exposant du premier $e$ est surchargé et lu « $2\omega$ » avec doute ; la division par $1-e^\omega$ qui y mène donne $2(1+e^\omega)$, et c'est ce que portent ses deux citations françaises (vues 210 et 224).} La première, $\sin\frac12\omega=0$, soit $\omega=2i\pi$ (§ 49), fait vibrer chaque moitié comme une verge appuyée à ses deux bouts, avec des sons $4,16,36,64,\dots$, deux octaves au-dessus du cas IV. Au § 50, Euler en calcule les coefficients et trouve que tous s'annulent sauf $\gamma$ et $\gamma'$, avec $\gamma'=-\gamma$. Le signe est faux : pour $\omega=2i\pi$, l'égalité des pentes en $L$ exige $\gamma'=\gamma$, et le mode est le seul sinus $y=\sin\omega u$ tout le long de la verge ; avec $\gamma'=-\gamma$, les deux portions feraient en $L$ un angle. Le $-1$ vient de ce que ces coefficients sont des formes $0/0$ en $\sin\frac12\omega=0$, auxquelles Euler donne une valeur en multipliant tout par $\sin\frac12\omega$.\footnote{Vues 200 et 201, pages (156) et (157). Vérifié : le vecteur $(\gamma,\gamma')=(1,1)$, les six autres coefficients nuls, annule les huit équations pour $\lambda=\frac12$, $\omega=2\pi$ ; le vecteur $(1,-1)$ laisse un résidu dans l'équation des pentes. Dans $\beta'$, la copie écrit deux fois l'exposant $\frac12\omega$ ; sa citation française de la vue 212 porte $e^{\frac12\omega}-e^{-\frac12\omega}$. Un mot court est biffé d'une rature épaisse devant « $\sin\frac12\omega=0$ » et n'est pas lu.}

La seconde famille (§ 51), $\tan\frac12\omega=\tanh\frac12\omega$, est l'équation du cas V pour la demi-longueur : chaque moitié vibre « comme si elle était simplement appuyée en $E$ et en $F$, et entièrement encastrée en $L$ », avec $\omega\approx\frac{5\pi}{2},\frac{9\pi}{2},\dots$ (exactement $7{,}8532$, puis $14{,}1372$) et des sons deux octaves au-dessus du cas V. Euler conclut qu'il est « tout à fait digne de remarque » que les deux moitiés puissent vibrer de deux manières, l'une qui relève du cas IV et l'autre du cas V.\footnote{Vues 201 et 202. La copie se lit « altero, qui cum casu superiore quinto, altero vero, qui cum casu quinto congruit », de part et d'autre du changement de page ; pour la première manière, le sens demande le cas IV ; sa citation française (vue 228) écrit « l'une qui se rapporte au cas IV et l'autre au cas V ».} Les deux familles sont de vrais modes de la verge sur trois appuis : la première est formée des modes antisymétriques par rapport au milieu, la seconde des modes symétriques, dont la pente est nulle en $L$ par symétrie et qui font travailler le stylet. La seconde solution d'Euler est juste ; c'est elle que les « Remarques » vont contester.

\subsection*{La verge courbe et l'anneau (§ 52 à 55)}

\pagerange{202}{206}
Le § 52 étend tout ce qui précède à une verge naturellement courbe, en prenant pour axe la courbe naturelle elle-même et pour $y$ le déplacement normal : même équation, mêmes six cas, mêmes sons.\footnote{Vues 202 et 203, pages (158) et (159). La phrase « dum ante nobis erat linea recta cum figura curvae congruens » est ajoutée d'une plume plus fine au-dessus de la ligne, avec un renvoi ; « hance » est lu avec doute et peut n'être que « hanc » ; « supra curvaturam naturam » est écrit ainsi.} Ce n'est vrai qu'à la limite d'une courbure faible devant celle des ondulations : dans une verge courbe, déplacement normal et déplacement tangentiel sont liés, et l'équation d'une verge droite ne vaut plus.

Les §§ 53 à 55 traitent de la verge fermée sur elle-même, « de sonis annulorum elasticorum ».\footnote{Vues 203, 205 et 206, pages (159) à (161). Le chiffre de la figure, empâté, est lu « 10 » d'après la suite des figures. À la vue 206, un mot est biffé sous une rature épaisse avant « percipietur », écrit au-dessus, et n'est pas lu.} Euler pose $y=0$ au point $A$ pris pour origine et demande qu'elle s'annule encore en $s=a,2a,3a,\dots$, $a$ étant le périmètre : les termes exponentiels disparaissent, $\delta=0$, et $\sin i\omega=0$ pour tout $i$, d'où $\omega=i\pi$ et des sons comme $1,4,9,16,\dots$ Il en conclut que de tels anneaux, comme les disques et les « bassins en forme de cloche », rendent des sons plus purs que les cordes, et que leurs sons sont en raison inverse du carré du périmètre. Dans son propre modèle, la condition d'un anneau n'est pas que $y$ s'annule en $A$ à chaque tour, mais que $y$ et ses dérivées reprennent leurs valeurs : $\sin(i\pi s/a)$ change de signe après un tour quand $i$ est impair, ce qui ferait un angle en $A$. Seuls $\omega=2j\pi$ conviennent, et les sons du modèle sont alors comme $4,16,36,\dots$, c'est-à-dire comme $1,4,9,\dots$ entre eux. La loi actuelle de l'anneau mince est autre, et ce n'est pas une loi de carrés ; elle est énoncée une fois, avec l'anneau de Sophie Germain et les observations de Chladni, à la sous-section « L'anneau et la lame courbe, contre Euler et avec Chladni » du long mémoire (vues 574 à 592).

\section*{« Remarques sur le mémoire d'Euler » : la seconde solution contestée}

\pagerange{208}{334}
Après la copie viennent, de sa main pour les brouillons, plusieurs états d'une critique du § 47 d'Euler.\footnote{La transcription des vues 210 à 316 et 320 à 389 cite ces feuillets comme les siens : la première personne du féminin y est constante (« afin d'être autorisée », vues 240 et 332). La mise au net des vues 320 à 389 est d'une main régulière, que les transcriptions attribuent diversement ; la lecture ne tranche pas.} On lit ici la mise au net comme texte principal jusqu'à son § 7, et les brouillons en notes ; la suite, à partir de son § 8, est la section suivante.

\subsection*{Deux titres, deux brouillons, une mise au net}

\pagerange{208}{220}
Un feuillet de titre, « Remarques sur la lame elastique », ouvre l'ensemble.\footnote{Vue 208, n\textsuperscript{o} 101 de la série à l'encre ; rien d'autre sur le feuillet, verso blanc.} Suivent deux rédactions. La première, « Remarque sur le mémoire d'Euler : Investigatio motuum … in act. Acad. pétrop ann. 1779 P.I p. 103 et seq. », tient quatre feuillets (vues 210 à 216) : elle résume le § 47 à la troisième personne, conteste la seconde solution et commence la transformation de l'équation finale.\footnote{Vues 210 à 216, n\textsuperscript{os} 102 à 105 de la série à l'encre ; ses feuillets sont numérotés 2, 3, 4 en haut à gauche, le premier chiffre étant sous une tache. Le renvoi « (2d N\textsuperscript{o} 35) » de la vue 210 est lu avec doute ; le mot « pu » de la vue 212, qui ne fait pas sens, aussi. « Premier » et « second » brouillon sont ici des noms d'ordre de reliure : rien, sur les feuillets, ne dit lequel fut écrit d'abord.} La seconde, « Remarques sur le memoire d'Euler … », repart du titre avec sa propre pagination, 1 à 42 (vues 218 à 308) : une ouverture historique, une traduction française des §§ 47 à 51, puis des remarques numérotées, qui se poursuivent jusqu'au problème de la lame libre et aux plaques carrées. Enfin un feuillet de titre, « Memoire presenté et examiné / sur la lame » (vue 318), précède une mise au net paginée (1) à (35) (vues 320 à 389), qui reprend l'argument dans le même ordre, en plus bref.\footnote{Vue 318, n\textsuperscript{o} 155 ; « examiné » est lu avec doute, sa fin étant repassée. Vue 320, page (1), n\textsuperscript{o} 156 ; les pages (2) à (11) sont aux vues 322 à 340, chacune au recto d'un feuillet dont le verso est blanc, les numéros entre parenthèses barrés d'un trait. Le titre y est écrit « laminæ e virgæ elasticæ » ; la transcription des vues 320 et suivantes lit la fraction de la longueur $x$, par une lettre proche d'un $\lambda$ : on écrit $\lambda$ partout, comme dans la copie d'Euler.}

L'ouverture de la seconde rédaction situe la question : Euler avait déjà traité ces matières dans la \emph{Methodus inveniendi curvas maximi minimive proprietate gaudentes}, « imprimée en 1744 » ; le mémoire de 1779 les met dans un ordre plus méthodique, sans « aucun changement essentiel dans sa méthode », et « l'addition la plus importante » est l'examen d'un point d'appui placé ailleurs qu'aux extrémités. Elle annonce qu'elle rapportera cette partie du mémoire « en y joignant les problèmes relatifs aux cas qu'il a négligé de traiter » et, ajoute la marge, « les reflexions que l'examen plus aprofondi de cette matière m'a fait naitre ».\footnote{Vue 218, page 1, n\textsuperscript{o} 106. Le premier mot, une abréviation de deux lettres, est couvert d'une rature et n'est pas lu. L'addition marginale n'a pas de signe d'insertion dans la ligne ; la transcription la place d'après le sens. La page citée après « l'auteur » est lue « p. 151 » ou « 131 » ; le § 47 commence à la page (151) de la copie latine (vue 193), ce qui s'accorde avec la première lecture.}

\subsection*{Le § 47 d'Euler, en français}

\pagerange{218}{228}
La seconde rédaction traduit alors les §§ 47 à 51, chaque ligne ouverte d'un guillemet, au plus près du latin de la copie : l'énoncé, les huit équations, l'équation finale, le cas $\lambda=\frac12$ et ses deux facteurs, les coefficients du § 50 avec $\gamma=-\gamma'$, et la seconde solution du § 51.\footnote{Vues 218 à 228, pages 1 à 6, n\textsuperscript{os} 106 à 111. La citation se ferme sur un guillemet à la troisième ligne de la vue 228. Écarts relevés par la transcription : le facteur du temps est écrit $c\sin(\zeta+t\surd\frac{g}{k})$ (vue 220), comme dans la copie latine de la vue 193 ; la valeur de $\gamma$ est écrite sans son dénominateur $\sin\lambda\omega$ et le second bout est noté $E$ au lieu de $F$ (vue 222) ; un verbe manque après « la question est réduite » et une formule est biffée d'un zigzag après « réduite a celle ci » (vue 224) ; « $s=\frac{u}{a}$ » est écrit pour $u=s/a$ (vue 226). Aucun de ces écarts ne touche l'argument.} La mise au net remplace cette traduction par un résumé à la troisième personne (vue 320), et la première rédaction faisait de même (vues 210 et 212) : l'auteur, après les six cas, suppose la lame appuyée en outre en un point intermédiaire « de façon que la communication entre ces différentes parties ne soit pas interrompue », choisit le cas le plus simple, obtient l'équation finale, renonce à la résoudre, examine $\lambda=\frac12$, et termine sur les deux manières de vibrer.

\subsection*{L'objection : la condition $\sin\omega = 0$ (§ 1)}

\pagerange{320}{322}
Sa « première réflexion » est que la seconde solution « sort de l'état de la question ».\footnote{Vue 322, page (2), n\textsuperscript{o} 157, son paragraphe I ; la marge porte « fig. 7 » (le premier mot lu avec doute). À la dernière ligne, « puisse est comprise » est écrit pour « puisse être comprise ». Même paragraphe au brouillon, vue 228 (page 6), marqué « I. », avec « le respect dû a l'opinion d'Euler » et un renvoi « (… cas IV §. 35 \&c.) » dont la première lettre n'est pas lue.} Les mouvements de la lame soumise au stylet\footnote{Le feuillet écrit « stilet » ; on écrit stylet, comme pour le « stylus » de la copie latine.} en $L$, dit-elle, doivent être au nombre de tous les mouvements réguliers dont la même lame est capable quand ses extrémités sont appuyées — Euler n'a-t-il pas employé pour les trouver quatre équations tirées de l'état des extrémités ? Or ces mouvements exigent $\sin\omega=0$ (cas IV, § 35). La première solution s'y accorde, $\sin\frac12\omega=0$ entraînant $\sin\omega=0$ ; la seconde non, car $\tan\frac12\omega=\tanh\frac12\omega$ donne $\sin\omega=\tanh\omega$, qui ne s'annule que pour $\omega=0$.\footnote{Le calcul $\sin\omega=\tanh\omega$ est celui des cinq formules marginales de la vue 214, dans la première rédaction : de $\tan\frac12\omega=(e^\omega-1)/(e^\omega+1)$ elle tire $\cos^2\frac12\omega=(e^\omega+1)^2/2(e^{2\omega}+1)$, puis $\sin\omega=(e^{2\omega}-1)/(e^{2\omega}+1)$ (un facteur 2, sous une tache, est lu avec doute). Le calcul est juste.} Le respect que l'on doit aux opinions d'Euler la décide pourtant à examiner le cas général.

Le calcul est juste ; la prémisse ne l'est pas. Elle suppose que le stylet ne fait que choisir parmi les mouvements déjà possibles de la lame sans stylet. La première rédaction le dit en toutes lettres : l'effet des points d'appui « est considéré comme purement exclutif », il ne peut « faire naître des mouvemens inconciliables avec l'état des extrèmités » et « se borne a déterminer quelques uns des mouvemens possibles déja connus en rendant les autres impossibles ».\footnote{Vue 214, feuillet 3 de la première rédaction. La même rédaction ajoute (vue 212) que « la supposition d'un point fixe au milieu … équivaut a sa séparation en deux moitiés entièrement indépendantes », et (vue 216) qu'excepté $\lambda=\frac12$ les deux portions séparées par un appui fixe ne peuvent rendre le même ton ; elle annonce enfin (vue 216) la transformation de l'équation finale reprise au § 5 de la mise au net.} C'est supposer que le stylet n'exerce aucune force. Or il en exerce une en général : c'est l'effort $F$ qu'Euler fait entrer dans les équations I et II de son § 8 précisément pour une verge « fixée comme au moyen d'un stylet ». Un mode de la lame avec stylet n'est un mode de la lame sans stylet que si cette force est nulle, c'est-à-dire si $y'''$ est continu en $L$ ; pour les autres modes, rien n'oblige $\sin\omega=0$. Pour $\lambda=\frac12$, la première famille ne charge pas le stylet, la seconde le charge : c'est exactement la différence qu'elle a vue, mais elle en tire la conclusion inverse de la vraie.

\subsection*{Le cas général : points de repos « sympathiques » et parties aliquotes (§ 2 à 5)}

\pagerange{322}{330}
Elle reprend l'équation du § 47 sous la forme
\[
0=2(1-e^{2\omega})\sin\lambda\omega\,\sin\mu\omega-\bigl(1-e^{2\mu\omega}\bigr)\bigl(e^{2\lambda\omega}-1\bigr)\sin\omega,\qquad \mu=1-\lambda,
\]
et, pour $\lambda=1/n$, la divise par $e^{2\lambda\omega}-1$ :\footnote{Vues 322 et 324, pages (2) et (3), son paragraphe 2. Au brouillon (vue 230, page 7, « 2. »), le second terme est précédé d'un « $-$ » : c'est une faute de signe, que la mise au net corrige. Vérifié numériquement pour $n=3,4,5$ : la forme de la mise au net est exactement l'opposé de l'équation d'Euler divisée par $e^{2\lambda\omega}-1$.}
\[
2\bigl(1+e^{2\lambda\omega}+e^{4\lambda\omega}+\dots+e^{2\mu\omega}\bigr)\sin\lambda\omega\,\sin\mu\omega+\bigl(1-e^{2\mu\omega}\bigr)\sin\omega=0 .
\]
Elle observe que « par la nature de la question » $\lambda$ est toujours une fraction et que, les deux extrémités étant appuyées, cette fraction est « partie aliquote de l'unité, ou multiple d'une telle partie ».\footnote{C'est une hypothèse, non une conséquence : le stylet peut être posé en tout point, et c'est précisément ce qui est en question. Au brouillon (vue 230), les mots « ou multiple d'une telle partie » sont ajoutés sous la ligne.} Par les développements classiques de $\sin n\theta$ en puissances de $\sin\theta$ — $\sin\omega=\frac1\lambda\sin\lambda\omega-\frac{1-\lambda^2}{6\lambda^3}\sin^3\lambda\omega+\dots$ pour $n$ impair, le même précédé de $\cos\lambda\omega$ et avec $1-4\lambda^2$, $1-16\lambda^2$ pour $n$ pair, développements qui sont justes —, $\sin\lambda\omega$ est facteur de l'équation. D'où la solution $\omega=i\pi/\lambda$ : la lame se partage en $1/\lambda$ parties égales qui vibrent comme autant de lames appuyées à leurs deux bouts, « tant dans le point où est placé le stilet, que dans les autres points de repos que l'on pourroit nommer sympatiques », et rendent des sons proportionnels à $1,4,9,16,\dots$ divisés par $\lambda^2$.\footnote{Vue 324. Le « $=0$ » qui suit « la condition $\sin\omega$ » à la dernière ligne est lu avec doute.} Tout cela est vrai, et c'est la bonne notion : ce sont les modes $\sin(i\pi u/\lambda)$ de la lame sans stylet dont un nœud tombe sur le stylet, et pour eux le stylet n'exerce aucune force.

Au § 3, elle généralise les coefficients d'Euler et trouve, comme lui, $\gamma=-\gamma'$ ; elle décrit alors les parties successives $EL$, $LL'$, $L'L''$, … par $y=\pm C\sin(\dots)\sin(\dots)$, le signe $+$ pour les arcs au-dessus de l'axe, le signe $-$ pour ceux qui sont au-dessous.\footnote{Vue 326, page (4), son paragraphe 3 ; première mise au net à la vue 232 (feuillet monté sur onglet, sans numéro de page, écrit sur son cinquième supérieur) et reprise à la vue 234 (page 8), avec renvoi marginal à la « fig. 7 ». L'exposant de $L$ dans « $L^{\text{vi}}E'$ » est lu avec doute ; le dernier facteur des équations de $y$ est écrit « sin » suivi de la lettre de la fraction et de $\omega$.} Le signe d'Euler est faux, on l'a vu (§ 50) ; sa description reste juste si l'abscisse de chaque partie est comptée depuis son propre bout gauche : le mode est un seul sinus, dont les arches alternent.

Au § 4, elle montre que le second facteur de l'équation ne peut s'annuler, si l'on impose $\sin\omega=0$, que pour $\omega=0$ : c'est exact, et c'est dire seulement que la seconde solution n'est pas un mode de la lame sans stylet. Elle annonce alors davantage : que, même sans cette condition, « il se trouve plusieurs valeurs » de $\lambda$ pour lesquelles l'équation n'a aucune solution du type du cas V, où chaque portion serait « fixée » au stylet et appuyée à son autre bout.\footnote{Vue 328, page (5), son paragraphe 4 ; brouillon aux vues 236 et 238 (pages 9 et 10, « 4. » et le début de « 5 »). Au brouillon, le premier terme de la première accolade est écrit $\frac2\lambda$ (lu avec doute) là où la mise au net et l'autre accolade ont $\frac1\lambda$.} Au § 5, elle récrit l'équation finale
\[
\sin\lambda\omega\,\bigl(\coth\mu\omega\,\sin\mu\omega-\cos\mu\omega\bigr)+\sin\mu\omega\,\bigl(\coth\lambda\omega\,\sin\lambda\omega-\cos\lambda\omega\bigr)=0,
\]
forme juste et symétrique en $\lambda$ et $\mu$ ;\footnote{Vue 330, page (6), son paragraphe 5 ; brouillons aux vues 216 et 238. Elle l'écrit avec les exponentielles, $\frac{e^{\mu\omega}+e^{-\mu\omega}}{e^{\mu\omega}-e^{-\mu\omega}}$ pour $\coth\mu\omega$. Divisée par $\sin\lambda\omega\sin\mu\omega$, c'est l'équation $\coth\lambda\omega-\cot\lambda\omega+\coth\mu\omega-\cot\mu\omega=0$ donnée plus haut. Vérifié numériquement.} pour que le stylet pût être tenu pour fixe, il faudrait que les deux parenthèses s'annulent ensemble, $\tan\lambda\omega=\tanh\lambda\omega$ et $\tan\mu\omega=\tanh\mu\omega$ : les deux portions vibreraient alors comme deux lames du cas V de même son. La caractérisation est exacte : ce sont les modes dont la pente est nulle au stylet. Avec les valeurs approchées d'Euler, $\omega\approx(4k+1)\pi/4$, les deux suites de sons, proportionnelles à $(4k+1)^2/\lambda^2$ et $(4j+1)^2/\mu^2$, n'ont aucun terme commun quand, $\lambda=1/\lambda'$, le nombre $\lambda'-1$ est pair ou de la forme $4n+3$ : l'arithmétique est juste, puisque l'égalité demande $(4k+1)(\lambda'-1)=4j+1$.\footnote{Au brouillon (vue 240, page 11), elle poursuit jusqu'à $\lambda'=6$, où « ce nombre est en effet égal au premier de la première suite ». Avec les racines exactes de $\tan z=\tanh z$, l'égalité n'est qu'approchée : $5z_1=19{,}6330$ contre $z_6=19{,}6350$. La mise au net ne reprend pas ce cas. Au brouillon, le numérateur de $(1-\lambda)^2$ est sous une tache, et deux chiffres sont lus 1 avec doute.} Mais aucun de ces faits ne touche l'existence des autres modes. Pour tout $\lambda$, l'équation a une infinité de racines : pour $\lambda=0{,}3$, $\omega=5{,}132$, $9{,}277$, $11{,}780$, … ; pour $\lambda=\frac13$, $\omega=5{,}335$, $3\pi$, $11{,}144$, …, où seul $3\pi$ est un mode « sympathique ».

\subsection*{La conclusion contre Euler, et Euler contre lui-même (§ 6 et 7)}

\pagerange{332}{334}
Elle conclut :
\begin{quote}
J'ai voulu entrer dans tous ces détails, àfin d'être autorisée a conclure qu'Euler a été induit en erreur, lorsqu'il a cru pouvoir admettre sa seconde solution…
\end{quote}
et elle ajoute que, sans sortir de ce sujet, Euler a lui-même « varié de sentiment » : dans la \emph{Methodus}, il aurait rejeté toutes les solutions de la lame libre qui donnent un nombre impair de nœuds, à cause du « mouvement giratoire » autour du nœud du milieu, et, dans le mémoire de 1779, il les admet et n'emploie la considération de ce mouvement qu'à prouver que le fondamental de la lame libre coupe l'axe en deux points (§ 21).\footnote{Vue 332, page (7), son paragraphe 6 ; renvois « V. methodus \&c. P. 298 § 83 » et « V. Memoire § 21 ». Ce qu'elle dit de la \emph{Methodus} est son rapport ; la lecture ne l'a pas vérifié. Le brouillon (vue 240, page 11, « 6 », très corrigé, et le feuillet de la vue 242, en tête « seconde P… », la fin non lue) porte le même argument, avec en plus Giordano Riccati : son mémoire « Delle Vibrazioni Sonore de cilindri », au premier volume des Mémoires de la Société italienne (p. 457), « s'étonne de cette opinion » et remarque que le calcul permet d'échapper à une conclusion contraire à l'expérience ; Riccati, dit-elle, ne paraît pas avoir connu le mémoire de 1779, « quoique le volume qui renferme le sien soit de 1782 ». Le brouillon écrit « P. 898 §. 83 » à côté du nom de Riccati et « le Scholie §. 81 » là où la mise au net écrit « P. 298 § 83 » pour la \emph{Methodus} et « § 21 » pour le mémoire ; le § 21 de la copie latine est bien le scholie 2 sur les deux nœuds (vue 164). Ces deux lectures du brouillon seraient à revoir sur l'image.} La conclusion ne suit pas : les détails qu'elle a réunis sont justes, mais ils portent sur les modes sans force au stylet, et la seconde solution d'Euler est un mode avec force. Sur le second point, le mémoire de 1779 a raison contre la \emph{Methodus} telle qu'elle la rapporte : les modes de la lame libre sont alternativement symétriques et antisymétriques, ces derniers ont un nombre impair de nœuds, et aucun ne comporte de rotation d'ensemble, leur quantité de mouvement et leur moment cinétique étant nuls.

Au § 7, elle revient à la première solution : elle est générale, elle montre comment les nœuds se placent selon $\lambda$, puisque $y=0$ encore aux points où $\sin2\lambda\omega$, $\sin3\lambda\omega$, … s'annulent, et le mouvement est le même quel que soit celui de ces points où l'on pose le stylet.\footnote{Vues 332 et 334, pages (7) et (8), son paragraphe 7. La condition « pourvu toutefois qu'en désignant par $\sin n\lambda\omega$ un des sinus de cette suite, $n$ ne soit pas facteur du dénominateur de la fraction » est écrite « diviseur » au brouillon (vue 244, page 12, alinéa marqué 7 en marge, « Je reviens a present a 12 »).} Elle en tire que « tout mouvement régulier, et par conséquent tout son mesurable, est impossible à obtenir de la lame dont les deux extrémités sont appuyées, lorsque le stilet est placé à tout autre point qu'à une division aliquote ». C'est faux : pour toute position du stylet, la lame a une suite infinie de mouvements réguliers ; ce qui est vrai, c'est que hors des positions rationnelles aucun d'eux n'est un mode de la lame sans stylet, tous faisant travailler le stylet. Elle note l'analogie avec la corde vibrante et regrette que ce cas des extrémités appuyées soit « le seul où de semblables analogies se fassent remarquer » ; puis elle annonce le problème qu'elle juge important : appliquer la méthode d'Euler à la lame libre à ses deux bouts, touchée par un stylet en un point quelconque.\footnote{Vue 334, page (8). La marge porte, dans un cadre, une « N.ta » : les vibrations transversales des plaques élastiques carrées « y sont comprises, puisque la largeur n'influe pas » (« traité » d'acoustique de Chladni, le mot lu avec doute). La largeur n'influe pas sur le son d'une lame (§ 46 d'Euler), mais une plaque carrée n'est pas une lame : voir la section suivante et celle du long mémoire. Au brouillon (vue 244), le regret porte sur ce que ce cas « n'est pas celui qui se presente le plus commodement dans l'experience \& la pratique » (« \& » lu avec doute).} Ce problème, son § 8, ouvre la section suivante.

\subsection*{Ce que dit la théorie actuelle de la poutre sur appui intermédiaire}

\pagerange{320}{334}
On énonce ici une fois, pour cette section et pour la suivante, ce que l'on sait aujourd'hui de la lame touchée par un stylet, dans le modèle même d'Euler (équation du § 4 sans effort axial).\footnote{Énoncés actuels, que les feuillets précèdent et ne connaissent pas ; ils ont été vérifiés numériquement pour cette lecture sur l'équation du § 47.}
\begin{enumerate}
\item \emph{Conditions au stylet.} En $L$ ($u=\lambda$) : $y=0$ ; $y$, $y'$, $y''$ continus ; $y'''$ saute d'une quantité proportionnelle à la force $R$ que le stylet exerce, $bc^4\,[y''']_L=-R$ à la convention de signe près. Ce sont les quatre conditions d'Euler au § 47, la force étant l'inconnue que son § 8 introduit.
\item \emph{Existence.} Pour toute position $0<\lambda<1$, la lame appuyée à ses bouts et touchée par le stylet a une suite infinie de modes, racines de $\coth\lambda\omega+\coth\mu\omega=\cot\lambda\omega+\cot\mu\omega$, $\mu=1-\lambda$ ; ajouter l'appui relève chaque fréquence sans dépasser la suivante de la lame sans stylet, dont les racines sont $\pi,2\pi,3\pi,\dots$\footnote{Pour $\lambda=0{,}3$ : $5{,}132\in[\pi,2\pi]$, $9{,}277\in[2\pi,3\pi]$, $11{,}780\in[3\pi,4\pi]$. C'est le théorème de séparation des valeurs propres pour une contrainte ajoutée, dû à Rayleigh et généralisé par Courant et Fischer ; les noms seuls sont donnés ici, les feuillets les précèdent.}
\item \emph{Modes sans force au stylet.} Un mode de la lame avec stylet n'exerce aucune force sur lui si et seulement si c'est un mode de la lame sans stylet dont un nœud tombe en $L$ : si $R=0$, $y'''$ est continu et $y$ vérifie l'équation d'un bout à l'autre ; réciproquement un tel mode satisfait toutes les conditions avec $R=0$. Les bouts étant appuyés, cela demande $\sin k\pi\lambda=0$ : $\lambda=p/q$ rationnel, $k$ multiple de $q$. Ce sont ses « points sympathiques ». Pour $\lambda$ irrationnel, il n'y en a aucun, et tous les modes chargent le stylet.
\item \emph{Le stylet au milieu.} Pour $\lambda=\frac12$ et des bouts de même nature, les modes se séparent en antisymétriques et symétriques par rapport au milieu. Un mode antisymétrique a $y=0$ et $y''=0$ au milieu, et $y'''$ y est continu : il ne charge pas le stylet, et chaque moitié vibre comme une lame de longueur $\frac12a$ gardant sa condition extérieure et appuyée au milieu. Un mode symétrique a $y'=0$ au milieu, et $y'''$ y saute : il charge le stylet, et chaque moitié vibre comme une lame gardant sa condition extérieure et encastrée au milieu. Bouts appuyés : moitiés appuyée--appuyée (cas IV, $\omega=2\pi,4\pi,\dots$) et appuyée--encastrée (cas V, $\omega=7{,}8532,\dots$) ; le fondamental est antisymétrique, et le premier mode symétrique est au-dessus dans le rapport $1{,}562$, très près de $25/16$. Bouts libres : moitiés libre--appuyée (cas II) et libre--encastrée (cas III). Bouts encastrés : moitiés encastrée--appuyée (cas V) et encastrée--encastrée (cas VI).
\item \emph{Limite.} Quand $\lambda\to0$, l'équation tend vers celle du cas V : un appui collé contre un bout appuyé l'encastre.
\end{enumerate}
Le verdict sur les « Remarques » se lit dans ces énoncés. Les deux solutions qu'Euler donne pour $\lambda=\frac12$ sont justes, et la seconde en particulier est la famille des modes symétriques, qui font travailler le stylet. L'objection repose sur une hypothèse que les feuillets ne disent qu'une fois, dans la première rédaction : que l'effet du stylet est « purement exclutif », c'est-à-dire qu'il n'exerce aucune force ; le § 8 d'Euler l'exclut. Ce qu'elle établit de juste en chemin — les points sympathiques aux positions rationnelles, les développements de $\sin n\theta$, la forme symétrique de l'équation, la rareté des modes à pente nulle au stylet — reste juste ; ce qu'elle en conclut — qu'aucun son mesurable n'existe hors des divisions aliquotes, et que la seconde solution d'Euler est à rejeter — est faux.

\section*{La lame touchée par un stylet : extrémités libres, extrémités fixées}

\pagerange{246}{396}
Après la seconde solution contestée, la critique du § 47 d'Euler passe à un problème qu'Euler n'a pas traité et qu'elle juge plus important : la lame libre, ou encastrée, à ses deux bouts, touchée par un stylet en un point quelconque de sa longueur. Le volume en garde deux états qui suivent le même ordre, paragraphe par paragraphe : un brouillon très repris et une mise au net. On suit ici la mise au net ; le brouillon est cité en note, rapproché par le contenu et non par les numéros.\footnote{Vues 246 à 316 : le brouillon, écrit au recto seul. Il est paginé de sa main en tête, les numéros soulignés : 13 à 20 aux vues 246 à 260, 21 à 29 aux vues 262 à 280, 30 à 38 aux vues 283 à 299, 39 à 45 aux vues 301 à 314 ; la vue 316 n'a pas de numéro. Il porte les n\textsuperscript{os} 120 à 154 de la série à l'encre. Ses alinéas sont numérotés 8, 10, 11, 12 dans la marge (le 9 manque ; une marque biffée ouvre la vue 256), puis 13 à 24. La vue 281 photographie de nouveau le feuillet de la vue 280, béquet relevé ; elle montre six lignes barrées d'une croix. La vue 302 reprend la vue 301. Vues 334 à 389 : la mise au net, pages (8) à (35) de sa pagination, entre parenthèses et barrées ; n\textsuperscript{os} 163 à 190 de la série à l'encre ; paragraphes 8 à 25, sans paragraphe 13. La vue 359 reprend la vue 358. Elle est d'une main régulière, que les transcriptions attribuent diversement ; le texte y parle au féminin (« je me suis assurée », vue 377). Correspondance par le contenu : le § 8 de la mise au net répond aux vues 246 à 254 du brouillon ; les § 9 et 10 aux vues 254 et 256 ; le § 11 aux vues 258 et 260 ; le § 12 aux vues 260 et 262 ; l'alinéa sans numéro de la vue 350 au § 13 du brouillon (vues 264 et 266) ; le § 14 aux vues 268 et 270 ; la remarque sur Chladni et Riccati à la vue 272 ; le § 15 aux vues 270 et 274 ; le § 16 aux vues 274 et 276 ; le § 17 aux vues 278 et 280 ; le § 18 aux vues 280 et 283 ; le § 19 aux vues 283 à 287 ; le § 20 aux vues 287 et 289 ; le § 21 aux vues 291 à 295 ; le § 22 à la vue 297 ; le § 23 à la vue 299 ; le § 24 aux vues 301 à 308 ; le § 25 aux vues 310 à 316.}

On garde les notations de la section précédente. L'ordonnée de la lame est $y=Y(u)\,\sin\bigl(\zeta+t\,\omega^2c\sqrt{2gb}/a^2\bigr)$, avec $u=s/a$, et
\[
Y(u)=\alpha e^{\omega u}+\beta e^{-\omega u}+\gamma\sin\omega u+\delta\cos\omega u
\]
sur la portion $EL$, $0\le u\le\lambda$, les coefficients accentués $\alpha',\beta',\gamma',\delta'$ sur la portion $LF$ ; les sons sont comme $\omega^2$. On pose $\mu=1-\lambda$ et, pour abréger,
\[
P(t)=\cosh t\,\sin t-\sinh t\,\cos t,\qquad C_{+}(t)=1+\cosh t\,\cos t,\qquad C_{-}(t)=1-\cosh t\,\cos t .
\]
L'équation $P(t)=0$, c'est-à-dire $\tan t=\tanh t$, est celle d'une lame de longueur réduite $t$ libre à un bout et appuyée à l'autre (cas II d'Euler), ou appuyée et encastrée (cas V) ; $C_{+}(t)=0$ celle d'une lame libre et encastrée (cas III) ; $C_{-}(t)=0$ celle d'une lame libre aux deux bouts (cas I) ou encastrée aux deux bouts (cas VI).\footnote{Le feuillet écrit « sin. », « cos. », « tang. » et des sommes d'exponentielles ; on écrit $\sin$, $\cos$, $\tan$, et $\cosh t$, $\sinh t$, $\tanh t$ pour $\frac12(e^t+e^{-t})$, $\frac12(e^t-e^{-t})$ et leur quotient ; ainsi son $\frac{e^{t}+e^{-t}}{e^{t}-e^{-t}}$ est $\coth t$. La fraction de la longueur où se tient le stylet est une petite lettre crochue, lue $\lambda$ aux vues 244 à 280 et 342 à 358, et $x$ (« une lettre proche de $\lambda$ ») aux vues 283 à 308, 320 à 340 et 361 à 389 ; on écrit partout $\lambda$. Le feuillet écrit $\left(\frac{ddy}{dx^2}\right)$ ; on écrit $Y''$, les dérivées étant prises en $u$. « Fixée » est son mot pour « encastrée ».}

Tout ce qui suit se juge sur les énoncés donnés à la fin de la section précédente (« Ce que dit la théorie actuelle de la poutre sur appui intermédiaire », vues 320 à 334) : au stylet, $y=0$, et $y$, $y'$, $y''$ sont continus tandis que $y'''$ saute de la force que le stylet exerce ; pour toute position du stylet la lame a une suite infinie de modes ; un mode n'exerce aucune force sur le stylet si et seulement si c'est un mode de la lame sans stylet dont un nœud tombe au stylet. Ses raisonnements, dans ces pages, reposent sur une hypothèse qu'elle énonce au § 20 : le mouvement de la lame touchée par le stylet doit être l'un des mouvements réguliers de la lame sans stylet. Ce qu'elle prouve sous cette hypothèse est vrai des modes sans force ; ce qu'elle en conclut pour tous les modes ne l'est pas.

\subsection*{Le problème de la lame libre aux deux bouts (§ 8 à 11)}

\pagerange{334}{348}
Elle pose le problème :
\begin{quote}
Si une lame élastique, libre à l'une et à l'autre de ses extrémités, est soumise à l'action d'un stylet dans un point quelconque de sa longueur : trouver tous les mouvemens de vibration dont elle est susceptible ?
\end{quote}
Elle reprend les quatre équations I à IV qu'Euler écrit au stylet dans son § 47, avec raison « puisqu'elles sont indépendantes de l'état des extrémités » : $Y(\lambda)=0$ des deux côtés, pente et courbure continues.\footnote{Vues 334 à 342, pages (8) à (12), son § 8, sous les titres soulignés « Problème » et « Solution ». Brouillon : vues 246 à 254, pages 13 à 17, alinéa 8 en marge. Le brouillon renvoie pour les conditions des bouts à « cas 1\textsuperscript{er} N\textsuperscript{o} 5 du memoire cité », la mise au net à « Memoire d'Euler § 19 », le 9 lu avec doute.} Les bouts libres donnent $Y''=Y'''=0$ en $u=0$ et en $u=1$ :
\[
\alpha-\beta=\gamma,\qquad \alpha+\beta=\delta,
\]
\[
\alpha'e^{\omega}+\beta'e^{-\omega}-\gamma'\sin\omega-\delta'\cos\omega=0,\qquad
\alpha'e^{\omega}-\beta'e^{-\omega}-\gamma'\cos\omega+\delta'\sin\omega=0 ,
\]
ses équations V à VIII ; elles sont exactes. L'élimination est conduite pas à pas : $\gamma'$ et $\delta'$ par VII et VIII, $\beta'$ par II, puis $\beta,\gamma,\delta$ par V, VI et I ; l'équation III donne une première valeur de $\alpha/\alpha'$, et l'équation IV, ajoutée à I et II et divisée par 2, soit $\alpha e^{\lambda\omega}+\beta e^{-\lambda\omega}=\alpha'e^{\lambda\omega}+\beta'e^{-\lambda\omega}$, une seconde. En les égalant, elle obtient l'équation finale, qui s'écrit
\[
C_{+}(\mu\omega)\,P(\lambda\omega)+C_{+}(\lambda\omega)\,P(\mu\omega)=0 .
\]
Elle est juste : c'est l'équation aux fréquences de la lame libre sur un appui intermédiaire.\footnote{Le feuillet (vue 342, et au brouillon vue 254) l'écrit $\{2+(e^{(1-\lambda)\omega}+e^{(\lambda-1)\omega})\cos(1-\lambda)\omega\}\{(e^{\lambda\omega}+e^{-\lambda\omega})\sin\lambda\omega-(e^{\lambda\omega}-e^{-\lambda\omega})\cos\lambda\omega\}+\{\dots\}\{\dots\}=0$, le second produit obtenu en échangeant $\lambda$ et $1-\lambda$ ; c'est quatre fois le premier membre écrit ici. Vérifié pour cette lecture : ses racines sont celles du déterminant $8\times8$ des conditions, calculé directement, soit $3{,}7502$ ; $7{,}8532$ ; $9{,}3882$ ; $14{,}1372$ pour $\lambda=\frac12$ et $4{,}4831$ ; $6{,}5842$ ; $10{,}5736$ ; $13{,}8915$ pour $\lambda=0{,}3$. Ses deux valeurs de $\alpha/\alpha'$ sont justes elles aussi. À la vue 342 une parenthèse et une accolade ne sont pas refermées, et le premier membre de la première égalité porte $\alpha$ là où l'on attendrait $\alpha'$ : accidents de copie sans effet sur la suite. Le terme de tête des accolades, 4, est écrit sur un autre chiffre à la vue 340 et lu avec doute aux vues 250 et 252 du brouillon ; le calcul demande 4. À la vue 336, le dernier terme d'une ligne est illisible ; le développement demande $\cos\omega\,\sin(1-\lambda)\omega$.} Divisée par $C_{+}(\lambda\omega)\,C_{+}(\mu\omega)$, elle devient
\[
\frac{P(\lambda\omega)}{C_{+}(\lambda\omega)}+\frac{P(\mu\omega)}{C_{+}(\mu\omega)}=0 :
\]
chaque terme est, au même facteur près, le couple qu'il faut appliquer en $L$ à l'une des portions, prise seule et libre à son autre bout, pour lui imposer une rotation donnée à la fréquence considérée ; l'équation dit que les deux couples se compensent, le stylet n'en transmettant aucun.\footnote{C'est la lecture qu'en donne aujourd'hui la méthode des raideurs dynamiques en rotation, bien postérieure aux feuillets, qui ne la connaissent pas. Le paragraphe 17 en énonce la structure dans ses propres termes.}

Pour $\lambda=\frac12$ (son § 9), l'équation se réduit à $2\,C_{+}(\tfrac12\omega)\,P(\tfrac12\omega)=0$ et a deux familles de solutions. La première, $\tan\frac12\omega=\tanh\frac12\omega$, fait vibrer chaque moitié « comme deux lames séparées de longueur $\frac12a$ » libres à un bout et appuyées à l'autre, et elle vérifie qu'elle s'accorde avec l'équation de la lame libre sans stylet : de $\cos\omega\cosh\omega=1$ avec $\sin\omega>0$, soit $\sin\omega=\tanh\omega$, on tire bien $\tan\frac12\omega=\sin\omega/(1+\cos\omega)=\tanh\frac12\omega$. Ce sont les modes antisymétriques de la lame libre, $\omega=7{,}8532$, $14{,}1372$, … , doubles des racines $3{,}9266$, $7{,}0686$, … du cas II ; ils ont un nombre impair de nœuds, celui du milieu sous le stylet, qui n'y travaille pas. La seconde famille, $C_{+}(\frac12\omega)=0$, soit $\cos\frac12\omega=-1/\cosh\frac12\omega$, ferait de chaque moitié une lame encastrée au stylet et libre à l'autre bout ; elle la rejette parce que cette condition ne peut s'accorder avec $\cos\omega\cosh\omega=1$ « que dans la seule supposition $\omega=0$ ».\footnote{Vues 342 et 344, pages (12) et (13), son § 9 ; brouillon, vues 254 et 256. L'incompatibilité est exacte : avec $c=\cosh\frac12\omega$, la condition donne $\cos\omega=2/c^2-1$, et l'égalité avec $1/(2c^2-1)$ revient à $(c^2-1)^2=0$. Le brouillon écrit qu'il faut rejeter cette solution « comme il me semble necessaire dele faire dans le cas examiné par Euler » ; la mise au net, « comme j'ai montré qu'il falloit le faire pour celle qui lui est analogue dans le cas examiné par Euler ».} L'incompatibilité est vraie, la conclusion non : ce sont les modes symétriques de la lame tenue en son milieu, $\omega=3{,}7502$, $9{,}3882$, $15{,}7095$, … , doubles des racines $1{,}8751$, $4{,}6941$, $7{,}8548$ du cas III, où la pente est nulle au milieu par symétrie et où le stylet exerce une force ; ils ne sont pas des modes de la lame libre, et c'est tout ce que dit l'incompatibilité. Le premier d'entre eux est même le son le plus grave de la lame ainsi tenue : $\omega^2\approx14{,}06$, contre $22{,}37$ pour le premier son de la lame libre.

Au § 10, elle note que l'analogie avec les bouts appuyés s'arrête là : pour eux, $\lambda=\frac12$ n'est qu'un cas particulier de la solution générale où $\lambda$ est une partie aliquote de l'unité, et « ces valeurs sont les seules pour lesquelles le mouvement de la lame soit possible ». Cette dernière clause est fausse, pour la raison donnée à la section précédente : pour toute position du stylet, la lame appuyée à ses bouts a une suite infinie de modes.\footnote{Vue 344, page (13), son § 10 ; brouillon, vue 256, alinéa 10 en marge, où « être une fraction quelconque aliquote de l'unité » est écrit au-dessus de « avoir toutes toutes les valeurs », biffé. À la vue 344, « $\sin(1-\lambda)\omega$ » est écrit sans « $=0$ ».}

Au § 11, elle démontre que, les bouts étant libres, les deux conditions $P(\lambda\omega)=0$ et $P(\mu\omega)=0$, qui feraient des deux portions deux lames séparées, libres au bout et appuyées au stylet, ne peuvent avoir lieu ensemble « lorsqu'on a égard à la condition qui résulte de l'état des extrémités de la lame primitive, que dans le seul cas où $\lambda=\frac12$ ». Sa transformation est exacte : en portant la première condition dans la seconde, on obtient
\[
\sin\lambda\omega\,\bigl[\cosh\bigl((1-2\lambda)\omega\bigr)\sin\omega-\sinh\omega\,\cos\omega\bigr]=0 .
\]
Le sinus ne peut s'annuler, puisque $P(t)=\mp\sinh t\neq0$ quand $\sin t=0$ et $t\neq0$ ; il reste $\tan\omega=\sinh\omega/\cosh\bigl((1-2\lambda)\omega\bigr)$. Avec l'équation de la lame libre, $\cos\omega=1/\cosh\omega$ et $\sin\omega=\pm\tanh\omega$, donc $\tan\omega=\pm\sinh\omega$, il faut $\cosh\bigl((1-2\lambda)\omega\bigr)=\pm1$ : le signe inférieur est impossible, le supérieur exige $\lambda=\frac12$.\footnote{Vues 346 et 348, pages (14) et (15), son § 11 ; brouillon, vues 258 et 260, alinéas 11 et 12 en marge. Vérifié pour cette lecture : l'identité qui mène au facteur $\sin\lambda\omega$ est exacte. Elle écarte $\sin\lambda\omega=0$ parce que la portion $\lambda a$ serait alors appuyée à ses deux bouts ; la raison donnée ici est plus directe. Au brouillon (vue 260), le signe négatif est dit répondre au cas « ou il y en a un nombre impair » : « impair » y est écrit deux fois, et le raisonnement demande « pair » ; la mise au net ne le répète pas. À la vue 348, l'exposant $(2\lambda-1)$ est écrit une fois sans $\omega$.} L'énoncé est donc juste tel qu'elle l'a écrit, avec sa restriction ; mais la restriction est l'hypothèse de la lame sans stylet. Sans elle, les deux portions vibrent comme deux lames libres et appuyées au stylet chaque fois que $\lambda\omega$ et $\mu\omega$ sont toutes deux racines de $\tan t=\tanh t$, soit pour $\lambda=r_i/(r_i+r_j)$ : avec $r_1=3{,}9266$ et $r_2=7{,}0686$, le stylet en $\lambda=0{,}35712$ et $\omega=10{,}99519$ donnent un vrai mode de la lame touchée, le couple au stylet nul et une petite force.\footnote{Vérifié pour cette lecture sur le déterminant $8\times8$ : il s'annule pour ces valeurs, $Y''(\lambda)=0$ des deux côtés, et $y'''$ saute. Ce n'est pas un mode de la lame libre : la troisième racine de $\cos\omega\cosh\omega=1$ est $10{,}99561$, et le nœud correspondant de ce mode de la lame libre est en $0{,}3558$. L'écart se resserre vite : pour $\lambda=r_2/(r_2+r_3)=0{,}40909$, la fréquence et la racine de la lame libre coïncident à cinq décimales et la force au stylet est plus de vingt fois plus faible.} La phrase qui clôt son paragraphe, « on est forcé de conclure que le cas où le stilet est appuyé au milieu … est le seul où cette lame puisse être considérée comme séparée en deux autres lames », perd la restriction et n'est vraie que des modes sans force.

\subsection*{Les extrémités fixées (§ 12)}

\pagerange{348}{350}
Pour les bouts encastrés, elle donne l'équation finale « au moyen de calculs que je supprime ici, à cause de leur parfaite analogie avec les précédens » ; elle l'écrit
\[
\Bigl(\frac{1}{\cosh\lambda\omega}-\cos\lambda\omega\Bigr)\bigl(\coth\mu\omega\,\sin\mu\omega-\cos\mu\omega\bigr)
+\Bigl(\frac{1}{\cosh\mu\omega}-\cos\mu\omega\Bigr)\bigl(\coth\lambda\omega\,\sin\lambda\omega-\cos\lambda\omega\bigr)=0 .
\]
L'équation juste est
\[
C_{-}(\mu\omega)\,P(\lambda\omega)+C_{-}(\lambda\omega)\,P(\mu\omega)=0,
\]
qui s'écrit dans sa forme en remplaçant chaque facteur $\coth t\,\sin t-\cos t$ par $\sin t-\tanh t\,\cos t$. La sienne revient à diviser les deux termes par des facteurs différents, $\cosh\lambda\omega\,\sinh\mu\omega$ et $\cosh\mu\omega\,\sinh\lambda\omega$ ; elle n'est donc pas équivalente, sauf pour $\lambda=\frac12$, et l'erreur est petite dès que $\tanh$ est voisin de 1.\footnote{Vues 348 et 350, pages (15) et (16), son § 12 ; brouillon, vues 260 et 262, où la même forme ouvre la vue 262. Vérifié pour cette lecture : pour $\lambda=0{,}3$, la première racine du déterminant est $6{,}3015$, celle de l'équation juste aussi, celle de la sienne $6{,}3138$ ; les suivantes s'accordent à trois décimales. La même forme reparaît pour les bouts libres dans le tableau du § 16.} Pour $\lambda=\frac12$, sa conclusion est juste : la première famille fait de chaque moitié une lame appuyée au stylet et encastrée à l'autre bout, $\omega=7{,}8532$, $14{,}1372$, … , qui sont les modes antisymétriques de la lame encastrée sans stylet. Elle passe l'autre famille sous silence : $C_{-}(\frac12\omega)=0$, les moitiés encastrées aux deux bouts, $\omega=9{,}4601$, … , qui sont les modes symétriques et chargent le stylet. Enfin elle affirme, comme au § 11, que les deux moitiés ne peuvent être des lames séparées appuyées au stylet pour aucune autre valeur de $\lambda$ ; même restriction, même verdict.

\subsection*{Pourquoi le milieu seul partage la lame (§ 13 et 14)}

\pagerange{350}{354}
Elle cherche alors à « se rendre compte de cette singularité ». Son explication part d'une prémisse : si plusieurs lames de longueurs différentes, libres ou encastrées, rendent le même son, la partie comprise entre un bout et le nœud le plus voisin est la même pour toutes ; et de même, pour des lames appuyées à un bout, la partie comprise entre le bout appuyé et le nœud le plus voisin. Si la lame touchée se partageait en deux lames appuyées au stylet et rendant le même son, les parties comprises entre le stylet et les nœuds voisins, de part et d'autre, seraient donc égales ; « il ne reste plus » qu'à prouver qu'excepté le cas d'un nœud au milieu, aucun point de repos n'a ses ordonnées égales à distances égales de part et d'autre.\footnote{Vues 350 et 352, pages (16) et (17), alinéa sans numéro ; au brouillon c'est le § 13 (vues 264 et 266, pages 22 et 23), avec deux notes de sa main dans la marge, la première annulée d'un grand trait oblique. Le brouillon justifie la prémisse : on pourrait supprimer deux à deux les parties voisines du milieu sans altérer le son, et sinon les sons des parties extrêmes seraient « infiniment differens » selon les longueurs, « résultat contraire a tout ce que l'on sait sur la théorie des corps sonores ». La mise au net énonce la prémisse sans cet argument.}

La prémisse n'est qu'approchée. Pour des lames libres à un bout et appuyées à l'autre qui rendent le même son, de longueurs $3{,}9266/k$, $7{,}0686/k$, $10{,}2102/k$, $13{,}3518/k$, la distance du bout appuyé au nœud le plus voisin vaut $2{,}889/k$, $3{,}156/k$, $3{,}141/k$, $3{,}1416/k$ : elle tend vers $\pi/k$, mais varie de près d'un dixième entre les deux premières lames.\footnote{Calcul de cette lecture : la lame appuyée en $0$, libre en $\ell$, a pour forme $\sin kx+b\sinh kx$ avec $b=\sin k\ell/\sinh k\ell$, et $k\ell$ racine de $\tan t=\tanh t$. Pour des lames libres aux deux bouts, la distance du bout au premier nœud vaut de même $1{,}060/k$, $1{,}038/k$, $1{,}038/k$, $1{,}038/k$ pour les quatre premiers sons.} Le terme hyperbolique, qui décroît loin des bouts, en est la cause.

La démonstration du § 14, elle, est juste. Elle part de l'ordonnée de la lame libre qu'elle tire d'Euler,
\[
y(u)=e^{\omega u}\pm e^{\omega(1-u)}+(1\mp e^{\omega})\sin\omega u+(1\pm e^{\omega})\cos\omega u ,
\]
et forme la somme des ordonnées à la distance $\lambda'$ de part et d'autre du point $\lambda$. On vérifie l'identité
\[
y(\lambda-\lambda')+y(\lambda+\lambda')=2\bigl(e^{\lambda\omega}\pm e^{\mu\omega}\bigr)\bigl(\cosh\lambda'\omega-\cos\lambda'\omega\bigr)+2\cos\lambda'\omega\;y(\lambda),
\]
le double signe étant celui de $y$. En un nœud, $y(\lambda)=0$ ; le second facteur est positif pour $\lambda'\neq0$ ; le premier ne s'annule qu'avec le signe moins et pour $\lambda=\frac12$. Ainsi, dans un mode de la lame libre, ou encastrée, où les coefficients $\gamma$ et $\delta$ changent seulement de signe, aucun nœud n'a ses ordonnées opposées à distances égales, sinon le nœud du milieu des modes à nombre impair de nœuds ; a fortiori, aucun nœud autre que celui-là n'est à égale distance de ses deux voisins. C'est vrai, et elle le démontre.\footnote{Vues 352 et 354, pages (17) et (18), son § 14 ; brouillon, vues 268 et 270, pages 24 et 25. Vérifié symboliquement pour cette lecture. Deux fautes de signe s'y compensent. Le feuillet écrit le facteur $e^{\lambda\omega}\mp e^{\mu\omega}$ (vues 352 et 354, et au brouillon) alors que la somme donne le même signe que dans $y$ ; et il attribue le signe supérieur au nombre impair de nœuds, alors que le § 19 (vue 365) attribue à ce cas, où $\sin\omega$ est positif, les signes inférieurs. Avec le signe de $y$ et la convention du § 19, le facteur est $e^{\lambda\omega}-e^{\mu\omega}$ pour un nombre impair de nœuds, et la conclusion est celle du feuillet. À la vue 354, le feuillet écrit $e^{\lambda\omega}=\pm e^{-\lambda\omega}$ là où le raisonnement attend $e^{\mu\omega}$. Le brouillon renvoie pour l'ordonnée au « § 21 » d'Euler (vue 352, lu avec doute), la dernière figure se confondant avec la parenthèse.} Ce qu'elle prouve porte sur les modes de la lame sans stylet ; joint à la prémisse, cela redonne la conclusion du § 11 pour les modes sans force, et rien pour les autres.

\subsection*{L'inégalité des intervalles entre les nœuds : Chladni et Riccati}

\pagerange{354}{356}
Elle en tire une remarque sur la lame libre : « Je ne crois pas que l'on ait encore fait cette remarque de l'inégalité des parties comprises entre deux nœuds de vibrations ». Chladni se contente de dire qu'une partie située à un bout est à peu près moitié d'une partie comprise entre deux nœuds, ce qui semble supposer ces dernières égales entre elles. Giordano Riccati, dans son mémoire sur les vibrations des cylindres, donne pour la lame libre des valeurs de l'angle d'où elle tire, lui ne l'ayant pas fait, que la partie située au bout est moins que la moitié de la suivante, que les parties croissent du bout vers le milieu, et que les différences sont beaucoup plus grandes près des bouts.\footnote{Vues 354 et 356, pages (18) et (19) ; brouillon, vue 272, page 25, les deux dernières lignes d'une plume plus fine, une mention lue « 2.\textsuperscript{eme} » avec doute en tête. Chladni : « traité d'acoustique pag 4 §§ 71 » au brouillon, « P. 94 § 71 » à la mise au net, le 94 lu avec doute (peut-être « g4 », pour « pag 4 ») ; Riccati : « Delle vibrazioni de' cilindri », au premier volume des Mémoires de la Société italienne, avec un renvoi « § XXX P 503 » où trois croix tiennent la place du numéro du paragraphe, aux deux états. Le brouillon dit que Riccati donne des valeurs de « $\omega$ », la mise au net de « $\lambda\omega$ ». Aucun de ces ouvrages n'a été consulté pour cette lecture. Le long mémoire revient à Riccati aux vues 427 à 429. Sa phrase porte sur l'état de la littérature ; la lecture ne l'examine pas.} Le calcul actuel des nœuds de la lame libre donne, en fractions de la longueur :
\begin{itemize}
\item premier son, $\omega=4{,}7300$ : bout $0{,}2242$, milieu $0{,}5517$ ;
\item deuxième son, $\omega=7{,}8532$ : $0{,}1321$, puis $0{,}3679$ ;
\item cinquième son, $\omega=17{,}2788$ : $0{,}0601$ ; $0{,}1664$ ; $0{,}1827$ ; $0{,}1818$ au milieu ;
\item huitième son, $\omega=26{,}7035$ : $0{,}0389$ ; $0{,}1076$ ; $0{,}1182$ ; $0{,}1176$ ; $0{,}1177$, pour $\pi/\omega=0{,}11765$.
\end{itemize}
En unités de $1/\omega$, dès le troisième son, la partie du bout vaut $1{,}038$, la suivante $2{,}874$, la troisième $3{,}156$, la quatrième $3{,}141$, les autres $\pi$ à moins d'un millième près. La partie du bout est donc environ $0{,}36$ de la suivante ($0{,}41$ au premier son) et le tiers d'une partie intérieure : moins que la moitié, comme elle le dit, et le « à peu près moitié » de Chladni est grossier. Les parties croissent du bout vers le milieu sur les trois premières, et les différences sont bien concentrées près des bouts ; mais la troisième partie dépasse $\pi/\omega$ d'un demi-centième et les suivantes oscillent autour de $\pi/\omega$ : « d'autant plus étendue qu'elle est plus voisine du milieu » n'est exact que pour les trois premières parties.\footnote{Calcul de cette lecture, par les racines de $\cos\omega\cosh\omega=1$ et les zéros de $\cosh\omega u+\cos\omega u-\sigma(\sinh\omega u+\sin\omega u)$, $\sigma=(\cosh\omega-\cos\omega)/(\sinh\omega-\sin\omega)$.}

\subsection*{L'analogie des équations finales (§ 15 à 18)}

\pagerange{356}{363}
Au § 15, elle oppose les bouts appuyés : dans les modes $\sin\lambda\omega=0$, les nœuds à $\lambda\pm\lambda'$ vérifient $\sin(\lambda-\lambda')\omega+\sin(\lambda+\lambda')\omega=2\sin\lambda\omega\cos\lambda'\omega=0$, les parties entre nœuds sont égales, et c'est pourquoi ce cas seul admet une « solution générale » de l'équation finale.\footnote{Vue 356, page (19), son § 15 ; brouillon, vues 270 et 274, où les arguments $(\lambda-\lambda')$ et $(\lambda+\lambda')$ sont écrits sans $\omega$ et le résultat $2\sin\lambda\cos\lambda'\omega$.} L'identité est exacte, et dans ces modes $y$ est proportionnel à $\sin\omega u$ ; mais ces modes ne sont pas les seuls de la lame appuyée touchée par le stylet.

Au § 16, elle range les équations finales des quatre cas où le stylet touche la lame, « afin que l'on puisse saisir plus facilement ces rapports ». Dans les notations de cette section :\footnote{Vue 358, page (20), son § 16 ; brouillon, vues 274 et 276 (les équations) et 278 (le dernier cas, avec en marge « voyez § 47 et suiv. du memoire d'Euler »). Vérifié pour cette lecture sur les déterminants : pour $\lambda=0{,}3$, première racine $4{,}4831$ pour les bouts libres, contre $4{,}4028$ par sa forme du tableau ; racines $5{,}1318$ ; $9{,}2769$ pour les bouts appuyés et $6{,}2055$ ; $10{,}1984$ ; $12{,}1248$ pour le dernier cas, égales à celles de ses formes. À la vue 358, l'accolade ouverte après $\sin\lambda\omega$ n'est pas refermée. Les renvois « § 5 », « § 8 », « § 12 » qui suivent ces équations semblent être ses propres paragraphes.}
\begin{itemize}
\item bouts appuyés : $\sin\lambda\omega\,\dfrac{P(\mu\omega)}{\sinh\mu\omega}+\sin\mu\omega\,\dfrac{P(\lambda\omega)}{\sinh\lambda\omega}=0$, ce qui est sa forme ; elle est exacte ;
\item bouts libres : sa forme, $\bigl(\frac{1}{\cosh\lambda\omega}+\cos\lambda\omega\bigr)\bigl(\coth\mu\omega\,\sin\mu\omega-\cos\mu\omega\bigr)+(\lambda\leftrightarrow\mu)=0$, n'est pas équivalente à l'équation juste $C_{+}(\mu\omega)P(\lambda\omega)+C_{+}(\lambda\omega)P(\mu\omega)=0$, qu'elle avait pourtant établie au § 8 ; même défaut qu'au § 12 ;
\item bouts encastrés : la forme du § 12, même verdict ;
\item un bout encastré, l'autre appuyé : $(\tan\lambda\omega-\tanh\lambda\omega)(\tan\mu\omega-\tanh\mu\omega)=2\tan\lambda\omega\,\tanh\lambda\omega\,\Bigl(1-\dfrac{1}{\cosh\mu\omega\,\cos\mu\omega}\Bigr)$, la portion $\lambda a$ étant du côté du bout appuyé ; elle est exacte.
\end{itemize}
Elle y joint les équations des six cas d'Euler sans stylet :
\[
\text{I : }\cos\omega\cosh\omega=1,\quad
\text{II : }\tan\omega=\tanh\omega,\quad
\text{III : }\cos\omega\cosh\omega=-1,
\]
\[
\text{IV : }\sin\omega=0,\quad
\text{V : }\tan\omega=\tanh\omega,\quad
\text{VI : }\cos\omega\cosh\omega=1,
\]
toutes exactes, et remarque au brouillon qu'on peut « changer » l'encastrement et la liberté d'un bout l'un en l'autre « sans que l'équation qui détermine la valeur de l'angle $\omega$ change de forme » : c'est vrai des équations aux fréquences (I et VI, II et V ; III est inchangé), non des formes des modes.\footnote{Le feuillet écrit $\cos\omega=\frac{2e^{\omega}}{e^{2\omega}+1}$ pour I et VI, $\tan\omega=\frac{e^{2\omega}-1}{e^{2\omega}+1}$ pour II et V, $\cos\omega=-\frac{2e^{\omega}}{e^{2\omega}+1}$ pour III. Les numéros d'Euler cités sont 23, 30, 35, 39, 43 pour les cas II à VI aux deux états ; pour le cas I, la mise au net écrit « § 5 » (vue 358, comme « cas I § 5 » à la vue 344), le brouillon « N\textsuperscript{o} 9 » (vue 276), et la mise au net elle-même « § 15 » au § 24 (vue 381) ; la lecture ne tranche pas.}

Au § 17, elle observe que, pour des bouts semblables, chaque équation finale est une somme de deux termes, dont chacun est le produit de la quantité qui s'annulerait si la portion $\lambda a$ vibrait seule, appuyée au stylet, par celle qui s'annulerait si la portion $\mu a$ vibrait seule, encastrée au stylet ; le second terme échange les deux portions. C'est exact des équations justes : $P(\lambda\omega)$ est la condition de la portion $\lambda a$ appuyée au stylet, $C_{\pm}(\mu\omega)$ celle de la portion $\mu a$ encastrée au stylet, et c'est la structure $P/C+P/C=0$ relevée au § 8. Pour des bouts dissemblables, un bout encastré et l'autre appuyé, elle lit de même, le stylet regardé comme appuyé, les conditions $\sin\lambda\omega=0$ et $\tan\mu\omega=\tanh\mu\omega$, et, le stylet regardé comme encastré, $\tan\lambda\omega=\tanh\lambda\omega$ et $\cos\mu\omega\cosh\mu\omega=1$ : c'est juste.\footnote{Vues 361 et 363, pages (21) et (22), son § 17 ; brouillon, vues 278 et 280, pages 28 et 29. Au brouillon (vue 280), « tang $\lambda$ » est écrit sans $\omega$. À la vue 361, le signe lu « § » devant « precedent » est incertain.}

Au § 18, elle constate que ces analogies donnent peu. Pour les bouts appuyés, l'équation finale montre que, si le stylet est en $\lambda a$, on a $\sin\lambda\omega=0$ et des points de repos en $2\lambda a$, $3\lambda a$, … : ce n'est vrai que de la famille aliquote. Pour les bouts libres ou encastrés, le seul résultat général serait que, s'il y a un nœud à la distance $\lambda a$ d'un bout, il y en a un autre à la même distance de l'autre bout, « parce que ces équations sont fonctions homogènes » de $\lambda$ et $1-\lambda$. La symétrie de l'équation en $\lambda$ et $\mu$ ne prouve que ceci : le stylet en $\lambda$ et le stylet en $1-\lambda$ donnent les mêmes sons, par une figure renversée. Elle ne place pas de nœud en $1-\lambda$ : pour $\lambda=0{,}3$ et les bouts libres, les trois premiers modes ont en $0{,}7$ une ordonnée égale à $-0{,}075$, $-0{,}628$ et $+0{,}460$ fois leur plus grande ordonnée. L'énoncé est vrai des modes sans force, qui sont ceux de la lame libre, symétriques ou antisymétriques ; il ne l'est pas des autres.\footnote{Vue 363, page (22), son § 18 ; brouillon, vues 280 et 283. À la vue 363, « $\sin xa=0$ » est écrit avec $a$ pour $\omega$. Au brouillon, un long trait oblique traverse l'alinéa 18 sans qu'on puisse dire s'il l'annule ; une note de marge le reprend (« la methode qui conduit a tous ces resultats est loin d'être aussi utile qu'on auroit pû d'abord l'esperer »), et le texte renvoie à un « § 9 » lu avec doute. Vérifié pour cette lecture sur le déterminant, à $\omega=4{,}4831$, $6{,}5842$, $10{,}5736$.}

\subsection*{Tout, sans les équations finales (§ 19)}

\pagerange{365}{367}
« Il me semble que l'on peut parvenir à tous ces résultats, sans avoir recours aux équations finales », par les seules conditions des bouts et la supposition $y=0$ aux nœuds. Pour les bouts appuyés, $y$ est proportionnel à $\sin\lambda\omega$ (§ 35 d'Euler), et $\sin\lambda\omega=0$ entraîne $\sin2\lambda\omega=0$, … , et, avec $\sin\omega=0$, $\sin\mu\omega=0$. Pour les bouts libres et un nombre impair de nœuds ($\sin\omega>0$, signes inférieurs), $y(\frac12)=0$ s'écrit $(1+e^{\omega})\sin\frac12\omega+(1-e^{\omega})\cos\frac12\omega=0$, c'est-à-dire $\tan\frac12\omega=\tanh\frac12\omega$, la condition d'une lame de longueur $\frac12a$ libre et appuyée. Puis, en multipliant cette relation par $2\cos\frac12(1-2\lambda)\omega$ et en la retranchant de $y(\lambda)=0$, elle obtient $y(1-\lambda)=0$ ; pour un nombre pair de nœuds, elle fait de même avec $(1+e^{\omega})\sin\frac12\omega-(1-e^{\omega})\cos\frac12\omega=0$, tirée de $\sin\omega=-\tanh\omega$, et le multiplicateur $2\sin\frac12(1-2\lambda)\omega$ ; les bouts encastrés se traitent de même. Tout est exact : c'est, calculé terme à terme, le fait que les modes d'une lame à bouts semblables sont symétriques ou antisymétriques par rapport au milieu, et que leurs nœuds sont donc placés symétriquement.\footnote{Vues 365 et 367, pages (23) et (24), son § 19 ; brouillon, vues 283 à 287, pages 30 à 32. Vérifié pour cette lecture. Au brouillon (vue 285), la conclusion est écrite une fois « lorsque l'on change $x$ en $x-1$ », et un signe y est illisible sous un pâté ; la mise au net écrit « $(1-x$ » sans parenthèse fermante et, à la vue 365, « $(1-x)\omega$ » pour $(1-x)a$. Elle renvoie à Euler « § 2(1) », la dernière figure douteuse, pour l'ordonnée de la lame libre, et « §§ 15 » (lu avec doute) pour ses conditions.} Ce qu'elle obtient ainsi, ce sont précisément les résultats vrais des § 11 à 18, et eux seuls : ceux qui concernent les modes sans force.

\subsection*{Le reproche : la continuité de la courbe et les huit coefficients (§ 20 à 23)}

\pagerange{369}{379}
Vient le reproche. Au § 47, Euler admet que le stylet laisse le mouvement se communiquer d'une portion à l'autre, mais « suppose cependant que l'action de ce stilet interromp la continuité de la courbe », les deux portions étant exprimées par des équations qui diffèrent par leurs coefficients. Or, dit-elle, le seul effet du stylet est de déterminer la formation de nœuds, et donc de donner lieu à des mouvements réguliers « compris au nombre de ceux qui ont été traités dans l'examen des six cas », qui s'exécutent suivant une courbe continue ; les coefficients $\alpha',\beta',\gamma',\delta'$ doivent être égaux chacun à son correspondant. « Cependant je n'ose admettre une opinion si contraire à celle d'Euler, sans la soumettre à un examen plus approfondi. »\footnote{Vues 369 et 371, pages (25) et (26), son § 20 ; brouillon, vues 287 et 289, pages 32 et 33. Au brouillon, le corps de la vue 289 est barré d'une grande croix et récrit dans la marge ; la première rédaction dit : « Je n'ose admettre ce resultat contraire a l'opinion d'Euler sur une hypothèse par laquelle il semble avoir ouvert la carrière a l'examen d'une foule de cas dont il donne un exemple comme pour guider dans des recherches semblables ». À la vue 369, trois mots courts sont biffés après « § 47 » et illisibles.} Le reproche décrit mal Euler. Ses équations I à IV rendent la courbe continue en $L$, avec sa pente et sa courbure ; seule la dérivée troisième y change, et c'est l'effet de la force que le stylet exerce. Que l'expression analytique change en $L$ n'est pas une rupture de la courbe. Quant à la prémisse, que le mouvement est l'un de ceux des six cas, c'est l'hypothèse de la force nulle.

Au § 21, elle énonce la proposition : toutes les fois que le stylet laisse subsister une hypothèse déterminée et semblable sur l'état des deux bouts, les coefficients des deux portions sont égaux chacun à son correspondant.\footnote{Vue 371, page (26), son § 21 ; brouillon, vue 291, page 34. L'énoncé est écrit à la vue 371 d'une écriture plus droite et plus ronde que le reste ; la transcription ne dit pas s'il s'agit d'une autre main. Après le premier signe $=$ de la valeur de $\alpha/\alpha'$, une fraction entière est biffée et illisible.} Elle prouve d'abord qu'une seule égalité entraîne les autres : si $\alpha=\alpha'$, l'équation $\alpha e^{\lambda\omega}+\beta e^{-\lambda\omega}=\alpha'e^{\lambda\omega}+\beta'e^{-\lambda\omega}$ donne $\beta=\beta'$, puis III et IV donnent $\gamma=\gamma'$ et $\delta=\delta'$. C'est exact, et cela revient à dire que $y'''$ est continu, c'est-à-dire que le stylet n'exerce aucune force. Elle veut ensuite prouver que $\alpha=\alpha'$ est nécessaire. Pour les bouts libres, la valeur de $\alpha/\alpha'$ tirée de I, II et IV est
\[
\frac{\alpha}{\alpha'}=\frac{P(\mu\omega)}{P(\lambda\omega)}\cdot\frac{\sin\lambda\omega-\cos\lambda\omega-e^{-\lambda\omega}}{e^{-\omega}(\sin\mu\omega+\cos\mu\omega)+e^{-\lambda\omega}},
\]
et il suffit, dit-elle, de montrer que l'égalité des deux membres à 1 « a lieu d'elle-même sous les conditions indiquées ». Elle la développe et la réduit, pour $\sin\omega$ positif puis négatif, et montre que ce qui reste est la somme des deux expressions $y(\lambda)$ et $y(1-\lambda)$ de l'ordonnée de la lame libre, multipliées par des exponentielles : quantité nulle, dit-elle, puisque $y=0$ en ces points.\footnote{Vues 371 à 375, pages (26) à (28) ; brouillon, vues 291 à 295, pages 34 à 36, où un dénominateur porte $e^{x\omega}$ sans signe à l'exposant et une ligne est récrite. Vérifié symboliquement pour cette lecture : sa valeur de $\alpha/\alpha'$ est exacte (à $\lambda=0{,}3$, $14{,}2197$ et $-80{,}3004$ aux deux premiers modes, comme les vecteurs propres du déterminant), et, avec $y(u)=e^{\omega u}-e^{\omega(1-u)}+(1+e^{\omega})\sin\omega u+(1-e^{\omega})\cos\omega u$, la différence des deux membres de son égalité s'écrit exactement $e^{-\lambda\omega}$ fois $y(\lambda)\bigl[e^{-\mu\omega}-\sin\mu\omega+\cos\mu\omega+e^{-\omega}y(1-\lambda)\bigr]+y(1-\lambda)\bigl[e^{-\lambda\omega}-\sin\lambda\omega+\cos\lambda\omega\bigr]$, sans reste. À la vue 375, le « 2 » d'une ligne est tracé gras.} Le calcul est juste, et il prouve exactement ceci : si le stylet est en un point où l'ordonnée de la lame libre s'annule, ainsi qu'au point symétrique, alors $\alpha=\alpha'$. Les « conditions indiquées » sont celles du § 19, $y=0$ aux nœuds de la lame libre : l'hypothèse est dans la prémisse. C'est la moitié facile de l'énoncé rappelé en tête de cette section, un mode de la lame libre avec un nœud au stylet est un mode sans force ; il ne prouve pas que toute solution des huit équations ait $\alpha=\alpha'$. Pour le premier mode symétrique de la lame tenue au milieu, $\omega=3{,}7502$, on trouve $\alpha/\alpha'=6{,}52$ ; pour $\lambda=0{,}3$, les valeurs ci-dessus.

Le § 22 traite de même les bouts encastrés : les deux expressions de $\alpha/\alpha'$ ne diffèrent que par le signe de $e^{-\lambda\omega}$, compensé par les changements de signe des équations des nœuds, et elle conclut encore $\alpha=\alpha'$. L'expression est exacte, et le verdict est le même.\footnote{Vue 377, page (29), son § 22 ; brouillon, vue 297, page 37, où le premier mot écrit « Lorsque les équations sont l'une et l'autre fixées » pour « les extrémités ». Vérifié pour cette lecture : à $\lambda=0{,}3$, $\alpha/\alpha'=59{,}111$ et $-696{,}137$ aux deux premiers modes, par sa formule comme par le déterminant.}

Au § 23, les bouts appuyés : la relation d'Euler $\alpha/\alpha'=\bigl(e^{\lambda\omega}-e^{(2-\lambda)\omega}\bigr)/\bigl(e^{\lambda\omega}-e^{-\lambda\omega}\bigr)$ ne mène pas à $\alpha=\alpha'$, mais, fait-elle observer, tous les coefficients sauf un sont nuls en vertu de l'état des bouts, de sorte que $\alpha/\alpha'$ est « essentiellement une quantité indéterminée » — ce qui, ajoute-t-elle, a sans doute empêché Euler d'apercevoir l'identité des deux séries. C'est juste pour les modes $\sin\lambda\omega=0$, où $y$ est proportionnel à $\sin\omega u$ d'un bout à l'autre. Elle cite ensuite les valeurs d'Euler pour $\lambda=\frac12$ et leur réduction, $\alpha=\alpha'=\beta=\beta'=\delta=\delta'=0$ et $\gamma=-\gamma'$, et dit le changement de signe « attribué au coefficient $\gamma'$, au lieu d'être considéré comme appartenant à l'ordonnée » prise de l'autre côté de l'axe.\footnote{Vues 377 et 379, pages (29) et (30), son § 23 ; brouillon, vue 299, page 38. Deux différences : le coefficient qui n'est pas nul est $\gamma$ à la mise au net, souligné, et $\delta$ au brouillon, où l'égalité réduite se lit $\delta=-\delta'$ ; les valeurs d'Euler sont renvoyées à son « § 48 » à la mise au net, à son § 50 au brouillon. Le relevé qui suit, où $\delta=0$, s'accorde avec $\gamma$. Un mot est biffé et illisible après « § 47 » à la vue 377.} Dans les équations I à IV, les deux portions sont rapportées à la même origine $E$ et à la même variable $u$ : la continuité de la pente exige $\gamma'=\gamma$, et l'alternance des signes d'un côté à l'autre du stylet est déjà portée par $\sin\omega u$. Le signe d'Euler ne vient pas d'une convention sur l'ordonnée ; il vient de ce que ses coefficients sont des formes $0/0$ en $\sin\frac12\omega=0$, dont les deux dénominateurs, $\sin\frac12\omega$ et $\sin\frac12\omega-\cos\frac12\omega\tan\omega$, s'annulent avec des signes opposés.\footnote{Vérifié pour cette lecture : le rapport des deux valeurs d'Euler tend vers $-1$ quand $\omega\to2\pi$, de part et d'autre. La section précédente (§ 50 d'Euler) établit que $\gamma'=\gamma$.} Elle conclut que l'examen du stylet ne donne rien de plus que $y=0$ en son point, joint aux conditions des bouts, que les deux premières des équations I à IV se réduisent à $y=0$ et les deux dernières à des identités dès qu'on admet $\alpha=\alpha'$, $\beta=\beta'$, … ; le brouillon ajoutait que « l'appareille des quatre équations trouvées par Euler est au moins inutile ». Ce n'est vrai que des modes sans force. Pour les autres, les équations III et IV sont précisément ce qui fixe les sons, et la seconde solution d'Euler en est un exemple.

\subsection*{Six cas réduits à quatre (§ 24)}

\pagerange{379}{385}
Au § 24, elle tire des moitiés de lame une réduction des six cas. La lame libre à ses deux bouts, avec un nombre impair de nœuds, se partage en deux lames libres et appuyées ; la lame encastrée, dans le même cas, en deux lames encastrées et appuyées. On peut donc conclure le cas II du cas I et le cas V du cas VI, et les six cas se réduisent à quatre « réellement indépendans les uns des autres », I, III, IV et VI.\footnote{Vues 379 à 383, pages (30) à (32), son § 24 ; brouillon, vues 301 à 306, pages 39 à 41, paragraphe 24 en marge. Les renvois à Euler sont § 15, § 30, § 35 et § 43 pour les cas I, III, IV et VI, § 23 et § 39 pour les cas II et V. La transcription du lot 20 a vu par transparence, à travers la vue 380, un béquet portant « (29) » sur le feuillet de la vue 379 ; il n'a pas été transcrit, et cette lecture n'en dit rien de plus.} Elle vérifie les coefficients. Au cas I, pour $\sin\omega>0$, $\alpha=1$, $\beta=-e^{\omega}$, $\gamma=1+e^{\omega}$, $\delta=1-e^{\omega}$, avec $\tan\frac12\omega=\tanh\frac12\omega$ ; au cas II, $\alpha=1$, $\beta=-e^{2\omega}$, $\gamma=1+e^{2\omega}$, $\delta=1-e^{2\omega}$, avec $\tan\omega=\tanh\omega$ : on passe de l'un à l'autre en changeant $\frac12\omega$ en $\omega$, la moitié de la lame libre devenant une lame entière. C'est exact.\footnote{Aux deux états (vues 301, 304 et 381), le feuillet écrit $\tan\frac12\omega=\frac{1-e^{\omega}}{1+e^{\omega}}$ et $\tan\omega=\frac{1-e^{2\omega}}{1+e^{2\omega}}$, avec le signe contraire de celui qu'elle écrit elle-même aux § 16 et 19 et que demandent les coefficients ; on rétablit le signe. À la vue 381, une phrase intervertit les chiffres des cas I et II ; à la vue 304 du brouillon, un groupe très repassé, lu avec doute « $A=-e^{\omega}$ », « $x=1$ », suit les coefficients.}

Entre le cas V et le cas VI, le même rapport semble manquer : au cas VI, $\alpha=1$, $\beta=-e^{\omega}$, $\gamma=-(1+e^{\omega})$, $\delta=-(1-e^{\omega})$, alors qu'au cas V Euler donne
\[
y\propto e^{\omega u}-e^{-\omega u}\mp\sqrt{2\bigl(e^{2\omega}+e^{-2\omega}\bigr)}\,\sin\omega u ,
\]
sans rapport apparent. Elle lève le paradoxe : Euler met l'origine au bout appuyé, alors que la demi-lame du cas VI doit l'avoir au bout encastré. En changeant $u$ en $1-u$ et en portant $\sin\omega=\pm(e^{\omega}-e^{-\omega})/\sqrt{2(e^{2\omega}+e^{-2\omega})}$ et $\cos\omega=\pm(e^{\omega}+e^{-\omega})/\sqrt{2(e^{2\omega}+e^{-2\omega})}$, elle obtient $y\propto e^{\omega u}-e^{(2-u)\omega}-(1+e^{2\omega})\sin\omega u-(1-e^{2\omega})\cos\omega u$, c'est-à-dire les coefficients du cas VI où $\omega$ tient la place de $\frac12\omega$. Le calcul est juste.\footnote{Vues 381 et 383 ; brouillon, vues 304 et 306. La transcription rend le radical $\sqrt{2}\,(e^{2\omega}+e^{-2\omega})$ ; le calcul demande le radical sur tout le produit, $\sqrt{2(e^{2\omega}+e^{-2\omega})}=2\sqrt{\cosh2\omega}$, et c'est ainsi qu'on l'écrit. C'est une lecture de cette édition, à vérifier sur l'image. Avec $\tan\omega=\tanh\omega$, $\sin\omega=\pm2\sinh\omega/\sqrt{2(e^{2\omega}+e^{-2\omega})}$, d'où le terme d'Euler $-2(\sinh\omega/\sin\omega)\sin\omega u$. À la vue 383, les deux valeurs portent le même numérateur $e^{\omega}-e^{-\omega}$ ; le brouillon (vue 306) donne $e^{\omega}+e^{-\omega}$ pour le cosinus, qui est juste. À la vue 381, le premier « cas VI » de l'alinéa désigne le cas V.} Elle confirme par les conditions d'Euler pour un bout appuyé (§ 13), $y=0$ et $y''=0$ : au milieu de la lame libre, pour un nombre impair de nœuds, $y''(\frac12)=-\bigl[(1+e^{\omega})\sin\frac12\omega+(1-e^{\omega})\cos\frac12\omega\bigr]=0$ en même temps que $y(\frac12)=0$. C'est exact.

L'énoncé actuel est celui-ci. Tout mode d'une lame libre et appuyée de longueur $\ell$ est la moitié d'un mode antisymétrique d'une lame libre de longueur $2\ell$, et réciproquement : le prolongement impair par-delà le bout appuyé garde $y=y''=0$ au raccord et rend $y'$ et $y'''$ continus. De même, la lame encastrée et appuyée est la moitié d'un mode antisymétrique de la lame encastrée. Les cas II et V sont donc contenus dans les cas I et VI. Le cas III ne l'est pas : le prolongement pair d'un mode, qui donne les modes symétriques, a au raccord $y'=0$ et $y'''=0$, non $y=y'=0$. C'est le bout que le long mémoire appelle « libre au sens nouveau » ; voir son § 3 (vues 414 à 432). La réduction de six à quatre est juste.

Elle ajoute pourquoi, sans nœud au milieu, les deux moitiés d'une lame libre ne peuvent être des lames entières soumises aux mêmes conditions : la moitié de la partie centrale ne peut égaler la partie du bout, par l'inégalité du § 14. C'est vrai, et la raison actuelle est plus simple : dans un mode symétrique, chaque moitié a au milieu $y'=0$ et $y'''=0$, et c'est une lame libre à un bout et « libre au sens nouveau » à l'autre, non une lame libre aux deux bouts.\footnote{Vue 385, page (33), premier alinéa, qui renvoie à « l'observation que j'ai faites, § 14 » (le 14 lu avec doute) ; brouillon, vue 308, page 42, qui s'arrête sur « extremités ». Au troisième son de la lame libre, par exemple, la demi-partie centrale vaut $0{,}1442$ de la longueur, la partie du bout $0{,}0944$.}

\subsection*{Vers la plaque carrée (§ 25)}

\pagerange{385}{389}
Le dernier paragraphe passe aux plaques, dont la théorie « semble devoir être beaucoup plus compliquée ». Il relève une singularité : alors que la lame libre ne peut être formée de deux demi-lames libres, plusieurs figures de plaques carrées libres de tous côtés peuvent être regardées comme composées de quatre plaques soumises aux mêmes conditions (Chladni, Traité d'acoustique, § 110). Les figures composées sont celles qu'aucune ligne nodale ne coupe diamétralement en deux parties égales, parce qu'une partie contiguë ne peut tenir lieu d'un bord libre que si elle est elle-même libre, et donc en mouvement. La figure 63 de Chladni, la plus simple de celles qui ont des lignes nodales en différents sens, deux lignes qui se croisent au centre d'après son dessin de la vue 391, répétée quatre fois, donne la figure 68 a, comme l'auteur le remarque ; répétée neuf fois, la figure 75 ; seize fois, la figure 82. Elle l'a vérifié en y traçant des lignes ponctuées.\footnote{Vues 385 et 387, pages (33) et (34), son § 25, le 5 en forme de S ; « 110 », « 63 » et « contigue » lus avec doute. Une note de marge précise que les seules figures composées dont il s'agit sont décomposables en quatre figures simples, celles qui en contiennent un nombre impair ayant des lignes nodales diamétrales. Brouillon, vues 310 et 312, pages 43 et 44 : la figure composée de quatre est la « fig 64. a », et « huit » est biffé devant « 16 fois » ; la phrase qui suit, « mais cependant M\textsuperscript{r} Chladni ne dit pas avoir obtenu ces deux dernieres figures par le rapprochement de 9 et 16 plaques portant la fig. 63 », est barrée. Les figures redessinées sont à la vue 391. Au § 8 (vue 334), une note encadrée disait déjà que les vibrations transversales des plaques carrées sont comprises dans le cas de la lame ; voir la section précédente.} Elle en tire un énoncé général : si une plaque carrée de côté $am$ est divisée en $m^2$ carrés de côté $a$ et qu'un stylet fait naître dans un carré d'angle une figure nodale, cette figure se répète dans les $m^2-1$ autres ; les sections de la surface par des plans perpendiculaires aux côtés coupent le plan de repos aux points $\alpha$, $\alpha+a$, $\alpha+2a$, … et $\alpha'$, $\alpha'+a$, … , et sont symétriques autour de ces points. Pour les membranes tendues, par le § 24 du mémoire de Biot au tome 4 des Mémoires de l'Institut, $\alpha=\alpha'=0$ ; pour les composées de la figure 63, $\alpha=\alpha'=\frac12a$.\footnote{Vues 387 et 389. À la vue 387, le troisième terme de la seconde suite est écrit « $2\alpha'+2a$ » pour $\alpha'+2a$. Le brouillon (vue 312) cite le « § 26 » de Biot, le 6 lu avec doute, là où la mise au net écrit « § 24 » ; il porte en marge un petit dessin d'un carré partagé en quatre et la remarque « pour les figures composées dela fig 63 $\alpha$ et $\alpha'=\frac12a$ ». Le mémoire de Biot n'a pas été consulté pour cette lecture.}

Ce que dit la théorie actuelle. Pour la membrane fixée à son bord, comme pour la plaque simplement appuyée sur ses quatre côtés, les modes $\sin\frac{p\pi x}{L}\sin\frac{q\pi y}{L}$ se répètent exactement dans $m^2$ carrés, dont les côtés sont des lignes nodales. C'est ce que dit $\alpha=\alpha'=0$.\footnote{Pour la plaque, c'est la solution dite de Navier, que les feuillets ne nomment pas ; l'égalité de ces modes avec ceux de la membrane ne vaut que pour des bords droits.} Pour la plaque libre, la répétition n'est pas exacte. Une ligne intérieure de part et d'autre de laquelle un mode se prolonge par réflexion est soit une ligne où $z$ et sa dérivée seconde normale s'annulent (réflexion impaire : une ligne d'appui, nodale), soit une ligne où les dérivées première et troisième normales s'annulent (réflexion paire : un axe de symétrie). Un bord libre n'est ni l'un ni l'autre. Dans les composées de la figure 63, où $\alpha=\frac12a$, les côtés des petits carrés sont des axes de symétrie : chaque petit carré d'angle y vibre comme une plaque dont deux côtés sont libres et deux sont des axes de symétrie, où la pente et la dérivée troisième normales s'annulent, non comme une plaque libre. L'accord des figures de Chladni avec cette composition est donc approché, ce que les dessins ne peuvent trancher ; les conditions de bord de la plaque libre sont énoncées dans la section du long mémoire (N\textsuperscript{o} 9, vues 433 à 438).\footnote{Le calcul des figures de la plaque carrée libre par approximations successives est celui de Ritz (1909) ; le nom seul est donné, les feuillets le précèdent.}

La théorie complète de la plaque carrée comprendrait, dit-elle, « quinze » cas selon que les côtés sont appuyés, libres ou fixés. Quinze est le nombre de façons de choisir combien de côtés sont dans chaque état. Si l'on distingue deux côtés de même état adjacents ou opposés, et qu'on identifie les dispositions échangées par une symétrie du carré, il y en a vingt et un.\footnote{Dénombrement de cette lecture : $3^4=81$ dispositions, réparties en 21 classes sous les huit symétries du carré, et en 24 sous les seules rotations. « quinze » est lu avec doute à la vue 389 (« quize » en apparence) ; le brouillon écrit « quinze » (vue 314) et « 15 » (vue 316).} Ces cas, ajoute-t-elle, auraient entre eux des relations qui permettraient de n'en calculer que quelques-uns. Ainsi une figure de la plaque libre à deux lignes diamétrales se coupant au centre pourrait peut-être se ramener à quatre plaques carrées ayant chacune deux côtés libres et deux appuyés. C'est exact pour les modes impairs par rapport aux deux médianes : une fonction impaire en $x$ a $z=0$ et $\partial^2z/\partial x^2=0$ sur la médiane, qui est ainsi une ligne d'appui.

Elle finit par une réflexion sur le mot « appuyé ». Pour la lame, il désignerait un état fixe, « entièrement équivalent aux nœuds de vibrations », puisque le mouvement doit se transmettre à travers le point d'appui et qu'on fait abstraction de la largeur. Pour la plaque, le mouvement pourrait passer par les points voisins, et l'irrégularité de certaines figures de Chladni lui fait croire que le mouvement est parfois plus gêné dans une moitié de la plaque que dans l'autre. La page s'arrête sur « dans la ».\footnote{Vue 389, page (35) ; « appuyé » est souligné. Les derniers mots sont serrés contre le bord droit, et aucune vue ne continue la phrase. Brouillon, vues 314 et 316 : la première page annule plusieurs passages par des hachures ; la seconde, sans numéro, reprend l'ouverture de la première et la poursuit autrement. Elle y trouve l'« immensité de cas » augmentée « d'une manière effrayante » par les degrés de pression que Chladni rapporte : une pression faible diffère peu de la liberté, et la fixité est la limite d'une échelle « composée d'un nombre infini de degré ». Un passage barré y ajoutait la position, la pression et la vivacité de l'archet.} Pour la lame, l'équivalence du stylet et d'un nœud est encore l'hypothèse de la force nulle : elle vaut pour les modes sans force et non pour les autres. Pour la plaque, un stylet impose $z=0$ en un point, qui n'est pas une ligne. Les modes de la plaque ainsi retenue ne sont en général pas ceux de la plaque libre, et ceux qui n'exercent aucune force sur le stylet sont ceux de la plaque libre dont une ligne nodale passe par ce point. Ses « degrés de pression » répondent aujourd'hui à un appui élastique, dont la raideur va de zéro, le bord libre, à l'infini, l'appui ; l'encastrement demande en plus qu'on empêche la rotation.

\subsection*{Deux feuillets de figures, deux calculs}

\pagerange{391}{396}
Entre la fin de cette mise au net et le long mémoire, quatre feuillets sans texte suivi.\footnote{Vues 391 à 396, n\textsuperscript{os} 191 à 194 de la série à l'encre, sans pagination de sa main. Les vues 393 et 395 sont des revers blancs ; la vue 392 n'est pas le revers de la vue 391, mais un feuillet plus petit.} La vue 391 porte six figures. En haut, une lame $EF$ et sa courbe de vibration, qui coupe l'axe en $L$ (« fig 1 ») ; puis la même lame avec une courbe à sept arcs et six nœuds $L$, $L'$, … , $L^{\text{VI}}$ (« fig. 1$'$ »). Le chiffre de ces légendes se lit 1 ou 7, et la mise au net cite en marge une « fig. 7 » dont les parties sont $EL$, $LL'$, … , $L^{\text{vi}}E'$ (vues 322 et 326) : c'est probablement elle.\footnote{Rapprochement de cette lecture, non écrit sur le feuillet.} Suivent les figures 63, 68 a, 75 et 82 de Chladni, redessinées avec les lignes ponctuées que le § 25 annonce, la figure 63 répétée quatre, neuf et seize fois.

Les vues 392 et 394 sont deux petits feuillets de calcul, sans titre, dans les notations de la lame libre. Le premier cherche si l'ordonnée peut s'annuler à la fois en $\lambda$ et en $\frac12\lambda$. De $y(\frac12\lambda)=0$ multipliée par $2\cos\frac12\lambda\omega$ et retranchée de $y(\lambda)=0$, il tire $2\cos\frac12\lambda\omega=e^{\frac12\lambda\omega}+e^{-\frac12\lambda\omega}$, « resultat absurde » pour $\lambda\omega\neq0$ puisque $\cos t<\cosh t$, et conclut que la seconde équation est inadmissible. Pour l'ordonnée qu'il écrit, celle de la lame libre, une étape est fausse. Le membre de droite doit être $e^{\lambda\omega}\pm e^{\omega-\lambda\omega}-1\mp e^{\omega}$, et non $e^{\lambda\omega}\pm e^{\omega-\lambda\omega}+1\pm e^{\omega}$ ; avec le bon signe, on n'aboutit pas à une absurdité. Le calcul du feuillet est juste en revanche pour l'ordonnée de la lame encastrée, où les termes trigonométriques changent de signe ; pour elle, la conclusion est démontrée.\footnote{Vérifié symboliquement pour cette lecture. Pour la lame libre, aucun des dix premiers modes n'a un nœud à une distance du bout double de celle d'un autre ; l'écart le plus petit, $0{,}023$ de la longueur, est au dixième mode. Le signe lu « \&c » après la dernière exponentielle est douteux ; le feuillet s'arrête sur « or cette ».} Le second feuillet suppose l'ordonnée la plus grande en $\frac12\lambda$, soit $dy/du=0$, et, avec $y(\lambda)=0$, obtient
\[
\bigl(1+e^{\lambda\omega}\bigr)\bigl(1\pm e^{(1-\lambda)\omega}\bigr)=2\bigl(e^{\frac12\lambda\omega}\mp e^{(1-\frac12\lambda)\omega}\bigr)\sin\tfrac12\lambda\omega ;
\]
les étapes sont justes, et le feuillet s'arrête là, sans conclure.\footnote{Vue 394, avec le cachet « Bibliothèque Impériale MSS. » ; un exposant de la dernière ligne est empâté. Le feuillet écrit $x$ pour la fraction ; on écrit $\lambda$.} Rien sur ces deux feuillets ne les rattache à un paragraphe de la critique.

La vue 396 porte six plaques carrées tracées sur un quadrillage pointillé : les figures 67 c, 72 a, 78 et 85 de Chladni, avec des lettres aux sommets et aux milieux des côtés, et deux figures sans numéro. Aucun texte de ces pages ne les cite.

\section*{Le mémoire sur la question proposée par la première classe de l'Institut}

\pagerange{398}{490}
Des vues 398 à 592 court un seul long mémoire sur les surfaces élastiques, paginé par
elle de (1) à (165), écrit au recto et au verso de feuillets qui sont ensuite montés
sur onglets ; c'est une mise au net corrigée, dont les brouillons de l'ouverture sont
ailleurs dans le volume (vues 593 à 626, section suivante). Cette section en lit les
pages (1) à (80) : l'équation du mouvement, ses intégrales particulières, les
conditions aux limites de la lame puis de la plaque, et l'explication des figures de
Chladni sur la plaque carrée. La comparaison des sons, qui commence à la page (81)
avec le N\textsuperscript{o} 16, est lue plus bas, dans la section « Le mémoire
sur la question de la première classe (suite) : la théorie devant l'expérience »
(vues 491 à 592), après les autres états de l'ouverture.\footnote{Vues 398 à 490. Pages (1) à (80) de sa pagination, en tête des
pages, les impaires souvent barrées d'un trait ; série à l'encre en haut à droite, un
nombre par feuillet, de 195 (vue 398) à 240 (vue 489). Sont reliés dans ce passage
sans en faire partie : un feuillet de calculs (vues 410 et 411, n\textsuperscript{o}
201), le petit feuillet sur la cire (vue 416, n\textsuperscript{o} 204, lu avec le
§ 6 du mémoire, qu'il corrige), un petit feuillet de calculs (vue 428, n\textsuperscript{o}
210), le fragment sur la mesure de la courbure (vue 465, n\textsuperscript{o} 228) et
le feuillet de calculs de la vue 483 (n\textsuperscript{o} 237), ces deux derniers lus
avec les feuillets épars sur la courbure. Les vues 430, 431, 466 et 484 sont des dos
blancs ou une seconde prise (431 répète 429). Que ce mémoire soit celui qu'elle
appelle ailleurs « mon mémoire pour 1813 » est une inférence de cette lecture,
exposée à la dernière sous-section de la section suivante.}\footnote{Le feuillet
écrit les dérivées avec un $d$ droit ($d^4z/dx^4$, $ddz/dx^2$, $d^2z/dxdy$) ; on écrit
$\partial^4 z/\partial x^4$, $\partial^2 z/\partial x^2$, $\partial^2 z/\partial x\partial y$,
et dans le commentaire $z_{xx}$, $z_{xy}$. On note
$\Delta = \partial^2/\partial x^2 + \partial^2/\partial y^2$ et $\Delta^2$ le
bilaplacien.}

\subsection*{Le titre, la question, l'épigraphe}

\pagerange{398}{398}
Le titre donne la question telle que l'Institut l'avait posée : « Donner la théorie
mathématique des vibrations des surfaces élastiques et la comparer à
l'expérience ». Il est suivi d'une épigraphe tirée du \emph{Novum Organum} de Bacon
(livre I, aphorisme XCII) : « Sed longe maximum progressibus scientiarum […] obstaculum
deprehenditur in desperatione hominum et suppositione impossibilis ».\footnote{Vue 398,
page (1), n\textsuperscript{o} 195 de la série à l'encre ; un brouillon très corrigé,
lignes entières barrées d'un trait épais et leçon nouvelle écrite dans l'interligne.
La citation latine est d'une plume plus fine ; le feuillet écrit « isdem » et
« impossibilitis », que la mise au net de la vue 593 écrit « iisdem » et
« impossibilis ». Traduction : « Mais de loin le plus grand obstacle aux progrès des
sciences, et à l'entreprise de tâches et de domaines nouveaux, se trouve dans le
désespoir des hommes et dans la supposition de l'impossible. » La référence porte
« tome 1 p. 298. edit in f\textsuperscript{o} », puis un mot qui commence par « lond »
(lu avec doute) et dont la fin est illisible ; la mise au net écrit « London
1740 ».}\footnote{Contexte, non tiré des feuillets : les pièces envoyées aux concours
de l'Institut étaient anonymes et se reconnaissaient à leur épigraphe, reprise sur un
billet cacheté qui portait le nom de l'auteur.}

Le préambule annonce la méthode. Traiter la question dans toute sa généralité
mènerait à des calculs très compliqués ; mais si l'on se borne aux surfaces planes, et
à leur mouvement « régulier et comparable à celui d'un pendule simple », on parvient
aisément à une équation dont les intégrales particulières s'appliquent sans peine.
Elle défend ce choix de deux manières : les expériences de Chladni ne portent que sur
ces mouvements réguliers, « les seuls en effet qui appartiennent à l'acoustique,
puisqu'eux seuls présentent des lois déterminables » ; et la supposition d'un
déplacement $z$ très petit se transporte au petit nombre de mouvements que Chladni a
observés sur des surfaces courbes, qu'une même équation « fort simple » réunit à
plusieurs autres corps sonores.\footnote{Le préambule s'est fait en plusieurs
couches ; l'ordre des corrections est celui que la page suggère. « examiner » remplace
« embrasser » ; « différentielle dont les intégrales particulières sont facilement
applicables » remplace « fort simple » ; la note de la marge, « puisqu'eux seuls
presentent des lois déterminables », remplace une addition barrée (« en tant qu'elle
s'occupe des lois relatives au son ») ; « je soumets » remplace « j'ai l'honneur de ».
Sont barrées, entre autres, la phrase « de l'autre, cet habile physicien n'a fait
qu'un petit nombre d'essais sur le son des surfaces courbes » (« n'a fait qu'un » lu
avec doute) et une phrase dont on ne lit que « surfaces d'une courbure quelconque en
substituant », « différentielles », « de la courbe » et « desorte que les valeurs de »,
la fin de chacun de ces groupes étant illisible. Ce que disait cette phrase barrée, un
autre état de l'ouverture le donne en entier (vues 595, 598 et 614, section suivante) :
transporter la théorie aux surfaces courbes en prenant pour coordonnées les lignes de
courbure principales. La lettre $z$ est soulignée.}

En termes actuels, le « mouvement régulier » est un mode propre : une solution de la
forme $z = Z(x,y)\sin(\Omega t + \zeta)$, où toute la surface oscille à une même
fréquence. Se borner aux modes est légitime pour l'acoustique, puisque l'équation est
linéaire et que tout mouvement est une superposition de modes ; elle le dit elle-même
plus loin, en écartant « la coexistence des sons » comme une question qui n'ajoute pas
de difficulté (vue 406).

\subsection*{§ 1. L'équation de la surface élastique vibrante (N\textsuperscript{os} 1 et 2)}

\pagerange{399}{404}
Le N\textsuperscript{o} 1 cherche la somme des moments des forces d'élasticité sur
toute la surface.\footnote{Vues 399 à 404. Pages (2) à (7) ; n\textsuperscript{os}
195 (verso écrit), 196, 197 et 198 de la série à l'encre. Mise au net peu corrigée :
« d'où il suit » devient « Ainsi je trouve » (vue 399). La lettre de l'épaisseur,
un jambage coiffé d'un trait, est lue $\tau$ avec doute à sa première occurrence (vue
402) ; elle a la forme d'un $\Gamma$ à la vue 413.} Elle prend quatre points voisins
de la surface, $(x,y,z)$, $(x, y+dy, z+dz)$, $(x+dx, y, z+dz)$ et
$(x+dx, y+dy, z+2dz+d^2z)$, fait passer un plan par le premier, le deuxième et le
quatrième, un autre par le premier, le troisième et le quatrième, appelle $e$ l'angle
extérieur de ces deux plans, $\beta$ leur petite ligne d'intersection et $E$ la force
qui, agissant le long de $\beta$, s'oppose à l'inflexion ; la somme des moments est
$\iint E\beta\,de$. Le calcul de $e$ par les coefficients $A, B, C$ et $A', B', C'$ des
deux plans est juste : le numérateur de $\sin^2 e$ vaut exactement
\[
dx^2dy^2\,(d^2z)^2\left[(2\,dz + d^2z)^2 + dx^2 + dy^2\right],
\]
ce qu'elle écrit, et, en négligeant $dz$ devant $dx$ et $dy$ pour une surface voisine
de son plan,
\[
\sin^2 e = \frac{(dx^2+dy^2)\,(d^2z)^2}{dx^2\,dy^2}, \qquad \beta^2 = dx^2 + dy^2 .
\]
Elle prend ensuite la force proportionnelle à la somme des courbures principales et à
un facteur de matière et d'épaisseur, « hypothèse que je crois être la plus
convenable » :
\[
E = b\,\tau^3\left(\frac{1}{r} + \frac{1}{r'}\right), \qquad
\frac{1}{r} + \frac{1}{r'} = -\,\Delta z ,
\]
$b$ dépendant de la rigidité naturelle, $\tau$ étant l'épaisseur, la seconde égalité
valant au premier ordre quand $\partial z/\partial x$ et $\partial z/\partial y$ sont
petits.\footnote{Le feuillet écrit $\frac{1}{(r)}+\frac{1}{(r')} = -\left(\frac{ddz}{dx^2}+\frac{ddz}{dy^2}\right)$.
La relation de la vue 399, $-\cos e = \frac{AA'+B'B'+CC'}{\ldots}$, porte $B'B'$ au
numérateur là où la ligne précédente a $BB'$ ; c'est un lapsus, le calcul de la vue
400 employant $BB'$. À la vue 401, la première ligne s'arrête au bord sur « $+d$ » :
le calcul demande $d^2z$. À l'avant-dernière ligne du calcul du numérateur, la
dernière accolade, $\{dx\,d^2z\}$, n'a pas son exposant 2, le bord de la ligne
touchant la marge.} Une intégration par parties, dont elle laisse de côté les
termes aux limites, donne alors pour la somme des moments
$-\iint b\tau^3\,\Delta^2 z\,\delta z\,dx\,dy$ ; en l'égalant à
$\iint \partial_t^2 z\,\tau\,dx\,dy\,\delta z$ (la masse de l'élément étant prise
proportionnelle à son épaisseur, la gravité pour unité), elle obtient
\[
\frac{\partial^2 z}{\partial t^2} + b\,\tau^2\,\Delta^2 z = 0 \qquad \text{(A)}.
\]
\footnote{Le feuillet écrit
$b\tau^2\left\{\frac{d^4z}{dx^4}+2\frac{d^4z}{dx^2dy^2}+\frac{d^4z}{dy^4}\right\}+\frac{ddz}{dt^2} = 0$,
vue 403. Un 4 est écrit sous l'exposant de $\tau^3$, contre l'accolade, sans qu'on
voie ce qu'il marque.}

C'est l'équation des plaques minces vibrantes, celle qu'on appelle aujourd'hui
équation de Germain--Lagrange : le bilaplacien de la flèche équilibre
l'accélération.\footnote{Contexte, non tiré des feuillets : la pièce qu'elle avait
envoyée au premier concours portait une équation fautive, que Lagrange, commissaire, a
corrigée en 1811 sous la forme qu'on lit ici ; le nom de « Germain--Lagrange » vient de
là.} La dépendance en épaisseur est celle de la théorie actuelle : une rigidité en
$\tau^3$ divisée par une masse en $\tau$ donne des fréquences proportionnelles à
$\tau$, comme le coefficient $D/(\rho h) = Eh^2/12\rho(1-\nu^2)$ de
Kirchhoff.\footnote{Kirchhoff, 1850 ; le feuillet est antérieur et ne le cite pas.
Ici $E$ est le module d'Young, $h$ l'épaisseur, $\rho$ la densité, $\nu$ le coefficient
de Poisson ; ce $E$ n'est pas la force $E$ du feuillet.} Mais la dérivation n'est pas
une démonstration. L'angle $e$ qu'elle calcule est le pli des deux plans le long de la
diagonale de l'élément : sa différence $d^2z$, prise avec le même $dz$ sur les deux
côtés, est la différence seconde mixte
$z(x+dx,y+dy) - z(x+dx,y) - z(x,y+dy) + z(x,y) \approx z_{xy}\,dx\,dy$, si bien que $e$
mesure la torsion de l'élément et non sa courbure moyenne, alors que la force $E$ est
supposée proportionnelle à $\Delta z$ ; le passage de
$\iint E\beta^2\,\delta(d^2z/ds)$ à $\iint dd.(E\beta^2)/ds\;\delta z$, avec
$ds = dx\,dy$ et « $dd.$ » lu comme un laplacien, est formel ; et la valeur de $E$ est
une hypothèse. Ce qui en ferait une démonstration est connu : se donner une énergie de
flexion par unité d'aire, quadratique dans les courbures — celle que ses propres
conditions aux limites supposent (N\textsuperscript{o} 9 plus bas) est
$\frac12 b\tau^3\iint(\Delta z)^2\,dx\,dy$ — et appliquer le principe de Hamilton, qui
donne à la fois (A) et les conditions naturelles au bord ; l'énergie elle-même se
justifie ensuite à partir de l'élasticité à trois dimensions sous les hypothèses de
Kirchhoff. L'énergie de Kirchhoff diffère de la sienne par un multiple de
$\iint(z_{xx}z_{yy} - z_{xy}^2)$, qui ne change pas l'équation intérieure et ne change
que les conditions au bord libre ; c'est le lagrangien nul que la lecture de
l'\emph{Exposition des principes} met en lumière.

Le N\textsuperscript{o} 2 passe au mouvement régulier. Si la surface oscille comme un
pendule simple de longueur $K$, la gravité et la seconde prises pour unités, on a
$\partial_t^2 z = -z/K$ et (A) devient
\[
K b\tau^2\,\Delta^2 z = z \qquad \text{(B)}.
\]
En posant $K b\tau^2 = h^4$, elle observe que (B) « peut s'abaisser au second degré » :
\[
h^4\Delta^2 - 1 = (h^2\Delta - 1)(h^2\Delta + 1),
\]
de sorte que toute solution de $h^2\Delta z = z$ ou de $h^2\Delta z = -z$ est solution de
(B), la seconde étant l'équation des mouvements réguliers des surfaces tendues —
l'équation qu'on appelle aujourd'hui de Helmholtz.\footnote{Le nom est postérieur au
feuillet. Le feuillet écrit la seconde équation $h^2\{\frac{d^2z}{dy^2}+\frac{d^2z}{dx^2}\}=-z$,
le signe moins serré contre $z$, peut-être ajouté ; la lettre $h$ y est tracée comme
un $b$ barré.} La réciproque, qu'elle n'énonce pas ici mais qu'elle emploie à la
vue 436, est vraie aussi : toute solution de (B) s'écrit $z = z' + z''$ avec
$z' = \frac12(z + h^2\Delta z)$ solution de la première et
$z'' = \frac12(z - h^2\Delta z)$ solution de la seconde.

Elle en conclut que « quelle que soit la figure que puisse prendre une surface tendue
renfermée dans des limites données, la même figure peut appartenir également à une
surface élastique renfermée dans les mêmes limites ».\footnote{La phrase précédente
porte plusieurs pluriels annulés d'une petite croix, et après « que » une lettre
barrée qui n'a pas été lue (vue 404) ; « cette équation n'a point d'intégrale générale
en termes finis » est dit de (B).} Pour l'équation intérieure, c'est exact. Pour un
corps donné, avec son bord, ce ne l'est que si le mode de la membrane satisfait aussi
les conditions de la plaque : c'est le cas d'un polygone à côtés droits appuyés, où
$z = 0$ sur un côté entraîne $\Delta z = -z/h^2 = 0$, et d'un bord à pente normale nulle,
où $\partial_n z = 0$ entraîne $\partial_n\Delta z = 0$ ; ce n'est pas le cas d'un bord
encastré ou libre, ni d'un bord appuyé courbe (sur un cercle, le moment de flexion fait
intervenir $\nu\,\partial_r z/r$). Même quand la figure est commune, le son ne l'est
pas : la membrane a des fréquences proportionnelles au nombre d'onde $k$, la plaque à
$k^2$.

\subsection*{§ 2. L'intégration : les intégrales (C) à (I) (N\textsuperscript{os} 3 à 6)}

\pagerange{405}{415}
Le N\textsuperscript{o} 3 sous-entend désormais le facteur du temps
$\sin(t\sqrt{1/K} + \zeta)$ et laisse de côté la superposition des modes, ce qui est
permis puisque l'équation est linéaire.\footnote{Vues 405 à 415. Pages (8) à (14) ;
n\textsuperscript{os} 199 à 203 de la série à l'encre. La pagination repasse deux fois
par (9) et (10) (vues 406 et 407, puis 408 et 409) ; les vues 408 et 409, puis 412 et 413, sont
deux mises au net successives de la seconde moitié de la page (10) de la vue 407, et
les pages (11) et (12) (vues 412 et 413) reprennent la matière des vues 408 et 409.
Entre les deux, le feuillet de calculs des vues 410 et 411 (n\textsuperscript{o} 201,
écrit en travers, sans pagination) dérive deux fois en $x$ une série de la forme de
(F), avec $\Pi - \Pi'$ pour paramètre, et essaie
$z = \sin\frac{mx}{A}\sin\frac{my}{A}\{\cos\sqrt{m^2+n^2}\frac{x}{A} + \cos\sqrt{m^2+n^2}\frac{y}{A}\}$ ;
un facteur noirci y reste illisible (vue 411).} Elle cherche alors des intégrales
particulières de (B) sous diverses conditions entre $x$, $y$ et $z$.

\emph{L'intégrale (C).} Sous la condition $\partial^2 z/\partial x\,\partial y = 0$, $z$
est la somme d'une fonction de $x$ et d'une fonction de $y$ :
\[
z = \mathcal{A}\,f(x) + B\,f(y), \qquad
f(s) = a\,e^{\omega s/A} + b\,e^{-\omega s/A} + c\sin\frac{\omega s}{A} + d\cos\frac{\omega s}{A},
\]
où $A$ est le côté de la plaque ; comme $f'''' = (\omega/A)^4 f$, (B) est satisfaite si
$Kb\tau^2(\omega/A)^4 = 1$, et le son, qu'elle mesure par $\sqrt{1/K}/\pi$, vaut
$\omega^2\tau\sqrt b/(\pi A^2)$ : c'est ce qu'elle écrit.\footnote{La fréquence est
$\sqrt{1/K}/2\pi$ dans ses unités ; sa mesure en diffère d'un facteur constant, sans
effet sur les rapports. Le coefficient de (C) est une capitale cursive bouclée,
distincte du $A$ des dénominateurs, rendue $\mathcal{A}$ par la transcription, qui
note qu'elle a aussi la forme d'un B cursif.} La fonction $f$ est celle de la lame
d'Euler, et c'est à cette intégrale, munie des coefficients d'Euler pour les bouts
libres, que servira tout le § 4. Elle est la seule des intégrales du § 2 qui ne
satisfasse pas l'équation des surfaces tendues.

\emph{L'intégrale (D).} Sous la condition
$\partial^2 z/\partial x^2 = (n^2/m^2)\,\partial^2 z/\partial y^2$, qu'elle présente
comme l'expression d'une périodicité, elle prend
\[
z = \mathcal{A}\sin\frac{\pi n x}{A}\,\sin\frac{\pi m y}{A}, \qquad
\frac{\sqrt{1/K}}{\pi} = \frac{\pi(m^2+n^2)\,\tau\sqrt b}{A^2},
\]
ce qui est juste : $\Delta z = -\pi^2(m^2+n^2)z/A^2$, et (B) donne ce son.\footnote{La
condition n'exprime pas à elle seule la périodicité : c'est une équation des ondes en
$(x,y)$, dont la solution générale est $\varphi(mx+ny) + \psi(mx-ny)$, et les produits
de sinus n'en sont qu'une famille particulière. Le second dénominateur de (D) est un
$A$ traversé d'un trait oblique, accent ou rature (vue 406).} C'est un mode de
membrane, et c'est exactement un mode de la plaque carrée de côté $A$ appuyée sur son
contour, quelle que soit la matière : sur un côté droit où $z = 0$, la condition
d'appui de Kirchhoff (moment nul) se réduit à $\Delta z = 0$, que ce produit
satisfait.

\emph{Les intégrales (E) et (F).} Elle écrit ensuite des séries en puissances de
$\cos u$, $u = \pi n x/A$, dont les coefficients dépendent d'un paramètre $\Pi$ :
\[
S_1 = \cos u + \frac{\Pi+3}{2\cdot 3}\cos^3 u + \frac{(\Pi+3)(\Pi+15)}{2\cdot3\cdot4\cdot5}\cos^5 u + \cdots,
\]
\[
S_2 = \frac{2}{\Pi} + \cos^2 u + \frac{\Pi+8}{3\cdot 4}\cos^4 u + \frac{(\Pi+8)(\Pi+24)}{3\cdot4\cdot5\cdot6}\cos^6 u + \cdots,
\]
et pose, pour (E), $z = \sin u\,\sin v\,(a'S_1 + b'S_2)_u\,(A'S_1 + B'S_2)_v$ avec
$v = \pi m y/A$ et le même $\Pi$ dans les deux facteurs, pour (F) la même forme avec
$u = \pi m x/A$ et $\Pi$ en $x$, $\Pi'$ en $y$.\footnote{Le même signe, un $\pi$ aux
jambages hauts, sert sur le feuillet pour la demi-circonférence et pour le paramètre ;
la distinction $\pi$ / $\Pi$ est du transcripteur. Le point entre deux chiffres est
son signe de produit : « $3.5$ » et « $3.8$ » valent 15 et 24. Les termes généraux
qu'elle donne à la vue 408,
$\frac{(\Pi+3)(\Pi+15)\cdots(\Pi+4n^2-1)}{2\cdot3\cdots(2n+1)}\cos^{2n+1}$ et
$\frac{(\Pi+8)(\Pi+24)\cdots(\Pi+(2n-1)^2-1)}{3\cdot4\cdots 2n}\cos^{2n}$, sont ceux
que donne la récurrence des coefficients ; la copie de la vue 412 écrit
$2\cdot3\cdot4\cdot5\cdots 2n$ au dénominateur des lignes en $b''$ et $B''$, ce qui est
un lapsus. Dans (E) (vue 407), le dernier terme de la ligne en $B'$ porte $nx$ là où
$my$ est attendu ; dans (F), plusieurs arguments des lignes en $y$ portent $mx$, et
les accents de $\Pi'$ y sont ajoutés sur des groupes barrés ; la mise au net de la vue
408, étiquetée (F$'$), écrit partout $my$ et met $\Pi$ en $x$, $\Pi'$ en $y$. On suit
cette mise au net.} Ce que sont ces séries se voit en posant $\Pi = 1 - p^2$ : ce sont
des développements de sinus et de cosinus d'arcs multiples,
\[
\sin u\;S_1(\cos u) = \frac{1}{p}\,\sin p\!\left(\frac{\pi}{2} - u\right), \qquad
\sin u\;S_2(\cos u) = \frac{2}{\Pi}\,\cos p\!\left(\frac{\pi}{2} - u\right)
\qquad (0 < u < \pi),
\]
identités qu'elle retrouvera elle-même, pour $p$ entier, dans les \emph{Leçons sur le
calcul des fonctions} de Lagrange (vue 477), et qui valent pour tout $p$.\footnote{Vérifié
pour cette lecture à trente décimales, pour $p = 1{,}5$, $2{,}7$ et $\sqrt{24}$ et
plusieurs valeurs de $u$.} Le facteur en $x$ de (E) est donc la solution générale de
$f'' = -p^2(\pi n/A)^2 f$, celui en $y$ de $g'' = -p^2(\pi m/A)^2 g$, et (E) est un
produit de deux fonctions trigonométriques : un mode de membrane, de son proportionnel
à $p^2(m^2+n^2) = (1-\Pi)(m^2+n^2)$. Elle écrit $(\Pi - 1)(m^2+n^2)$ ; elle corrige le
signe à la vue 476. De même (F) a les nombres d'onde $p\pi m/A$ et $q\pi m/A$
($\Pi' = 1 - q^2$), le son $m^2(p^2+q^2)$, qu'elle écrit $m^2(\Pi+\Pi'-2)$, et la
condition $\partial_x^2 z/\partial_y^2 z = (\Pi-1)/(\Pi'-1) = p^2/q^2$ qu'elle donne à
la vue 408.\footnote{La première rédaction (vue 407) porte au dénominateur de cette
condition $\Pi' - 1$ et au numérateur un groupe barré qui n'a pas été lu, et la
formule des sons y a une lettre noircie devant $\Pi+\Pi'-2$.} Une conséquence qu'elle
ne tire pas, et qui corrige ce qu'elle en dira aux vues 475 et 478 : les lignes
nodales de (E) et de (F) sont toujours des droites parallèles aux côtés, là où
s'annule l'un des deux facteurs, également espacées ; le « facteur périodique »
$\sin u$ n'en donne pas d'autres, car la série qu'il multiplie devient infinie quand
$\cos^2 u = 1$ et le produit y reste fini et non nul.

\emph{L'intégrale (G)} (N\textsuperscript{o} 5). En prenant $z = z'z''$ avec
$z' = \sin u\sin v$, $u = \pi m x/A$, $v = \pi m y/A$, elle forme
\[
z = \sin u\,\sin v\,\Big[(a'S_1 + b'S_2)_{u,\Pi}\,(A'S_1 + B'S_2)_{v,\Pi'}
+ (a''S_1 + b''S_2)_{u,\Pi''} + (A''S_1 + B''S_2)_{v,\Pi''}\Big],
\]
\[
\Pi + \Pi' = \Pi'',
\]
avec le son de (F). D'après ce qui précède, c'est une somme de trois produits de
fonctions trigonométriques, $f(x)g(y) + h(x)\sin v + \sin u\,k(y)$, dont les nombres
d'onde carrés valent $(p^2+q^2)$, $(r^2+1)$ et $(1+r^2)$ en unités $(\pi m/A)^2$, avec
$r^2 = 1 - \Pi''$ : sa condition $\Pi + \Pi' = \Pi''$ est exactement l'égalité de ces
trois valeurs propres, et (G) est bien une solution de (B). C'est, avec (C), la seule famille
du § 2 qui superpose des modes différents de même son, et qui donne par là des lignes
nodales courbes ; le N\textsuperscript{o} 15 s'en servira.\footnote{Sur la vue 412, le
groupe noyé sous une rature en boucle après « $\Pi + \Pi' =$ » n'est pas lu ; « ayent »
est écrit au-dessus d'un mot barré, lu « aient » avec doute.}

\emph{Les intégrales (H) et (I).} (H) multiplie le facteur
$\sin(\pi n x/A)\sin(\pi m y/A)$ par la somme d'une série en $x$ de paramètre
$2 - m^2$ et d'une série en $y$ de paramètre $2 - n^2$, avec le son $m^2n^2$. Elle ne satisfait
pas (B) en général : le premier terme a pour valeur propre
$(\pi/A)^2\,[n^2(1-\Pi) + m^2]$, le second $(\pi/A)^2\,[n^2 + m^2(1-\Pi')]$, et leur
égalité demande $n^2\Pi = m^2\Pi'$, ce que ses valeurs $\Pi = 2-m^2$, $\Pi' = 2-n^2$ ne
satisfont que pour $m = n$ ; le son serait alors $m^2n^2 + m^2 - n^2$, non
$m^2n^2$.\footnote{Les coefficients de (H) s'écrivent sur le feuillet « $(2-m^2+3$ »,
la parenthèse ouverte et non fermée ; le premier chiffre de ces groupes, chargé
d'encre, se lit 2 sans certitude (vue 409), et les constantes $+3$, $+8$, $+15$, $+24$
sont presque toutes repassées sur d'autres chiffres à la vue 412, où « $+24$ » est lu
sous une tache d'après les lignes voisines. La lecture n'a pas retrouvé la condition
dont ces coefficients procèdent.} Pour (I), le N\textsuperscript{o} 6 part d'une
« équation aux cordes vibrantes »
$\partial_t^2 z' = (K/n^2)\,\partial_x^2 z'$ pour
$z' = \sin(t\sqrt{1/K}+\zeta)\sin(\pi n x/A)$, relation que cette forme de $z'$ ne
donne pas (le calcul donne le coefficient $A^2/\pi^2 n^2 K$), et forme
$z = \sin(\pi n x/A)\{a'S_1 + b'S_2 + \mathfrak A\sin(\cdots y/A)\}$ avec le paramètre
$2 - m^2$ ; la surface n'est périodique que dans un sens.\footnote{L'argument du
dernier sinus se lit « $\pi\,may$ » avec doute à la vue 409 et « $\pi.mn.y$ » avec
doute à la vue 413 ; la majuscule ornée est rendue $\mathfrak A$. La formule des sons
de (I) est d'abord $\pi m^2 n^2\tau\sqrt b/A^2$, barrée à la vue 409 et remplacée par
« $= \pi.n^2(m^2+1)$ » sans le facteur $\tau\sqrt b/A^2$, puis de nouveau
$\pi m^2n^2\tau\sqrt b/A^2$, non barrée, à la vue 413.} Avec la lecture « $mn$ », le
terme en $\mathfrak A$ a le son $n^2(m^2+1)$ de la correction de la vue 409, mais la
série a le son $n^2(m^2-1)$ : les deux termes ne vibrent pas ensemble, et (I) telle
qu'elle est écrite n'est pas une solution ; elle le serait avec le paramètre $-m^2$
au lieu de $2 - m^2$. Elle n'emploiera (H) que pour la nommer et (I) que pour une
explication qu'elle dira « conjecturale » (vues 487 à 490).

Elle termine le § 2 par une remarque juste : toutes ces intégrales, sauf (C),
conviennent aussi aux surfaces tendues, et c'est ce qui les rend simples ; les
solutions propres à la surface élastique contiendraient aussi les exponentielles
$e^{\pm\omega}$, et l'angle $\omega$ ne s'obtiendrait plus que par approximation.

\subsection*{§ 3. Les limites analytiques de la lame (N\textsuperscript{os} 7 et 8)}

\pagerange{414}{432}
Le § 3, « Observations sur les propriétés analytiques des différens points des corps
sonores », pose un principe.\footnote{Vues 414 à 432. Pages (13) à (26) ;
n\textsuperscript{os} 203 à 211 de la série à l'encre. Le petit feuillet de la vue
416 (n\textsuperscript{o} 204), sur l'effet de la cire, interrompt la phrase entre la
page (14) et la page (15) ; il est lu avec le § 6 du mémoire, qu'il corrige. Le petit feuillet
de la vue 428 (n\textsuperscript{o} 210), des calculs sans un mot, calcule d'une part
les carrés $f^2$, $g^2$, $h^2$ de deux cordes consécutives d'une courbe et de la corde
qui les joint (l'angle de deux éléments par la loi des cosinus), d'autre part la
variation $\iint \Delta z\;\delta\,\partial_x^2 z\;dx\,dy$ intégrée par parties, ce qui
est exactement le calcul d'où sortent les conditions aux limites du N\textsuperscript{o}
9 ; la parenthèse de sa deuxième ligne porte $\frac{d^2z}{dx^2}$ deux fois, et deux
signes y restent illisibles, avec le coefficient 2 lu avec doute. La vue 431 répète la
vue 429.} Les points du contour d'un corps sonore satisfont, outre l'équation du
mouvement, des conditions propres aux limites ; aucun point, aucune ligne ne peut
être pris pour la limite d'une partie vibrante s'il ne satisfait ces conditions, et il
suffit qu'il les satisfasse pour qu'on puisse le considérer comme une « limite
analytique » des parties qu'il sépare. En termes actuels : ce sont les conditions
naturelles du problème variationnel, et la restriction d'un mode à une sous-région
bordée de lignes où ces conditions sont remplies est un mode de cette sous-région pour
les bords correspondants. Le principe est exact, et c'est l'idée qui porte tout le
reste du mémoire.

Pour la corde, la condition au bout se réduit à $R\,\partial_x z\,\delta z = 0$
(« formule P 332 de la 1\textsuperscript{re} édition de la Méc. Analy. ») : elle est
remplie par $z = 0$, aux nœuds, mais aussi par $\partial_x z = 0$, au milieu des
ventres. Ces derniers points sont des bouts libres au sens d'une corde dont
l'extrémité glisserait sans frottement ; elle note qu'on ne peut les réaliser sur la
corde, mais que Lagrange, dans le second volume des Mémoires de Turin, a tiré
$\partial_x z = 0$ de l'égalité de l'élasticité de l'air au bout ouvert d'un tuyau et
à l'extérieur. Tout cela est exact : le bout ouvert d'une flûte est un ventre de
déplacement, et la flûte ouverte aux deux bouts est, pour son premier mode, le
« renversement » de la corde — deux limites $\partial_x z = 0$ aux bouts et un nœud au
milieu, au lieu de deux nœuds aux bouts et d'une limite $\partial_x z = 0$ au milieu.
Pour les membranes, les conditions tirées du mémoire de Biot (« P. 87 de son mémoire
t. 4 de l'institut ») se réduisent de même à $z = 0$ ou à une dérivée normale nulle ;
toute ligne de repos y est donc une vraie limite, et c'est pourquoi l'analogie, que
suggère « le programme publié par la classe », entre les nappes vibrantes et les ondes
des cordes de Sauveur, et entre les courbes de repos et les nœuds, vaut pour les
surfaces tendues. Elle ne vaut plus pour les lames et les surfaces élastiques, où il y
a deux genres de lignes de repos : celles qui satisfont les conditions des limites, et
celles qui séparent des parties qui, isolées, ne pourraient pas prendre les courbures
qu'elles prennent dans le corps entier. C'est juste : une ligne nodale d'une plaque
porte en général un moment de flexion non nul, et la partie qu'elle borde ne vibre
alors ni comme une plaque appuyée ni comme une plaque libre sur ce bord. Elle dit la
remarque nouvelle, et la tire de la théorie d'Euler (\emph{Acta} de Pétersbourg pour
1779).\footnote{Le nom d'Euler est tracé « Lobr », son E en forme de L (vue 419) ;
« cordes », dans un groupe barré, est lu avec doute.}

Le N\textsuperscript{o} 8 prend les conditions aux bouts d'une lame données par
Lagrange dans la nouvelle édition de sa Mécanique, $K\,z'''\,\delta z = 0$ et
$K\,z''\,\delta z' = 0$ (les primes désignent les dérivées selon la longueur), et en
tire quatre couples possibles : $(z = 0,\ z' = 0)$, $(z = 0,\ z'' = 0)$,
$(z' = 0,\ z''' = 0)$, $(z'' = 0,\ z''' = 0)$.\footnote{« P. 197 » : la page de
Lagrange est chargée d'encre, 197 lu avec doute. La lettre devant les deux équations
est un $K$, lu avec doute pour la seconde. Le feuillet écrit « $\frac{d^3z}{dx^3} = 0$ »
et ne met pas « $= 0$ » après $\frac{dz}{dx}$ et $\frac{d^2z}{dx^2}$ dans l'énumération
d'Euler.} Euler n'en a examiné que trois (« P. 116 du mémoire cité ») : le bout fixe
(encastré), le bout appuyé (articulé) et le bout libre $(z'', z''')$. Elle soutient que
le quatrième couple donne lui aussi des bouts libres. Il faut le dire autrement : un
bout où $z' = 0$ et $z''' = 0$ est un bout \emph{guidé} (encastrement glissant) ;
l'effort tranchant y est nul, mais un couple est nécessaire pour y maintenir la pente
nulle, puisque le moment $z''$ n'y est pas nul, et le seul bout physiquement libre est
celui où $z'' = z''' = 0$. Il n'en reste pas moins une limite analytique, et c'est ce
qui compte pour elle : au milieu d'un mode symétrique, $z' = 0$ et $z''' = 0$ ont lieu
d'eux-mêmes, et la ligne de symétrie est exactement un bout guidé de chaque
demi-lame. Elle écrit avoir obtenu ces sons sur une lame que le jeu d'appuis rendait
propre à cette manière de vibrer, mais « la lame […] s'est brisée trop tôt pour que je
pusse constater et mesurer ce résultat ».\footnote{Le mot barré devant « constater »
est lu « mesurer » avec doute ; « parvenue » est au féminin (vue 420).}

Elle invoque à l'appui les fourches métalliques, que Chladni dit suivre la loi de la
lame appuyée à ses deux bouts bien que leurs bouts soient physiquement libres (« V.
traité de Chladni N.o 88 P. 115 ») : leurs bouts satisferaient donc au couple
$(z', z''')$.\footnote{Vue 421. « métalliques » est écrit seul dans la marge ;
« satisfont » est écrit sur « satisfait » et lu avec doute ; une addition interlinéaire
annonce un effet de la courbure « que j'expliquerai dans la suite » ; les dernières
lignes, sur la facilité de la lame droite à se ployer au second cas, sont barrées. Ce
que dit Chladni n'est pas vérifié ici.} L'inférence est cohérente : une lame guidée aux
deux bouts a les modes $\cos(k\pi u)$ et les sons en $k^2$, exactement ceux de la lame
appuyée aux deux bouts. Ce n'est pas l'explication qu'on donne aujourd'hui du
diapason, dont on traite chaque branche comme une lame encastrée à la base et libre au
bout. Les quatre cas nouveaux qu'elle renonce à calculer (vue 422) se calculent en
effet « presque sans calcul » : guidée et guidée, $\cos k\pi u$ ; guidée et appuyée,
$\cos(k+\frac12)\pi u$ ; guidée et libre, ou guidée et encastrée,
$\tan\omega + \tanh\omega = 0$, soit $\omega = 2{,}365$, $5{,}498$, $8{,}639$,
moitiés des racines des modes symétriques de la lame libre ($4{,}730$, $10{,}996$,
$17{,}279$).

Elle reprend alors la lame libre d'Euler (« l'équation P. 120 »), dans la variable
$u = s/a$ :
\[
z = e^{u\omega} \mp e^{(1-u)\omega} + (1 \pm e^{\omega})\sin u\omega + (1 \mp e^{\omega})\cos u\omega ,
\]
les signes inférieurs pour un nombre pair de points de repos, les supérieurs pour un
nombre impair.\footnote{Vues 422 à 424. Le feuillet change $y$ en $z$ et écrit « u » ;
la lecture garde $u$. Le signe du second crochet de la seconde équation de la vue 422
est effacé sous une tache. À la vue 423, dans la dernière paire d'équations, le
second crochet est écrit $(1-e^{\omega})$ devant $\sin u\omega$ là où les deux
équations dont elle est tirée portent $(1+e^{\omega})$ : c'est un lapsus. « suivant la
remarque d'Euler N.o 47 » remplace « P. 157 », barré, lu avec doute.} Ces fonctions
vérifient bien $z'' = z''' = 0$ en $u = 0$ (vérifié pour cette lecture), et les deux
bouts sont libres quand $\cos\omega\cosh\omega = 1$, dont les racines sont
$\omega = 4{,}7300$, $7{,}8532$, $10{,}9956$, $14{,}1372$, $17{,}2788$ ; la première
famille est symétrique par rapport au milieu, la seconde antisymétrique. Ses
démonstrations sont exactes. Dans un mode symétrique, aucun point ne satisfait
$z = z'' = 0$, car la somme des deux équations donnerait
$e^{u\omega} + e^{(1-u)\omega} = 0$ ; le milieu satisfait $z' = z''' = 0$, et la lame
« vibre comme deux lames séparées de moitié moins longues », libres à un bout et
guidées au milieu. Dans un mode antisymétrique, $z = z'' = 0$ exige
$e^{u\omega} = e^{(1-u)\omega}$, donc $u = \frac12$ : le milieu est un appui au sens
d'Euler et chaque moitié est une lame libre à un bout, appuyée à l'autre, dont la
condition $\tanh\frac{\omega}{2} = \tan\frac{\omega}{2}$ est celle du second cas
d'Euler pour une longueur moitié. Le second cas est donc contenu dans le
premier.\footnote{Les cas sont ceux d'Euler, dans son ordre : I, les deux bouts
libres ; II, un bout libre et l'autre appuyé ; III, un bout fixé (encastré) et l'autre
libre ; IV, les deux bouts appuyés ; V, un bout fixé et l'autre appuyé ; VI, les deux
bouts fixés. Le feuillet écrit la condition
$\frac{e^{\omega}-1}{e^{\omega}+1} = \text{tang}\,\frac12\omega$.}

Pour le troisième cas (encastré et libre), elle affirme qu'aucun point intérieur ne
satisfait l'un des quatre couples ; c'est exact pour les cinq premiers modes, vérifiés
pour cette lecture, bien que loin de l'encastrement, pour les modes élevés, les termes
hyperboliques deviennent si petits que certains points en approchent. Pour le
quatrième (deux bouts appuyés), où la lame se ploie comme la corde, tout nœud satisfait
$(z, z'')$ et tout milieu de ventre $(z', z''')$ : avec $n-1$ nœuds, cela fait $n+1$
points d'appui, bouts compris, et $n$ limites guidées, en tout $2n+1$ limites, comme
elle le dit. Le cinquième est contenu dans le sixième comme le second dans le premier ;
et dans le sixième la seule limite intérieure est le milieu, appuyé pour un nombre
impair de nœuds, guidé pour un nombre pair. Tout cela est juste ; la lecture du
reproche fait à Euler l'énonce aussi, sous le nom de six cas réduits à
quatre.\footnote{Vues 425 à 427. « remarque », barré après « la », est une lecture du
sens plus que de la lettre ; « d'extrêmités » (vue 426) est lu avec doute. Les
ordinaux « v\textsuperscript{eme} » et « vi\textsuperscript{eme} » désignent les cas
V et VI.}

Elle en conclut que les points de repos des lames ne doivent pas être assimilés aux
nœuds des cordes, et que ni Giordano Riccati, qui a calculé les angles de la lame
libre (« Delle vibrazioni de cilindri », Mémoires de la Société italienne, t. 1), ni
Chladni n'ont vu que les intervalles entre les nœuds sont inégaux. La partie comprise
entre le bout et le premier nœud est moindre que la moitié de l'intervalle suivant, et
les intervalles croissent vers le milieu, les différences étant bien plus grandes
près des bouts. Le calcul, fait une fois pour les premiers modes de la lame libre dans
la lecture de la critique d'Euler (« L'inégalité des intervalles entre les
nœuds : Chladni et Riccati », vues 354 à 356), le confirme pour le segment de
bout, qui vaut environ $0{,}36$ de l'intervalle suivant ($0{,}41$ au premier son),
et pour les premiers intervalles ; au-delà du troisième, les intervalles oscillent
autour de $\pi/\omega$ au lieu de croître jusqu'au milieu.\footnote{Vue 429. Le feuillet dit aussi que « la partie
comprise entre le second nœud et le troisième doit être moins grande que celle qui est
limitée par les deux premiers nœuds » : cette comparaison est inversée, par rapport au
calcul (le rapport du premier intervalle intérieur au suivant vaut $0{,}91$) comme par
rapport à sa propre conclusion, qu'une partie est « d'autant plus étendue qu'elle est
plus éloignée de l'extrêmité ». La fin de « confirme », barrée, est illisible. Elle
écrit « Ricati », « cilindri ».} Elle dit s'être assurée que l'expérience le confirme ;
la comparaison des valeurs de Riccati et des observations de Chladni est faite à
la même sous-section.

Les vues 429 à 432 reviennent sur l'exemple du § 47 d'Euler, la lame appuyée en un
point intermédiaire. Elle écrit que l'égalité qu'Euler établit entre les valeurs de
$z'$ sur les deux branches ne vaut que pour des bouts appuyés et ne peut subsister
quand les branches ont des courbures différentes. C'est inexact : en un appui
intermédiaire, la lame reste continue, et $z$, $z'$ et $z''$ sont continus quelles
que soient les conditions aux bouts ; seul $z'''$ saute, du montant de la réaction de
l'appui (voir « Ce que dit la théorie actuelle de la poutre sur appui
intermédiaire », vues 320 à 334). Le reste de la remarque est juste : un obstacle placé en un point de repos,
qui laisse passer le mouvement, ne fait que choisir un mode parmi ceux que la lame
possède, puisqu'en un nœud il n'exerce aucune force ; et pour des bouts libres ou
fixés, le point où Chladni tient la lame n'est pas un appui au sens d'Euler
($z = z'' = 0$), mais un simple nœud. Le mot « libre » et le mot « appuyé » sont donc,
dit-elle, bien vagues quand on les applique aux points de la lame ; la théorie de
l'appui intermédiaire est lue avec le mémoire d'Euler et le reproche qu'elle lui
fait.\footnote{Vue 432. Un mot barré après « qui », noirci, n'est pas lu. Vue 433,
début : un signe barré devant « remarqué » n'est pas lu.}

\subsection*{Les limites de la plaque : côtés, coins, plaque ronde (N\textsuperscript{o} 9)}

\pagerange{433}{438}
Le N\textsuperscript{o} 9 prend, dans la formule du N\textsuperscript{o} 1, les termes
qui ne sont affectés que d'un seul signe d'intégration.\footnote{Vues 433 à 438. Pages
(27) à (32) ; n\textsuperscript{os} 212 à 214 de la série à l'encre. À la vue 433, la
lettre après « pour » est lue $E$ avec doute, et la lettre grecque sous le carré,
$\beta$, avec doute. Son signe $S$ est une somme le long du contour ; le feuillet écrit
la première condition avec « $dx\,dy$ » sous $dz$, et les suivantes avec un facteur
$\frac{1}{dy}$.} Sur un côté parallèle à l'axe des $y$, les conditions aux limites
sont
\[
\oint \Delta z\;\delta\!\left(\frac{\partial z}{\partial x}\right) dy = 0, \qquad
\oint \frac{\partial \Delta z}{\partial x}\;\delta z\; dy = 0,
\]
ce sont exactement les conditions naturelles de l'énergie
$\frac12 b\tau^3\iint(\Delta z)^2$, dont la variation par parties est calculée sur le
petit feuillet de la vue 428. Elle en tire quatre couples pour un tel côté :
\[
(z = 0,\ z_x = 0), \quad (z = 0,\ \Delta z = 0), \quad
(z_x = 0,\ \partial_x\Delta z = 0), \quad (\Delta z = 0,\ \partial_x\Delta z = 0),
\]
et les quatre analogues en $y$, « l'analogie est frappante » avec la lame : les deux
premiers conviennent aux points de repos du contour, les deux derniers aux points
libres.

Voici, pour tout le volume, ce qu'en dit la théorie actuelle. Dans la théorie des
plaques de Kirchhoff, l'énergie de flexion est
\[
U = \frac{D}{2}\iint\Big[(\Delta z)^2 - 2(1-\nu)\big(z_{xx}z_{yy} - z_{xy}^2\big)\Big]\,dx\,dy,
\qquad D = \frac{E h^3}{12(1-\nu^2)},
\]
où $\nu$ est le coefficient de Poisson ; l'équation intérieure est la même que la
sienne. Sur un côté droit $x = \text{const}$, les conditions sont : pour un bord
encastré, $z = 0$ et $z_x = 0$ ; pour un bord appuyé, $z = 0$ et le moment nul,
$z_{xx} + \nu z_{yy} = 0$, qui se réduit à $z_{xx} = 0$, c'est-à-dire à $\Delta z = 0$,
puisque $z_{yy} = 0$ le long du côté ; pour un bord guidé, $z_x = 0$ et l'effort
tranchant nul, $z_{xxx} + (2-\nu)z_{xyy} = 0$, qui se réduit à $z_{xxx} = 0$,
c'est-à-dire à $\partial_x\Delta z = 0$, puisque $z_{xyy} = 0$ le long du côté ; pour
un bord libre enfin,
\[
z_{xx} + \nu\,z_{yy} = 0, \qquad z_{xxx} + (2-\nu)\,z_{xyy} = 0,
\]
et, en un coin où se rencontrent deux bords libres, la condition supplémentaire
$z_{xy} = 0$ (la force concentrée au coin vaut $2D(1-\nu)z_{xy}$).\footnote{Kirchhoff,
1850 ; le feuillet est antérieur et ne le cite pas.} Ses trois premiers couples
coïncident donc avec ceux de Kirchhoff quel que soit $\nu$ ; le quatrième, son bord
libre $(\Delta z = 0,\ \partial_x\Delta z = 0)$, n'est celui de Kirchhoff que pour
$\nu = 1$, valeur qu'aucune matière n'atteint (verre et laiton ont $\nu$ entre $0{,}2$
et $0{,}35$ environ). La différence des deux énergies est le terme en
$z_{xx}z_{yy} - z_{xy}^2$, un lagrangien nul qui ne touche que le bord libre ; la
lecture de l'\emph{Exposition des principes} en donne le sens géométrique (la courbure
de Gauss).

Deux de ses remarques sur ces couples demandent une correction. Elle dit que
$z = 0$ et $z_x = 0$ ne font plus un point fixe, parce qu'il faudrait encore
$z_y = 0$ pour que le plan tangent tombe dans le plan des $x$ et $y$ : c'est vrai d'un
point isolé, mais le long d'un côté entier $z = 0$ entraîne $z_y = 0$, et le couple
est bien l'encastrement.\footnote{Vue 434 : « appartiendroient » est lu avec doute. La
mise au net de la vue 596 écrit « comme fixé ».} Aux angles, qu'elle appelle
\emph{points de double limite}, elle compte quatre conditions, celles des deux côtés
qui s'y rencontrent. Ce ne sont pas des conditions nouvelles : pour une solution
régulière jusqu'au coin, elles découlent de celles des deux côtés. La condition
propre au coin, dans la théorie de Kirchhoff, est $z_{xy} = 0$ en un coin libre, et
elle manque dans la sienne parce qu'elle y est multipliée par $1 - \nu = 0$. Sa phrase
selon laquelle le couple $(z = 0,\ \Delta z = 0)$ ne peut appartenir qu'aux points de
double limite dit seulement que ce couple figure dans les deux listes. Les « points
d'origine » (vue 436), où $z$, $z_x$, $\Delta z$ et $\partial_x\Delta z$ s'annulent
ensemble, sont les points qu'on peut prendre pour origine par symétrie : par exemple
les croisements des lignes nodales de (D), comme elle le fera à la vue
472.\footnote{Vue 435 : dans le second couple, $\frac{dz}{dy}$ est écrit sans
« $=0$ » ; une petite croix barrée devant « car » n'appelle rien.}

Pour la plaque ronde, elle pose $x = r\cos\varphi$, $y = r\sin\varphi$, écrit
$z = z' + z''$ avec $z'$ et $z''$ solutions des deux équations du second ordre, de
sorte que $\Delta z$ est proportionnel à $z' - z''$ — ce qui est juste —, puis
transporte au bord $r = a$ les conditions écrites pour des côtés $x = \text{const}$ et
$y = \text{const}$. Dans ce transport, elle ne garde de $\partial_x$ que sa partie
angulaire, $-(\sin\varphi/a)\,\partial_\varphi$, et laisse tomber
$\cos\varphi\,\partial_r$ ; ses conclusions — des conditions remplies par
$\tan\varphi = 0$ ou $\cot\varphi = 0$, c'est-à-dire aux bouts des axes, et quatre
couples en $z' \pm z''$ et en leurs dérivées angulaires — en sont des artefacts. Sur
le cercle, les conditions naturelles de son énergie s'écrivent en chaque point de la
même façon, avec la dérivée normale : $(z,\ \partial_r z)$, $(z,\ \Delta z)$,
$(\partial_r z,\ \partial_r\Delta z)$, $(\Delta z,\ \partial_r\Delta z)$ ; aucun point
du bord n'y est privilégié, ce que la symétrie de la figure imposait. Chez Kirchhoff,
le moment radial au bord est
$z_{rr} + \nu\,(z_r/r + z_{\varphi\varphi}/r^2)$, qui ne se réduit pas à
$\Delta z$ ; ce qui en résulte pour le cercle appuyé est dit avec la plaque
circulaire du § 7 (vues 548 à 574).\footnote{Vues 436 et 437. Le groupe barré après « on peut » est lu
« $r=a$ » avec doute. Les points de suspension, « $=\ldots$ », sont d'elle : elle laisse
en blanc le facteur devant l'accolade. À la vue 437, dans la dernière des quatre
équations, le second terme est écrit dans l'ordre inverse de celui de la deuxième.}

Elle clôt ce paragraphe sur ce qui en fait l'usage : les lignes où ces conditions sont
remplies découpent la surface en surfaces partielles, dont les bords peuvent être
« dans un état différent de ceux de la plaque entière », sans que l'équation de la
plaque cesse de s'appliquer à chacune. C'est exact, et c'est ce que le § 4 va
exploiter.

\subsection*{§ 4. La plaque carrée libre et l'intégrale (C) (N\textsuperscript{os} 10 à 12)}

\pagerange{439}{468}
Le § 4, « Comparaison de la théorie à l'expérience, dans le cas des plaques
quarrées », part de l'intégrale (C) munie des coefficients qu'Euler a déterminés pour
la lame libre (« P. 118 et suiv. de son m. »).\footnote{Vues 439 à 468. Pages (33) à
(60) ; n\textsuperscript{os} 215 à 229 de la série à l'encre, 228 étant le fragment
sur la courbure (vue 465), étranger au mémoire et relié entre les pages (58) et (59).
Le titre du § 4 porte un mot barré illisible. Les feuillets, plus petits, sont montés
bas sur leurs feuilles à partir de la vue 443 ; la main est celle de la mise au net,
de plus en plus corrigée, avec des notes de sa main dans les marges et des pieds de
page annulés en rédactions superposées (vues 444, 452, 457, 459).} Pour la plaque de
côté $A$, l'origine en un coin,
\[
z = \mathcal{A}\,f(x) + B\,f(y), \qquad
f(s) = e^{\omega s/A} \pm e^{(A-s)\omega/A} + (1 \mp e^{\omega})\sin\frac{\omega s}{A} + (1 \pm e^{\omega})\cos\frac{\omega s}{A},
\]
les signes supérieurs pour un nombre pair de nœuds de la lame (fonction symétrique),
les inférieurs pour un nombre impair (antisymétrique). Si $\mathcal{A}$ ou $B$ est nul,
la plaque montre autant de lignes nodales droites que la lame a de nœuds. Si
$\mathcal{A} = -B$, $z$ s'annule sur la diagonale $x = y$ ; pour $f$ symétrique, aussi
sur la seconde diagonale $y = A - x$. Sa vérification est juste :
\[
f(x) - f(A-x) = 2\sin\frac{(A-2x)\omega}{2A}\left[(1+e^{\omega})\sin\frac{\omega}{2} - (1-e^{\omega})\cos\frac{\omega}{2}\right],
\]
nul en vertu de $\tan\frac{\omega}{2} = (1-e^{\omega})/(1+e^{\omega})$, qui est la
condition des modes symétriques. Pour $f$ antisymétrique, il n'y a que la diagonale
$x = y$.\footnote{Vues 439 et 440. Devant « d'abord », « supp » est barré et « fesant »
lu avec doute ; après « auroit », quelques lettres noircies ne se lisent pas ; « est »
et un mot illisible sont barrés devant « faut ». Dans (C$'$), le facteur $\mathcal{A}$
n'est pas répété devant la seconde accolade. La formule réduite de la vue 440 a le
signe global opposé à celui qu'on écrit ici, ce qui ne change rien.} Quels que soient
$\mathcal{A}$ et $B$, la figure passe par les points où $f(x) = 0$ et $f(y) = 0$ à la
fois : $n^2$ « points de repos invariables » quand la lame a $n$ nœuds, communs à
toutes les figures de même $\omega$. C'est immédiat, et exact dans (C).

\emph{Ce que (C) est pour la plaque réelle.} (C) satisfait l'équation intérieure. Sur
un côté $x = 0$, l'effort tranchant de Kirchhoff, $z_{xxx} + (2-\nu)z_{xyy}$, vaut
$\mathcal{A}f'''(0) = 0$, et la condition de coin $z_{xy} = 0$ est remplie partout ;
mais le moment vaut $z_{xx} + \nu z_{yy} = \nu B f''(y)$, qui n'est pas nul. (C) est
donc un mode exact de la plaque carrée libre pour $\nu = 0$, et pour $\nu = 0$
seulement.\footnote{Un calcul de Ritz fait pour cette lecture (fonctions de Legendre,
théorie de Kirchhoff) redonne, pour $\nu = 0$, un mode double de la plaque carrée
libre exactement à la fréquence de la lame libre de même longueur.} Avec ses propres
conditions, qui sont celles de $\nu = 1$, la condition $\Delta z = 0$ au bord n'est pas
remplie non plus, puisque $\Delta z = B f''(y)$ sur ce côté. Pour une matière réelle,
les deux fonctions $f(x)$ et $f(y)$, de même son dans (C), sont couplées par le terme
en $\nu$ : le son double se sépare en deux, la figure des deux diagonales
($\mathcal{A} = -B$) et la figure à courbe fermée ($\mathcal{A} = B$) n'ont plus le même
son. Le même calcul de Ritz donne entre elles un rapport de fréquences de $1{,}19$ pour
$\nu = 0{,}25$ et de $1{,}27$ pour $\nu = 0{,}33$, soit une tierce mineure à une tierce
majeure. Ce que (C) donne des figures est donc qualitativement juste ; ce qu'elle en
tire sur l'égalité des sons ne l'est que pour $\nu = 0$. Les combinaisons de fonctions
de la lame qu'elle forme ici sont celles dont partira Ritz pour sa théorie de la plaque
carrée libre.\footnote{Ritz, 1909 ; le feuillet est antérieur et ne le cite pas.}

Avant de passer aux cas, elle fait une remarque juste (vues 441 à 443). Si les
parties de la lame entre deux nœuds étaient égales et symétriques — l'hypothèse
« dont on n'avoit pas encore remarqué l'inexactitude » —, $f$ serait un cosinus et
$f(x) = f(y)$ donnerait des droites parallèles et perpendiculaires aux diagonales. Avec
la vraie fonction d'Euler, $|f|$ est plus grand aux bouts de la lame qu'en aucun autre
point (le rapport du bout au plus grand maximum intérieur vaut $1{,}02$ à $1{,}11$ pour
les six premiers modes), et les droites sont remplacées par des courbes qui se
détachent du bord.\footnote{Vues 441 à 443. « fictives » et « consécutives », dont les
fins sont surchargées, sont lus avec doute, ainsi que « on » (vue 442).}

Le N\textsuperscript{o} 11 décrit, valeur de $\omega$ par valeur de $\omega$, la figure
$\mathcal{A} = -B$. Pour la première valeur (deux nœuds), $f$ est monotone sur chaque
moitié, $f(x) = f(y)$ équivaut à $y = x$ ou $y = A - x$, et la figure se réduit aux deux
diagonales, qui passent par les $4$ points invariables. Pour les valeurs suivantes,
elle trace les lignes « fictives » des suppositions $\mathcal{A} = 0$ et $B = 0$ en
pointillé, marque leurs croisements $i$, et suit la courbe nodale d'intervalle en
intervalle en comparant les plus grandes valeurs de $f(x)$ et de $f(y)$ dans chacun :
une diagonale et une courbe fermée pour la deuxième valeur (trois nœuds, $9$ points),
les deux diagonales et une courbe fermée pour la troisième ($16$ points), une
diagonale et deux courbes fermées dans chaque moitié pour la quatrième ($25$ points),
les deux diagonales et quatre courbes fermées pour la cinquième ($36$
points).\footnote{Vues 443 à 456. Les figures dessinées sur les pages (vues 445, 448,
450) sont décrites par la transcription ; à la vue 450, une seconde figure est annulée
de longs traits obliques. Sa note marginale de la vue 445, « N.\textsuperscript{ta}
j'ai refait cette fig parceque j'avois laissé entre la première ligne ponctuée et le
bord de la surface la moitié de l'intervalle entre deux lignes fictives cequi ne doit
pas être », lit « refait » avec doute : c'est la même inégalité des intervalles
qu'elle avait relevée sur la lame. Le pied de la vue 444, sur la numérotation des
figures de Chladni, est barré en deux rédactions enchevêtrées. Les renvois « (V pl.
jointe a ce M. fig. 67 c) » et « (V. la pl. … fig 72 a) » portent une lettre et un
groupe douteux ou illisibles ; la planche « jointe à ce mémoire » n'est pas dans le
volume. Les vues 451 et 452 (la règle des trois cas, $<$, $=$, $>$) sont très
récrites ; « sera » et un mot commençant par « tour » y restent douteux ou illisibles.}
Que la figure passe par les $n^2$ points invariables, qu'elle présente chaque fois
comme une confirmation (« comme le veut la théorie »), est automatique : ces points
sont sur $f(x) = f(y)$ par construction. Sa règle pour le tracé — la courbe tourne sa
concavité vers les lignes fictives parallèles à l'axe des $y$ quand, dans l'intervalle,
le maximum de $f(y)$ est moindre que celui de $f(x)$, vers les autres dans le cas
contraire, et coupe l'intervalle en quatre branches en cas d'égalité — est une bonne
description qualitative des courbes de niveau $f(x) = f(y)$ ; la lecture ne l'a pas
vérifiée figure par figure.\footnote{Vues 456 et 457 : les conditions que doit remplir
la courbe fermée des cas suivants sont très raturées ; « 12 » points invariables, au
lieu de 8 pour la troisième valeur, est lu avec doute, ainsi que « lorsqu'il s'agit » ;
la fin d'une ligne barrée et un mot noirci sont illisibles.}

Le N\textsuperscript{o} 12 ajoute que, pour une même $\omega$, tout autre rapport de
$\mathcal{A}$ à $B$ donne une autre figure, faite de lignes sinueuses autour des lignes
fictives et passant par les mêmes points invariables — c'est exact dans (C), et c'est
avec la réserve de $\nu$ qu'il faut le lire. Elle reprend la \emph{flexion} de Chladni
(« V. N.o 109 de son traité d'acoustique »), en lui reprochant de faire passer l'axe
d'une ligne sinueuse par ses points les plus hauts ou les plus bas au lieu de ses
points moyens, compte une flexion pour la première valeur, une et demie pour la
deuxième, deux, deux et demie et trois pour les suivantes, et rapporte à ses valeurs
de $\omega$ les figures 64, 65, 67, 67 b, 67 c, 72 a, 72 b, 73 a, 73 b, 78 et 85 de
Chladni : la figure 64 est celle des deux diagonales, la figure 65 celle où les lignes
infléchies « se rejoignent très près du bord de la plaque pour former une courbe
rentrante », la figure 67 c celle qu'elle a décrite pour la deuxième valeur, 72 a pour
la troisième.\footnote{Vues 457 à 460. « tirer », dans sa note sur la définition de
Chladni, et « que j'ai exposée » (vue 458) sont lus avec doute ; une parenthèse barrée
sur la renumérotation de deux figures par Chladni est lue en partie avec doute, sa fin
illisible (vue 459) ; « § 3 » (vue 460) est lu avec doute. Le renvoi à « la table
p. 150 de Chladni », où il indique les sons de chaque figure, est ajouté dans
l'interligne (vue 458).} Elle note enfin que la figure « primitive » ($B = 0$) ne se
réalise que pour la deuxième valeur (figure 67), et que les figures 78 et 85 ne
s'accordent qu'au premier coup d'œil ; elle les concilie par les limites analytiques :
les parties qu'une telle ligne sépare peuvent vibrer comme si elles étaient isolées,
pourvu qu'elles rendent le même son, et Chladni, qui touche plusieurs points à la fois,
force chaque quart de la plaque à l'une ou l'autre des figures possibles (vues 460 à
463).\footnote{Vues 461 à 463. « tel », « naturel », « en eff », « dans », « -sent »,
la lettre $r$ d'une note marginale, « 85 » (peut-être 89), « lignes », « tasse » et
« 64 » sont lus avec doute ; la figure où elle a marqué les points $a\,n\,m\,y$,
$g\,k\,h\,d$, $c$ et $s$ n'est pas sur ces feuillets.} Ce dernier point demande une
précision. Pour $f$ symétrique (nombre pair de nœuds), les médianes $x = A/2$ et
$y = A/2$ satisfont exactement $z_x = 0$ et $\partial_x\Delta z = 0$ : ce sont des
bords guidés, et la plaque se partage en quatre carrés égaux. Pour $f$ antisymétrique,
$z = B f(y)$ sur la médiane, qui n'est une limite que pour la partie
$\mathcal{A}f(x)$ ; c'est ce que concède sa restriction « en n'ayant égard qu'à une
seule des variables ». La phrase de la vue 461 donne cette propriété « dans le cas d'un
nombre impair » ; pour la plaque, elle n'est exacte que lorsque $f$ est symétrique,
c'est-à-dire, dans la convention de ses vues 422 et 439, pour un nombre pair de nœuds
de la lame. La valeur de $\omega$ de son expérience, qui suit, est de ce cas.

\emph{Son expérience} (vues 463 et 464). Sur une plaque carrée qui donnait la figure 64
de Chladni, elle fait couper le quart $BCDE$ en réservant l'endroit $c$ où
s'appliquaient les doigts, replace les doigts en $c$, tire le son de l'archet en $F$ :
« Le son est resté le même que sur la plaque complette, et la figure est demeurée aussi
nette ». Elle y voit la confirmation, qu'elle n'avait pu obtenir sur la lame, des bouts
libres par $z' = 0$ et $z''' = 0$.\footnote{Vue 464. « on sait » est lu avec doute ;
une croix et quelques lettres, lues « m. » avec doute, semblent un appel sans texte ;
la « fig. A » (lue aussi « fig. 1 ») qu'elle renvoie « ci-joint » n'est pas dans le
volume ; « (V. § 3 N.o 8) » est lu sur un signe de paragraphe tracé comme un 8.} Le
raisonnement ne tient pas tel quel : couper la plaque le long des médianes fait de ces
lignes des bords libres, où moment et effort tranchant sont nuls, alors que sur la
plaque entière la médiane est une ligne guidée qui porte un moment de flexion — dans
sa théorie comme dans celle de Kirchhoff, puisque $\Delta z$ et
$z_{xx} + \nu z_{yy}$ n'y sont pas nuls. La plaque amputée n'a donc pas exactement le
mode de la plaque entière. Un calcul de Ritz fait pour cette lecture, dans la théorie de
Kirchhoff, pour $\nu = 0{,}25$ et $0{,}33$, donne pour la plaque en équerre qui reste
un mode très voisin de celui des deux diagonales — la diagonale qui traverse le quart
enlevé reste nodale par symétrie, l'autre ne l'est plus qu'à peu près —, mais un son
plus grave d'un peu plus d'un ton majeur : rapport des fréquences $0{,}884$ au plus
($213$ cents), valeur qui ne peut que baisser si l'on raffine le calcul. Si la
comparaison qu'elle rapporte est exacte, l'accord qu'elle a entendu n'est pas celui
que prédit la théorie actuelle ; ni la figure « ci-jointe » ni une mesure du son ne
sont sur les feuillets pour en juger.

Elle ajoute (vues 464 à 468) que (C) ne s'applique dans toute sa généralité qu'aux
rectangles, les seules surfaces où une variable reste constante le long de chaque côté
— c'est juste —, et que les diagonales nodales sont aussi des limites analytiques,
puisqu'elles satisfont $z = 0$ et $\Delta z = 0$. C'est exact dans (C), et c'est même la
condition d'appui de Kirchhoff pour tout $\nu$, puisque sur une droite où $z = 0$ la
dérivée seconde le long de la droite est nulle et que le moment se réduit alors à
$\Delta z$. Elle en tire des triangles rectangles isocèles dont les mouvements sont
contenus dans ceux de la plaque carrée, en changeant $\omega$ en $\frac12\omega$ selon le
rapport des côtés. Le premier — les côtés de l'angle droit appuyés, l'hypoténuse libre
au sens d'Euler, pour un nombre pair de nœuds — est l'un des quatre triangles que
découpent les deux diagonales, l'hypoténuse étant un côté de la plaque ; le deuxième —
l'hypoténuse appuyée, un côté guidé, l'autre libre au sens d'Euler — est bordé par une
diagonale, une médiane et un demi-côté ; la lecture ne retrouve pas le troisième parmi
les triangles que découpent diagonales et médianes. Ces triangles héritent de (C)
l'approximation de ses côtés libres : les bords appuyés et guidés sont exacts, les
bords libres ne le sont pas.\footnote{Vues 467 et 468. La condition de la vue 467,
« $d.\{\frac{ddz}{dx^2}+\frac{ddz}{dy^2}\}=0$ », est écrite sans dénominateur à la
première occurrence et avec $dx$ à la seconde. Le mot barré après « pour » (vue 468)
est illisible, et le mot barré « avantage » lu avec doute.}

\subsection*{Surfaces périodiques et intégrales en série (N\textsuperscript{os} 13 à 15)}

\pagerange{469}{490}
Elle salue d'abord « l'étonnante sagacité de M\textsuperscript{r} Chladni » qui,
sans connaître la théorie, avait jugé équivalentes à de simples droites parallèles des
figures en apparence si bizarres ; elle-même n'y avait vu d'abord qu'une
classification, et c'est la constance de $\omega$, sous un rapport variable de
$\mathcal{A}$ à $B$, qui explique pour elle l'uniformité des sons.\footnote{Vues 469
à 490. Pages (61) à (80) ; n\textsuperscript{os} 230 à 240 de la série à l'encre, 237
étant le petit feuillet de calculs de la vue 483, étranger au mémoire (variation d'une
intégrale double de $ds$ divisée par le carré de la somme des courbures), lu avec les
feuillets épars sur la courbure. « homage » est barré et récrit « hommage » (vue 469).
Les pages (61) à (63) de la seconde rédaction de la comparaison (vues 600 à 602)
redisent presque mot pour mot les vues 486 à 490 ; les différences sont notées plus
bas.} L'explication est juste dans (C), avec la réserve déjà faite : pour une plaque
réelle, les figures d'une même $\omega$ ne rendent pas exactement le même son.

\emph{Le N\textsuperscript{o} 13, « Surfaces périodiques ».} L'intégrale (D),
$z = \sin(\pi n x/A)\sin(\pi m y/A)$, l'origine en un coin, rend la plaque appuyée sur
tout son contour : aux quatre côtés, $z = 0$ et $\Delta z = 0$. Elle nomme \emph{ligne
d'appui} toute ligne où ces deux conditions ont lieu, rappelle que le mot « appuyé »
n'a pour elle qu'un sens analytique, et observe qu'une période complète est bordée
par quatre lignes d'appui et partagée en quarts par deux lignes où
$z_x = 0$, $\partial_x\Delta z = 0$ (ou les analogues en $y$). Tout cela est exact.
Elle en conclut que l'équation de la plaque carrée appuyée sur ses quatre côtés
« appartiendroit également, et sous les mêmes conditions, à une surface quarrée dont
les quatre côtés seroient libres », les quarts de période étant assemblés dans un cas
par leurs côtés libres, dans l'autre par leurs côtés appuyés, comme la corde et la
flûte ouverte (vues 470 et 471). C'est vrai si « libre » s'entend au sens nouveau
qu'elle a introduit : la plaque carrée à bords guidés ($z_n = 0$,
$\partial_n\Delta z = 0$) a exactement pour modes les produits de cosinus, quelle que
soit la matière. C'est faux pour une plaque physiquement libre : sur un côté guidé,
$\Delta z = -k^2 z$ et $z_{xx} + \nu z_{yy}$ ne sont pas nuls, et ni sa condition de
bord libre ni celle de Kirchhoff n'est remplie.\footnote{Vue 470 : dans la dernière
ligne de formules, $\frac{dz}{dy}$ est écrit sans « $=0$ » et la seconde condition finit
sur « $=$ » au bord du feuillet. « dans un grand nombre de cas » et « dans les cas dont
j'ai parlé » sont ajoutés dans l'interligne.} Il en va de même de ce qu'elle déduit de
(D) pour les lignes nodales (vue 471) : qu'elles sont également espacées, que la
distance de la dernière au côté « libre » est la moitié de leur intervalle, que chaque
intervalle est partagé par une limite guidée — tout cela est exact pour la plaque
guidée, et seulement approché pour une plaque libre, où, comme sur la lame libre
qu'elle a elle-même analysée, le segment de bord vaut à peu près $0{,}36$ de
l'intervalle suivant et non $0{,}5$.\footnote{Vue 471. Trois lignes barrées lettre par
lettre de petites croix annonçaient l'examen d'une intégrale dont l'étiquette ne se lit
pas ; « présenté », « l'observation », « s'applique » et « satisfera » sont lus avec
doute.} Les croisements des lignes qui terminent les périodes sont des points
d'origine (vue 472) : en un tel point, $z$, $z_x$, $\Delta z$ et $\partial_x\Delta z$
s'annulent, ce qu'on vérifie sur le produit de sinus. Et (D) n'admet pas d'autre ligne
nodale que ces droites, puisqu'un produit ne s'annule que si l'un de ses facteurs
s'annule.

Elle rapporte à (D) les figures 63, 66 a, 68 a, 71 a, 74 a, 75 a, 80 a, 81 a et 82 de
Chladni, $m$ et $n$ comptant les lignes nodales parallèles à chacun des axes, avec
des sons proportionnels à $m^2 + n^2$ :
\[
\begin{array}{lccccccccc}
\text{figure} & 63 & 66\,a & 74\,a & 68\,a & 71\,a & 75\,a & 80\,a & 81\,a & 82\\
(m,n) & (1,1) & (2,1) & (4,1) & (2,2) & (2,3) & (3,3) & (3,4) & (2,5) & (4,4)\\
m^2+n^2 & 2 & 5 & 17 & 8 & 13 & 18 & 25 & 29 & 32
\end{array}
\]
Les sommes sont justes ; les figures où $n = 1$ sont les « vibrations tournantes » de
Chladni (« V. T. d'acoustique N.o 84 … 87 et N.o 98 »), qui ne diffèrent des autres,
dit-elle, que par une seule période dans un sens, « composée de deux demies périodes
réunies dans la ligne de repos ».\footnote{Vues 472 à 475. L'étiquette de l'intégrale
est une capitale empâtée lue « D » ; « cette » (vue 473) et « duquel », « telle » (vue
475) sont lus avec doute ; la fin d'une ligne barrée de la vue 473 est illisible. Son 4
est tracé comme un « h ».} Pour la plaque libre de Kirchhoff, ces figures existent à
peu près, mais les sons ne suivent pas exactement la loi $m^2 + n^2$ : par exemple le
rapport des figures 63 et 66 a, que (D) fait égal à $5/2$, vaut $2{,}56$ pour
$\nu = 0{,}25$ et $2{,}60$ pour $\nu = 0{,}33$ (calcul de Ritz fait pour cette lecture).
La confrontation avec les sons de Chladni commence au N\textsuperscript{o} 16.

\emph{Le N\textsuperscript{o} 14 : (E) et (F).} Les constantes $a'$, $b'$, $A'$, $B'$
étant libres, elle y voit plusieurs figures pour un même son, et des sons de même forme
que ceux de (D) : $p^2(m^2+n^2)$ pour (E), $m^2(p^2+q^2)$ pour (F), en posant
$\Pi = 1 - p^2$, $\Pi' = 1 - q^2$ et en changeant le signe de $\Pi - 1$ « qui vient de
$\sqrt{(\Pi-1)^2}$ ». C'est juste, et c'est la correction du signe relevé au § 2.
Elle montre ensuite (vue 477) qu'avec $b' = B' = 0$ ou $a' = A' = 0$, « suivant que $p$
est pair ou impair », le facteur en série, après le changement de $x$ en
$\frac12 A - x$ et une multiplication par $p$ (ou par $(1-p^2)/2$), n'est autre que
l'expression de $\sin(\pi n p x/A)$ (ou de $\cos$) donnée « P. 143 des leçons sur le
calcul des fonctions in.8 an 1806 », de sorte que (E) se réduit à (D). C'est exact : la
série de $a'$ se termine quand $p$ est pair, celle de $b'$ quand $p$ est impair, et
ce sont les développements de $\sin p\theta$ et de $\cos p\theta$ en puissances de
$\sin\theta$ rappelés plus haut.\footnote{Le changement de $x$ en $\frac12 A - x$
échange sinus et cosinus quand $n = 1$ ; pour $n$ quelconque il faut changer $x$ en
$A/2n - x$. Le dernier alinéa de la vue 477 porte nettement « l'intégrale (E) », mais
les constantes $\Pi$, $\Pi'$ et le son $m^2(\Pi+\Pi'-2)$ sont ceux de (F). « fait »,
« soit » et « éd. » sont lus avec doute ; le second facteur « $16-p^2$ » se lit aussi
« $96-p^2$ ».} Mais sa conclusion (vue 478), que dans tout autre cas les séries ne se
terminent pas et qu'« il faudroit donc calculer à part chaque cas pour déterminer la
figure qui lui convient », est à corriger : terminées ou non, ces séries sont des sinus
et des cosinus, et les figures de (E) et de (F) sont toujours des réseaux de droites
parallèles aux côtés.

\emph{Le N\textsuperscript{o} 15 : (G).} Avec $a' = b' = A' = B' = 0$, $z''$ est la
somme de deux fonctions séparées de $x$ et de $y$ ; pour des côtés libres au sens
nouveau il faut encore $a'' = A'' = 0$, « car sans cette condition, elle ne
conviendroit qu'aux plaques dont les côtés seroient appuyés ». Avec
$\Pi'' = 1 - 9$ elle obtient, pour $u = \pi m x/A$ et $v = \pi m y/A$,
\[
z = \sin u\,\sin v\left\{b\cos^2 u + B\cos^2 v - \frac{b+B}{4}\right\}
= \frac{b}{4}\,\sin 3u\,\sin v + \frac{B}{4}\,\sin u\,\sin 3v ,
\]
la seconde forme, qu'elle n'écrit pas, se vérifiant par $\sin 3u = \sin u\,(4\cos^2 u - 1)$.
C'est la superposition des deux modes de même son $(3,1)$ et $(1,3)$, celle qui donne
les figures courbes les plus connues des membranes et des plaques appuyées, et son son,
$-m^2(\Pi''-2) = 10\,m^2$, vaut $10$ pour $m = 1$ et $40$ pour $m = 2$, comme elle
l'écrit.\footnote{Vue 479. « $\Pi'' = 1 - 9$ » est écrit ainsi ; la page porte « 71 »
sans parenthèses, barré. Dans la ligne barrée sur la figure 70, un mot court est
illisible. À la vue 481, le signe entre $\Pi''$ et $2$ est noyé sous une tache ; elle
ajoute en marge « dorenavant je changerai le signe des $\Pi$ », la dernière lettre lue
avec doute.} Elle y rapporte les figures 69, 70, 87 a, 87 b, 88 a et 88 b de Chladni :
69 pour $m = 1$, $b = -B$, où le facteur $\cos^2 u - \cos^2 v$ donne les deux
diagonales ; 70 pour une autre relation entre $b$ et $B$, qu'elle n'a pu mesurer. Sa
discussion de la vue 480 est exacte : une ligne nodale aboutit au côté où
$\cos v = 0$ quand $\cos^2 u = (b+B)/4b$, et au côté où $\cos u = 0$ quand
$\cos^2 v = (b+B)/4B$ ; en posant $r = B/b$, ces deux conditions sont possibles
ensemble si et seulement si $\frac13 \le r \le 3$ ou $r = -1$ ; l'une d'elles est
impossible, et il n'y a que deux lignes sinueuses, quand $b$ et $B$ sont de signes
différents sans que leur somme soit nulle, ou quand l'un des coefficients est moindre
que le tiers de l'autre, ce qu'elle dit exactement ; si $b$ ou $B$ est nul, les lignes
sont droites, sur $\cos(\pi x/A) = \frac12$, au tiers de l'intervalle, puisque
$\sin 3u = 0$.\footnote{Vue 480. Le long passage qui traite de ces cas est une addition
d'une plume plus fine, dans l'interligne, au bas de la page puis dans la marge ; un
groupe barré, « B par exemple sera $>$ que 3b », y est lu avec doute, et un mot d'une
ligne barrée est illisible. La lecture « $\frac12$ » des fractions mal formées est celle
de la transcription, d'après « d'un tiers de l'intervalle ».}

Pour la figure 87 b ($m = 2$, $b = -B$), les quatre lignes du facteur périodique et
les lignes diagonales, dont deux passent chacune par deux points d'origine, font six
diagonales, dont quatre forment un carré ; pour 87 a et 88 a, quatre lignes sinueuses
à deux flexions complètes ; pour 88 b, elle incline à la relation $b = B$, qui donne
$\cos(2\pi x/A) = \cos(2\pi y/A) = \sqrt{1/2}$ aux bouts des lignes, et le son $40$
(vues 481 à 485). Avec $\Pi'' = 1 - 25$, (G) devient
\[
z = \sin u\,\sin v\left\{b\Big(\cos^2 u - \frac{4}{3}\cos^4 u\Big) + B\Big(\cos^2 v - \frac43\cos^4 v\Big) - \frac{b+B}{12}\right\}
\]
\[
= -\frac{b}{12}\,\sin 5u\,\sin v - \frac{B}{12}\,\sin u\,\sin 5v ,
\]
les modes $(5,1)$ et $(1,5)$, de son $26$ ; pour $m = 1$ et $b = -B$, le facteur
$(\cos^2 u - \cos^2 v)\big(1 - \frac43(\cos^2 u + \cos^2 v)\big)$ donne les deux
diagonales et quatre lignes courbes qui aboutissent au bord là où
$\cos = \pm\sqrt3/2$ ; elle y rapporte la figure 79 a. Les deux identités ont été
vérifiées pour cette lecture ; ce sont, pour la plaque appuyée ou guidée, des figures
exactes, et pour la plaque libre des figures approchées.\footnote{Vues 481 à 485.
Deux passages barrés de la vue 481 sont repris aussitôt presque mot pour mot ; à la
vue 482, « particulières » est coupé au bord, sa fin passant sous le montage, et « semblement » corrigé en
« semblablement » ; la lettre de « (G) » peut se lire G ou 6 (vue 485). La
transcription relève, comme remarque de transcripteur, que les nombres $10$, $40$ et
$26$ s'accordent avec $\Pi'' = 1 - 9$ pour $m = 1$ et $m = 2$ et $\Pi'' = 1 - 25$ pour
$m = 1$.} Elle renonce à poursuivre (G) et loue encore Chladni d'avoir deviné, parmi
les distorsions des figures, celles qui tiennent aux imperfections de la plaque (vue
486).\footnote{Vue 486. Le mot barré sous « deviné » et le groupe barré après « s'est
bien gardé de » sont illisibles ; le bas de la page est une récriture en interligne,
où « senti » et « la presentoit » sont lus avec doute dans les lignes barrées.}

\emph{(H) et (I), et les figures 79 b et 86.} Pour (H), elle note que les figures
auraient des lignes droites et des sons suivant la loi de la lame appuyée (ou guidée),
et que ses séries ne se terminent jamais ; elle la laisse. (I) n'est périodique que
dans un sens, ses séries ne se terminent pas non plus ; elle croit pourtant pouvoir y
rapporter les figures 79 b et 86, avec $b' = 0$ et $m = 5$, puis $m = 6$ :
\[
z = \sin\frac{\pi x}{A}\left\{a'\cos\frac{\pi x}{A}\Big(1 - \frac{10}{3}\cos^2\frac{\pi x}{A} + \frac43\cos^4\frac{\pi x}{A} + \cdots\Big) + \mathcal{A}\sin\frac{5\pi y}{A}\right\}.
\]
Les coefficients sont bien ceux de $S_1$ avec $\Pi = 2 - 25 = -23$ (et, pour $m = 6$,
$-31/(2\cdot3)$, $31\cdot19/(2\cdot3\cdot4\cdot5)$ avec $\Pi = -34$) ; mais le premier
terme a alors pour carré du nombre d'onde $24\,(\pi/A)^2$, et le terme en
$\mathcal{A}$ en a $25\,(\pi/A)^2$ ou $26\,(\pi/A)^2$ selon qu'on le lit hors de
l'accolade ou dedans : les deux termes ne rendent pas le même son, et cette expression
n'est pas une solution de (B). Elle-même tient cette application pour « conjecturale »,
« et il faudroit avoir recours au calcul pour s'assurer qu'elle satisfait en effet à la
description exacte de ces figures » ; les sons $25$ et $36$ qu'elle assigne sont ceux
du seul terme en $\mathcal{A}$.\footnote{Vues 487 à 490. À la vue 488, « $n = 4$ » est
écrit au-dessus de la ligne, alors que l'argument est $\pi x/A$ ; la seconde rédaction
(vue 601) écrit « $n = 1$ ». « Donnés » et « Avant », barrés, et la lettre $b'$ de
« entre $b'$ et $\mathcal{A}$ » (le sens voudrait $a'$, que donne la vue 601) sont lus
avec doute ; un signe du second facteur est noyé sous une tache. À la vue 489,
« $n' = 1$ » est lu avec doute et la première valeur de $z$, barrée, a un numérateur
illisible. À la vue 487, (I) est dite périodique « dans le sens des $y$ », alors que
les vues 409 et 413 disaient « dans le sens des $x$ ». Dans la seconde rédaction, les
passages sur (H), étiquetée « (++) », et sur les figures 79 b et 86 sont traversés de
longs traits obliques qui paraissent les annuler (vues 600 à 602), et une note de la
marge de la vue 600 ramène les sons de ces figures à ceux des intégrales
périodiques.}

Elle conclut la description des figures en observant qu'elle a rendu compte de plus
de la moitié de celles que Chladni a observées sur les plaques carrées, que les autres
ne sont « ni plus bizares ni plus compliquées », et qu'elle suivra pour leurs sons
les équivalences de Chladni lui-même ; la comparaison des sons, qui suit, est lue
plus bas, dans la section « Le mémoire sur la question de la première classe
(suite) : la théorie devant l'expérience » (vues 491 à 592).\footnote{Vue 490, page « 80 » sans
parenthèses. La note marginale, appelée par une croix, ajoute la ressource
d'« employer à la fois ces diverses intégrales comme je le dirai en parlant des plaques
rectangulaires », son premier mot, « soit », paraissant barré.}

\section*{Les autres états de l'ouverture}

\pagerange{593}{626}
L'ouverture du mémoire existe dans le volume en trois états, auxquels s'ajoute la
suite du troisième, et un brouillon plus tardif sur les surfaces cylindriques.
Comparés ligne à ligne, les trois états se rangent : le troisième, dont la suite est
aux vues 614 à 618, porte en clair le texte que les deux autres portent barré ; le
brouillon des vues 398 et 399 et la mise au net des vues 593 à 595 en donnent la
rédaction corrigée. Cet ordre est une inférence de cette lecture, tirée des seuls
textes.\footnote{Comme dans la section précédente, le feuillet écrit les dérivées avec
un $d$ droit ; on écrit $\partial$.}

\subsection*{Une seconde mise au net : pages (1) et (2), (23) et (24)}

\pagerange{593}{597}
Une écriture plus grande et très régulière copie le titre, la question, l'épigraphe
et le début du mémoire.\footnote{Vues 593 à 597. N\textsuperscript{o} 292 de la série
à l'encre pour le feuillet des vues 593 et 595 (la vue 594 le photographie une seconde
fois), 293 pour celui des vues 596 et 597 ; pagination (2) à la vue 595, (23) et (24)
aux vues 596 et 597, aucune à la vue 593. Le titre est d'une main calligraphiée ; la
référence de l'épigraphe et le titre du paragraphe sont soulignés.} Le texte gardé est
celui du brouillon de la vue 398 une fois corrigé, presque mot pour mot : « examiner »,
« une équation différentielle dont les intégrales particulières sont facilement
applicables », « d'après cette remarque », « la théorie que je soumets », « puisqu'eux
seuls présentent des lois déterminables », « le petit nombre de mouvemens que cet
habile physicien a observé dans les surfaces courbes ; car, comme je le dirai dans la
suite, ces mouvemens se trouvent compris dans une équation fort simple qui convient à
plusieurs genres de corps sonores ». L'épigraphe y est corrigée (« iisdem »,
« impossibilis ») et la référence complétée (« edit. in f.o London 1740 »). Mais cette
mise au net a d'abord été écrite sur l'ancien texte : entre ses lignes gardées courent
des lignes plus anciennes, hachurées de petites croix, qui disent que Chladni « n'a
fait qu'un petit nombre d'essais sur le son des surfaces courbes ; et d'ailleurs les
géomètres sentiront aisément que les simplifications auxquelles donne lieu la
supposition $z$ très petit pourroient se transporter aux surfaces d'une courbure
quelconque en substituant aux différentielles des deux ordonnées rectangulaires $x$ et
$y$ les différentielles des deux courbes qui mesurent les courbures principales de ces
surfaces ; desorte que les valeurs de $z$ exprimeroient encore, comme dans le cas du
plan, la quantité dont chaque point des mêmes surfaces s'éleveroit pendant leur
mouvement au dessus de leur situation naturelle ».\footnote{Vues 593 et 595. « nombre »
et « d'essais » ne se lisent que sous les hachures, d'où le doute. Le mot barré après
« au dessus de » est « leur », avec des lettres récrites par-dessus, peut-être
« sa ».} Ce que la mise au net abandonne, c'est donc un programme : étendre la théorie
aux surfaces courbes en prenant pour coordonnées les lignes de courbure. Elle le
remplace par une affirmation plus modeste, que les mouvements observés sur des surfaces
courbes relèvent d'une équation simple commune à plusieurs corps sonores. Le programme
abandonné n'aurait pas suffi : dans la théorie actuelle des coques, la courbure ne
change pas seulement les coordonnées, elle couple la flexion à l'extension de la
surface moyenne, et c'est ce couplage qui domine le son des surfaces courbes (voir plus
bas, les surfaces cylindriques).

Le paragraphe premier suit, sans numéro (« §. Recherche de l'équation différentielle
de la surface élastique vibrante »), et son N\textsuperscript{o} 1 redit celui de la
vue 399, avec la correction « Ainsi je trouve » déjà faite et le double signe
d'intégration écrit $SS$ ; il omet, dans la définition des coefficients d'un plan,
« sur les plans coordonnés verticaux », et s'arrête sur « le cosinus de l'angle ».

Les pages (23) et (24) (vues 596 et 597) redisent les pages (28) et (29) du mémoire
(vues 433 à 435) : la plaque rectangulaire aux côtés parallèles aux axes, les quatre
couples d'équations de chaque côté, les points de double limite. Le texte est le même,
à trois écarts près : la mise au net écrit $\delta x$ au lieu de $\delta z$ dans la
seconde condition des côtés parallèles à l'axe des $y$, ce qui est un lapsus (la même
condition porte $\delta z$ à la vue 597) ; elle écrit « comme fixé » pour « comme
fixe » ; et elle laisse sans « $= 0$ » deux conditions que la vue 434 complète en
partie. Qu'elle arrive à la page (23) là où le mémoire en est à la page (28), dans une
écriture plus grande, montre qu'elle a laissé de côté une partie de ce qui précède ;
les feuillets ne disent pas laquelle.\footnote{Vues 596 et 597. « en repos » est écrit
d'une encre plus lourde par-dessus d'autres lettres ; le « 0 » de « $dy = 0$ » pourrait
être lu « a ». La vue 597 s'arrête sur « on ait $dy = 0$, » au bas du feuillet.}

\subsection*{Un troisième état}

\pagerange{598}{598}
Un brouillon sans pagination, très corrigé, reprend l'ouverture sous le titre
« Mémoire sur cette question proposée par la première classe de l'institut », sans
épigraphe.\footnote{Vue 598, n\textsuperscript{o} 294 de la série à l'encre. « fort »
est taché d'un pâté et lu avec doute, peut-être biffé ; le mot qui suit « de l'autre »,
dans une addition, se lit « habile » avec doute ; les deux lignes « la théorie …
question. » sont hachurées de petites croix, « je crois » lu sous les hachures. La vue
599, dos blanc, n'est pas transcrite.} Il écrit « embrasser d'abord cette question dans
toute sa généralité », annonce « une équation différentielle fort simple », après avoir
barré « du quatrième degré » et une phrase commencée, « à la vérité cette équation
n'est pas intégrable en » ; il espère que la classe trouvera « la théorie que j'ai
l'honneur de soumettre à son jugement » aussi générale que le demande l'état de la
question ; il définit les mouvements réguliers comme ceux où « des sons perceptibles
accompagnent des figures constantes » ; et il s'arrête, aux deux tiers du feuillet, sur
« les géomètres sentiront aisément que la simplification à laquelle donne lieu la
supposition de $z$ très petit pourroit se transporter aux surfaces d'une courbure
quelconque en substituant aux. » Ce sont, mot pour mot, les leçons que le brouillon
de la vue 398 porte barrées sous ses corrections (« embrasser », « cette », « fort
simple », « j'ai l'honneur de », « surfaces d'une courbure quelconque en substituant »)
et que la mise au net des vues 593 à 595 porte hachurées. Ce troisième état est donc,
selon toute apparence, le plus ancien des trois : c'est l'inférence de cette lecture,
que la reliure, qui le place après la mise au net, ne contredit ni ne confirme.

\subsection*{Un brouillon adressé à « la classe »}

\pagerange{614}{618}
Trois feuillets d'un autre papier, d'une main rapide et corrigée, sans
pagination, s'adressent à « la classe ».\footnote{Vues 614, 616 et 618 ;
n\textsuperscript{os} 301, 302 et 303 de la série à l'encre. Les vues 615, 617 et 619
sont leurs dos blancs. Ils sont reliés après la seconde rédaction de la comparaison
des sons (vues 600 à 613), dont ils ne sont pas.} Le premier commence au milieu d'une
phrase : « […] differentielles des deux ordonnées rectangulaires $x$ et $y$ les
differentielles des deux courbures qui mesurent les courbures principales de la
surface desorte que les valeurs de $z$ [exprime]roient, comme dans le cas du plan, la
quantité dont chaque point de cette surface s'éleveroit pendant le mouvement, audessus
de la situation naturelle ». C'est exactement la fin de la phrase que la vue 598 laisse
sur « en substituant aux. », et que la mise au net des vues 593 et 595 porte hachurée
dans une rédaction presque identique : la vue 614 continue la vue 598. La transcription
ne fait pas ce rapprochement ; il est de cette lecture, et il repose sur le texte
seul.\footnote{Vue 614. Le premier mot, barré, n'a pas été lu ; « courbures », dont
la fin est surchargée, est lu avec doute ; « qui mesurent les courbures principales »
est écrit au-dessus de « pareillement rectangulaires », barré, et « roient »
au-dessus de « n'exprimeroit encore », barré, sans que le verbe corrigé soit
récrit. « comme on valevoir » (va le voir) et « delà » sont lus avec doute ; un mot
barré après « des mêmes cas » est illisible.}

La suite est un résumé anticipé du mémoire. L'équation des surfaces élastiques est du
quatrième degré, n'a pas d'intégrale générale en termes finis, mais s'abaisse au second
degré, et l'une des deux équations est celle des surfaces tendues : toute figure
d'une surface tendue renfermée dans des limites données peut appartenir à une surface
élastique renfermée dans les mêmes limites — c'est, presque dans les mêmes mots, le
N\textsuperscript{o} 2 du mémoire (vue 404), dont la lecture a dit plus haut dans quelle
mesure il est exact. Les géomètres, poursuit la vue 616, se sont peu occupés de ces
intégrales, et Biot est le seul, croit-elle, à avoir démontré l'existence des lignes
nodales parallèles aux côtés des membranes rectangulaires ; elle montrera que c'est une
suite de la périodicité, et comparera ses intégrales à l'expérience. Puis vient la
« singularité remarquable » : les cas où les surfaces élastiques prennent des figures
qui pourraient appartenir à des surfaces tendues sont les seuls où les nappes
vibrantes répondent aux ondes des cordes de Sauveur, et les courbes de repos à leurs
nœuds, analogie « indiquée par les termes mêmes du programme publié par la classe » —
c'est la remarque du N\textsuperscript{o} 7 du mémoire (vues 418 et 419). La vue 618
annonce le plan (l'équation, puis des intégrales particulières comparées à
l'expérience) et s'arrête : « je m'estimerai heureux si la classe juge que j'ai reussi
a ebaucher une théorie que des ».\footnote{Vues 616 et 618. « ni en general », « ce
savant » et « d'après les le plus » sont lus avec doute sous les ratures ; « et encore
moins » est écrit au-dessus d'un groupe barré, « moins » lui-même traversé d'un trait
léger. « aucun essais », « présenteroit », « correspondroit » sont écrits ainsi.} Dans
cet état, l'ouverture annonçait donc les résultats du § 1 et du § 3 ; dans la rédaction
des vues 398 et 593, ils ont pris place dans le corps du mémoire, et le programme des
surfaces courbes a disparu.

Les participes de l'auteur y sont au masculin, « parvenu » et « heureux » (vue 618),
comme « heureux » dans la phrase finale du mémoire (vue 592) ; le mémoire lui-même dit
« je crois être parvenue » (vue 420), et le brouillon suivant « privée » et « résolue »
(vue 620). Les feuillets le montrent ; ils ne disent pas pourquoi.

\subsection*{Les surfaces cylindriques, et « mon mémoire pour 1813 »}

\pagerange{620}{626}
Quatre pages paginées par elle, 1 (annulé) à 4, d'une main rapide et très corrigée,
reprennent le sujet pour des surfaces cylindriques.\footnote{Vues 620, 622, 624 et
626 ; n\textsuperscript{os} 304 à 307 de la série à l'encre. À la vue 620, à côté du
304, un petit chiffre noirci, peut-être 1, barré et souligné ; puis 2, 3 et 4 soulignés
en tête des vues 622 à 626, le 4 barré. Les dos blancs (621, 623, 625, 627) ne sont
pas transcrits.} Le plus simple des cas « après ceux dont je viens de parler » est la
surface cylindrique à base circulaire de rayon $R$, pour laquelle son équation (e)
devient
\[
\kappa b^{2}\left\{\Delta^2 z'' \pm \frac{1}{R^2}\,\Delta z''\right\} = z'' ,
\]
et, pour une portion d'arc $A$ qui répond à l'arc $a$ du cercle de rayon $1$,
$R = A/a$ et $1/R^2 = a^2/A^2$, ce qui est juste. L'équation (e) n'est pas dans le
volume : le long mémoire désigne ses équations par des majuscules, de (A) à (I).
« Le genre de surface auquel appartiendroit » cette équation « n'a encore, je crois,
été l'objet d'aucune expérience ».\footnote{Vue 620. Le coefficient se lit $\kappa$
(ou un $k$ à boucle) et $b$ avec un 2 en exposant, sous lequel un second petit signe
n'a pas été lu, dans les deux équations. Le signe devant $\frac{1}{R^2}$ est noirci
d'une tache dans la première, $+$ dans la seconde. Le feuillet écrit les dérivées
$\frac{d^4z''}{dx^4}$, $\frac{d^2z''}{dx^2}$ ; le laplacien est de cette lecture.}
Suit une confession : « Je me suis résolue beaucoup trop tard à reprendre le sujet qui
fait l'objet du présent mémoire », le temps lui manquant à la fois pour faire exécuter
toutes les expériences et pour chercher les intégrales qui en rendraient raison ;
« je me trouvé privée ici de tout secours étranger ».\footnote{« résolue » est lu avec
doute, serré après « avisé », barré ; deux lignes barrées commençaient « Le tems m'a
manqué pour effectuer toutes celles que j'avois projettés ».}

Les expériences (vues 622 à 626) portent sur des plaques de verre carrées, taillées
puis cambrées selon diverses courbures, essayées avant et après. Elle voit « avec
étonnement » que la première figure de Chladni, deux droites qui se croisent au centre,
se forme de la même manière sur la surface cylindrique et y rend le même ton ; sur une
quarantaine de surfaces, sept ou huit fois un demi-ton d'écart, la surface courbée plus
aiguë, qu'elle attribue à la haute température nécessaire pour ployer le verre, « le
peu d'activité des ateliers » l'ayant empêchée de varier les essais. Sur une quarantaine
de plaques carrées planes, l'intervalle entre les tons des figures de deux droites
perpendiculaires et de trois droites dont l'une est perpendiculaire aux deux autres est
le plus constant ; sur les surfaces cylindriques il varie avec la courbure. Sur les
plaques carrées, un nombre donné de lignes nodales est toujours accompagné du même ton,
principe qui a guidé la classification de Chladni « et la théorie a pleinement
confirmé, comme je l'ai dit ailleurs » ; sur les surfaces cylindriques, des figures
d'un même nombre de lignes, parallèles les unes aux côtés courbes, les autres aux côtés
droits, rendent des sons différents. Enfin, la figure des deux diagonales, si facile
sur les plaques carrées, « suivant la théorie que j'ai developpée N.o 11 de mon memoire
pour 1813 », est équivalente à celle de deux droites parallèles aux côtés, que Chladni
n'a jamais vue se former sur les plaques carrées.\footnote{Vues 622 à 626. Plusieurs
paragraphes de la vue 622 sont annulés de grands traits obliques, dont celui qui
décrit des portions cylindriques d'arc $A$ ; « ployés » et « cambrer » sont lus avec
doute. À la vue 626, trois lignes barrées mot à mot disaient que la figure des deux
diagonales « ne peut se former sur les surface cylindriques il suffit d'une courbure a
peine sensible pour que cette figure ne paroisse plus ». La note marginale de la vue
624, appelée semble-t-il par un astérisque, est celle sur la variation de l'intervalle
avec la courbure.}

Ces observations sont celles que prévoit la théorie actuelle des coques minces, et
l'équation (e) n'en contient pas la cause. Sur une surface cylindrique peu profonde,
de génératrices parallèles à $x$, une flexion qui ne courbe pas les génératrices
($z_{xx} = 0$) ne change pas la courbure de Gauss au premier ordre et se fait sans
étendre la surface moyenne ; une flexion qui les courbe l'étend, et la rigidité
d'extension, bien plus grande que celle de flexion, en élève fortement le son. La
torsion $z \approx xy$ de la première figure est du premier genre : son ton ne dépend
presque pas de la courbure, comme elle l'a vu. Les figures de lignes parallèles aux
côtés droits (flexion $z = g(y)$, du premier genre) et aux côtés courbes (flexion
$z = f(x)$, du second) rendent des sons différents, comme elle l'a vu. Et la figure des
diagonales, qui demande que $f(x)$ et $f(y)$ vibrent ensemble au même son, disparaît dès
que la courbure les sépare — ce que disaient ses lignes barrées de la vue 626. La
variation d'extension que l'équation (e) ignore est ce qu'on appelle aujourd'hui la
distinction entre flexion inextensionnelle et flexion avec extension.\footnote{Rayleigh
a étudié les vibrations inextensionnelles des surfaces courbes dans les années 1880 ;
le feuillet est antérieur et ne le cite pas. Sur les plaques planes, l'intervalle
qu'elle trouve le plus constant ne dépend, dans la théorie de Kirchhoff, que de $\nu$,
comme tous les rapports de sons d'une plaque carrée libre ; il vaut $2{,}56$ pour
$\nu = 0{,}25$ (calcul de Ritz fait pour cette lecture).} Quant à la remarque de
Chladni, qu'on ne voit jamais deux droites parallèles aux côtés sur une plaque carrée,
la théorie actuelle l'explique par le même couplage en $\nu$ qui sépare, sur la plaque
plane, les figures $\mathcal{A} = -B$ et $\mathcal{A} = B$ : le mode pur $f(x)$ n'existe
que pour $\nu = 0$.

Le « N.o 11 de mon memoire pour 1813 » est, dans le long mémoire des vues 398 à 592,
exactement la première valeur de $\omega$ du N\textsuperscript{o} 11 (vues 443 et 444),
où la supposition $\mathcal{A} = -B$ donne les deux diagonales, de même son que la
figure $B = 0$ de deux droites parallèles à un côté. Cette coïncidence fait de ce long
mémoire, selon toute apparence, celui qu'elle appelle son mémoire pour 1813, et de ces
quatre pages une reprise postérieure du sujet. Un brouillon d'introduction du volume
dit d'ailleurs que c'est « dans le mémoire couronné qui fut présenté en 1815 qu'on
trouve les premiers essais sur la théorie des surfaces courbes », et qu'« à raison de
sa plus grande simplicité » on y avait traité « le cas des surfaces cylindriques »
(vue 103) : ces quatre pages en seraient un brouillon. Ce sont deux inférences de cette
lecture ; aucun feuillet ne date le long mémoire, et ces pages ne se nomment
pas.\footnote{Contexte, non tiré des feuillets : le prix sur les surfaces
élastiques a été mis au concours à plusieurs reprises, avec des échéances en 1811, 1813
et 1815. Ce que les feuillets du volume disent de cette histoire, c'est l'Avertissement
de l'\emph{Exposition des principes} (vues 3 à 9) qui le rapporte, et les brouillons de l'introduction du mémoire sur la courbure (vues 93 à 114),
d'où vient la citation de la vue 103.}

\section*{Le mémoire sur la question de la première classe (suite) : la théorie devant l'expérience}

\pagerange{491}{592}
Cette section reprend le long mémoire sur les surfaces élastiques là où la
précédente l'a laissé, à la fin de son N\textsuperscript{o} 15 : les figures
nodales de la plaque carrée sont expliquées par les intégrales (C) à (I), et il
reste à comparer les sons.\footnote{Vues 491 à 592 : ses pages (81) à (165), la
fin de la copie de travail mise au net et corrigée qui commence à la vue 398
(page (1)) ; écrites recto et verso, le numéro de la série à l'encre (241 à 291)
sur un côté de chaque feuillet. Sont reliés parmi ces pages et ne font pas partie
du texte : un brouillon sur la chaleur (vue 549, traité avec la chaleur), un
feuillet blanc (vues 551--552), les petits feuillets de calcul des vues 553, 555,
561 et 563--564, et la note de la vue 575 (traitée à part ci-dessous). Les vues
550, 554, 556, 558, 562 et 576 sont des dos blancs, la vue 559 une seconde prise
de la vue 557. Le feuillet écrit $d^4z/dx^4$, $ddz/dx^2$ ; on écrit
$\partial^4 z/\partial x^4$, $\partial^2 z/\partial x^2$, ou $z_{xx}$ dans le
commentaire. Que ce mémoire soit son « mémoire pour 1813 », jugé en janvier 1814,
est une inférence de cette lecture (voir la fin de la section « Les autres
états de l'ouverture », vues 620 à 626) ; elle
trouve ici un appui : l'Exposition des principes (vue 38) renvoie pour
l'expérience de la cire au « M. pour janvier 1814 », et le § 6 ci-dessous
contient cette expérience, jusqu'au rapport de dix à un (vue 546).}
Elle a trois parties : la plaque carrée (N\textsuperscript{o} 16), les plaques
rectangulaires (§ 5), l'effet d'une charge de cire (§ 6) ; puis la plaque
circulaire (§ 7), qui la mène à l'anneau et à ses propres expériences sur des
lames de verre courbes.

Deux conventions servent partout. Le « nombre auquel le son est proportionnel »
est, à un facteur près, la fréquence : pour une surface
$z = \sin(m\pi x/a)\,\sin(n\pi y/b)$, l'équation
$\partial_t^2 z + b\tau^2\,\Delta^2 z = 0$ donne la pulsation
\[
\omega_{mn} = \pi^2\tau\sqrt{b}\,\Bigl(\frac{m^2}{a^2}+\frac{n^2}{b^2}\Bigr),
\]
et sur le carré de côté $A$ ses nombres sont $m^2+n^2$.\footnote{Le $b$ de
$b\tau^2$ est sa « rigidité naturelle », à ne pas confondre avec le côté $b$ de
la formule ; le feuillet écrit les sons $\sqrt{1/K}\,/\pi = \pi(m^2+n^2)\tau\sqrt{b}/A^2$
(vue 406, intégrale (D)).} Les intervalles sont des rapports de fréquences ; on
les mesure aussi en cents (l'octave vaut 1200 cents, le demi-ton tempéré
100).\footnote{Le cent est une mesure moderne, étrangère aux feuillets. Pour
les noms d'intervalles qu'elle emploie : ton majeur $9/8$ (204 cents), ton
mineur $10/9$ (182), demi-ton majeur $16/15$ (112), demi-ton mineur $25/24$
(71), comma syntonique $81/80$ (21,5) ; son « comma » est $128/125$ (41
cents), qu'on nomme aujourd'hui petite diésis.} Les notes de Chladni sont
nommées comme sur le feuillet, le chiffre donnant l'octave (sol 1,
si$^{\flat}$3) ; pour les confronter à ses rapports, on les compte en
demi-tons tempérés, ce qui néglige les qualificatifs « + aigu », « $-$ grave »
de la table de Chladni. La notation $m|n$ de Chladni désigne une figure de $m$
lignes nodales droites parallèles à un côté et de $n$ parallèles à
l'autre.\footnote{Vue 497 : « $m/n$ signifie […] qu'il devra y avoir $m$
lignes droites nodales parallelles à un des axes et $n$ de pareilles lignes
parallelles à l'autre axe ». Le feuillet écrit $m/n$ ou $m|n$ ; on écrit
$m|n$.}

\subsection*{N\textsuperscript{o} 16. La comparaison des sons de la plaque carrée}

\pagerange{491}{507}
Elle commence par fixer la tolérance. Chladni dit lui-même que dans des
expériences bien faites l'erreur n'excède jamais un demi-ton ; deux sons
mesurés peuvent donc s'écarter d'un ton au plus du rapport théorique, et cet
écart suppose que les deux erreurs soient maximales et de sens
contraires.\footnote{Vue 491, sa page (81), n\textsuperscript{o} 241 de la série
à l'encre : fin de la phrase de la vue 490, puis « N.o 16 \emph{comparaison
des sons} ». La citation est entre ses guillemets, avec le renvoi « p 154 de
son livre », ajouté d'une plume fine (le « p » peut être un §) ; un renvoi
« (V § 54.) » qui la terminait est barré.} Le raisonnement est juste : si
chaque son est connu à un demi-ton près, leur rapport l'est à un ton près.

\emph{Les vibrations « tournantes ».} Ce sont les figures qui ne contiennent,
dans le sens d'un des axes, que deux demi-périodes : une seule ligne nodale
de ce côté, soit les figures $1|n$.\footnote{Vue 492, sa page (82). Le mot
barré devant « demies periodes » est lu « deux » avec doute. Les renvois
« N.o 13 » et « N.o 15 » sont aux paragraphes de la première partie du mémoire,
d'où viennent les nombres.} Ses nombres sont $2, 5, 10, 17, 26$ pour les
figures 63, 66\,a, 69--70, 74 et 79\,a de Chladni, c'est-à-dire $1+n^2$ pour
$n = 1, \dots, 5$ ; et $25$, $36$ pour 79\,b et 86, qu'elle rapporte, par
conjecture, à l'intégrale (I). Rapportés à la figure 63 et comparés aux notes
de la table de Chladni (p. 152), ils donnent :
\[
\begin{array}{lcccl}
\text{fig.} & \text{nombre} & \text{rapport à 63} & \text{cents} & \text{Chladni (écart)}\\
63 & 2 & 1 & 0 & \text{sol 1}\\
66\,a & 5 & 5/2 & 1586 & \text{si 2}\ (+14)\\
69 & 10 & 5 & 2786 & \text{si 3}\ (+14)\\
70 & 10 & 5 & 2786 & \text{ut}^{\sharp}4\ (+214)\\
74 & 17 & 17/2 & 3705 & \text{si}^{\flat}4\ (+195)\\
79\,a & 26 & 13 & 4441 & \text{fa}^{\sharp}5\ (+259)\\
79\,b & 25 & 25/2 & 4373 & \text{fa}^{\sharp}5\ (+327)\\
86 & 36 & 18 & 5004 & \text{ut}\,6\ (+296)
\end{array}
\]
(l'écart, en cents, est celui de la note de Chladni au rapport théorique ;
positif, le son observé est plus aigu). Ses conclusions sont exactes. De 63 à
66\,a, une octave, un ton majeur et un ton mineur font $2\cdot\frac98\cdot\frac{10}{9}
= \frac52$ : l'accord est parfait, comme de 63 à 69 ($5$). La figure 70, que
la théorie met à l'unisson de 69, en diffère d'un ton, ce qu'elle déclare ne pas
comprendre, en relevant que Chladni les tient pourtant pour
équivalentes.\footnote{Vue 493, sa page (83), n\textsuperscript{o} 242. Le
renvoi « (V. N.o 15) » est lu avec doute pour le 1.} De 63 à 74, l'écart
($75/68$ entre ses bornes, 170 cents) est entre le demi-ton et le ton, dans la
tolérance ; de 63 à 79\,a il est d'environ un ton et quart ($15/13$, 248 cents) :
le son observé est « trop aigu d'environ 1 ton », hors de la tolérance ;
79\,b et 86 s'en éloignent davantage.\footnote{Vues 493--494, sa page (84). Pour
79\,b elle ajoute que l'écart des nombres $26/25$ (68 cents) est « beaucoup
moindre qu'un demi ton » : c'est vrai du demi-ton majeur (112 cents), non du
demi-ton mineur (71 cents), qu'il égale presque. Pour 86 elle trouve
$\frac{64}{3}$ contre $18$ et dit le son « plus aigu d'un ton mineur, et
$\frac12$ demi ton majeur que ne le veut la théorie » ; le rapport
$\frac{64/3}{18} = \frac{32}{27} = \frac{10}{9}\cdot\frac{16}{15}$ est
exactement un ton mineur et un demi-ton majeur entier (294 cents). La seconde
rédaction (vue 606) écrit « d'un ton et $\frac12$ ton majeur ».} Elle rappelle
que ses explications de 79\,b et 86 ne sont qu'« hypothetique[s] ».

\emph{Les figures à périodes complètes.} Rapportées à la figure 68\,a ($2|2$,
nombre $8$), les figures 71, 75, 80\,a, 81\,a, 82, 87--88 et 79 ont les nombres
$13, 18, 25, 29, 32, 40, 26$.\footnote{Vues 495--497, ses pages (85), 86 et (87),
n\textsuperscript{os} 243 et 244. Dans « $\frac{25}{8}$ » (vue 495) le 5 est
surchargé d'un chiffre qui le fait lire 6.} Par rapport à la note de 68
(si$^{\flat}$3), les écarts à la table de Chladni sont faibles :
\[
\begin{array}{lccccc}
\text{fig.} & 71 & 75 & 80 & 81 & 82\\
\text{rapport} & 13/8 & 9/4 & 25/8 & 29/8 & 4\\
\text{Chladni} & \text{fa}^{\sharp}4 & \text{ut}\,5 & \text{fa}^{\sharp}5 & \text{sol}^{\sharp}5 & \text{si}^{\flat}5\\
\text{écart} & -41 & -4 & +27 & -30 & 0
\end{array}
\qquad
\begin{array}{lcc}
\text{fig.} & 87\text{--}88 & 79\\
\text{rapport} & 5 & 13/4\\
\text{Chladni} & \text{ré}\,6 & \text{fa}^{\sharp}5\\
\text{écart} & +14 & -41
\end{array}
\]
Elle vérifie chaque cas par un rapport juste, et ses vérifications tiennent :
$16/5$ contre $25/8$ pour 80, écart $128/125$, son « comma » ; $16/5$ contre
$13/4$ pour 79, écart $65/64$, moindre que ce comma ; pour 87--88, deux octaves et
une tierce majeure donnent $5$.\footnote{Vue 496 : elle écrit pour 87--88
« d'une octave et de deux tons », ce qui ferait $2\cdot\frac{81}{64}$ et non
$5$ ; de si$^{\flat}$3 à ré 6 il y a deux octaves et une tierce majeure, et la
seconde rédaction (vue 607) écrit « de deux octaves et de deux tons ». Pour 81,
elle dit que $\frac{29}{8}$ surpasse $\frac{32}{9}$ « de plus d'un comma » :
l'écart est $\frac{261}{256}$ (33 cents), plus grand que le comma syntonique
$81/80$, plus petit que son comma $128/125$ des vues 497 et 499. Vue 497 : « le
comma $\frac{128}{125}$ », en toutes lettres (la seconde rédaction, vue 607,
de même).}

\emph{Les cas sans figure.} Pour les cas dont Chladni n'a pas donné de figure,
dans sa notation $m|n$, les nombres sont $m^2+n^2$ : $4|2$, $5|3$, $6|3$, $5|4$,
$6|4$, $5|5$, $6|5$ donnent $20, 34, 45, 41, 52, 50, 61$, et, rapportés à $2|2$
(8), des écarts à la table de $+14$, $-5$, $+10$, $+71$, $+59$, $+27$ et $+83$
cents. Elle choisit, quand Chladni donne deux notes, la plus favorable, et le dit.
Ses vérifications sont justes : pour $6|3$, $\frac{45}{8}$ contre $\frac{50}{9}$
diffèrent d'un comma syntonique exactement ; pour $5|5$, $\frac{25}{4}$ contre
$\frac{32}{5}$, de $128/125$ ; pour $6|5$, $\frac{61}{8}$ contre $8$, de $64/61$
(83 cents), moins qu'un demi-ton majeur ; pour $5|4$, après correction,
$\frac{41}{8}$ est au-dessous de $\frac{16}{3}$ de $128/123$ (69 cents), à
peine moins qu'un demi-ton mineur, comme elle l'écrit.\footnote{Vues 497--499,
ses pages (87) à (89), n\textsuperscript{o} 245. À la vue 499, pour $5|4$, elle
avait d'abord écrit que l'intervalle observé est « $> \frac{128}{25}$ » et que la
différence était « d'une quantité absolument inappréciable a l'oreille » ; elle
barre les deux et écrit « $<\frac{16}{3}$ », « au-dessous d'un demi ton
mineur » ($\frac{128}{25}$ est lu avec doute). La seconde rédaction (vue 609)
garde le premier état. Dans la table (vue 498), la ligne de $2|2$ est ajoutée
au-dessus des autres, et $6|5$ n'a pas de chiffre d'octave.}

Sa conclusion, qu'on ne pouvait « espérer plus d'accord entre la théorie et
l'experience », est vraie de ces comparaisons : entre figures à périodes
complètes, aucun écart ne dépasse 41 cents, et, pour les cas sans figure pris
au plus favorable, 83 cents ; Chladni lui-même tenait pour équivalentes des
figures dont les sons diffèrent d'un demi-ton, d'un ton ou d'une
tierce.\footnote{Vue 500, sa page (90) : « Il étoit impossible d'espérer plus
d'accord entre la théorie et l'experience », et « (V. latable p. 152) ».}

\emph{La difficulté.} Mais les deux familles ne s'accordent pas entre elles.
Rapportées aux figures de l'intégrale (C), dont les sons sont ceux de la lame
libre d'Euler, toutes les figures à périodes complètes sonnent environ un ton
trop grave ; et les vibrations tournantes, rapportées à ces mêmes sons, sont
plus graves encore, de un à deux tons et demi, d'autant moins qu'elles ont plus
de lignes nodales ; les plaques rectangulaires montrent le même
écart.\footnote{Vues 500 à 504, ses pages (90) à (94), n\textsuperscript{os}
246 et 247. Exemples de la vue 501 : $0|4$ (fig. 72--73), sol$^{\sharp}$3, de
nombre $\frac{25}{4}$ « d'après la théorie d'Euler », et $4|3$ (fig. 80),
fa$^{\sharp}$5, de nombre $25$, devraient être à deux octaves ; il y a deux
octaves moins un ton. De la vue 502 : 63 et 68, dans le rapport $2:8$, devraient
être à deux octaves, et sont à deux octaves et un ton et demi (de sol 1 à
si$^{\flat}$3, 27 demi-tons). L'arithmétique est juste. Une réserve sur le
premier exemple : par la règle que le feuillet applique ailleurs à la lame
libre (un nombre $(k-\frac12)^2$ pour $k$ lignes nodales, $\frac{25}{4}$ pour
trois lignes à la vue 508, $\frac{81}{4}$ pour cinq à la vue 520, conformes aux
racines d'Euler : $(\omega/\pi)^2 = 2{,}27$ ; $6{,}25$ ; $12{,}25$ ; $20{,}25$
pour 2, 3, 4, 5 nœuds), quatre lignes donneraient $\frac{49}{4}$, et le rapport
$25 : \frac{49}{4}$, une octave à peine, renverserait le sens de l'écart. Les
figures 72 appartiennent à sa « troisième valeur de l'angle $\omega$ », à
quatre nœuds sur la lame (vues 443--447). Comment le feuillet fait correspondre
$0|4$ à $\frac{25}{4}$, cette lecture ne le reconstitue pas.} Elle juge cet
écart trop constant pour être accidentel.

Elle en propose une cause, en la donnant pour « une simple conjecture » : les
quarts de période (parties où la surface n'a qu'une demi-période dans chaque
sens) et les demi-périodes des surfaces libres. Leur influence serait maximale
quand il n'y a que des quarts de période, moindre quand s'y ajoutent des
demi-périodes, minimale avec des périodes complètes ; de sorte que l'accord ne
serait parfait que pour une plaque appuyée de tous côtés, que Chladni n'a pas
essayée.\footnote{Vues 504 à 506, ses pages (94) à (96), n\textsuperscript{o}
248. Deux lignes barrées ouvrent la vue 505 : « si il s'agissoit de surfaces
appuyées de tout côté les mêmes équations, seroient applicables ». Elle ajoute
qu'une lame libre par $z_{xx}=0$, $z_{xxx}=0$, « s'il avoit été essayé avec
succès », comparée à la lame appuyée, aurait décidé la question (« m'avoit »
est lu avec doute ; le sens demande « m'auroit »).} Elle écarte l'analogie de
la corde et de la flûte, la théorie des instruments à vent n'étant pas
rigoureuse, et compare sa difficulté à celle de la vitesse du son, que
« l'ingénieuse hypothèse » de Laplace, introduite dans le calcul par Poisson, a
résolue ; elle ajoute que dans le cas d'une flûte on ne saurait séparer
l'effet de la chaleur dégagée de celui des demi-périodes.\footnote{Vues 506--507,
ses pages (96) et (97), n\textsuperscript{o} 249 : « M\textsuperscript{r} De la
Place », « M\textsuperscript{r} Poisson » — ce feuillet nomme Poisson. Il
s'agit de la correction de Laplace à la vitesse du son de Newton par la chaleur
dégagée dans la compression (contexte, non tiré des feuillets). La page finit :
« Je vais m'occuper àprésent des plaques rectangulaires differentes du
quarré. »}

Ce que dit la théorie actuelle. Ses nombres $m^2+n^2$ sont exactement ceux de
la plaque rectangulaire simplement appuyée, dont les modes sont
$\sin(m\pi x/a)\sin(n\pi y/b)$ : pour une telle plaque, sa conclusion
conjecturale — l'accord parfait avec les surfaces périodiques — est
vraie.\footnote{Le résultat pour la plaque appuyée est dû à Navier (contexte,
non tiré des feuillets) ; il tient à ce que, sur un bord droit, $z = 0$ et
$\Delta z = 0$ équivalent aux conditions d'appui de Kirchhoff.} Sur une plaque
libre, ni les surfaces périodiques ni l'intégrale (C) ne satisfont les conditions
du bord libre, hors, pour (C), le cas limite $\nu = 0$ (voir, dans la première
partie du mémoire, les conditions de Kirchhoff au N\textsuperscript{o} 9, vues
433 à 438, et ce qu'elles donnent pour (C) au § 4) ; les fréquences exactes ne sont données par aucune des deux
familles, et les figures ne leur ressemblent qu'approximativement. L'écart
qu'elle observe est de cet ordre : pour la lame libre, le nombre exact de $k$
nœuds est voisin de $(k-\frac12)^2$ et non de $k^2$, et le bord libre abaisse de
même les sons des figures de la plaque par rapport à ceux de la plaque appuyée.
Son explication par les quarts de période est la sienne ; ce n'est pas la cause
qu'on donne aujourd'hui.

\subsection*{§ 5. Les plaques rectangulaires (N\textsuperscript{os} 17 à 19)}

\pagerange{508}{537}
\emph{L'intégrale (C) sur le rectangle.} L'intégrale (C) veut la même valeur de
$\omega$ dans ses deux fonctions, de $x$ et de $y$. Sur un rectangle, cela
n'est possible en général qu'avec $B = 0$ ou $\mathcal{A} = 0$ : des lignes
nodales droites et parallèles ; les figures « à distorsions » n'apparaissent que
si deux valeurs de $\omega$ compensent l'inégalité des côtés.\footnote{Vue 508,
sa page (98) : le titre souligné « §.5 Comparaison dela théorie à l'experience
dans le cas des plaques rectangulaires differentes du quarré », puis
« N.o 17 ». Le coefficient bouclé est rendu $\mathcal{A}$ comme dans la
première partie ; l'ajout « ou $\mathcal{A}=0$ » est au-dessus de la ligne.}
Son exemple est juste : sur la plaque dont les côtés sont comme $5:7$, trois
lignes de la lame libre ($\frac{25}{4}$) sur le côté 5 et quatre
($\frac{49}{4}$) sur le côté 7 donnent le même son, à la précision des nombres
d'Euler, puisque
$\frac{25/4}{5^2} = \frac{49/4}{7^2} = \frac14$ ; Chladni y voit, en effet, des
figures « intermédiaires » (Traité d'acoustique, § 136).\footnote{Vues 508--509,
ses pages (98) et (99), n\textsuperscript{o} 250. Le feuillet dit les trois
lignes « transversales au côté qui est comme 7 » et les quatre « transversales
au côté qui est comme 5 » ; l'égalité demande que les trois lignes partagent le
côté 5 et les quatre le côté 7. Les dénominateurs de $\frac{25}{4}$ et de
$\frac{49}{4}$ sont surchargés.} Sur la plaque $3:1$, les cas $5|0$ et $0|2$
(fig. 77, § 122) sont de même à l'unisson, à très peu près :
$\frac{9}{2}\cdot\frac13 = \frac32\cdot 1$ (calcul de cette lecture ; les
racines d'Euler donnent $1{,}506$ au lieu de $\frac32$ pour deux nœuds, soit 6
cents d'écart).\footnote{Vue 510, sa page (100). Pour la plaque $3:5$ (§ 121),
le feuillet dit que « 2 lignes nodales parallelles dans un des sens doit donner le
même son que 4 pareilles lignes dans l'autre sens » ; le premier chiffre peut se
lire 3, le second est surchargé et lu 4. Le calcul demande 2 et 3 :
$\frac{3/2}{3} = \frac{5/2}{5}$, ce que ni $2$ et $4$ ni $3$ et $4$ ne
donnent. À vérifier sur l'image. Une note marginale remplace un passage barré
sur les figures de la « pl. 6 » et les fig. 170 ; un paragraphe barré annonçait
le détail de la fig. 163\,b.} Elle suit ensuite, sans la décrire en entier, la
fig. 163\,b : la règle des valeurs de $f(x)$ et $f(y)$ à l'origine, de la forme
$2(1+e^{\omega})$ ou $2(1-e^{\omega})$ selon la parité des nombres de lignes,
la supposition $\mathcal{A}(1-e^{\omega}) = -(1+e^{\omega})B$ au lieu de
$\mathcal{A} = -B$, et le passage de la courbe par les douze points de repos
invariables ; elle en tire la règle générale : si $\mathcal{A} = 0$ donne $m$
lignes droites et $B = 0$ en donne $n$, la figure est faite de $m$ lignes
sinueuses à $n/2$ flexions ou de $n$ lignes à $m/2$ flexions.\footnote{Vues 511
à 514, ses pages (101) à (104), n\textsuperscript{os} 251 et 252. Le « b » de
« 163\,b » et de « 175\,b » est souligné et peut se lire 6. Ces règles tiennent
à l'intégrale (C) (voir la première partie du mémoire) et en partagent la
limite : (C) ne satisfait pas les conditions du bord libre, hors le cas limite
$\nu = 0$.}

\emph{Les surfaces périodiques ; la formule du son.} Pour un rectangle de côtés
$s:t$ ($s>t$), elle multiplie les ordonnées parallèles au petit côté par
$s/t$ :
\[
z = \dots\,\sin\frac{\pi t m x}{tA}\,\sin\frac{\pi s n y}{tA},
\qquad \text{son} \propto \frac{t^2m^2+s^2n^2}{t^2},
\]
avec $m$ lignes parallèles à l'axe des $y$ et $n$ à l'axe des
$x$.\footnote{Vues 514--515, ses pages (104) et (105), n\textsuperscript{o}
253 : « N.o 18 \emph{surfaces periodiques} ». La lettre qui précède $ny$ est
écrite sur une autre et lue $s$ ; les deux dénominateurs sont $tA$.} La plaque
va donc de $0$ à $A$ en $x$ (le grand côté) et de $0$ à $tA/s$ en $y$ ; la
formule est celle de Navier, $A^2\bigl((m/a)^2+(n/b)^2\bigr)$ avec $a = A$,
$b = tA/s$. Dans ses propres termes :
\[
\text{son} \propto \frac{t^2\,T^2 + s^2\,L^2}{t^2},
\]
où $T$ est le nombre des lignes \emph{transversales} (parallèles au petit côté,
qui partagent le grand) et $L$ celui des lignes \emph{longitudinales} ; c'est
exact pour la plaque appuyée, approché pour la plaque libre (voir la difficulté
du N\textsuperscript{o} 16).\footnote{Les lettres changent d'une page à
l'autre, non la règle. Aux vues 515 à 520, $m$ compte les lignes transversales
et $n$ les longitudinales : $\frac{t^2m^2+s^2n^2}{t^2}$, lecture du lot 26,
conforme aux calculs des vues 518 et 520. Aux vues 524 à 534, elle note « avec
M.r Chladni, par $n/m$ la figure qui seroit composée de $n$ lignes transversales
et de $m$ lignes longitudinales », et la formule devient
$\frac{t^2n^2+s^2m^2}{t^2}$, lecture du lot 27, conforme à « $s=2$ $t=1$ …
$n^2+4m^2$ » de la vue 534. Les deux lectures sont donc justes, chacune à sa
page. La seule formule discordante est la spécialisation de la vue 524,
« $\frac{64m^2+81n^2}{64}$ », qui garde les lettres des vues 515--520 dans la
phrase des vues 524--534 ; le tableau qui la suit (vue 525) calcule $a|b$ par
$64a^2+81b^2$, le premier chiffre étant le nombre transversal : il suit la
règle. Le passage barré de la vue 533 donne d'ailleurs les deux formules, l'une
pour « $n$ lignes transversales $m$ lignes longitudinales », l'autre pour
« $m$ lignes transversales et $n$ lignes longitudinales ».} Les sommes de
solutions particulières sont admissibles, dit-elle, quand leurs termes ont le
même son : c'est exact, une figure stationnaire composée demande des modes de
même fréquence ; une ligne nodale commune à deux termes est alors une
« véritable ligne invariable de repos », et le mot « libre » n'a plus aucun de
ses deux sens.\footnote{Vues 516 à 518, ses pages (106) à (108),
n\textsuperscript{o} 254. « pourra » (vue 517) est lu avec doute.}

Ses exemples. Sur la plaque $6:5$, $4|1$ et $2|3$ ont les nombres
$\frac{16\cdot25+36}{25} = \frac{4\cdot109}{25}$ et
$\frac{4\cdot25+9\cdot36}{25} = \frac{4\cdot106}{25}$ ; leur intervalle
$109/106$ (48 cents) dépasse à peine son comma : les deux cas peuvent se
mêler, et la ligne longitudinale médiane, commune aux deux ($n = 1$ et $n = 3$
sont impairs), est la ligne invariable de la fig. 157.\footnote{Vues 518--519, ses
pages (108) et (109), n\textsuperscript{o} 255. Le feuillet écrit « $16.25+96$ »
là où le calcul demande $36$ ($s^2n^2$) ; la transcription garde 96 (son 3 à
queue ressemble au 9). La note marginale nomme le cas « $4|1$ » (lu avec doute)
et remarque que $106 < 109$ va dans le sens des vibrations tournantes, plus
graves.} Sur la plaque $5:4$, Chladni donne le même son à $5|0$ et $1|4$
(§ 115, fig. 158). Le feuillet décrit $5|0$ comme cinq lignes parallèles au grand
côté, et le terme périodique de $1|4$ comme quatre lignes parallèles au petit
côté et une dans l'autre sens, la médiane, commune aux deux termes et donc
invariable ; il calcule $5|0$ par le nombre d'Euler $\frac{81}{4}\cdot\frac{25}{16}
= \frac{2025}{64}$, conforme à cette description, et $1|4$ par
$\frac{16+25\cdot16}{16} = 26$, d'où le rapport $2025:1664$, environ une tierce
(340 cents). Mais $26$ est le nombre d'une ligne transversale et de quatre
longitudinales ; quatre transversales et une longitudinale, la figure que le
feuillet décrit, donnent $\frac{16\cdot16+25}{16} = \frac{281}{16}$, et un
rapport d'environ $1{,}80$ (1019 cents).\footnote{Vues 520--521, ses pages (110)
et (111), n\textsuperscript{o} 256. Le feuillet porte une fraction
intermédiaire barrée, dont la transcription lit le numérateur « 281 » : c'est
exactement $16\cdot16+25$. Le premier facteur du numérateur gardé est
illisible ; $26$ demande 1. Entre sa description et son calcul, la lecture
tient pour la description, que le chiffre barré confirme ; elle ne tranche pas
ce que Chladni désigne par $1|4$.} Avec 26, elle trouve un écart de deux tons
environ, qu'elle rapproche de la gravité connue des vibrations tournantes pour
expliquer que les deux termes se combinent en figures « compliquées par leur
influence simultanée » ; avec $\frac{281}{16}$, l'écart est d'environ une
septième (1019 cents), et l'explication ne suffit plus.

Sur la plaque $2:1$, $2|1$ et $3|0$ (fig. 173) n'ont pas de ligne invariable
commune, et n'en montrent pas ; $5|1$ et $1|3$ (fig. 174) ont chacun deux
lignes médianes, qui se retrouvent, droites, dans les figures ; elle écarte les
figures de plaques dont Chladni a « un peu alteré » les
proportions.\footnote{Vues 521 à 523, ses pages (111) à (113),
n\textsuperscript{o} 257. Le groupe « 18. » après « fig. 173 » est lu avec
doute.}

\emph{N\textsuperscript{o} 19. Les sons de la plaque $9:8$.} Elle laisse de
côté les vibrations tournantes et calcule les dix-huit cas $2|2$ à $6|3$ par
$\frac{64T^2+81L^2}{64}$ : $\frac{145}{16}$, $\frac{985}{64}$, $\frac{97}{4}$,
$\frac{2281}{64}$, $\frac{225}{16}$, $\frac{1305}{64}$, $\frac{117}{4}$,
$\frac{2601}{64}$, $\frac{337}{16}$, $\frac{1753}{64}$, $\frac{145}{4}$,
$\frac{3049}{64}$, $\frac{481}{16}$, $\frac{2329}{64}$, $\frac{181}{4}$,
$\frac{3625}{64}$, $\frac{657}{16}$, $\frac{3033}{64}$ — tous justes — puis
leurs rapports au nombre $8$ de la figure 68 de la plaque carrée.\footnote{Vues
524 à 526, ses pages (114) à (116), n\textsuperscript{o} 258. Les lettres de la
formule générale de la vue 524 sont lues avec doute (voir la note sur les
lettres, plus haut).} Chladni remarquant (p. 161) que toutes ses plaques
rectangulaires sonnent environ un demi-ton plus haut que la plaque carrée, elle
suppose que le verre de la plaque carrée différait, et prend pour 68
« si$^{\flat}$3 + aigu » au lieu de « si$^{\flat}$3 aigu ».\footnote{Vues
527--528, ses pages (117) et (118), n\textsuperscript{o} 259 : « quelque
circonstance physique inappréciable, telle qu'une différence de rigidité du
verre ou de son epaisseur ». La table de Chladni (p. 160) est recopiée à la vue
528.} Elle compare alors, cas par cas, le rapport théorique à l'intervalle
entendu, écrit comme un rapport juste, et conclut : « parmi les 18 comparaisons
[…] il ne se trouve qu'un seul cas » ($5|5$) où l'écart excède un
demi-ton.\footnote{Vues 529 à 533, ses pages (119) à (123), n\textsuperscript{os}
260 à 262. Ses produits en croix sont justes : 960 et 985, 97 et 96, 2281 et
2304, 2025 et 2048, 1011 et 1024, 15777 et 16384, 15245 et 16384, 481 et 512,
2329 et 2304, 3625 et 4096, 5913 et 6400, 15165 et 16384.}

Ce jugement est à reprendre. Comptées en demi-tons tempérés depuis si$^{\flat}$3,
les notes de Chladni sont au-dessus des rapports théoriques dans presque tous
les cas, de 13 à 211 cents ; la correction « + aigu », que le feuillet ne chiffre
pas, réduit d'autant les écarts positifs. Ses propres rapports donnent plus d'un
cas hors du demi-ton : $5|5$ (212 cents, plus qu'un ton majeur), $4|5$ (125),
$6|2$ (137) et $6|3$ (134), bien que pour les trois derniers l'intervalle entendu
ne soit qu'une borne (« un peu moindre que »). Trois de ses phrases sont
inexactes : pour $3|3$ l'écart $4096/3915$ (78 cents) n'est pas « bien moindre
qu'un demi-ton mineur » (71 cents) ; pour $4|3$, $16384/15777$ (65 cents) est
moins, et non « un peu plus », qu'un demi-ton ; pour $5|5$, $4096/3625$ passe le
ton majeur.\footnote{Vue 529 : la vérification écrite, « $\frac{3.1305}{8.512} <
\frac{25}{24}$, $\frac{783}{512} < \frac53$, $2349 < 2560$ », compare au
demi-ton un rapport plus petit que 1, ce qui est toujours vrai ; le rapport à
comparer est l'inverse, et $4096\cdot24 = 98304 > 3915\cdot25 = 97875$. Sa
première rédaction, barrée, disait « $\frac{4096}{3915}$ est $<\frac{16}{15}$ »,
ce qui est juste. Vue 530 : pour $3|5$ le feuillet écrit « $\frac{2601}{512}$
est en effet $>$ que $\frac{16}{3}$ » ; il est plus petit, comme le suppose
l'écart $\frac{8192}{7803}$ qu'elle donne (84 cents). Vue 532 : pour $5|5$,
« l'excès […] \emph{depres} d'un ton mineur » (« depres » lu avec doute) ; pour
$6|2$, « donné par l'experience » pour « par la théorie ».} Comptés en demi-tons tempérés, avec la note la plus
favorable quand Chladni en donne deux, quinze cas sur dix-huit restent à moins
d'un demi-ton majeur de la note avant même sa correction, dans la tolérance
qu'elle s'est donnée ; $5|5$, $6|2$ et $6|3$ n'y sont pas.

\emph{La plaque $2:1$.} Par $T^2+4L^2$ elle trouve pour $2|2$, $2|3$, $3|2$,
$3|3$, $4|2$, $4|3$, $5|2$, $5|3$ les nombres $20, 40, 25, 45, 32, 52, 41,
61$, justes, et les compare à la table de Chladni (p. 168). Six cas sont à moins
d'un demi-ton ; $2|2$ et $2|3$, dont l'intervalle mutuel est exactement l'octave
des deux côtés, sont trop hauts dans la théorie d'un peu plus d'un
ton.\footnote{Vues 534 à 536, ses pages (124) à (126), n\textsuperscript{o} 263.
En demi-tons tempérés depuis si$^{\flat}$3, les écarts sont $-186$, $-186$,
$+27$, $-90$, $+100$, $-41$, $+71$, $+83$ cents. Pour $4|2$ (si 5) elle écrit
que l'intervalle est « un peu plus grande que 2 octaves », « à fort peu près
exprimée par 4 » : de si$^{\flat}$3 à si 5 il y a deux octaves et un demi-ton.
Le cas $4|3$ de la dernière table est lu d'après la première (tache). Un
paragraphe citant le § 124 de Chladni (sons de la plaque carrée « à peu près dans
le rapport des nombres 6, 15, 30 \&c. », tendant vers la suite des nombres
naturels quand un diamètre diminue) est annulé d'une grande croix.}

Elle note enfin que les figures des équations périodiques se formeraient aussi
sur une membrane tendue — renversées en partie si les côtés sont libres,
identiques si la plaque est appuyée — avec des sons proportionnels aux « simples
nombres » là où la plaque a leurs carrés, comme Euler l'a remarqué de la lame
appuyée et de la corde : c'est exact, la membrane ayant $\omega \propto k$ et
la plaque $\omega \propto k^2$ pour le même nombre d'onde
$k$.\footnote{Vue 537, sa page (127), n\textsuperscript{o} 264. « Je » est lu
avec doute. Elle dit avoir voulu expérimenter sur des plaques appuyées sur un ou
plusieurs côtés, sans y parvenir : c'est l'occasion du § 6.}

\subsection*{§ 6. L'effet d'un corps mou : la cire (N\textsuperscript{o} 20)}

\pagerange{538}{547}
Voulant réaliser un appui en serrant un bord d'une plaque carrée dans une
rainure de bois, elle attendait un son plus aigu et l'obtint plus grave : le
poids du bois l'emportait. Elle charge alors un ou plusieurs côtés de
plaques de verre de cire collante, uniformément répartie, et
trouve l'effet équivalent à un prolongement fictif de la plaque, d'une étendue
égale en poids à la cire : même son, lignes nodales déplacées en
conséquence.\footnote{Vues 538--539, ses pages (128) et (129),
n\textsuperscript{o} 265. Le titre souligné « De l'effet de l'application d'un
corps mou, aux côtés d'une plaque rectangulaire élastique vibrante » est précédé
d'un signe lu « § 6 » avec doute ; « N.o 20 » suit.}

\emph{Les expériences.} Plaque carrée de 106 mm de côté, pesant 1 once 18
grains ; le son de la figure à deux lignes nodales croisées ($1|1$), un peu
plus bas que sol 1 (l'octave 1 étant celle du la du diapason). Chaque dose de
cire, 1 gros 2 grains, est le huitième de ce poids ; elle vaut donc
$106/8 \approx 13$ mm de plaque.\footnote{Vues 539 à 542, ses pages (129) à
(132), n\textsuperscript{os} 265 et 266. En grains (1 once = 8 gros = 576
grains), la plaque pèse 594 grains et la dose 74, le huitième étant 74,25 : la
dose est bien le huitième. La vue 540 écrit pour la première dose « 1 once 2 gros
de cire », que la vue 541 (« une egale quantité, savoir 1 gros 2 grains ») et la
vue 542 contredisent : à vérifier sur l'image.} Une dose sur un côté : le son
tombe entre mi 1 et fa 1, plus d'un ton plus bas ; la ligne nodale parallèle au
côté chargé est à 45 mm de lui et à 58 mm du côté opposé, contre $46\frac12$ et
$59\frac12$ calculés — l'écart est la largeur de la ligne et de la cire, environ
3 mm. Deux doses sur deux côtés contigus : environ deux tons majeurs. Deux doses
sur deux côtés opposés : ré 1 exactement, et la figure ne bouge pas, la charge
étant symétrique.\footnote{Vues 541 à 544, ses pages (131) à (133),
n\textsuperscript{o} 267. « un plus grave » (vue 541) est écrit ainsi. Elle
recommande de laver souvent la plaque à l'esprit-de-vin, la moindre trace de
cire empêchant les figures de se former.}

\emph{Le calcul, et sa correction.} Elle traite la plaque chargée comme un
rectangle $119:106 \approx 9:8$ et écrit son nombre $\frac{9^2+8^2}{64} =
\frac{145}{64}$ contre $2$ pour la plaque carrée, d'où un rapport $\frac{145}{128}$,
voisin du ton majeur. Ce calcul mêle deux unités : $\frac{145}{64}$ est rapporté
au grand côté du rectangle (119 mm), $2$ au côté du carré (106 mm) ; tel quel,
il ferait la plaque agrandie plus aiguë. Le rapport juste est
\[
\frac{2/8^2}{1/9^2+1/8^2} = \frac{2\cdot 9^2}{9^2+8^2} = \frac{162}{145}
\quad (192 \text{ cents}),
\]
un ton mineur à peu près, la plaque chargée étant la plus grave. Le feuillet de
la vue 416 donne exactement cette correction.\footnote{Vues 542--543, ses pages
(132) et (132), n\textsuperscript{o} 267. Vue 416, petit feuillet déchiré relié
au début du mémoire, d'une plume plus fine : « Il regne dans tous le passage
relatif aux effets de l'application dela cire une erreur » ; le poids de la
plaque chargée étant à celui de la plaque non chargée comme $N$ à $n$, « on a
pris le son proportionnel au nombre $\frac{N^2+n^2}{n^2}$ parconséquent ce son
seroit $>$ aigu que $\frac{2n^2}{n^2}$ […] le contraire doit avoir lieu et le
premier son est proportionnel au nombre $\frac{N^2+n^2}{N^2}$ », et l'intervalle
« que j'ai exprimé par $\frac{N^2+n^2}{2n^2}$ doit l'être par
$\frac{2N^2}{N^2+n^2}$ qui heureusement ne different pas sensiblement ». Avec
$N=9$, $n=8$, ce sont $\frac{145}{128}$ et $\frac{162}{145}$ : le feuillet
corrige ce § 6, non le N\textsuperscript{o} 15 de l'Exposition des principes
(vues 38--40), qui décrit la même expérience sans chiffres.} Pour une dose, la
correction ne change presque rien (216 contre 192 cents). Pour deux côtés
contigus, le rapport $81/64$ (deux tons majeurs) est juste, la plaque restant
carrée. Pour deux côtés opposés, en revanche, la même erreur donne
$\frac{10^2+8^2}{8^2} = \frac{41}{16}$ et un intervalle $\frac{41}{32}$ (429
cents), « un peu plus que deux tons majeurs », comme l'expérience ; corrigé, il
devient $\frac{50}{41}$ (344 cents), moins d'une tierce majeure, alors que de
ré 1 à sol 1 elle entend près d'une quarte. Le feuillet de la vue 416 ne relève
pas ce cas, où la correction n'est pas négligeable.\footnote{Vue 544, sa page
(133). « à » (« l'addition à deux des côtés opposés ») paraît corrigé sur « au ».
Avec les longueurs mêmes (106 et 132 mm) le rapport corrigé est $1{,}216$, 338
cents.}

\emph{La lame chargée.} Une lame de verre de 20 cm sur 23 mm, pesant 4 gros 12
grains, montre une ligne nodale longitudinale et deux transversales, l'espace
entre celles-ci double de celui qui les sépare des bouts. Chargée à un bout de
1 gros 3 grains, elle montre les deux lignes transversales à 12 mm du bout chargé
et 121 mm l'une de l'autre, 61 mm restant jusqu'à l'autre bout : une portion dix
fois moindre que l'autre rend le même son.\footnote{Vues 545--546, ses pages
(13?) et (135), n\textsuperscript{o} 268 ; le dernier chiffre de la page est
récrit, sans doute 4 sur 2, et deux pages consécutives portent (132). Le « 2 »
des lignes transversales est lu avec doute ; « les deux lignes transversales »
l'appuie.} Elle ne l'explique pas sur la page ; sa règle le fait : la cire pèse
75 grains, le quart des 300 de la lame, et vaut 50 mm ; sur une lame fictive de
250 mm partagée en $\frac14$, $\frac12$, $\frac14$, les segments seraient
$12{,}5$, $125$ et $62{,}5$ mm, contre 12, 121 et 61 observés (6 mm étant pris
par la cire et les lignes). Elle conclut qu'on peut, par une matière pesante et
non élastique, faire rendre à une portion aussi petite qu'on veut le son d'une
portion donnée, et que l'analyse permet la substitution à toute distance de
l'origine, les termes qui doivent s'annuler aux limites étant multipliés par un
facteur de rigidité.\footnote{Vues 546--547, ses pages (135) et (136),
n\textsuperscript{o} 269. « paroît », « lagene » (« la gêne »), « equivalante »
et « satisfera » sont lus avec doute. L'Exposition des principes (vues 38--40)
explique le même phénomène par son hypothèse sur la force élastique ; voir « Le
bord, les lignes nodales, la plaque chargée » (vues 31 à 41).}

Ce que dit la théorie actuelle. L'équivalence « égale en poids » est exacte au
premier ordre : à un bord libre, le moment de flexion est nul, si bien qu'une
bande ajoutée de largeur $\delta$ apporte l'énergie cinétique de sa masse et une
énergie de flexion d'ordre supérieur en $\delta$ ; au premier ordre, prolonger
et charger d'un poids égal abaissent la fréquence de la même quantité (quotient
de Rayleigh).\footnote{Le nom est postérieur aux feuillets (Rayleigh, 1877), et
le feuillet ne donne pas cet argument.} Pour des charges fortes, la masse
ajoutée n'équivaut plus à une longueur : elle ajoute de l'inertie sans
raideur, et l'équivalence n'est qu'approchée ; la limite d'une charge
infinie est un bord immobile, où la ligne nodale voisine vient se confondre avec
le bord, ce qui justifie qualitativement sa « portion aussi petite qu'on
voudra ».\footnote{Calcul fait pour cette lecture sur la lame libre aux deux
bouts, prise pour modèle (éléments finis, 200 éléments). Pour une masse ajoutée à
un bout égale à 1\,\% de celle de la lame, la fréquence du premier mode avec la
masse et celle de la lame prolongée de 1\,\% diffèrent de 0{,}09\,\% ; à 10\,\%,
de 6\,\% ; pour une masse égale à celle de la lame, la fréquence avec la masse est
2,9 fois celle de la lame doublée. Quand la masse passe de 1\,\% à dix fois celle
de la lame, la ligne nodale voisine du bout chargé passe, dans le même modèle, de
$0{,}78$ à $0{,}99$ de la longueur, et la fréquence tend vers une limite finie,
alors que celle d'une lame réellement prolongée baisserait sans limite.
L'Exposition des principes (N\textsuperscript{o} 15, vues 38 à 40) rapporte le
même effet, sans chiffres.}

\subsection*{§ 7. La plaque circulaire (N\textsuperscript{os} 21 à 23)}

\pagerange{548}{574}
Les applications faites jusque-là ne regardent que les plaques
rectangulaires ; avant de terminer, elle cherche une intégrale particulière
« immédiatement applicable » à la plaque circulaire.\footnote{Vue 548, sa page
(137), n\textsuperscript{o} 269 : le titre souligné « § 7. Recherche d'une
intégrale particulière immédiatement applicable au cas d'une plaque elastique
circulaire », puis « N.o 21 » (le 2 pourrait être un 3). La page s'arrête sur
« on pourra toujours faire » ; la suite est la page (138) de la vue 557. Entre
les deux sont reliés le brouillon sur la chaleur de la vue 549, un feuillet
blanc, et deux petits feuillets de calcul (vues 553, 555).}

\emph{L'équation en coordonnées polaires.} Avec $x = r\cos\varphi$,
$y = r\sin\varphi$, elle calcule toutes les dérivées jusqu'au quatrième ordre
et trouve
\[
\Delta^2 z = \partial_r^4 z + \frac{2}{r}\partial_r^3 z - \frac{1}{r^2}\partial_r^2 z
+ \frac{1}{r^3}\partial_r z + \frac{2}{r^2}\partial_r^2\partial_\varphi^2 z
- \frac{2}{r^3}\partial_r\partial_\varphi^2 z + \frac{4}{r^4}\partial_\varphi^2 z
+ \frac{1}{r^4}\partial_\varphi^4 z,
\]
ce qui est juste.\footnote{Vues 557 et 560, ses pages (138), barrée, et (139),
n\textsuperscript{o} 275 (vue 559 : seconde prise de la vue 557). Vérifié par
calcul symbolique. Le dernier dénominateur $d\varphi^4$ de la vue 557 est lu
avec doute. Les petits feuillets reliés avant ces pages sont des essais : à la
vue 553, en haut (n\textsuperscript{o} 273), des dérivées en $r$, $\varphi$ et
un troisième angle $\delta$, où le dernier terme de $\partial z/\partial y$ porte
$dz/dr$ là où l'on attend $dz/d\delta$ ; en bas (n\textsuperscript{o} 272), la
condition du contour prise « en prenant le rayon constant $r = a$ », où
$\partial z/\partial x$ se réduit à $-\frac{\sin\varphi}{a}\partial_\varphi z$ —
c'est la seule dérivée tangentielle, la dérivée normale au contour étant
$\partial_r$ — et la remarque que $\partial^2 z/\partial x\partial y = 0$ donne
$z = F(x)+f(y)$. À la vue 555 (n\textsuperscript{o} 274), des « cas
particuliers de l'intégrale (E) » et leurs dérivées, avec trois signes
illisibles.} Posant $b\tau^2 = e^4/a^4$, l'équation (B) du mouvement régulier
devient $\frac{a^4}{e^4}z = \Delta^2 z$ ; comme seules des dérivées paires en
$\varphi$ y figurent, elle pose $z = u(r)\sin(\beta\varphi + B)$, où $\beta$ sera
le nombre des diamètres nodaux, et obtient
\[
\frac{a^4}{e^4}\,u = u'''' + \frac{2}{r}u''' - \frac{1+2\beta^2}{r^2}u''
+ \frac{1+2\beta^2}{r^3}u' + \frac{\beta^4-4\beta^2}{r^4}u \qquad (\text{B}')
\]
puis, avec $u = r^\beta s$ et $n = 1+2\beta$ (le $n$ de ces pages, sans rapport
avec celui des plaques rectangulaires),
\[
\frac{u^4}{e^4}\,s = s'''' + n\Bigl\{\frac{2}{r}s''' + (2\beta-1)\Bigl(\frac{s''}{r^2} - \frac{s'}{r^3}\Bigr)\Bigr\}
\qquad (\text{B}''),
\]
toutes deux justes.\footnote{Vues 560 et 565, ses pages (139) et (140),
n\textsuperscript{o} 278. Le feuillet emploie $u$ à la fois pour la fonction de
$r$ et, dans $u^4/e^4$, pour le paramètre de fréquence ; on garde ses lettres.
Le dernier terme du développement de $d^4u/dr^4$ est écrit sans le facteur
$r^\beta$. La lettre lue $\tau$ ressemble ici à un $r$ ; le $e$ pourrait être un
$c$. Vérifié par calcul symbolique.}

\emph{La solution : les fonctions de Bessel.} En cherchant $s$ en série de
puissances croissantes de $r$, elle trouve, avec $g = u^2/e^2$,
\[
s = A\Bigl\{1 - \frac{g r^2}{2(n+1)} + \frac{g^2 r^4}{2\cdot4\,(n+1)(n+3)} - \dots\Bigr\}
+ A'\Bigl\{1 + \frac{g r^2}{2(n+1)} + \frac{g^2 r^4}{2\cdot4\,(n+1)(n+3)} + \dots\Bigr\}.
\]
Comme $(n+1)(n+3)\cdots(n+2j-1) = 2^j(\beta+1)\cdots(\beta+j)$, les deux séries
multipliées par $r^\beta$ sont, à un facteur près, $J_\beta(kr)$ et $I_\beta(kr)$,
$k = u/e$ : la fonction de Bessel de première espèce et la fonction de Bessel
modifiée. La solution est donc $z = \bigl(A\,J_\beta(kr) + A'\,I_\beta(kr)\bigr)
\sin(\beta\varphi + B)$ à des constantes près, et
$\Delta z = k^2\bigl(-A\,J_\beta + A'\,I_\beta\bigr)\sin(\beta\varphi+B)$, comme
elle l'écrit.\footnote{Vues 565--566, ses pages (140) et (141). Coefficients
vérifiés jusqu'au cinquième terme pour $\beta = 0, \dots, 3$, et la série
vérifiée dans (B$''$). Le feuillet ne nomme pas ces fonctions ; le nom est celui
de Bessel (1824), et si ce mémoire est celui de 1813 — inférence, voir plus
haut —, il le précède. Le premier dénominateur $2(n+1)$ est lu avec doute (le
chiffre peut passer pour un 8). Elle note qu'on ne gagnerait rien à chercher
$s$ sous la forme de sinus et cosinus de $\frac{u}{e}r$ multipliés par des
séries, qui ramèneraient à la même forme.} Elle ne garde que deux constantes :
c'est juste pour une plaque pleine, les deux autres solutions ($Y_\beta$,
$K_\beta$) étant infinies au centre ; pour un anneau ou une plaque percée,
comme celle de la vue 589, il les faudrait.

\emph{Les couples de conditions au contour.} Elle reprend les quatre couples de
conditions de son N\textsuperscript{o} 9 et les écrit en $z'$ et $z''$, les deux
parties de la solution qui répondent aux deux facteurs de l'équation
(N\textsuperscript{o} 2).\footnote{Vues 566--567, ses pages (141) et (142),
n\textsuperscript{o} 279. Dans le premier couple, deux $z$ n'ont pas leur accent
sur la page ; les huit équations de la vue 567 sont régularisées par la
transcription.}

— Le premier couple (points appuyés) est satisfait par
$\sin(\beta\varphi+B) = 0$ : $2\beta$ points du contour, extrémités de
$\beta$ diamètres qui sont des lignes d'appui ; les secteurs égaux se superposent,
et la dépendance en $\varphi$ a la forme de l'équation des cordes. Pour
$\beta = 0$, on ne peut prendre $B = 0$, sans quoi la plaque ne bougerait pas.
Ou bien par $A' = 0$ et $J_\beta(ka) = 0$ : la plaque appuyée sur tout son
contour, dont l'équation est celle du tambour d'Euler, les figures celles du
tambour et les sons les carrés des siens.\footnote{Vues 567 à 569, ses pages
(142) à (144), n\textsuperscript{o} 280 : « que Euler adonnée dans le tome X de
pétersbourg. (V. N.o 13 du m. \emph{De motu vibratorio tympanorum}) ». C'est
dans ce mémoire d'Euler sur la membrane circulaire que la série de Bessel
apparaît pour la première fois (contexte, non tiré des feuillets). Dans « entre
$t$ et $\varphi$ », le $t$ est lu avec doute, de même « uu » et « aux ».}
Dans sa théorie, c'est exact : si $z = 0$ sur le cercle et $\Delta z = 0$, le
mode est celui de la membrane, et pour un même nombre d'onde $k$ la plaque a
$\omega\propto k^2$ là où la membrane a $\omega\propto k$. Dans la théorie de
Kirchhoff, c'est faux pour un matériau réel : sur le cercle $z = 0$, le moment
radial vaut $-D\bigl(\partial_r^2 z + \nu\,\partial_r z/a\bigr) =
D(1-\nu)\,\partial_r z/a$ pour un mode de membrane, et ne s'annule que si
$\nu = 1$ ; l'artifice de Navier ne vaut que pour les bords
droits.\footnote{Son énergie $\iint(\Delta z)^2$ est celle de Kirchhoff avec
$\nu = 1$ (voir la première partie du mémoire et la section de l'Exposition des
principes). Kirchhoff (1850) est postérieur aux feuillets.}

— Le second couple (contour libre) : elle l'écrit $\Delta z = 0$ et
$\partial_x \Delta z = 0$, « ou, cequi est ici lamême chose » $\partial_y\Delta z
= 0$, et le satisfait par la seule équation
$A\,J_\beta(ka) - A'\,I_\beta(ka) = 0$ (à des constantes près) ; comme on peut
toujours choisir $A$ et $A'$ pour qu'elle ait lieu, elle conclut que les
conditions du contour « auront lieu d'elles-mêmes » et ne peuvent servir à
trouver $u$, qui « devr[a] être donné d'ailleurs ».\footnote{Vues 569--570, ses
pages (144) et (145), n\textsuperscript{o} 280 ; « N.o 22 » ouvre l'alinéa. Le signe
$\neq$ et la lettre $\alpha$ qui terminent les deux premières séries de la vue
570 sont lus avec doute. À la vue 570, la série alternée est donnée pour $A$, à
la vue 571 pour $A'$ ; c'est la seconde attribution qui annule la combinaison.
Le dernier dénominateur de la vue 569 porte un 4 là où les autres ont $(n+3)$.}
Ceci est faux. Au contour, la dérivée à prendre est la dérivée normale
$\partial_r$, et les conditions du bord libre de sa théorie sont $\Delta z = 0$
et $\partial_r\Delta z = 0$ : deux équations en $A$, $A'$,
\[
-A\,J_\beta(ka) + A'\,I_\beta(ka) = 0, \qquad -A\,J'_\beta(ka) + A'\,I'_\beta(ka) = 0,
\]
dont le déterminant $J_\beta(ka)I'_\beta(ka) - I_\beta(ka)J'_\beta(ka) = 0$ fixe
les valeurs de $u$. C'est l'équation même de la plaque encastrée
($z = \partial_r z = 0$) : $ka = 3{,}196$ ; $4{,}611$ ; $5{,}906$ ; $7{,}144$ pour
$\beta = 0, 1, 2, 3$. Dans sa théorie, la plaque circulaire libre a donc les sons
de la plaque encastrée ; et toute déformation harmonique ($\Delta z = 0$, par
exemple $r^\beta\sin\beta\varphi$) n'y coûte aucune énergie.\footnote{Pour tout
contour, si $w$ est un mode encastré, $\Delta w$ satisfait $\Delta(\Delta w) = 0$
et $\partial_n\Delta(\Delta w) = 0$ au bord, puisque $\Delta^2 w = k^4 w$ : les
modes libres de l'énergie $\iint(\Delta z)^2$ sont les laplaciens des modes
encastrés. Racines calculées numériquement. La nullité de l'énergie pour
$z$ harmonique est celle que la section de l'Exposition discute (lagrangien
nul).} La plaque libre de Kirchhoff ($\nu\approx0{,}3$) a une autre équation ;
ses modes à diamètres seuls ont des fréquences dans les rapports
$1 : 2{,}32 : 4{,}08 : 6{,}25$ pour $\beta = 2$ à $5$, voisins des carrés
$1 : 2{,}25 : 4 : 6{,}25$ que Chladni observe.\footnote{Premières racines
$ka = 2{,}315$ ; $3{,}527$ ; $4{,}673$ ; $5{,}788$ pour $\nu = 0{,}3$, calculées
par les conditions $M_r = 0$, $V_r = 0$ (voir la première partie du mémoire pour
leur énoncé).}

\emph{L'ajout sur les lignes annulaires.} Un petit feuillet relié avant la page
(140), « addition a la p. 140 », traite le cas sans diamètres, où toutes les
lignes nodales sont annulaires. Elle y écrit les conditions du contour avec la
dérivée radiale, $\Delta z\,\delta(\partial_r z) - \partial_r(\Delta z)\,\delta z
= 0$ (multipliées par $\cos\varphi$ et par $\sin\varphi$), ce qui est juste ; puis,
« à cause de » $\Delta z = \frac{e^2}{a^2}z$ et de sa dérivée, elle les dit
satisfaites par $z = 0$ ou par $\partial_r z = 0$ : toutes les lignes de repos
satisferaient aux conditions des extrémités sans être des lignes d'appui « dans le
sens analytique attribué a ce mot par Euler », de même toute ligne circulaire où
la tangente est parallèle au plan de la plaque ; et « les extremités réelles et
physiques sont assujetties a la condition $\frac{dz}{dr} = 0$ ». Au verso, avec
$g = u^2/e^2$, elle choisit $A$ et $A'$ proportionnels à $I_0'(R)$ et à
$-J_0'(R)$, pour que $\partial_r z$ s'annule au bord quel que soit $g$, et
donne pour lignes annulaires les racines $r<R$ de $A\,J_0(r)+A'\,I_0(r) = 0$ (à
des facteurs près) ; ces séries sont justes.\footnote{Vues 563--564,
n\textsuperscript{o} 277, sans pagination ; le 4 de « p. 140 » a la forme d'un
b. « analytique » est lu avec doute (la vue porte « ahotytique » ou approchant).
Au verso, deux $r$ et un dénominateur sont lus avec doute, et les exposants du
premier terme de la dernière ligne, surchargés, ne se lisent pas ; le calcul
s'arrête sur un signe $=$ au bord. Le petit feuillet de la vue 561
(n\textsuperscript{o} 276) tire de $\frac1r z_r + z_{rr} = \frac{u^2}{e^2}z$
l'équation du quatrième ordre de la plaque à lignes annulaires, « comprise dans
celle qui a été donnée 139 » : c'est juste, à un signe près dans une ligne
intermédiaire ($+\frac{1}{r^3}z_r$ où le calcul demande $-$), que la réduction
corrige ; l'exposant $u^4/e^4$ est lu avec doute.} L'argument est faux :
$\Delta z = \frac{e^2}{a^2}z$ n'est vrai que de l'une des deux parties de la
solution, non de $z = A\,J_0 + A'\,I_0$ ; et $\partial_r z = 0$ n'est pas une
condition du bord libre, c'est la moitié de l'encastrement. Les conditions de
sa théorie restent $\Delta z = \partial_r\Delta z = 0$, qui déterminent $u$.

Les lignes nodales circulaires sont de simples lignes de repos, non des lignes
d'appui : une ligne d'appui circulaire exigerait $A' = 0$, ce qui ne convient
qu'au contour appuyé — ce qui est exact dans sa théorie. Diamètres et cercles
peuvent exister ensemble ou séparément.\footnote{Vues 570--571, ses pages (145)
et (146), n\textsuperscript{o} 281. La note marginale « + qui seront des lignes
d'appui » est appelée après « diamétrales ». À la vue 570 elle compare le cas où
l'on néglige la variation du rayon à celui des plaques rectangulaires où une
seule coordonnée varie.} Le troisième couple (contour fixé) donne
$A\,J_\beta + A'\,I_\beta = 0$, qui est la moitié des conditions de
l'encastrement ; le quatrième donne aussi $\cos(\beta\varphi+B) = 0$ : des
droites qui partagent chaque secteur en deux secteurs égaux, chacun appuyé d'un
côté et « libre » de l'autre au sens de $\partial_\varphi z = 0$,
$\partial_\varphi^3 z = 0$ — c'est exact, ce sont les axes de symétrie du
mode.\footnote{Vue 572, sa page (147).}

\emph{N\textsuperscript{o} 23. Les figures de Chladni, et l'anneau.} La
difficulté de trouver $u$ l'empêche de comparer les sons ; elle range les
figures de Chladni par nombre de diamètres : 104 et 109 pour $\beta = 0$ ;
105, 110\,a, 116\,a pour $\beta = 1$ ; 99, 106, 111 pour $\beta = 2$ ; 100,
107 pour $\beta = 3$ ; 101\,a, 108 pour $\beta = 4$ ; 102\,a pour
$\beta = 5$.\footnote{Vue 573, sa page (148), n\textsuperscript{o} 282. La
lettre qui suit « N. » dans « (V. N.o 2) » est empâtée.} Faisant ensuite
$r = a$ dans l'équation et effaçant les dérivées en $r$, elle obtient
$\frac{u^4}{e^4}z = \frac{1}{a^4}\partial_\varphi^4 z$, qu'elle donne pour
l'équation de l'anneau élastique, cas particulier de celle de la plaque, et qui
ne diffère, dit-elle, de celle d'Euler, à la fin de son mémoire sur la lame, qu'en
ce que les arcs y remplacent les ordonnées. Ce n'est pas une dérivation :
poser $r = a$ dans l'équation de la plaque ne donne pas celle d'un anneau, qui
est une poutre courbe avec ses propres équations ; l'équation obtenue est celle
d'une lame droite appuyée refermée sur elle-même, qui est l'hypothèse d'Euler
(§ 53 à 55 ; voir la section sur la copie d'Euler). Il s'ensuit, sur la page, que
les sons sont comme les carrés de $2, 4, 6, \dots$ (le nombre des nœuds), et
qu'Euler, qui donne la suite $4, 9, 16, 25, \dots$, n'a pas vu que l'anneau ne
peut avoir qu'un nombre pair de nœuds. Ce dernier point est exact : un mode de
flexion de l'anneau en $\cos\beta\varphi$ a $2\beta$ nœuds.\footnote{Vue 574, sa
page (149). La phrase se poursuit à la vue 577, page (150), par-dessus la note de
la vue 575 : « qu'un nombre pair de nœuds dans l'anneau équivaut à un nombre
impair de nœuds dans la lame droite ». Elle ajoute (vue 577) que, sur une plaque
libre n'ayant que des diamètres, $u$ restant indéterminé, « rien n'empêche donc
de faire $\frac{u}{e} = \beta$ », d'où des sons comme $\beta^2$ : c'est un
choix, non une conséquence de l'équation, et l'on a vu que $u$ n'est pas
indéterminé.}

\subsection*{L'anneau et la lame courbe, contre Euler et avec Chladni (N\textsuperscript{os} 24 et 25)}

\pagerange{574}{592}
\emph{L'objection de Chladni.} Chladni observe (Traité d'acoustique, p. 118)
que la théorie de l'anneau d'Euler n'est pas confirmée : les sons de l'anneau
partagé en 4, 6, 8 parties sont comme les carrés de 3, 5, 7, et non de 4, 6, 8.
Elle ne doute pas que le cas observé soit celui des extrémités appuyées, les
parties étant égales (§ 89 du Traité) ; elle attribue l'écart à une circonstance
hors du calcul. Chladni ayant vu, sur une bande de fer ployée (Traité, section
VI, p. 113, fig. 37 à 42 ; la fourche), les sons baisser et les nœuds se
rapprocher dans les parties courbes, elle soutient que la courbure prise
analytiquement ne change rien aux formules — les instruments à vent en
témoignent — mais que la bande pliée a perdu de son élasticité : dans l'anneau,
uniformément ; de là des sons plus graves, d'autant moins que l'arc entre deux
nœuds est petit, et un rapport $(2\beta-1)^2 : (2\beta)^2$ avec la théorie
d'Euler, qui tend vers 1 quand $\beta$ croît.\footnote{Vues 577 à 581, ses
pages (150) à (154), n\textsuperscript{os} 284 à 286. « Pour concilier des
expériences sur l'exactitude desquelles on ne peut élever aucun doute, avec une
théorie donnée par un des géometres dont l'autorité a le plus de poids ». Le
numéro de figure après « Dans la fig. » (vue 578) est deux petites lettres lues
« ee ». À la vue 581, elle écrit que les nombres de Chladni sont à ceux d'Euler
« $::2\beta-1:2\beta$ » et nomme « différence » le rapport
$\frac{(2\beta)^2}{(2\beta-1)^2}$.}

Ce que dit la théorie actuelle. Chladni a raison, et la loi carrée — celle
d'Euler comme la sienne — est fausse pour l'anneau. Les modes de flexion d'un
anneau mince dans son plan, à $2n$ nœuds, ont des fréquences proportionnelles à
\[
\frac{n(n^2-1)}{\sqrt{n^2+1}}, \qquad n = 2, 3, 4, \dots,
\]
soit $1 : 2{,}83 : 5{,}42 : 8{,}77$ pour 4, 6, 8, 10 nœuds ; les carrés de
Chladni donnent $1 : 2{,}78 : 5{,}44 : 9$, les siens
$1 : 2{,}25 : 4 : 6{,}25$.\footnote{La loi est celle de Hoppe (1871), postérieure
aux feuillets. Pour les vibrations hors du plan de l'anneau, la loi diffère
un peu mais donne des rapports voisins ($1 : 2{,}87 : 5{,}54$ pour une section
circulaire, $\nu = 0{,}3$). Rapports calculés.} La cause n'est pas une
élasticité diminuée : la courbure lie le déplacement radial au déplacement
tangentiel (l'anneau fléchit sans s'allonger), et elle entre dans les équations,
contrairement à ce que le feuillet affirme. Deux de ses remarques coïncident
pourtant avec la loi actuelle : la courbure abaisse les sons, et d'autant moins
que l'arc entre deux nœuds est court — le rapport de la loi de l'anneau à celle
de la lame droite, $\frac{n^2-1}{n\sqrt{n^2+1}}$, vaut $0{,}67$ ; $0{,}84$ ;
$0{,}91$ pour $n = 2, 3, 4$ et tend vers 1 ; et l'anneau n'a qu'un nombre pair de
nœuds.

\emph{Ses expériences sur le verre (N\textsuperscript{o} 25).} Pour écarter
l'élasticité variable du métal plié, elle fait tailler dans le verre des lames
courbes. Des fourches (deux branches droites réunies par un arc) vibrent comme la
lame droite libre, les nœuds étant plus espacés vers le milieu : le contraire de
la fourche de fer de Chladni, qu'elle explique par l'élasticité
variable.\footnote{Vues 581 à 585, ses pages (154) à (158),
n\textsuperscript{os} 286 à 288. Elle récuse aussi l'argument du § 54 d'Euler
(« V. N.o 54 de son M. sur la lame »), qui, prenant $s = a$, $s = 2a$, $s = 3a$,
exclut tout autre cas que les extrémités appuyées : ces suppositions, sauf la
première, sont « purement mathématiques ». Vue 584 : la mesure doit se prendre
sur le contour extérieur, la largeur n'influant pas sur le son, selon elle. La
condition de la seconde espèce de « libre » est écrite « $\frac{dz}{dx}=0$ et
$\frac{d^3z}{dx^3}=0$ ».} Un anneau de verre de 83 mm de diamètre et 9 mm de
large donne trois figures : trois nœuds inégaux, si 1 ; quatre nœuds égaux,
mi 2 ; six nœuds égaux, fa$^{\sharp}$3. Elle rapporte les deux dernières aux
extrémités appuyées, de sons $16 : 36$, soit $\frac94$, une octave et un ton,
« précisément comme l'indique l'expérience » ; la première aux extrémités
fixées, deux nœuds intérieurs de la lame encastrée ayant le nombre
$\frac{49}{4}$, et $16 : \frac{49}{4} = \frac{64}{49}$ (462 cents) étant voisin
de la quarte (498 cents) de si 1 à mi 2.\footnote{Vues 585--586, ses pages (158)
et (159). « La fi » est une amorce reprise à la ligne suivante ; le chiffre
d'octave de fa$^{\sharp}$3 est lu 3 avec doute (il ressemble à un 9). Vue 586 :
« est $\frac{49}{4}:\frac{64}{49}$ doit donc exprimer » — un mot semble manquer.
Le nombre $\frac{49}{4}$ est bien celui d'Euler pour la lame encastrée aux deux
bouts à deux nœuds intérieurs, $(10{,}996/\pi)^2 = 12{,}25$.} L'arithmétique est
juste ; mais la loi actuelle de l'anneau met entre quatre et six nœuds un
rapport voisin de $2{,}8$ (1800 cents), non $\frac94$ (1404 cents). Tel que la
transcription le lit, son anneau de 9 mm contredirait cette loi d'une tierce
majeure environ ; la lecture ne tranche pas entre une note mal lue et un autre
mode de vibration. Elle conclut, « conformément à la doctrine d'Euler », que la
courbure n'influe pas sur les manières de vibrer, que les écarts de Chladni
viennent de la variation de l'élasticité, et qu'Euler a eu tort d'exclure pour
l'anneau les extrémités fixées ; elle n'a pas vu de mode qu'elle puisse rapporter
aux extrémités libres, sans l'exclure pour autant.\footnote{Vue 587, sa page
(160), n\textsuperscript{o} 289.} La première conclusion est fausse, pour la raison
dite plus haut.

Les anneaux de 20 mm de large ne donnent pas le cas des extrémités fixées, mais
bien les deux figures des extrémités appuyées, et l'un d'eux une figure à deux
nœuds, que Chladni n'a pas observée sur l'anneau métallique.\footnote{Vues
588--589, ses pages (161) et (162), n\textsuperscript{o} 290. La largeur des
autres anneaux, « 35 » mm, est lue avec doute. La note marginale précise qu'il
ne s'agit que des cas où l'anneau vibre comme une simple lame. La vue 589 donne
les sons : « près du sol plus bas de 2 octaves que l'octave 1 » pour les deux
nœuds, sol$^{\sharp}$1 pour « la figure qui […] se composoit de deux nœuds »,
si 3 pour six nœuds ; le second « deux » est attendu « quatre » par la suite et
par la théorie : à vérifier sur l'image. Tels qu'ils sont lus, de
sol$^{\sharp}$1 à si 3 il y a deux octaves et une tierce mineure (environ
$4{,}8$), loin de $\frac94$ ; avec si 2 on aurait environ $2{,}4$. Elle écrit
que ces sons « ne s'en éloignent pourtant pas » plus que les causes physiques ne
le permettent.} Dans la théorie de l'anneau mince, deux nœuds ($n = 1$) sont un
mouvement rigide ; une figure à deux nœuds sur un anneau aussi large relève de la
plaque annulaire. Enfin une plaque circulaire et un anneau de même diamètre,
percé d'un trou de 11 mm, donnent les mêmes figures et les mêmes sons : quatre
diamètres, ut$^{\sharp}$2 ; six, mi 3 — une octave et une tierce mineure
(1500 cents), un demi-ton de plus que $\frac94$, écart qu'elle juge suffisamment
petit, d'autres essais ayant donné l'intervalle juste ou trop grand d'un ton.
Ici la plaque de Kirchhoff donne 1458 cents entre deux et trois diamètres :
l'observation est juste, et plus près de la théorie actuelle que de la
sienne.\footnote{Vues 589--590, ses pages (162) et (163). Après « 11 », un signe
barré n'est pas lu. Avec un trou de 11 mm sur 83, la plaque percée diffère peu de
la plaque pleine ; le calcul est fait pour la plaque pleine, $\nu = 0{,}3$.}

\emph{La conclusion du mémoire.} De ce que l'équation de l'anneau
$z = \dots\,\mathcal{A}\sin(\beta\varphi+B)$, qui est celle de la corde, convient
aussi à la plaque circulaire, elle étend la même forme, en négligeant la
variation du facteur radial, à la lame droite, au tambour circulaire et à la
cloche hémisphérique, et en conclut que dans ceux de ces corps qui sont
élastiques les sons sont comme $\beta^2$ ; Chladni trouve en effet, pour les
plaques circulaires à diamètres seuls, les carrés de 2, 3, 4, 5, 6 (table p. 185)
et, pour la cloche d'harmonica, des sons comme les carrés de la même
suite.\footnote{Vues 590--591, ses pages (163) et (164), n\textsuperscript{o}
291. « au tambour circulaire » est ajouté sous la ligne ; « 185 » est lu avec
doute (le 5 peut être un 9) ; les deux tables sont sans doute celles du Traité
d'acoustique, ce que la transcription donne comme inférence.} Pour la plaque, la
loi observée est approximativement celle de la théorie actuelle, on l'a vu, mais
elle ne découle pas de son équation, où le facteur radial ne peut être négligé ;
pour l'anneau, elle est fausse ; pour la cloche, le feuillet ne la démontre
pas.

Le mémoire finit sur ces mots : « Je termine enfin ce long mémoire dans lequel je
n'ai pourtant examiné qu'une partie des différens cas que l'expérience présente
et des cas plus nombreux encore que la théorie montre possibles », et le vœu
d'avoir rempli « au gré de mes juges les conditions du programme proposé par la
classe ».\footnote{Vue 592, sa page (165), le premier tiers du feuillet. Trois
lignes et demie barrées : « la classe juge que je ne me suis pas tout à fait
égaré mépris sur la route qu'il falloit tenir », et « dans la carrière ouverte
par cet illustre auteur ». « heureux » est au masculin sur la page.}

\subsection*{Une note sur Fourier, Poisson et Navier}

\pagerange{575}{575}
Reliée entre les pages (149) et (150), sans interrompre la phrase, une note d'un
autre papier, à la première personne, sans doute plus tardive que le
mémoire : sa santé l'a
empêchée d'étudier « le mémoire que Monsieur fourier a eu la complaisance de me
prêter » ; elle veut pourtant lui communiquer une observation. Dans le
« N.o d'aout, le dernier qui ait paru », Poisson répond à Navier que
« l'hypothese n'est point admissible » ; il semble ne connaître « cette
hypothese que j'ai donné en 1811 » que par la citation de Navier. Or Poisson
tient l'hypothèse pour inapplicable à la plaque circulaire dont le contour est
soumis à des forces normales ; elle a donc cherché, au § VIII de son mémoire,
p. 189, où il traite de la plaque circulaire, une équation qu'on ne pourrait
déduire de l'hypothèse, et trouve p. 190 $\Delta z = \varphi$ et p. 191
$\Delta\varphi = \varphi_{rr} + \frac1r\varphi_r$ : ces deux équations
composent l'opérateur $\Delta\Delta$, celui même de l'équation que son
hypothèse donne.\footnote{Vue 575, n\textsuperscript{o} 283, sans pagination ; écriture
plus rapide et plus lourde que celle des pages qui l'entourent ; dos blanc (vue
576). Le numéro de page « 438 » est lu avec doute (« b 38 »), ainsi que
« traiteroit ». Une ligne serrée dans l'interligne : « il dira qu'il a entendu tel
ou tel restriction ». « absurde » est barré devant « inadmissible ». À la page
191, le feuillet écrit $\frac{d\varphi}{dy^2}$ sans exposant au $d$. Que le
périodique soit les \emph{Annales de chimie et de physique} et l'échange celui de
Poisson et de Navier en 1828 est l'inférence de la transcription, que rien sur
le feuillet n'établit. Que la note soit plus tardive que le mémoire est
l'inférence de cette lecture : elle cite l'hypothèse « de 1811 » et un numéro
« le dernier qui ait paru » d'un périodique où Poisson répond à Navier.
Sur l'entrée de 1811, voir la note de contexte au § 1 du long mémoire (vues
399 à 404). Le feuillet ne dit pas quel mémoire de Fourier lui a été prêté, et
cette lecture ne l'identifie pas.} La note s'arrête sur ces deux
équations, sans en tirer la conclusion par écrit.

\section*{Une seconde rédaction de la comparaison avec Chladni}

\pagerange{600}{613}
Les pages (61) à (72), de sa main régulière de mise au net, reprennent presque
mot pour mot les pages (76) à (93) du long mémoire (vues 486 à 503) : la fin de
l'intégrale (G), les intégrales (H) et (I), les figures 79\,b et 86, le
N\textsuperscript{o} 16 et sa difficulté.\footnote{Vues 600 à 613, ses pages (61)
à (72), n\textsuperscript{os} 295 à 300 ; les vues 611 et 612 sont une seconde
prise des vues 609 et 610. Le premier chiffre de (61) est pris dans un pâté ;
(62) à la vue 601 l'appuie. La page (72) s'arrête au milieu d'une phrase ; la
suite n'est pas dans le volume à cette place. Les vues 486 à 490 sont traitées
dans la première partie du mémoire.} Ce n'est pas une simple copie : ces pages
portent des retraits et des corrections qui leur sont propres.

— Les intégrales (H) et (I) sont désignées par deux petites croix,
« (++) » ; deux alinéas sur (H) sont traversés d'un long trait oblique, et
l'alinéa sur (I) est récrit (« renferme une série qui ne peut jamais se
terminer ») ; une note marginale ajoute que les nombres
$(m^2+n^2)(\Pi-1)(\Pi+\Pi'-2)$ se réduisent, pour des valeurs convenables de
$\Pi$, $\Pi'$, à la forme de ceux des intégrales périodiques.\footnote{Vue 600.
Le $\Pi$ est la lettre que la première partie rend ainsi, distincte du
$\pi$ ; ici elle ressemble à « rr ». Le numéro « 796 » (pour 79\,b) est lu avec
doute ; plusieurs mots de la note sont perdus dans la reliure, et l'ordre de ses
derniers mots n'est pas assuré.}

— Les passages sur les figures 79\,b et 86 (vues 601, 602, 605, 606), et sur 79
dans la comparaison à 68 (vue 607), sont traversés de longs traits obliques, et
les lignes correspondantes des tableaux sont barrées (vues 604, 606) : ici elle
retire ces figures, que la première rédaction garde.\footnote{Vues 601 à 607. Au
second facteur de la vue 601, le coefficient de $\cos^4$ est $\frac{4}{3}$ dans
la formule et $\frac{14}{3}$ quatre lignes plus bas ; le calcul
($\Pi = 2-m^2 = -23$, $\frac{(-20)(-8)}{2\cdot3\cdot4\cdot5} = \frac43$) donne
$\frac43$, comme la vue 488. Pour $m = 6$, la vue 601 lit « 31, 31.19, 91.19 »
là où la vue 489 a « 31, 31.19, 31.19 » ; le troisième numérateur est
$(-31)(-19)(1)$, précédé d'un signe moins que ni l'une ni l'autre ne porte. Le
trait des vues 604--606 ne barre aucun mot en particulier.}

— Vue 607 : « de deux octaves et de deux tons » pour 87--88, où la vue 496 dit
« d'une octave » ; seul « deux octaves » est juste. Les qualificatifs des notes
de Chladni diffèrent : « + grave » ici (vues 606, 608) là où la première
rédaction écrit « $-$ aigu » (vues 495, 498).

— Vue 609, pour $5|4$ : « $>$ que $\frac{128}{25}$ ; et en effet […]
$\frac{41}{8}$ est $>$ que $\frac{128}{25}$ ; mais d'une quantité absolument
inappréciable a l'oreille » — exact ($1025/1024$, 2 cents), et c'est mot pour
mot le premier état, barré, de la vue 499. Vue 613, au contraire, la figure 79
« se rapproche […] d'avantage du rapport donné par [la] théorie puisqu'il n'en
differe plus que de l'intervalle d'un ton et $\frac12$ ton », correction écrite
sur un texte barré qui est celui de la vue 502 (79\,b et 86, ut 6 contre ut 5,
36 contre 18).

Chaque rédaction porte ainsi une correction dont le premier état est dans
l'autre. Ces feuillets ne disent pas laquelle des deux est la première, et cette
lecture ne le décide pas.\footnote{Vues 609, 610 et 613. La difficulté du ton
trop grave (vue 610) est donnée dans les mêmes termes qu'aux vues 500--501,
avec l'exemple de $0|4$ et $4|3$ ; la vue 613 s'arrête sur « dans la
comparaison au son des figures ».}

\section*{Le mouvement de la chaleur dans la sphère, le cône et le cylindre}

\pagerange{549}{716}
Une quarantaine de pages du volume ne parlent ni d'élasticité ni de nombres,
mais de la propagation de la chaleur dans les solides. Elles sont reliées en
trois endroits : une feuille isolée au milieu du long mémoire sur les surfaces
élastiques (vue 549), une suite continue de feuillets de calcul et de rédaction
(vues 630 à 701), et deux feuillets pris entre les notes d'arithmétique, juste
après le manuscrit C (vues 712 à 716).\footnote{Vues 549, 630 à 646, 652 à
701, 712 à 716. Série à l'encre : 270 (vue 549), 309 à 344 (vues 630 à 700),
350 à 352 (vues 712 à 716). La vue 650 (feuillet 319), sur la somme des
courbures principales de la sphère, est reliée au milieu de ces calculs mais
n'en fait pas partie : elle est lue avec les feuillets sur la courbure
(section « Feuillets épars sur la courbure »). Les feuillets 712 à 716 sont
reliés entre le manuscrit C (vues 708 à 710) et les notes sur les résidus
quadratiques (vues 718 à 730), sans rapport de sujet avec eux. Aucun de ces
feuillets n'est daté. Ils citent la « théorie de la chaleur » de Fourier par
pages et par articles ; si c'est la \emph{Théorie analytique de la chaleur},
parue en 1822 (contexte, non tiré des feuillets), ils sont postérieurs à cette
date : inférence de cette lecture, que la lecture ne vérifie pas contre
l'ouvrage. Une note reliée dans le long mémoire sur les surfaces élastiques
(vue 575) parle d'un mémoire que Fourier lui a prêté, et de Poisson ; elle est
lue avec ce mémoire (« Une note sur Fourier, Poisson et Navier »).} Le problème est celui que Fourier avait résolu pour la sphère dont
chaque couche concentrique est également échauffée, et pour le cylindre dont
chaque couche cylindrique l'est aussi : elle le reprend lorsque le solide a
été « échauffé d'une manière entièrement arbitraire », c'est-à-dire lorsque la
température dépend aussi des angles. Elle cite Fourier par pages (110, 141,
277, 348) et par articles (104, 111, 119, 120, 245, 306 à 308), nomme Laplace
et Poisson, et compare sa manière à celle de Poisson.

Les notations sont celles de la théorie qu'elle cite, et les siennes : $v$ la
température (comptée à partir de celle de l'air), $t$ le temps, $K$ la
conductibilité intérieure, $CD$ le produit de la capacité de chaleur et de la
densité — elle dit elle-même que $CD\,r^2\sin\theta\,d\theta\,d\varphi\,dr$
« exprime combien il faut de chaleur pour élever de 0 à 1 la molécule » —, $h$
la conductibilité extérieure, $k = K/CD$.\footnote{La lettre lue $D$ est, sur
tous ces feuillets, une capitale en forme de dôme posé sur une barre ; les
transcriptions la lisent $D$ d'après la suite (A), (B), (C), (D) des équations,
et notent qu'on pourrait y voir un $\Omega$ (vues 549, 636, 646, 652, 676).
Elle écrit tantôt $C.D$, tantôt $c.D$ ; on écrit $CD$. Tout au long, les
dérivées partielles sont écrites avec un $d$ droit : le feuillet écrit
$\frac{d^2v}{dr^2}$, $\frac{dv}{d\varphi}$ ; on écrit
$\frac{\partial^2 v}{\partial r^2}$, $\frac{\partial v}{\partial \varphi}$.
Ses « sin », « cos », « tang », « cot », avec ou sans point, sont rendus
$\sin$, $\cos$, $\tan$, $\cot$ ; son point de multiplication est omis.} Dans
tout solide homogène, l'équation du mouvement varié de la chaleur est
\[
\frac{\partial v}{\partial t} = k\,\Delta v, \qquad
\Delta v = \frac{\partial^2 v}{\partial x^2} + \frac{\partial^2 v}{\partial y^2}
+ \frac{\partial^2 v}{\partial z^2},
\]
et à la surface, où la chaleur passe dans l'air, la condition est
\[
\frac{\partial v}{\partial n} + \frac{h}{K}\,v = 0 ,
\]
$n$ étant la normale extérieure.\footnote{Elle écrit la condition de surface
sous la forme générale $m\frac{dv}{dx} + n\frac{dv}{dy} + p\frac{dv}{dz} +
\frac{h}{k}vq = 0$, $(m,n,p)$ étant un vecteur normal et $q$ sa longueur, et la
cite comme l'équation (B) de la « page 141 » de la théorie de la chaleur (vue
662). Le rapport y est écrit $\frac{h}{k}$ avec un petit $k$ aux vues 549, 662,
672 et 676, et $\frac{h}{K}$ avec une capitale dans la copie (vue 693) ; à la
vue 678 elle convient d'appeler $k$ le rapport $K/CD$ et $h$ le rapport
$\frac{h}{k}$ de la condition de surface. Pour la condition de Newton à la
surface, le rapport est celui de la conductibilité extérieure à la
conductibilité intérieure : on écrit partout $h/K$.} Tout le travail de ces
feuillets consiste à écrire $\Delta$ dans des coordonnées adaptées à la figure
— sphériques $r, \theta, \varphi$ pour la sphère et le cône, polaires $r,
\varphi$ pour le cylindre —, à retrouver la même équation par le bilan de
chaleur d'une molécule (« en traitant la question d'une manière spéciale »,
comme Fourier), puis à intégrer. L'angle $\theta$ est compté depuis le pôle
(elle le dit « l'arc compris entre le pôle et un point quelconque de la
sphère », vue 652), $\varphi$ est la longitude ; on note
\[
\Delta_S v = \frac{1}{\sin\theta}\frac{\partial}{\partial\theta}\Big(\sin\theta\,
\frac{\partial v}{\partial\theta}\Big) + \frac{1}{\sin^2\theta}\,
\frac{\partial^2 v}{\partial\varphi^2}
\]
la partie angulaire du laplacien.\footnote{Aux vues 663 et 664 elle appelle
« latitude » ce même angle $\theta$ compté depuis le pôle ; c'est la
colatitude. La notation $\Delta_S$ est de la lecture.}

\subsection*{Un brouillon sur le cylindre, égaré dans le mémoire sur les plaques}

\pagerange{549}{549}
La feuille de la vue 549, reliée entre deux pages du long mémoire sur les
surfaces élastiques, est un brouillon de l'article 3 de la pièce sur le
cylindre (vues 676 à 678) : mêmes phrases, dans le même ordre, avec plus de
ratures.\footnote{Vue 549. Feuille de brouillon d'un autre papier que la mise au
net qui l'entoure, main rapide, deux encres, très raturée ; cachet de la
Bibliothèque impériale dans la marge ; série à l'encre 270. La transcription
ne décide pas si l'argument est le sien ou des notes sur le mémoire d'un autre ;
la comparaison avec les vues 676 à 678, qui reprennent ces phrases presque mot
pour mot dans une pièce qui est la sienne, est l'argument de la présente
lecture pour y voir son brouillon (inférence de cette lecture). En tête, d'une
encre pâle, une ligne isolée, interrompue, « $z = \dots\{e^{\frac{\omega
x}{a}} - e^{\frac{(a-\omega)}{a}x} + (1$ », qui rappelle les fonctions de la
lame élastique et n'a pas de rapport visible avec la suite. Dans la marge
gauche, des essais de plume (« tranche », « on voit »).} Après avoir « effacé
ce qui se détruit », la condition de surface se réduit à
\[
\frac{\partial v}{\partial r} + \frac{h}{K}\,v = 0 \qquad \text{(D)} ;
\]
on doit parvenir aux équations (B) et (D) « en traitant la question d'une
manière spéciale » ; le mouvement de la chaleur a lieu à la fois dans deux
directions ; le flux entre les couches de rayons $r$ et $r + dr$ fournit dans
l'équation (C) les deux termes $\frac{K}{CD}\big(\frac{\partial^2 v}{\partial
r^2} + \frac1r \frac{\partial v}{\partial r}\big)$, et le flux entre les
sections placées aux distances $r\varphi$ et $r\varphi + r\,d\varphi$ fournit le
troisième, $\frac{K}{CD\,r^2}\frac{\partial^2 v}{\partial\varphi^2}$.\footnote{Le
feuillet écrit ce troisième terme $\frac{K}{CDrr}\frac{dv}{d\varphi}$ ; un « 2 »
isolé au-dessus de la ligne semble l'exposant oublié de la dérivée seconde,
que le calcul demande. Le diviseur de la réduction, ajouté au-dessus de la
ligne, est lu avec doute « $+\frac{dv}{dr}$ ». Un mot de l'addition sur la
« couche » du cylindre est illisible, ainsi qu'un mot barré et un petit mot
abrégé au-dessus de « même ».} Ces termes sont exacts. Elle renvoie pour les
deux flux à « l'analyse employée N\textsuperscript{o} 19 » et
« N\textsuperscript{o} 14 » ; la rédaction des vues 676 à 678 dit
« N\textsuperscript{o} 119 (T. de la chaleur) » et
« N\textsuperscript{o} 104 ».\footnote{« 19 » et « 14 » sont lus avec doute,
le second dans une addition interlinéaire et dans une ligne barrée. La
lecture ne décide pas si ce sont des abréviations des mêmes articles ou des
renvois à une autre numérotation.}

\subsection*{Calculs : le cône, la sphère, le cylindre, l'anneau}

\pagerange{630}{646}
Ces feuillets de calcul ne sont pas reliés dans l'ordre de sa pagination, et
la lecture ne le rétablit que là où une phrase passe d'une page à
l'autre.\footnote{Vues 630 à 646. Série à l'encre 309 à 317. Sa pagination,
petite, souvent barrée : 7 (vue 632), 8 (vue 635), 7 (vue 636), 6 (vue 639) ;
aucune aux vues 630, 640, 643 à 646. Deux pages portent 7. La phrase qui finit
la vue 639 (« au lieu d'être comprise entre les deux calotte sphériques ») se
poursuit en tête de la vue 636, dont le verso (vue 637) achève la note
marginale annoncée par « Retournez » : on lit donc 639, 636, 637. La vue 635
s'arrête sur « l'an » coupé et n'a pas de suite reconnue. Sur la vue 639, un
« 312 » barré : ce feuillet a été numéroté deux fois.}

\textbf{Le cône en variable $\sin\theta$ (vue 630).} Sous le titre « cône »,
elle pose $x = r\sin\theta\cos\varphi$, $y = r\sin\theta\sin\varphi$, $z =
r\cos\theta$, et annonce qu'il faut différentier « non par rapport à l'arc
$\theta$, mais par rapport au sinus de cet angle », parce que l'arc $\theta$
« qui mesure la calotte sphérique est remplacé ici par $\sin\theta$ qui mesure
la base du cône ».\footnote{Vue 630, sans pagination de sa main. Le feuillet
écrit $z = \cos\theta$, sans le facteur $r$ que la différentielle $dz = dr\cos
\theta - r\sin\theta\,d\theta$ de la ligne suivante suppose. Un groupe noirci
en fin de ligne est illisible ; le chiffre de l'addition qui le remplace est lu
2, peut-être 3. La page s'arrête au bas du feuillet sur les termes en
$\frac{dv}{d\varphi}$ de $\frac{d^2v}{dx^2}$.} Avec $\sigma = \sin\theta$, ses
dérivées $\frac{\partial\sigma}{\partial x} = \frac{\cos^2\theta\cos\varphi}{r}$,
$\frac{\partial\sigma}{\partial y} = \frac{\cos^2\theta\sin\varphi}{r}$,
$\frac{\partial\sigma}{\partial z} = -\frac{\sin\theta\cos\theta}{r}$ sont
exactes, et le développement de $\frac{\partial^2 v}{\partial x^2}$ l'est
aussi, terme à terme.\footnote{La lecture a refait le calcul (dossier de
vérification). Les deux exposants que la transcription donne avec doute — le
$4$ de $\cos^4\theta$ devant $\frac{\partial^2 v}{\partial\sigma^2}$ et le $2$
du dernier terme de l'accolade de $\frac{\partial v}{\partial\sigma}$ — sont
ceux que le calcul demande.} Le calcul complet, repris au net aux vues 712 et
716, aboutit à
\[
\Delta v = \frac{\partial^2 v}{\partial r^2} + \frac2r\frac{\partial v}{\partial r}
+ \frac{\cos^2\theta}{r^2}\frac{\partial^2 v}{\partial\sigma^2}
+ \frac{1}{r^2}\Big(\frac{\cos^2\theta}{\sin\theta} - \sin\theta\Big)
\frac{\partial v}{\partial\sigma}
+ \frac{1}{r^2\sin^2\theta}\frac{\partial^2 v}{\partial\varphi^2},
\]
qui est exact. Ce n'est pas une autre équation que celle de la sphère : c'est
le même laplacien, où l'on a pris $\sigma = \sin\theta$ pour variable au lieu
de $\theta$ (avec $\partial_\theta = \cos\theta\,\partial_\sigma$). L'équation
intérieure de la chaleur est la même dans tout solide ; ce qui distingue la
sphère, le cône terminé par une calotte et le cône terminé par une base plane,
ce sont la surface et la condition qu'on y écrit, non l'équation.

\textbf{Le cône opposé (vue 632).} Les plans menés par le sommet
perpendiculairement aux côtés du cône enveloppent un second cône, « opposé » ;
si $A$ est l'angle du côté avec l'axe et $A'$ celui du côté du cône opposé,
elle note en marge « $A + A' = \frac\pi2$ », ce qui est exact : le cône polaire
d'un cône de révolution de demi-angle $A$ a pour demi-angle $\frac\pi2 -
A$.\footnote{Vue 632, pagination 7 barrée, série à l'encre 310 ; de petits
dessins de fleurs et de feuilles dans la marge gauche. Les angles des
intersections avec les plans $zx$ et $xy$ sont écrits
$\frac{A\,v}{\frac{\pi}{2}}$ et $\frac\pi2 - \frac{A\,v}{\frac\pi2}$ ; ce que
désigne la lettre lue $v$ n'est pas dit sur la page, et la lecture ne peut
juger ces expressions. Sont lus avec doute un $y$, un $z$, « cone » et
« fera ». La page s'arrête aux deux tiers.}

\textbf{Le cylindre, l'anneau, la plaque (vue 635).} Dans le cylindre dont la
base est dans le plan des $yz$, l'élément d'arc est celui du cercle de base ;
« pour avoir le cas de l'anneau il suffit donc de faire $x$ constant dans
l'équation du cylindre ». Dans le cône, l'élément d'arc du plan des $yz$ est
incliné sur la base ; à la limite du cône « infiniment surbaissé », qui est la
plaque circulaire, il devient l'élément du rayon.\footnote{Vue 635, sa
pagination 8, sans numéro de la série à l'encre sur cette face (le feuillet
porte 311 sur l'autre face, blanche, vue 634). « surbaissé » et la lettre $r$
sont lus avec doute. La page s'arrête sur « $A$ étant l'an ».} Supprimer la
variation le long de l'axe donne l'équation de la plaque ; l'anneau mince de
Fourier (l'« armille ») demande en outre qu'on supprime la variation le long du
rayon, ce qu'elle dit à la vue 691 (« le cas de l'armille correspond à celui
où on ferait le rayon constant »). Faire $x$ constant seul ne suffit pas.

\textbf{La molécule de la sphère et celle du cône (vues 639, 636, 637).} La
molécule de la sphère a pour côtés $dr$, $r\sin\theta\,d\varphi$, $r\,d\theta$,
pour faces $r^2\sin\theta\,d\varphi\,d\theta$, $r\sin\theta\,d\varphi\,dr$,
$r\,d\theta\,dr$ ; en retranchant la chaleur qui entre par trois faces de celle
qui sort par les faces opposées, elle obtient
\[
\frac{\partial v}{\partial t} = k\Big(\frac{\partial^2 v}{\partial r^2}
+ \frac2r\frac{\partial v}{\partial r} + \frac{1}{r^2}\frac{\partial^2
v}{\partial\theta^2} + \frac{\cos\theta}{r^2\sin\theta}\frac{\partial
v}{\partial\theta} + \frac{1}{r^2\sin^2\theta}\frac{\partial^2
v}{\partial\varphi^2}\Big),
\]
qui est l'équation exacte.\footnote{Vue 639, sa pagination 6 barrée. Le
feuillet écrit $\frac{dv}{dr}$ au premier membre de la deuxième ligne, là où le
calcul attend $\frac{dv}{dt}$ ; le numérateur de la première fraction est une
tache (le sens attend 1) ; le dernier terme porte un 2 au-dessus du $d$ du
numérateur et aucun exposant au dénominateur ; un dénominateur de la troisième
accolade s'arrête sur « $r\sin\theta$ », sans $d\varphi$.} Puis, « au lieu de
considérer le solide conique qui s'appuie sur la calotte sphérique », elle
veut « s'occuper du simple cône » dont la base est plane : sa molécule est
comprise entre deux plans et non entre deux calottes, et elle lui donne pour
côtés $dr$, $r\sin\theta\,d\varphi$, $r\cos\theta\,d\theta$, le flux à
travers la face $r\sin\theta\,d\varphi\,dr$ étant multiplié par $\cos\theta$
« parce que la direction » du flux « fait l'angle $\theta$ avec la
perpendiculaire au plan de cette face ». Le bilan donne
\[
\frac{\partial v}{\partial t} = k\Big(\frac{1}{r^2}\frac{\partial}{\partial
r}\Big(r^2\frac{\partial v}{\partial r}\Big) + \frac{1}{r^2\sin\theta\cos\theta}
\frac{\partial}{\partial\theta}\Big(\sin\theta\cos\theta\,\frac{\partial
v}{\partial\sigma}\Big) + \frac{1}{r^2\sin^2\theta}\frac{\partial^2
v}{\partial\varphi^2}\Big),
\]
dont le terme du milieu vaut $\frac{\cos\theta}{r^2}\frac{\partial^2
v}{\partial\sigma^2} + \frac{1}{r^2}\big(\frac{\cos\theta}{\sin\theta} -
\frac{\sin\theta}{\cos\theta}\big)\frac{\partial v}{\partial\sigma}$.\footnote{Vue
636, sa pagination 7 barrée, série à l'encre 312. En développant, le feuillet
écrit $\frac{\cos^2\theta}{r^2}$ devant $\frac{d^2v}{(d\,\sin\theta)^2}$ ; la
dérivation de son propre premier membre donne $\frac{\cos\theta}{r^2}$, et
c'est avec $\cos\theta$ que la remarque qui suit est exacte. La variable de la
seconde différentiation (« la seconde de même par rapport à … ») est
illisible ; le sens demande $\sin\theta$. Une parenthèse fermante sans
ouvrante et un $\cos\theta$ barré dans la deuxième accolade ; « obliquement »
et « et sort » sont lus avec doute.} Elle constate que cette équation diffère
de celle qu'elle a « trouvée d'une autre manière » (le changement de variable
de la vue 630) par un facteur $\cos\theta$ sur les deux termes en $\sigma$, et
elle l'explique : la molécule traversée obliquement n'a pas la « solidité » de
la molécule traversée perpendiculairement, les deux sont « $::\cos\theta : 1$ »,
et en corrigeant de ce rapport « les équations seront identiques ». La
conclusion est juste, et la raison peut se dire plus simplement : le côté de
la molécule le long du méridien est $r\,d\theta = r\,d\sigma/\cos\theta$, non
$r\cos\theta\,d\theta$, et son volume est $r^2\sin\theta\,dr\,d\theta\,d\varphi$
; avec ce volume, le bilan redonne exactement l'équation de la vue 630. Sa
correction d'« obliquité » rétablit ce volume ; il n'y a pas, pour le cône,
d'équation intérieure différente de celle de la sphère.\footnote{Vue 637, verso
de la vue 636, sans numéro : la note marginale du recto s'y achève (« ne mesure
plus la quantité de chaleur accumulée dans la molécule »), et le texte finit
sur « seront identiques ». Le « 8 » visible en tête des vues 637 et 639 est sur
une bande de papier qui dépasse au-dessus et paraît être le haut de la vue 635.}

\textbf{« Recherches sur les cones » (vue 640).} La page s'ouvre sur
$\big(\frac{\partial z}{\partial x}\big)^2$ et $\frac{\partial^2 z}{\partial
y^2}$ en $\theta, \varphi$, le calcul même de la feuille sur la courbure de la
sphère reliée plus loin (vue 650).\footnote{Vue 640, série à l'encre 314, sans
pagination de sa main. Elle écrit $v$ dans deux numérateurs de
$\frac{d^2z}{dy^2}$ où les autres portent $z$ ; un facteur couvert d'une tache
dans la dernière accolade est illisible. Pour la vue 650, voir la section
« Feuillets épars sur la courbure ».} Sous le titre souligné, elle cherche la
face de l'élément du cône parallèle à la base : une portion de secteur entre
deux cercles concentriques de rayons $r\sin\theta$ et $r'\sin(\theta -
d\theta)$, avec $r' = r\cos\theta/\cos(\theta - d\theta)$ pour rester dans le
même plan. La largeur est $r\,d\theta/\cos\theta$, exacte au premier ordre ;
l'aire est donc $r^2\sin\theta\,d\theta\,d\varphi/\cos\theta$.\footnote{Le
feuillet donne « $\frac{r^2}{2}\frac{d\varphi\,d\theta\sin\theta}{\cos\theta}$
pour la surface cherchée » : un facteur $\frac12$ de trop, qu'on s'explique si
l'on prend la moitié de $\rho_1^2 - \rho_2^2$ sans le facteur 2 de
$\rho_1^2 - \rho_2^2 \approx 2\rho\,d\rho$ (explication de la lecture). Le
numérateur lu $dv$ à la dernière ligne est empâté.} La page s'arrête sur
l'énumération des autres faces, avant toute équation.

\textbf{Un axe incliné (vues 643, 644).} Le petit feuillet 315 change de
variables dans $v$ avec $x = r\sin A\cot\theta\cos\varphi$, $y = r\sin
A\cot\theta\sin\varphi$, $dz = \sin A\,dr$ : $z$ est proportionnel à $r$, et
$\cot\theta$ est le rapport de la distance à l'axe à la hauteur $z$ ; ce sont
des coordonnées adaptées à des sections planes parallèles à la base. Ses
dérivées $\frac{\partial r}{\partial z} = \frac{1}{\sin A}$,
$\frac{\partial\theta}{\partial x} = -\frac{\cos\varphi\sin^2\theta}{r\sin A}$,
$\frac{\partial\theta}{\partial z} = \frac{\sin\theta\cos\theta}{r\sin A}$ sont
exactes ; celles de $\varphi$ ont le signe opposé à celui que donnent les
formules de $x$ et $y$, parce que son $dx$ porte $+r\sin A\cot\theta\sin\varphi
\,d\varphi$ là où il faut $-$ (et inversement pour $dy$).\footnote{Vue 643, petit
feuillet monté ; le nombre 315 est sur l'autre face, blanche (vue 642). Le
signe entre les deux termes de la première expression est couvert d'une tache
; elle écrit une fois $\frac{dy}{d\varphi}$ pour $\frac{dv}{d\varphi}$ ;
l'exposant d'un $\cos\theta$ est lu 3 avec doute ; la dernière ligne se perd
sous la garde. Les deux premières lignes, en $\cos^2(A-\theta)$ et
$\tan(A-\theta)$, relèvent, semble-t-il, d'un autre choix de l'angle ; la
lecture ne les a pas vérifiées.} Le feuillet 316 (vue 644) poursuit avec un angle $\theta'$
tel que $\sin\theta' = \cos\theta$ et un rayon $r'$, une « hypothèse »
$\frac{1}{\cos^2\theta}$ entourée d'un trait, des réductions partielles, et
s'arrête sur « Il faut encore ajouter à … » et une ligne barrée : aucune
équation n'est menée à son terme.\footnote{Vue 644, série à l'encre 316. Huit
signes ou groupes y sont illisibles (taches, exposants, un dénominateur barré) ;
« hypothese » (première occurrence) est lu avec doute ; un mot écrit de bas en
haut entre les colonnes est lu avec doute « Somme ». Un grand 4 isolé et un
essai de plume au bas de la feuille.}

\textbf{Le cylindre échauffé arbitrairement, avec son axe (vue 645).} Au dos du
feuillet 316, sous le titre « cylindre » : si le cylindre a été « échauffé
arbitrairement », on ne peut plus supposer $\frac{\partial^2 v}{\partial z^2} =
0$, et l'équation pour $u$ (avec $v = e^{-mt}u$) devient
\[
\frac{m}{k}\,u + \frac{\partial^2 u}{\partial r^2} + \frac1r\frac{\partial
u}{\partial r} + \frac{1}{r^2}\frac{\partial^2 u}{\partial\varphi^2} +
\frac{\partial^2 u}{\partial z^2} = 0,
\]
ce qui est exact. Elle propose la solution particulière $u = \sin\delta\varphi
\;r^\beta s\,\sin\frac{\nu z}{\rho}$ avec $\beta = \sqrt{\delta^2 +
\nu^2/\rho^2}$. Ce choix ne convient pas : en substituant, le facteur $\sin
\frac{\nu z}{\rho}$ ne modifie pas l'ordre de la fonction du rayon mais la
constante, et l'on a
\[
u = \sin\delta\varphi\;\sin\frac{\nu z}{\rho}\;r^{\delta}\,s(r),
\]
où $s$ est la série de la pièce sur le cylindre (voir plus bas) d'indice
$n = 2\delta + 1$, avec $g = \frac{m}{k} - \frac{\nu^2}{\rho^2}$ au lieu de
$g = \frac mk$.\footnote{Avec $\beta^2 = \delta^2 + \nu^2/\rho^2$, il reste
dans l'équation de $s$ le terme $\frac{\nu^2}{\rho^2 r^2}s$, qui ne se détruit
pas (calcul refait dans le dossier de vérification) ; la formule du feuillet
n'est d'ailleurs pas homogène si $\rho$ est une longueur. Vue 645, dos du
feuillet 316, sans numéro. « entierement » est écrit au-dessus
d'« arbitrairement », non barré. La lettre lue $\nu$ est un $v$ bouclé ;
celle lue $\rho$ est suivie d'une tache. La phrase s'arrête sur « on
satisfera à l'équation ». Elle renvoie à « p. 3 » et « p. 4 » : la valeur $u =
\sin\beta\varphi\,r^\beta s$ est à la page 4 du brouillon du cylindre (vue 678,
chiffre lu avec doute), ce qui s'accorde avec « la valeur de $u$ p. 4 » ;
« p. 3 » n'est pas identifié (inférence de cette lecture).}

\textbf{La molécule de la sphère, au net (vue 646).} « On parviendra à la même
équation » par la molécule comprise entre deux cônes de sommet au centre, deux
plans méridiens et deux calottes : côtés $r\,d\theta$, $r\sin\theta\,d\varphi$,
$dr$, faces $r^2\sin\theta\,d\theta\,d\varphi$, $r\,d\theta\,dr$,
$r\sin\theta\,d\varphi\,dr$, volume $r^2\sin\theta\,d\theta\,d\varphi\,dr$ ; la
chaleur accumulée pendant $dt$ est
$\frac{K}{r^2}\big(\partial_r(r^2\partial_r v) + \frac{1}{\sin\theta}
\partial_\theta(\sin\theta\,\partial_\theta v) + \frac{1}{\sin^2\theta}
\partial_\varphi^2 v\big)\,r^2\sin\theta\,d\theta\,d\varphi\,dr\,dt$, et en
divisant par $CD$ fois le volume on a l'accroissement de température. Tout est
exact ; le second membre de « $\frac{dv}{dt} =$ » est laissé en
blanc.\footnote{Vue 646, série à l'encre 317. Un signe entre $d\varphi$ et $dr$
et un petit mot noirci sont illisibles, ainsi qu'un mot parmi les essais de
plume de la marge (« S », « Sphere ») ; « des » est lu avec doute.}

\subsection*{« … dans une sphère solide échauffée d'une manière entièrement arbitraire »}

\pagerange{652}{670}
La pièce qui porte ce titre se lit dans l'ordre des vues 652, 654, 660, puis
662 à 670 ; une « Note » (vue 656) et un feuillet blanc (vue 658) sont reliés
entre ses deuxième et troisième pages.\footnote{Vues 652 à 670. Série à
l'encre 320 à 329. Les marques de page sont minces : un « 1 » suivi d'une
petite croix (vue 652), un chiffre barré lu 2 avec doute (vue 654), rien à la
vue 660, puis de petits chiffres barrés lus 4 (douteux), 5 (douteux, au verso
663), 6, 7, rien, 9 aux vues 662 à 670. Que la vue 660 soit sa page 3 est une
inférence : la Note de la vue 656 renvoie à « l'équation (A) p. 3 de la sphère
échauffée d'une manière arbitraire », et (A) est l'équation qui clôt la vue
660. La suite 1 à 9 est proposée par la transcription du lot 34 et n'est pas
une lecture.}

\textbf{L'équation (A) (vues 652, 654, 660).} Pour appliquer l'équation générale
« au cas d'une figure sphérique », elle prend le centre pour origine et,
« en adoptant les dénominations qui appartiennent à la théorie de la terre »,
l'équateur pour plan des $xy$, $\theta$ l'arc du pôle au point, $\varphi$ la
longitude comptée depuis le point du parallèle où $x$ est
maximal.\footnote{Vue 652. La correction de la première phrase est restée
incomplète (« Prenons centre de la sphère origine »).} Elle tire de $dx, dy,
dz$ les combinaisons $dr$, $r\,d\theta$, $r\sin\theta\,d\varphi$, puis les neuf
dérivées de $r, \theta, \varphi$ par rapport à $x, y, z$, puis
$\frac{\partial^2 v}{\partial x^2}$, $\frac{\partial^2 v}{\partial y^2}$,
$\frac{\partial^2 v}{\partial z^2}$ et leur somme ; tout est exact, et
l'équation générale devient
\[
\frac{\partial v}{\partial t} = k\Big(\frac{\partial^2 v}{\partial r^2}
+ \frac2r\frac{\partial v}{\partial r} + \frac{1}{r^2}\frac{\partial^2
v}{\partial\theta^2} + \frac{\cos\theta}{r^2\sin\theta}\frac{\partial
v}{\partial\theta} + \frac{1}{r^2\sin^2\theta}\frac{\partial^2
v}{\partial\varphi^2}\Big) \qquad \text{(A)},
\]
c'est-à-dire $\partial_t v = k\big(\partial_r^2 v + \frac2r\partial_r v +
\frac{1}{r^2}\Delta_S v\big)$, la forme actuelle du laplacien en coordonnées
sphériques.\footnote{Vues 654 et 660. À la vue 654, dans le coefficient de
$r\,d\theta$, le feuillet écrit « $-1$ » là où l'on attend
$-\cos\theta\sin\theta$ ; le résultat, $dr$, est juste. À la vue 660, le
coefficient 2 d'un terme de $\frac{d^2v}{dz^2}$ et le numérateur 1 du dernier
terme de la somme sont écrits sur des taches et lus avec doute ; le calcul
confirme 2 et 1. L'exposant de $\sin\theta$ au dernier dénominateur de
l'accolade de $\frac{dv}{d\theta}$ dans $\frac{d^2v}{dy^2}$ ne se lit pas ; la
lettre de renvoi (A) est récrite et empâtée.}

\textbf{La condition de surface (vues 662, 663).} Dans l'équation (B) de la
surface, la sphère donne $m = x$, $n = y$, $p = z$, $q = r$, et, avec les
valeurs de $\frac{\partial v}{\partial x}$, $\frac{\partial v}{\partial y}$,
$\frac{\partial v}{\partial z}$ « trouvées page 2 », tout se réduit à
$\frac{\partial v}{\partial r} + \frac hK v = 0$ : la même forme que dans le cas
d'une chaleur uniformément répandue dans chaque couche, à cette différence
près que $v$ varie d'un point à l'autre de la surface. Elle en donne la raison
: à la surface, « le mouvement ne peut s'exécuter que dans une seule
direction », que l'intérieur soit traité à une, deux ou trois dimensions. C'est
exact : la condition de Newton ne fait intervenir que la dérivée
normale.\footnote{Vues 662 et 663, série à l'encre 325, écrit des deux côtés.
Elle renvoie aux pages 141 et 110 de la théorie de la chaleur. « page 2 » est lu
avec doute, le chiffre ayant la forme de son $z$ ; il désigne la vue 654, où
sont ces valeurs. Un mot d'une addition très fine (« qui s'… sur ce parallele »)
est illisible.}

\textbf{La même équation, « d'une manière spéciale » (vues 663, 664).} Comme
Fourier à l'article 111 pour la sphère échauffée par couches, elle considère
l'élément compris entre deux cônes de sommet au centre (colatitudes $\theta$ et
$\theta + d\theta$), deux plans méridiens ($\varphi$ et $\varphi + d\varphi$) et
deux sphères ($r$ et $r + dr$), calcule les trois flux et leurs différences, et
retrouve une équation « identiquement la même » que (A) :
\[
\frac{\partial v}{\partial t} = \frac{K}{r^2 CD}\Big(\frac{\partial}{\partial
r}\Big(r^2\frac{\partial v}{\partial r}\Big) + \frac{1}{\sin\theta}
\frac{\partial}{\partial\theta}\Big(\sin\theta\,\frac{\partial
v}{\partial\theta}\Big) + \frac{1}{\sin^2\theta}\frac{\partial^2
v}{\partial\varphi^2}\Big).
\]
L'équation est exacte ; l'élément de surface qu'elle écrit en chemin ne l'est
pas. Elle donne à l'élément de la sphère perpendiculaire à $r$ l'aire
$r^2\sin^2\theta\,d\theta\,d\varphi$, et le justifie dans une note : la partie
du demi-fuseau comprise entre l'équateur et le parallèle de $45^\circ$ serait à
la surface entière « $::\frac12 r^2 : r^2$ ». L'élément est
$r^2\sin\theta\,d\theta\,d\varphi$ ; la zone comprise entre l'équateur et le
parallèle de latitude $45^\circ$ est à l'hémisphère dans le rapport $\sin
45^\circ \approx 0{,}707$, non $\frac12$. Le rapport $\frac12$ est celui des
projections de ces surfaces sur le plan de l'équateur : c'est la projection,
non la zone, qui suit $\sin^2$.\footnote{Vue 664, série à l'encre 326, petit 6
barré. L'exposant 2 de $\sin^2\theta$ est le même à chaque occurrence et la
note le confirme ; la transcription garde ce qui est écrit et signale l'erreur.
Plusieurs lettres, après $K$ et devant $\sin\theta$ dans les expressions de
flux, sont noircies jusqu'à la tache et ne se lisent pas ; un astérisque suivi
de deux parenthèses vides est un renvoi resté sans texte ; le $r^2$ du dernier
rapport de la note est lu avec doute. Elle écrit les rayons des bases des deux
cônes $r\sin\theta$ et $r(\sin\theta + d\theta)$ ; le second est
$r\sin(\theta + d\theta)$.} L'erreur ne passe pas dans l'équation finale, écrite
avec $\sin\theta$.

\textbf{L'intégrale de Fourier, multipliée par des coefficients en $\theta$ et
$\varphi$ (vue 666).} Si $\theta$ et $\varphi$ sont constants, (A) se réduit
à l'équation de Fourier
\[
\frac{\partial v}{\partial t} = k\Big(\frac{\partial^2 v}{\partial r^2} +
\frac2r\frac{\partial v}{\partial r}\Big) \qquad \text{(A)}',
\]
dont l'intégrale, qu'elle cite de la page 348 de la théorie de la chaleur, est
\[
v = \frac1r\Big(a_1 e^{-kn_1^2 t}\sin n_1 r + a_2 e^{-kn_2^2 t}\sin n_2 r +
\cdots\Big),
\]
les $n_i$ étant les racines de $nR\cos nR = \big(1 - \frac{hR}{K}\big)\sin nR$
pour une sphère de rayon $R$.\footnote{Vue 666, série à l'encre 327, petit 7
barré. Le feuillet écrit l'exposant $-Kn_1^2t$ ; l'équation (A)$'$ demande
$-\frac{K}{CD}n_1^2t$. L'équation des $n_i$ n'est pas sur le feuillet ; elle
résulte de la condition de surface appliquée à $\frac{\sin nr}{r}$ (calcul de
la lecture).} Elle remarque que cette intégrale satisfait encore à la condition
de surface si on la multiplie par des coefficients indépendants de $r$ — c'est
exact, la condition ne portant que sur $\partial_r$ — et propose de chercher
des coefficients en $\theta$ et $\varphi$ qui satisfassent en même temps aux
deux termes angulaires de (A). Le programme ne peut aboutir que dans un cas.
Si $v = w(r,t)\,Y(\theta,\varphi)$, (A) donne $Y\big(\partial_t w - k(\partial_r^2
w + \frac2r\partial_r w)\big) = \frac{k w}{r^2}\Delta_S Y$ : avec le $w$ de
Fourier, il faut $\Delta_S Y = 0$. Dans le cas général, $Y$ est une harmonique
sphérique de degré $l$, $\Delta_S Y = -l(l+1)Y$, et la fonction du rayon
change : $\frac{\sin nr}{r}$ est remplacée par la fonction de Bessel sphérique
$j_l(nr)$, et l'équation des $n$ par $nR\,j_l'(nR) + \frac{hR}{K}\,j_l(nR) =
0$.\footnote{Les harmoniques sphériques et les fonctions de Bessel sphériques
sont nommées par la lecture ; le feuillet ne les nomme pas et n'y conduit pas.
Pour $l = 0$, $j_0(nr) = \frac{\sin nr}{nr}$ : c'est le cas de Fourier.}
L'intégrale de Fourier, multipliée par des coefficients angulaires, ne résout
donc pas (A), sauf si ces coefficients annulent $\Delta_S$.

\textbf{« Un cas très remarquable » (vues 668, 670).} C'est précisément ce cas
qu'elle isole : celui où les températures d'une même couche diffèrent d'un
point à l'autre, mais où chacune ne varie qu'« à raison du mouvement de la
chaleur dans le sens des $r$ ». Il faut pour cela que le coefficient de $t$ dans
l'exposant soit indépendant des coefficients en $\theta$ et $\varphi$, et il
suffit que
\[
\frac{1}{\sin\theta}\frac{\partial}{\partial\theta}\Big(\sin\theta\,
\frac{\partial v}{\partial\theta}\Big) + \frac{1}{\sin^2\theta}\,
\frac{\partial^2 v}{\partial\varphi^2} = 0 ,
\]
car (A) se réduit alors à (A)$'$. Son exemple est $v = \frac1r\big(a_1
e^{-kn_1^2t}\sin n_1 r + \cdots\big)\cot\theta\sin\varphi$, et elle vérifie que
$\cot\theta\sin\varphi$ annule les deux termes angulaires ; la vérification
est juste.\footnote{Vue 668, série à l'encre 328 ; vue 670, série à l'encre 329,
petit 9 barré. Une première ligne barrée, « si on prend $2t^2 - 1 = m^2$ »
(le 2 lu avec doute), et une première vérification barrée portent sur
$\cot\theta\sin m\varphi$ ; deux lettres noircies devant $\theta$ et $\varphi$
dans l'exemple sont illisibles, peut-être le $m$ effacé. Or $\Delta_S(\cot\theta
\sin m\varphi) = (1 - m^2)\frac{\cos\theta}{\sin^3\theta}\sin m\varphi$ :
l'essai ne vaut que pour $m = 1$, et c'est $m = 1$ qu'elle garde. La famille
générale des solutions séparées est $\tan^m\frac\theta2\sin m\varphi$,
$\cot^m\frac\theta2\sin m\varphi$ (et de même avec $\cos m\varphi$), et
$\log\tan\frac\theta2$ ; on a $\cot\theta = \frac12\big(\cot\frac\theta2 -
\tan\frac\theta2\big)$. Calculs refaits dans le dossier de vérification.} Elle
en donne le sens physique : la chaleur qu'une molécule reçoit de sa voisine sur
le même parallèle est « exactement compensée par le froid » qui lui vient de sa
voisine sur le même méridien. C'est bien ce que dit $\Delta_S v = 0$, localement.

Mais la fonction $\cot\theta\sin\varphi$ est infinie sur l'axe, en $\theta =
0$ et $\theta = \pi$ : en coordonnées cartésiennes, c'est $\frac{yz}{x^2 +
y^2}$. Une température d'une sphère pleine doit rester finie et régulière dans
tout le solide, et la seule fonction régulière sur toute la sphère qui annule
$\Delta_S$ est la constante. Le « cas très remarquable », pour une température
admissible, se réduit donc au cas de Fourier, $v$ ne dépendant que de $r$ ; la
solution de la vue 670 est une solution de (A) hors de l'axe, non une
température de la sphère entière.\footnote{Au sens des distributions,
$\Delta\frac{yz}{x^2+y^2} = 2\pi z\,\partial_y\delta(x,y)$ : la fonction
représente une ligne de sources et de puits de chaleur le long de l'axe, dont
l'intensité croît avec $z$ ; c'est cette ligne qui entretient la compensation
qu'elle décrit. (Calcul de la lecture : $\frac{y}{x^2+y^2} = \partial_y\log
\rho$ et $\Delta_2\log\rho = 2\pi\delta$.) Ce qui reste constant pendant le
refroidissement, dans ce cas, est le rapport des températures de deux points
d'une même couche, non leur différence, qui décroît avec le facteur radial.} À
la vue 698, elle écrira elle-même que l'annulation des termes angulaires
exprime « la constance des températures » d'une couche, ce qui, énoncé
correctement, est le fait vrai (voir la dernière sous-section). La page 670
s'arrête, sous une tache d'eau, sur la définition de trois températures $v$,
$v'$, $v''$ de molécules voisines.\footnote{Les mots « est $v'$ », « $v''$ »,
« troisième » sont lus avec doute sous la tache ; la page s'arrête sur « sur
le même » et le verso (vue 671) est blanc.}

\textbf{La Note : le cylindre comme limite de la sphère (vue 656).} Dans (A),
elle fait $r$ infini, $\theta$ infiniment petit, et pose $r\sin\theta = s$,
quantité finie ; le terme $\frac2r\frac{\partial v}{\partial r}$ disparaît, et
il reste
\[
\frac{\partial v}{\partial t} = k\Big(\frac{\partial^2 v}{\partial r^2} +
\frac{\partial^2 v}{\partial s^2} + \frac1s\frac{\partial v}{\partial s} +
\frac{1}{s^2}\frac{\partial^2 v}{\partial\varphi^2}\Big),
\]
« l'équation du cylindre », $r$ tenant lieu de la coordonnée le long de
l'axe.\footnote{Vue 656, série à l'encre 322 soulignée, « Note » en tête. La
valeur de $\cos\theta$ est couverte d'une tache ; « $\theta$ étant infiniment
petit » est ajouté dans l'interligne. Un paragraphe et son équation, où elle
substituait $\frac{dv}{d(\sin\theta)}$ à $\frac{dv}{d\theta}$, sont annulés par
deux grands traits en X. « que $x$ » est une seule lettre lue avec doute ;
« ba », en fin de ligne, est un mot commencé et laissé ; la lettre lue $c$ dans
« $c = r + x$ » pourrait être un $r$ mal formé.} Le passage à la limite est
juste, et il peut se dire exactement : si l'on prend pour variables $r$, $s =
r\sin\theta$ et $\varphi$, le laplacien s'écrit, sans approximation,
\[
\Delta v = \frac{\partial^2 v}{\partial r^2} + \frac{\partial^2 v}{\partial s^2}
+ \frac1s\frac{\partial v}{\partial s} + \frac{1}{s^2}\frac{\partial^2
v}{\partial\varphi^2} + \frac{2s}{r}\frac{\partial^2 v}{\partial r\,\partial s}
+ \frac2r\frac{\partial v}{\partial r},
\]
et les deux derniers termes tendent vers zéro quand $r$ croît à $s$ fixé.
Elle juge cette manière « plus simple que celle employée par M\textsuperscript{r}
Poisson car je ne vois pas la nécessité de mettre $r - a$ à la place de $r$ »,
sinon « pour mieux faire sentir » que les ordonnées ont pour origine le point
où $x = 0$ ; cette substitution, dit-elle, « peut avoir lieu après coup ».
C'est exact quant à la forme de l'équation ($\partial_r = \partial_x$ si $r =
a + x$) ; la substitution sert à donner un sens aux points situés à distance
finie, ce qu'elle dit elle-même. Elle termine : on s'étonnerait de voir le
cylindre contenu dans la sphère si l'on ne voyait que le solide conique qui
s'appuie sur la calotte « devient un cylindre » quand le rayon est
infini.\footnote{La lecture n'identifie pas le travail de Poisson dont il est
question ; le feuillet ne le cite pas autrement.}

\subsection*{« … dans un cylindre solide » : brouillon et copie}

\pagerange{672}{695}
La pièce existe en deux états : un brouillon complet de dix pages, articles 1
à 6 (vues 672 à 691), et une copie posée dont le volume ne garde que les pages
1--2 et 7--8 (vues 692 à 695) ; ses pages 3 à 6 ne sont pas dans le
volume.\footnote{Vues 672 à 691, série à l'encre 330 à 339 ; sa pagination,
petite et barrée : rien (672), 2 (674), rien (676), 4 douteux (678), 5 douteux
(680), puis 6 à 10 (682 à 690), que le lot 35 lit nettement et qui confirment
les précédents. Vues 692 à 695, série à l'encre 340 et 341 : 692 sans page,
693 « 2 », 694 « 7 » barré, 695 « 8 ». Le brouillon est d'une écriture plus
droite que la pièce sur la sphère, presque sans rature aux vues 672 à 684, très
corrigé entre les lignes aux vues 686 à 688. Rien sur les feuillets ne dit
lequel des deux états fut écrit le premier ; la copie est la plus nette.} Le
problème est celui du cylindre infini « dont toutes les tranches parallèles à
la base » ont été « échauffées d'une manière entièrement arbitraire, mais
semblable pour chacune d'elles », puis exposé à un courant d'air plus froid.

\textbf{Les équations (A) à (D) (articles 1 à 3, vues 672 à 678).} Les tranches
étant échauffées de même, la chaleur ne se meut pas le long de l'axe,
$\frac{\partial^2 v}{\partial z^2} = 0$, et l'équation se réduit à
$\partial_t v = k(\partial_x^2 v + \partial_y^2 v)$, (A), avec la condition de
surface (B). En coordonnées polaires $x = r\cos\varphi$, $y = r\sin\varphi$,
elle calcule les dérivées et trouve
\[
\frac{\partial v}{\partial t} = k\Big(\frac{\partial^2 v}{\partial r^2} +
\frac1r\frac{\partial v}{\partial r} + \frac{1}{r^2}\frac{\partial^2
v}{\partial\varphi^2}\Big) \quad \text{(C)}, \qquad \frac{\partial v}{\partial
r} + \frac hK v = 0 \ \text{pour}\ r = R \quad \text{(D)},
\]
toutes deux exactes.\footnote{Vues 672 à 676. Dans $\frac{d^2v}{dx^2}$ et
$\frac{d^2v}{dy^2}$, le premier terme est écrit avec $r$ seul au dénominateur,
ce qui est juste ($\frac{\sin^2\varphi}{r}\frac{dv}{dr}$). Une ligne
commencée « En substituant … » est barrée, deux mots y sont lus avec doute ; le
$v$ d'un $\frac{dv}{d\varphi}$ ressemble à un $\varphi$. Le rayon $R$ n'est nommé
qu'à la vue 682.} Elle les retrouve « d'une manière spéciale », par
« l'analyse employée N\textsuperscript{o} 119 (T. de la chaleur) » pour les
deux termes en $r$ et « N\textsuperscript{o} 104 » pour le troisième, et note
que (D) est « absolument la même » que celle du N\textsuperscript{o} 120 pour le
cylindre uniformément échauffé, à ceci près que $v$ varie ici d'un point à
l'autre de la surface.\footnote{Vues 676 et 678. C'est le texte dont la vue 549
est le brouillon (voir plus haut).}

\textbf{L'intégration (article 4, vues 678 à 682).} Elle pose $v = e^{-mt}u$,
d'où $\frac mk u + \partial_r^2 u + \frac1r\partial_r u + \frac{1}{r^2}
\partial_\varphi^2 u = 0$, puis $u = \sin\beta\varphi\;r^\beta s$, « $\beta$
est un nombre entier et $s$ une fonction de $r$ seule », d'où, avec $g =
\frac mk$ et $n = 2\beta + 1$,
\begin{gather*}
\frac{d^2 s}{dr^2} + \frac{n}{r}\frac{ds}{dr} + g\,s = 0, \\
s = 1 - \frac{g r^2}{2(n+1)} + \frac{g^2 r^4}{2\cdot4\,(n+1)(n+3)} -
\frac{g^3 r^6}{2\cdot4\cdot6\,(n+1)(n+3)(n+5)} + \cdots
\end{gather*}
La série est exacte, et elle vérifie qu'elle satisfait à l'équation en
additionnant terme à terme.\footnote{Vues 678 et 680. Dans la seconde des deux
équations à intégrer (vue 678), $v$ manque après $\frac hk$. La condition sur
$s$ est obtenue en divisant par le facteur commun $\sin\beta\varphi\;r^\beta$,
qu'elle omet sans le dire. Au verso de la vue 678 (vue 679), neuf lignes
barrées de deux grands traits, qui ne suivent pas le recto : ce sont, mot pour
mot, les lignes qui ouvrent la vue 684 ; elle les a écrites là d'abord, puis
barrées.} C'est, au facteur près, une fonction de Bessel : $s =
\beta!\,\big(\frac{2}{\sqrt g}\big)^{\beta} r^{-\beta} J_\beta(\sqrt g\,r)$, de
sorte que $u = \text{const}\cdot J_\beta(\sqrt g\,r)\sin\beta\varphi$ ; on peut
mettre $\cos\beta\varphi$ au lieu de $\sin\beta\varphi$, et $e^{-gkt}(A\sin
\beta\varphi + B\cos\beta\varphi)\,r^\beta s$ est une valeur particulière de
$v$.\footnote{Le nom de Bessel n'est pas sur ces feuillets ; il est attaché à
ces fonctions par l'usage ultérieur (contexte, non tiré des feuillets). Les
feuillets ne disent pas si elle connaissait le travail de Bessel, et la lecture
ne le suppose pas. L'identification de la série avec $J_\beta$ a été vérifiée
terme à terme dans le dossier de vérification ($2\cdot4\cdots2j\,(n+1)(n+3)
\cdots(n+2j-1) = 4^j j!\,(\beta+1)(\beta+2)\cdots(\beta+j)$).}

La condition (D) pour $r = R$ donne l'« équation déterminée » en $g$. Le
feuillet l'écrit
\begin{align*}
& b\Big(1 - \frac{gR^2}{2(n+1)} + \frac{g^2R^4}{2\cdot4\,(n+1)(n+3)} - \cdots\Big) \\
& \qquad = \frac{2gR}{2(n+1)} - \frac{4g^2R^3}{2\cdot4\,(n+1)(n+3)} +
\frac{6g^3R^5}{2\cdot4\cdot6\,(n+1)(n+3)(n+5)} - \cdots,
\end{align*}
c'est-à-dire $b\,s(R) + s'(R) = 0$. Pour $v = r^\beta s$, la condition (D)
s'écrit $s'(R) + \big(\frac hK + \frac\beta R\big)s(R) = 0$ : il faut donc $b =
\frac hK + \frac\beta R$, ce que ces pages ne disent pas. En termes de
fonctions de Bessel, avec $x = \sqrt g\,R$ et $H = \frac{hR}{K}$, l'équation
est $x J_\beta'(x) + H J_\beta(x) = 0$.\footnote{Vue 682, série à l'encre 335,
sa page 6. Le troisième terme du second membre porte « $R^6$ » ; le calcul
demande $R^5$. La valeur de $b$ est une inférence de la lecture ; ces pages
disent seulement que « $b$ et $R$ sont des quantités déterminées ». « a » au
début de la première ligne est repassé sur une autre lettre.} Elle affirme,
« en se conformant à ce qui est démontré N\textsuperscript{o} 306 et suivants »,
que cette équation a une infinité de racines réelles, pour chaque ordre
$\beta$ ; c'est vrai (pour $h > 0$ les racines $g$ sont réelles et
positives).\footnote{« T. 2 », après « suivans », est lu avec doute. La preuve
actuelle de ce fait passe par le caractère autoadjoint du problème (théorie de
Sturm--Liouville, postérieure au feuillet et qu'il ne cite pas).}

Elle compose enfin l'intégrale générale en additionnant les solutions
particulières :
\[
v = B^0\big(a_1e^{-g_1kt}s_1 + a_2e^{-g_2kt}s_2 + \cdots\big) +
(A'\sin\varphi + B'\cos\varphi)\big(a'_1e^{-g'_1kt}s'_1 + \cdots\big) + \cdots
\]
jusqu'à l'ordre $\beta$, les exposants $'$, $2$, $\dots$, $\beta$ étant des
indices d'ordre et non des puissances.\footnote{Vues 682 et 684. L'exposant de
$B$ dans la première ligne est un pâté ; « 0 » est offert d'après la vue 684.
Dans cette intégrale, les $s$ ne portent pas le facteur $r^\beta$ des valeurs
particulières ; la vue 684 appelle $u$ les fonctions correspondantes, qui le
portent.} Sous cette forme, l'intégrale n'est pas la plus générale : elle
impose le même rapport $A^\beta : B^\beta$ à toutes les racines d'un même
ordre. L'intégrale générale est
\[
v = \sum_{\beta \geq 0}\ \sum_{i \geq 1}\big(A_{\beta,i}\sin\beta\varphi +
B_{\beta,i}\cos\beta\varphi\big)\,e^{-g_{\beta,i}kt}\,J_\beta\big(\sqrt{g_{\beta,i}}\,
r\big),
\]
avec deux constantes indépendantes par racine, que l'état initial détermine par
l'orthogonalité des fonctions propres ; ce calcul des coefficients n'est pas
sur ces feuillets. Elle dit pourtant que son intégrale « servira à résoudre dans
toute son étendue la question proposée » : c'est trop dire.\footnote{Vue 682.
La lecture corrige ici la forme ; la conclusion des vues 688 à 691 (voir plus
bas) ne dépend que de la première racine de chaque ordre et n'en est pas
affectée à la limite.}

\textbf{La réalité des racines (article 5, vues 684, 686).} Elle pose $\theta =
\frac{gR^2}{2(n+1)}$ (ici $\theta$ n'est plus un angle), et la série devient $y
= 1 - \theta + \frac{(n+1)^2\theta^2}{1\cdot2\,(n+1)(n+3)} - \cdots$, qui
satisfait à
\[
y + \frac{dy}{d\theta} + \frac{2\theta}{n+1}\frac{d^2y}{d\theta^2} = 0,
\qquad
\frac{d^iy}{d\theta^i} + \frac{n+2i+1}{n+1}\frac{d^{i+1}y}{d\theta^{i+1}} +
\frac{2\theta}{n+1}\frac{d^{i+2}y}{d\theta^{i+2}} = 0 ;
\]
ces équations sont exactes.\footnote{Vue 684, série à l'encre 336, sa page 7.
Les dénominateurs des valeurs remplacées par $\theta$ sont lus « $2^2$ »,
« $2.3$ », « $2.5$ » pour $n = 1, 3, 5$ ; la règle qu'elle énonce, $2(n+1)$,
donne $2\cdot2$, $2\cdot4$, $2\cdot6$. L'exposant « 3 » de $(n+1)$ dans la
troisième équation est très pâle ; le calcul demande $(n+1)$ sans exposant. La
lettre devant $\theta$ dans la marge est lue avec doute $f$. Le mot « art »
de « l'art 308 » est lu d'après le tracé.} Puis, « en appliquant ici les
remarques qui font l'objet de l'art. 308 », elle conclut : là où
$y^{(i+1)}$ s'annule pour $\theta > 0$, $y^{(i)}$ et $y^{(i+2)}$ ont des signes
opposés, « et on en conclura que les différentes valeurs de $g$ … sont toutes
réelles ». Ce qui est écrit est vrai et c'est le pas essentiel du raisonnement
qu'elle emprunte ; tel quel, il ne constitue pas une démonstration : il
faudrait en tirer le compte des racines réelles de $y$, puis de la combinaison
$b\,s + s'$, qui n'est pas $y$.\footnote{Vue 686, série à l'encre 337, sa page
8. Les mots « une valeur positive qui » sont barrés après la fraction.}

\textbf{L'ordre des racines, et l'état final (article 6, vues 686 à 691).} Le
facteur $A^\beta\sin\beta\varphi + B^\beta\cos\beta\varphi$ est « dû à la
différence de température entre les divers points d'une même couche ». Elle
envisage « $\beta$ infini, ou nul » comme « deux manières différentes de
présenter le même état » ; la lecture ne trouve pas de contenu mathématique
précis à ce passage, $\beta$ étant un entier et chaque ordre figurant dans la
somme.\footnote{Vue 686, très corrigée entre les lignes : « a la condition » et
« l'équation » sont en marge, « de chacune des couches », « existeroit »,
« tous les points. qui sont a une distance égale de l'axe du cylindre » dans
l'interligne ; « indiqueroit » est lu avec doute (écrit « indequeroit ») ; un
petit mot barré avant $v$ est lu avec doute « la ».} Elle en tire que
« plus le nombre $\beta$ est grand, plus aussi les différentes valeurs de
$g^\beta$ doivent être grandes » ; la conclusion est vraie (à rang égal, les
racines croissent avec l'ordre), la raison donnée — sinon $\beta$ infini
réduirait $s$ à l'unité — n'en est pas une preuve. Pour $H = 1$, par exemple,
les premières racines $x = \sqrt g\,R$ sont $1{,}256$ (ordre 0), $2{,}405$
(ordre 1), $3{,}518$ (ordre 2), $4{,}613$ (ordre 3).

Puis elle suppose (« si on admet ») que toutes les racines d'un ordre sont
plus grandes que celles de l'ordre précédent et plus petites que celles de
l'ordre suivant. Cette hypothèse est fausse : pour $H = 1$, la deuxième racine
d'ordre 0, $4{,}080$, dépasse la première d'ordre 1 et la première d'ordre 2.
Mais la conséquence qu'elle en tire n'en demande qu'une partie, qui est vraie :
la plus petite racine est $g_{0,1}$ (ordre 0), la suivante $g_{1,1}$ (ordre 1),
et toutes les autres sont plus grandes que $g_{1,1}$.\footnote{Vérifié
numériquement pour $H = hR/K$ de $0{,}01$ à $10\,000$ (dossier de
vérification) : $x_{0,1} < x_{1,1} < \min(x_{0,2}, x_{2,1})$ à chaque fois.
Pour $H = 1$ : $x_{0,1} = 1{,}256$, $x_{1,1} = 2{,}405$, $x_{2,1} = 3{,}518$,
$x_{0,2} = 4{,}080$. Vue 688, série à l'encre 338, sa page 9 : « si on admet »
remplace « En admettant donc » ; un petit mot ajouté au bord du feuillet est
illisible ; « sensiblement » est lu avec doute.} Pour $t$ assez grand, il reste
donc sensiblement
\[
v \approx a\,e^{-g_{0,1}kt}J_0\big(\sqrt{g_{0,1}}\,r\big) + (A'\sin\varphi +
B'\cos\varphi)\,e^{-g_{1,1}kt}J_1\big(\sqrt{g_{1,1}}\,r\big),
\]
les termes négligés décroissant plus vite.\footnote{Le feuillet garde toute la
série de l'ordre 0 et toute celle de l'ordre 1, sans le facteur $r$ de l'ordre
1 ; seules comptent, à la limite, leurs premières racines. L'état décrit est
atteint asymptotiquement, non à partir d'une « certaine valeur de $t$ »
précise.} Ses conséquences sont exactes dans cet état : deux points d'une même
couche aux extrémités d'un diamètre ont des facteurs angulaires opposés, et
« la demi-somme des températures des points opposés est égale à la
température moyenne de la couche » ; le maximum et le minimum sur une couche
sont aux deux bouts du diamètre $\tan\varphi_1 = A'/B'$ ; aux deux points
placés à $\frac\pi2$ de ceux-là, le facteur angulaire est nul, et « le
refroidissement s'opère » sur ce rayon « absolument de la même manière que si
ce rayon appartenait au cylindre » échauffé uniformément par
couches.\footnote{Vue 690, série à l'encre 339, sa page 10. La parenthèse
ouverte dans la première des deux égalités du bas n'est pas fermée.} Elle
rapproche ce résultat de l'« armille » de Fourier, « art 245 p 277 de la
Théorie de la chaleur » : chaque tranche se comporte comme une plaque ronde « qui
serait soustraite au refroidissement dû à l'action de l'air » — sur ses faces,
la tranche ne perdant de chaleur que par son bord —, et l'armille est le cas
où le rayon est constant. L'analogie est juste : dans l'anneau aussi,
l'harmonique d'ordre 1 survit seule, et les points diamétralement opposés ont
pour demi-somme la moyenne.\footnote{Vue 691, verso de la vue 690, cinq lignes ;
« est correspond » sur le feuillet. Le mémoire sur la courbure invoque la même
loi de l'anneau, en citant la page 277 de la même théorie (vues 55 à 63 ; voir
la section « Mémoire sur la courbure des surfaces »).}

\textbf{La copie (vues 692 à 695).} Le titre et les articles 1 à 3 sont
récrits presque mot pour mot (vues 692, 693), puis les pages 7 et 8 reprennent
le texte des vues 686 à 691 depuis « manières différentes de présenter » jusqu'à
la fin.\footnote{Vues 692 à 695. Variantes : $\frac hK$ avec une capitale (vue
693) pour $\frac hk$ ; « forcé » pour « mené », « Ainsi il y a » pour « Il y a
donc », les $s$ indicés ; « art 245 » sans page ni titre ; « correspond » pour
« est correspond ». À la vue 695, « $A'\cos\varphi - A'\sin\varphi = 0$ » : la
seconde lettre est écrite $A'$ où le brouillon (vue 690) porte $B'$, qui est
juste. Au premier numérateur de la vue 692, la lettre après $d$ est lue $v$
avec doute. La page 2 s'arrête sur « Pour y parvenir il » ; la vue 694 ne la
continue pas.}

\subsection*{Sphère, cône et cylindre : le dernier état}

\pagerange{696}{716}
Un nouveau titre, « Équations du mouvement de la chaleur dans la sphère, le
cône et le cylindre solides échauffés d'une manière entièrement arbitraire »,
ouvre un brouillon très raturé, paginé 1.\footnote{Vue 696, série à l'encre 342,
« 1 » barré au-dessus ; le dos (vue 697) est blanc. Les additions sont dans
l'interligne au-dessus des mots barrés ; un mot commencé après « les trois
équations. », un mot d'une addition fine et un mot barré sous une tache sont
illisibles ; « donnent », l'addition « qu'on cherche la … qu'elle et » et
« détermine » sont lus avec doute. Il n'est pas sûr, pour quelques mots de la
fin, que le trait qui les traverse soit une biffure.} Elle y dit ce qu'elle
cherche. On peut, « comme l'ont fait M\textsuperscript{rs} de la Place et
Poisson », transformer l'équation générale en y mettant « à la place des
coordonnées rectangulaires d'autres coordonnées appropriées à la figure du
corps » ; des transformations analogues donnent le cône et le cylindre ; mais
en suivant ce procédé il serait difficile « d'apercevoir la liaison qui existe
entre les trois équations ». Au contraire, en examinant chaque question
« d'une manière spéciale », « conformément à la méthode employée par
M\textsuperscript{r} Fourier », « on s'aperçoit aisément que la molécule de la
sphère appartient en même temps à un solide conique qui a son sommet au centre
de la sphère » et qui est terminé par une calotte. La page s'arrête sur la
différence entre cette molécule et celle du cône de même axe.

Ce que ces pages et les précédentes établissent est juste, à condition de le
dire ainsi : les trois équations sont une seule, $\partial_t v = k\,\Delta v$,
écrite en trois systèmes de coordonnées ; la sphère et le cône de sommet au
centre ont la même équation intérieure et diffèrent par la surface ; le
cylindre s'obtient de la sphère par le passage à la limite de la Note (vue 656
et vue 712), et la plaque du cylindre en supprimant le mouvement le long de
l'axe. La « liaison » qu'elle cherche est celle des formes d'un même opérateur.

\textbf{La sphère échauffée uniformément (vue 698).} Lorsque chaque couche est
à température uniforme, les termes angulaires sont nuls, et « il ne passe
aucune quantité de chaleur » par les côtés des quadrilatères que méridiens et
parallèles découpent sur une couche. Elle écrit la condition
$\frac{1}{r^2}\frac{\partial^2 v}{\partial\theta^2} + \frac{1}{r^2\sin^2\theta}
\frac{\partial^2 v}{\partial\varphi^2} = 0$ et la dit exprimer « la constance
des températures des différents quadrilatères ». Il manque le terme
$\frac{\cos\theta}{r^2\sin\theta}\frac{\partial v}{\partial\theta}$ : la
condition juste est $\Delta_S v = 0$, et, pour une température régulière sur
toute la couche, elle équivaut bien à la constance de $v$ sur la couche. C'est
ce fait qui écarte l'exemple de la vue 670.\footnote{Vue 698, série à l'encre
343, sans pagination ; le dos (vue 699) est blanc. Sans le terme en
$\frac{\partial v}{\partial\theta}$, l'équation écrite n'exprime pas cette
constance : $v = \cos\theta$ n'est pas constant, et $\Delta_S\cos\theta =
-2\cos\theta$. Dans le bilan qui suit, elle prend pour côtés du quadrilatère
l'arc $r\sin\theta\,d\theta$ du méridien et l'arc $r\,d\varphi$ du parallèle ;
ce sont $r\,d\theta$ et $r\sin\theta\,d\varphi$. Le numérateur de la fraction
$\frac{dv}{dt}$ écrite après « pendant le même tems » est lu $v$ avec doute ;
la page s'arrête sur « $\theta - d\theta$ : ».}

\textbf{L'hémisphère en quadrilatères (vue 701).} Au dos d'un petit feuillet,
elle partage la surface hémisphérique par une infinité de méridiens et de
parallèles, note que les quadrilatères d'une zone sont plus grands que ceux de
la zone suivante mais que la différence est négligeable pour des quadrilatères
infiniment petits, place un foyer de chaleur $f$ sur un côté méridien d'un des
parallélogrammes qui s'appuient sur l'équateur, et nomme $f, f', f'', \dots$ les
températures de la première rangée, $a, a', a'', \dots$ celles de la seconde.
La page s'arrête là : un modèle discret du même bilan, sans
équation.\footnote{Vue 701, verso d'un feuillet dont l'autre face (vue 700)
porte le nombre 344 et est blanche ; le feuillet 345 (vues 702, 703) est blanc.
« Si », « seront » et un mot barré (« fait ») sont lus avec doute ; un mot après
« négligeable à la » est illisible. La page passe de « quadrilatères » à
« parallélogrammes » pour les mêmes éléments.}

\textbf{Le laplacien en $r, \sin\theta, \varphi$, et ses limites (vues 712,
713, 716).} Le feuillet 350, écrit d'une main rapide et relié après le manuscrit
C, reprend au net le calcul du cône de la vue 630 : $\frac{\partial^2
v}{\partial y^2}$, la somme $\frac{\partial^2 v}{\partial x^2} +
\frac{\partial^2 v}{\partial y^2}$, $\frac{\partial^2 v}{\partial z^2}$, et le
total, qui est l'expression exacte donnée plus haut. Une « N.ota » explique un
pas : la dérivée de $\cos\theta$ par rapport à $\sin\theta$ est
$-\frac{\sin\theta}{\sqrt{1 - \sin^2\theta}}$, et multipliée par
$\frac{\partial\sin\theta}{\partial z} = -\frac{\sin\theta\cos\theta}{r}$ elle
donne $\frac{\sin^2\theta}{r}$ ; c'est exact.\footnote{Vue 712, série à l'encre
350. Elle écrit une fois $d^2u$ pour $d^2v$. Deux signes, couverts de pâtés,
sont offerts comme $-$ ; un $\varphi$ au bord du feuillet est lu avec doute. Le
calcul de la lecture confirme le développement de $\frac{d^2v}{dz^2}$, en
particulier le coefficient $\frac{2\sin\theta\cos^2\theta - \sin^3\theta}{r^2}$
de $\frac{dv}{d\sin\theta}$.} Puis trois cas particuliers.
\begin{itemize}
\item \emph{$r$ infini.} « L'équation appartient au cylindre » : avec $\theta =
0$, elle se réduit à $\partial_r^2 v + \partial_s^2 v + \frac1s\partial_s v +
\frac{1}{s^2}\partial_\varphi^2 v$, $s = r\sin\theta$ ; « cette équation est en
effet celle que M\textsuperscript{r} Poisson a donnée en considérant le cas
général », celui du cylindre échauffé arbitrairement. Le passage est celui de
la Note (vue 656), et il est juste au sens dit plus haut.
\item \emph{Sans mouvement le long de l'axe.} En supprimant $\partial_r^2 v$,
« l'équation est celle que j'avois trouvée » : celle du cylindre des vues 672 à
695, qui n'est autre que celle de la plaque circulaire, « abstraction faite de
la déperdition de chaleur due dans la plaque isolée au contact de l'air ».
C'est exact.\footnote{Sur le feuillet, « celle que j'avois trouvé » finit une
ligne et « j'avois remarqué, dès lors, » ouvre la suivante. Ce renvoi à une
équation « trouvée » auparavant, et la remarque sur la plaque, qui est la fin de
la pièce sur le cylindre (vues 691 et 695), font penser que ce feuillet est
postérieur à cette pièce : inférence de cette lecture.}
\item \emph{$\theta = 90^\circ$.} « L'axe du cylindre deviendra le rayon de la
plaque » : dans le plan de l'équateur, le second membre se réduit à
$\partial_r^2 v + \frac2r\partial_r v + \frac{1}{r^2}\partial_\varphi^2 v -
\frac{\sin\theta}{r^2}\frac{\partial v}{\partial\sigma}$. Au verso, pour
l'identifier à la plaque, elle admet que dans le dernier terme « $\sin\theta$
devienne égal au rayon du cercle dont le rayon est $r$ », d'où
$\frac{\sin\theta}{r^2}\frac{\partial v}{\partial\sigma} = \frac1r
\frac{\partial v}{\partial r}$, et l'équation devient $\partial_r^2 v +
\frac1r\partial_r v + \frac{1}{r^2}\partial_\varphi^2 v$, celle de la plaque.
Le résultat est exact, et la justification se fait ainsi : pour une température
de la plaque, $v$ ne dépend de $\sigma$ que par la distance à l'axe $\rho =
r\sigma$, donc $\partial_\sigma v = r\,\partial_\rho v$, et à $\theta =
90^\circ$, $\frac{\sin\theta}{r^2}\partial_\sigma v = \frac1r\partial_\rho v$
; c'est le sens exact de son « $\sin\theta$ égal au rayon ».\footnote{Vue 713,
verso du feuillet 350, sans numéro. Elle y ajoute que l'expression
$\frac{\sin\theta}{r^2}\frac{dv}{d\sin\theta}$ « peut être censée n'avoir pas
de valeur fixe quoique l'usage l'attribue toujours au cercle dont le rayon est
l'unité ». Elle écrit $\frac{d^2v}{dx}$, sans exposant au dénominateur, dans la
dernière équation. Calcul de vérification : pour $v = V(r\sin\theta, \varphi)$,
le laplacien sphérique vaut exactement, à $\theta = \frac\pi2$,
$V_{\rho\rho} + \frac1\rho V_\rho + \frac{1}{\rho^2}V_{\varphi\varphi}$.}
\end{itemize}
Le verso s'achève sur la molécule du cône, par ses trois faces, le deuxième
flux étant multiplié par $\cos\theta$ « à cause de l'obliquité » : le même
raisonnement qu'aux vues 636 et 637, qui s'arrête ici sur la troisième quantité
de chaleur.\footnote{Vue 713. Le dénominateur de la deuxième quantité,
$\frac{dv}{r\sin\theta}$, ne porte pas de différentielle.}

Le feuillet 352, sous le titre « Cone », donne le détail du même changement de
variables ; une première somme, où les dérivées de $\cos\theta$ et de
$\cos^2\theta$ n'avaient pas été prises, y est corrigée sur la page même
(« en faisant varier par rapport à $\sin\theta$, $\cos\theta$ et
$\cos^2\theta$ »), et le total final est l'expression exacte ; le bas du
feuillet redonne, pour $r$ infini, l'équation du cylindre.\footnote{Vue 716,
série à l'encre 352 ; le revers (vue 717) n'est pas transcrit : la
transcription le dit couvert d'autres calculs dans sa note de la vue 716, blanc
dans l'en-tête du lot 36 et dans sa note sur la numérotation. Dans la première somme, le coefficient de $\frac{dv}{dr}$ est
$\frac{\cos^2\theta + 1}{r}$ et celui de $\frac{dv}{d\sin\theta}$ porte $+$
au lieu de $-$ ; la somme corrigée a $\frac2r$ et le signe juste. Elle écrit
$\frac{d^2v}{dy}$ sans exposant, et un terme de $\frac{d^2v}{dy^2}$ sans
dénominateur, avec $d\sin\theta\,dr$ là où la ligne correspondante de
$\frac{d^2v}{dx^2}$ porte $d\sin\theta\,d\varphi$. Quatre signes couverts de
pâtés sont offerts d'après les lignes voisines ; un numérateur barré et un
signe récrit sont illisibles. Rien n'ordonne ce feuillet par rapport au
feuillet 350.}

\section*{Le manuscrit C (ff. 348r--349r, vues 708 à 710)}

\pagerange{708}{710}
Trois pages d'une mise au net, d'une main posée et régulière, presque sans rature, portent deux énoncés sur des équations de Fermat à exposants pairs, chacun suivi de sa « Démonstration »\footnote{Vues 708 à 710. Un feuillet écrit recto (vue 708, n\textsuperscript{o} 348 de la série à l'encre) et verso (vue 709, sans numéro), puis le recto du feuillet suivant (vue 710, n\textsuperscript{o} 349), écrit aux trois quarts ; son revers, vue 711, est blanc. Aucun numéro de page ni date de la main qui écrit. Le nom de « manuscrit C » et la cote ff. 348r--349r sont ceux de Laubenbacher et Pengelley (2010, § 1.3) ; la correspondance entre cette cote et les vues vient de la série à l'encre lue au lot 36, qui est tracée à la plume et que les transcriptions ne donnent donc pas comme foliotation. Le jugement de ces auteurs sur la démonstration n'est ni rapporté ni vérifié ici : ce qui suit est la vérification propre de cette lecture. Les feuillets 350 à 352 (vues 712 à 716), reliés juste après, portent des calculs sur la chaleur et sont lus avec les autres pièces sur la chaleur.}. Le premier porte sur $x^{2n}+y^{2n}=2z^{2n}$ pour $n$ impair et s'étend à $x^{2n}+y^{2n}=2z^{2(n+2d)}$ ; le second, qui en est déduit, est le théorème de Fermat pour les exposants pairs : $Z^{2n}=X^{2n}+Y^{2n}$ n'a pas de solution en entiers non nuls pour $n>1$. On restitue d'abord l'argument pas à pas, en disant à chaque pas s'il tient ; la dernière sous-section fait le bilan.

\subsection*{Le théorème sur $2z^{2n} = x^{2n} + y^{2n}$}

\pagerange{708}{709}
L'énoncé est : « $n$ étant un nombre impair plus grand que l'unité, il est impossible de satisfaire en nombres entiers à l'équation $2z^{2n}=y^{2n}+x^{2n}$ »\footnote{Vue 708. Le feuillet écrit, dans l'énoncé, $2z^{2n}=y^{2n}+z^{2n}$ : la dernière lettre a la forme du $z$ initial ; quatre lignes plus bas, la même équation s'écrit $2z^{2n}=y^{2n}+x^{2n}$, qu'on suit. « Théorème », « Démonstration » et le $n$ initial sont soulignés.}. Tel quel, il est faux : $x=y=z=1$ satisfait l'équation pour tout $n$. Ce qui est vrai est ceci : pour $n$ impair $>1$, les seules solutions en entiers premiers entre eux sont $x^2=y^2=z^2=1$. C'est une conséquence du théorème de Darmon et Merel\footnote{Darmon et Merel (1997) : pour tout nombre premier $\ell\ge 3$, l'équation $a^\ell+b^\ell=2c^\ell$ n'a pas de solution en entiers premiers entre eux avec $|abc|>1$ ; c'était une conjecture de Dénes. On l'applique à un facteur premier $\ell$ de $n$, avec $a=x^{2n/\ell}$, $b=y^{2n/\ell}$, $c=z^{2n/\ell}$ ; le cas $a=0$ donnerait $b^\ell=2c^\ell$, impossible. Ce résultat est de beaucoup postérieur aux feuillets, qui ne le citent évidemment pas.}. Pour $n=1$, au contraire, les solutions sont en nombre infini : $1^2+7^2=2\cdot 5^2$, $7^2+17^2=2\cdot 13^2$. On verra que la démonstration manque précisément le cas $x^2=y^2=1$.

\emph{Hypothèses.} La page suppose « comme à l'ordinaire » $x$, $y$, $z$ premiers entre eux. Il s'ensuit, ce qu'elle ne dit pas, que $x$, $y$ et $z$ sont impairs\footnote{Si l'un de $x$, $y$ était pair, l'autre serait impair et $x^{2n}+y^{2n}$ serait impair ; tous deux impairs, $x^{2n}+y^{2n}\equiv 2 \pmod 8$, donc $z^{2n}$ est impair.}. Elle observe que tout diviseur de $z^{2n}$ est une somme de deux carrés : c'est exact, car tout facteur premier impair de $(x^n)^2+(y^n)^2$, avec $x$ et $y$ premiers entre eux, est de la forme $4k+1$, et un tel nombre est somme de deux carrés (théorème des deux carrés de Fermat et Euler).

\emph{Première paramétrisation.} Puisque $(x^n)^2+(y^n)^2=2(z^n)^2$, la page pose
\[ z^n=P^2+Q^2,\qquad x^n=P^2-Q^2-2PQ,\qquad y^n=P^2-Q^2+2PQ. \]
C'est la paramétrisation classique des solutions primitives de $X^2+Y^2=2Z^2$, aux signes près (le signe de $x$ est libre, puisque seul $x^{2n}$ intervient) ; l'identité $(P^2-Q^2-2PQ)^2+(P^2-Q^2+2PQ)^2=2(P^2+Q^2)^2$ a été vérifiée.

\emph{La factorisation.} Pour $n$ impair, $x^{2n}+y^{2n}=(x^2+y^2)\,\Phi$ avec
\[ \Phi=x^{2n-2}-x^{2n-4}y^2+x^{2n-6}y^4-\dots+y^{2n-2}. \]
La page écrit $z=z'z''$, $x^2+y^2=2z'^{2n}$ et $\Phi=z''^{2n}$\footnote{Le feuillet écrit $z''^n=x^{2n-2}-x^{2n-4}y^2+\&c$ ; à la vue 710, en reprenant la démonstration, il parle des « valeurs de $F$, $H$ et $z''^{2n}$ ». Seul l'exposant $2n$ s'accorde avec $z=z'z''$ et $x^2+y^2=2z'^{2n}$ ; on le suit.}. Ce partage suppose $x^2+y^2$ et $\Phi$ premiers entre eux au facteur $2$ près ($\Phi$ est impair, somme de $n$ termes impairs). Or $\Phi\equiv n\,y^{2n-2}\pmod{x^2+y^2}$, de sorte que leur plus grand commun diviseur divise $n$. Il vaut $1$ lorsque aucun facteur premier de $n$ n'est de la forme $4k+1$, puisque les facteurs premiers impairs de $x^2+y^2$ sont de cette forme ; il peut valoir plus sinon : pour $n=5$, $x=1$, $y=3$, on a $x^2+y^2=10$ et $\Phi=1-9+81-729+6561=5905=5\cdot 1181$. La page ne traite pas ce cas : c'est une première lacune.

Avec $x^n+y^n=(x+y)H$ et $y^n-x^n=(y-x)F$\footnote{Le feuillet définit $H=x^{n-1}-x^{n-2}y+\&c$ et $F=x^{n-1}+x^{n-2}y+\&c$ ; on écrit les sommes complètes.}, et la seconde paramétrisation, appliquée à $x^2+y^2=2(z'^n)^2$,
\[ z'^n=p^2+q^2,\qquad x=p^2-q^2-2pq,\qquad y=p^2-q^2+2pq, \]
on a $\tfrac12(x+y)=p^2-q^2$, $\tfrac14(y-x)=pq$, et de même $\tfrac12(x^n+y^n)=P^2-Q^2$, $\tfrac14(y^n-x^n)=PQ$ ; d'où le système (1) de la page :
\[ P^2+Q^2=z''^n(p^2+q^2),\qquad P^2-Q^2=H(p^2-q^2),\qquad PQ=F\,pq. \]
Tout est exact. On note, ce que la page ne dit pas, que $p$ et $q$ sont premiers entre eux et de parités différentes, puisque $x$ et $y$ sont premiers entre eux et impairs.

\emph{Deuxième décomposition.} La page pose $z''^n=U^2+V^2$ et tire de la première équation de (1)
\[ P=Uq\pm Vp,\qquad Q=Up\mp Vq. \]
Cela ne découle pas de la seule égalité de deux sommes de deux carrés, un nombre en ayant en général plusieurs décompositions. On peut le justifier par les entiers de Gauss, que la page n'emploie pas\footnote{Contexte, non tiré des feuillets : Gauss introduit ces entiers en 1832.} : on a $(p+iq)^2=\tfrac{1-i}{2}(x+iy)$ et $(P+iQ)^2=\tfrac{1-i}{2}(x^n+iy^n)$ ; or $x^n+iy^n$ est divisible par $x+iy$ si $n\equiv 1\pmod 4$ et par $x-iy$ si $n\equiv 3\pmod 4$. Le quotient de $P+iQ$ par $p+iq$ (ou par $p-iq$) est donc un entier de Gauss, puisque son carré en est un, et sa norme est $z''^n$ : c'est, aux unités et à l'ordre des lettres près, le $U+iV$ de la page. En substituant, on obtient le système (2), dont on a vérifié les deux lignes :
\[ \bigl[U(p+q)\pm V(p-q)\bigr]\bigl[U(q-p)\pm V(p+q)\bigr]=H(p^2-q^2),\qquad (U^2-V^2)pq\pm(p^2-q^2)UV=F\,pq. \]
La première ligne, développée, s'écrit $(p^2-q^2)(V^2-U^2)\pm 4pq\,UV=H(p^2-q^2)$ : comme $p^2-q^2$ est premier avec $4pq$, il divise $UV$. La seconde montre que $pq$ divise $(p^2-q^2)UV$, donc $UV$. Ainsi $pq(p^2-q^2)$ divise $UV$, comme le dit la page.

\emph{Troisième décomposition.} La page pose $z''^{n-1}=U'^2+V'^2$, $z''=u^2+v^2$, $U=U'u\pm V'v$, $V=U'v\mp V'u$ (même remarque que plus haut), d'où l'identité, exacte,
\[ UV=(U'^2-V'^2)\,uv\pm(v^2-u^2)\,U'V'. \]
Elle poursuit : « Les valeurs de $VU$ données par l'un et l'autre signes devront également satisfaire » ; la somme $2(U'^2-V'^2)uv$ et la différence $2(v^2-u^2)U'V'$ seraient donc toutes deux divisibles par $N=pq(p^2-q^2)$. C'est la deuxième lacune : le signe est fixé par la décomposition choisie, une seule des deux valeurs est $UV$, et rien n'oblige l'autre à être divisible par $N$. Si on l'accorde, la conclusion de la page, $N\mid 2(v^2-u^2)uv$, est juste\footnote{Vue 709, qui continue la phrase de la vue 708 ; même main, sans rature. La page s'arrête aux trois quarts du feuillet, sur l'équation finale suivie d'un trait.} : $U'$ et $V'$ étant premiers entre eux, $U'^2-V'^2$ et $U'V'$ le sont aussi, et chaque puissance d'un nombre premier qui divise exactement $N$ divise alors $2uv$ ou $2(v^2-u^2)$.

\emph{L'alternative.} De $N\mid 2(v^2-u^2)uv$, la page tire : ou bien $|2(v^2-u^2)uv|\ge N$, ou bien $2(v^2-u^2)uv=0$\footnote{Le feuillet écrit $2(v^2-u^2)vu>pq(p^2-q^2)$ ; un multiple non nul de $N$ peut lui être égal.}.

Dans le premier cas, comme $y-x=4pq$ et $y+x=2(p^2-q^2)$, on a $N=\tfrac18(y^2-x^2)$, d'où $32(v^2-u^2)uv\ge 2(y^2-x^2)$ : c'est exact. La page affirme ensuite $2(y^2-x^2)>(x+y)^2$ ; mais
\[ 2(y^2-x^2)-(x+y)^2=(x+y)(y-3x), \]
qui n'est positif que dans certains cas. Elle borne ensuite l'autre membre : supposant implicitement $|y|>|x|$, elle tire de l'équation $z<|y|$, soit $z'^2(u^2+v^2)^2<y^2$, a fortiori $2z'^2(u^2+v^2)uv<y^2$, puis de « $z'$ étant impair et la somme de deux quarrés on a $z'>4$ » l'inégalité $32(u^2+v^2)uv<y^2$\footnote{Le feuillet porte ici $(v^2+u^2)$, avec le signe $+$, là où la supposition examinée porte $(v^2-u^2)$, et la transcription le signale. Le signe $+$ est celui que donne la borne, qui porte sur $z''=u^2+v^2$ ; le retour à $v^2-u^2\le v^2+u^2$ est sous-entendu.}. Pour contredire $32(v^2-u^2)uv>(x+y)^2$, il faudrait $(x+y)^2\ge y^2$, c'est-à-dire $x$ et $y$ de même signe ; et pour $x$, $y$ positifs l'inégalité précédente demande $y>3x$. La chaîne ne conclut donc que si $x$ et $y$ ont même signe et $|y|>3|x|$ : c'est la troisième lacune. Avec $p=3$, $q=2$, la paramétrisation de la page donne $x=-7$, $y=17$ ($x^2+y^2=2\cdot 13^2$) : on a bien $2(y^2-x^2)=480>(x+y)^2=100$, mais $(x+y)^2<y^2=289$ ; avec $x=7$, $y=17$, c'est la première inégalité qui tombe ($480<576$). De plus, « $z'>4$ » est faux pour $z'=1$, c'est-à-dire pour $x^2+y^2=2$, qui est justement le cas trivial ; pour $z'>1$, les facteurs premiers de $z'$ sont de la forme $4k+1$ et $z'\ge 5$, ce qui suffit à l'inégalité.

Dans le second cas, $v^2=u^2$ est impossible, parce que $z''=u^2+v^2$ divise $z$, qui est impair\footnote{La page écarte ce cas autrement : $F$ et $H$ auraient alors un diviseur commun, donc $x$ et $y$ aussi ; elle ne détaille pas ce pas.}. Si $v=0$ (ou $u=0$), $u$ divise $U$ et $V$, donc $P$ et $Q$, donc $u^2$ divise $x^n$ et $y^n$ : il faut $u=1$. Alors $z''=1$, $U=1$, $V=0$, $P=p$, $Q=q$ (aux signes près), et $x^n=x$, $y^n=y$. La page en conclut $n=1$ : « Ce cas est donc le seul dans lequel il soit possible de satisfaire à l'équation ». Mais $x^n=x$ et $y^n=y$ ont lieu aussi pour $x^2=y^2=1$, quel que soit $n$ : c'est la solution triviale que l'énoncé oublie.

\subsection*{L'extension, et le théorème de Fermat pour les exposants pairs}

\pagerange{710}{710}
La page reprend la même démonstration pour $2z^{2n}=y^{2m}+x^{2m}$ avec $m=n-2d$ : il suffit, dit-elle, de mettre $m$ au lieu de $n$ dans $F$, $H$ et $\Phi$, ce qui ne change pas la forme des systèmes (1) et (2)\footnote{Vue 710, n\textsuperscript{o} 349 de la série à l'encre, même main que les vues 708 et 709, écrite aux trois quarts ; le bas et le revers (vue 711) sont blancs. Le feuillet écrit $n-2d$ ; on abrège en $m$.}. Le premier cas de l'alternative tomberait « à plus forte raison » par $z'^{2n}(v^2+u^2)^{2n}<y^{2n-4d}$, qui découle bien de $2z^{2n}=x^{2m}+y^{2m}<2y^{2m}$\footnote{L'exposant de $(v^2+u^2)$ est écrit petit et serré ; la transcription le lit « 8n » avec doute, peut-être « 3n ». Le calcul demande $2n$, comme au passage parallèle de la vue 709 ; c'est une inférence de cette lecture.} ; le second donnerait $m=1$. En remettant $n$ pour $n-2d$, elle énonce : si $n$ est impair et plus grand que l'unité, quel que soit $d$, l'équation $2z^{2(n+2d)}=x^{2n}+y^{2n}$ est impossible en entiers\footnote{Le chiffre devant $d$ dans l'exposant est tracé d'un trait lourd, comme récrit ; la transcription le lit 2 avec doute, appuyée par la phrase qui précède (« $n+2d$ au lieu de $n$ »).}. Cet énoncé a lui aussi la solution triviale $x=y=z=1$, et sa démonstration hérite des quatre points relevés plus haut ; cette lecture ne cherche pas à établir le statut actuel de cette équation généralisée.

Vient le second théorème : « Pour toute valeur de $n$ plus grande que l'unité il est impossible de satisfaire en nombres entiers à l'équation $Z^{2n}=X^{2n}+Y^{2n}$ ». La réduction est la suivante. L'un de $X$, $Y$ est pair\footnote{La page le dit sans raison ; si tous deux étaient impairs, $X^{2n}+Y^{2n}\equiv 2\pmod 4$ ne serait pas un carré.} ; soit $X=2xx'$. Comme $Z$ et $Y$ sont impairs et premiers entre eux, $Z^n+Y^n$ et $Z^n-Y^n$ ont pour plus grand commun diviseur $2$ et pour produit $2^{2n}(xx')^{2n}$ ; la page écrit, avec un choix de signe convenable,
\[ Z^n\pm Y^n=2x'^{2n},\qquad Z^n\mp Y^n=2^{2n-1}x^{2n}, \]
d'où $Z^n=x'^{2n}+2^{2n-2}x^{2n}$ et $\pm Y^n=x'^{2n}-2^{2n-2}x^{2n}$. Elle passe de chacune de ces égalités entre puissances $n$-ièmes à une égalité entre les bases, en appelant $(x)$, $(x)'$, $(x)''$ des facteurs de $x$ :
\[ Z\mp Y=2^{2n-1}(x)^{2n},\qquad Z-x'^2=2^{2n-2}(x)'^{2n},\qquad x'^2\mp Y=2^{2n-2}(x)''^{2n}, \]
et, en ajoutant les deux dernières, $2(x)^{2n}=(x)'^{2n}+(x)''^{2n}$\footnote{Le feuillet écrit $2^{n-1}(x)^{2n}=2^{2n-2}\{(x)'^{2n}+(x)''^{2n}\}$ : l'exposant $n-1$ à gauche, là où la colonne de droite porte $2^{2n-1}$ ; la transcription le garde tel. Seul $2^{2n-1}$ donne l'équation finale que la page écrit ensuite. Les équations sont disposées sur deux colonnes, les trois lignes « $x=\dots(x)$, $=\dots(x)'$, $=\dots(x)''$ » à droite de $X=2xx'$.}, équation « dont l'impossibilité vient d'être démontrée ».

Ce pas est juste pour $n$ impair, sous la même réserve que plus haut. Chacune des trois différences $Z^n\mp Y^n$, $Z^n-(x'^2)^n$, $(x'^2)^n\mp Y^n$ se factorise par la somme ou la différence des bases, selon le signe, le cofacteur étant une somme de $n$ termes impairs, donc impair : toute la puissance de $2$ va à ce premier facteur, et le reste est une puissance $2n$-ième dès que la base et le cofacteur sont premiers entre eux, ce qui ne peut manquer qu'en un facteur premier de $n$. La solution obtenue n'est pas triviale : si $(x)'^2=(x)''^2=(x)^2$, on aurait $Z\pm Y=2x'^2$ et $Z^n\pm Y^n=2(x'^2)^n$, ce qui pour $n>1$ force $Z=Y$, donc $X=0$, dans un cas, et est impossible dans l'autre. Un facteur commun à deux des trois nombres divise le troisième, par l'équation même ; après division par leur plus grand commun diviseur, on retombe ainsi sur une solution primitive non triviale de la première équation. Pour $n$ pair, en revanche, les cofacteurs sont pairs et $Z+Y$ ne divise plus $Z^n+Y^n$ : la réduction de la page ne vaut pas, et d'ailleurs le premier théorème n'a été énoncé que pour $n$ impair. Le « pour toute valeur de $n$ » de l'énoncé n'est donc couvert, par les moyens de la page, que pour $n$ impair.

L'énoncé lui-même est vrai pour tout $n>1$ : si $n$ est pair, $4$ divise $2n$ et l'équation contient $X^4+Y^4=W^2$, que la descente infinie de Fermat exclut ; si $n$ est impair et $>1$, $2n$ a un facteur premier impair $\ell$, et l'équation contient le cas d'exposant $\ell$ du théorème de Fermat, démontré par Wiles\footnote{Wiles (1995), avec Taylor pour une étape. La descente pour l'exposant $4$ est de Fermat lui-même. Aucun des deux n'est cité sur la page, qui ne mentionne pas non plus l'exposant $4$.}.

\subsection*{Ce que la démonstration établit}

\pagerange{708}{710}
Sont justes, et ont été vérifiés par le calcul de cette lecture : les deux paramétrisations, les systèmes (1) et (2), l'identité donnant $UV$, la divisibilité de $UV$ par $pq(p^2-q^2)$, et la réduction du second théorème au premier pour $n$ impair, hors le cas où un facteur premier de $n$ s'introduit dans les plus grands communs diviseurs. Les décompositions $P=Uq\pm Vp$ et $U=U'u\pm V'v$ sont affirmées sans preuve, mais peuvent se justifier par les entiers de Gauss.

Ne sont pas établis, et la page ne démontre donc ni l'un ni l'autre théorème :
\begin{enumerate}
\item le partage $x^2+y^2=2z'^{2n}$, $\Phi=z''^{2n}$ lorsqu'un facteur premier de $n$ de la forme $4k+1$ divise $x^2+y^2$ (exemple : $n=5$, $x=1$, $y=3$) ;
\item la divisibilité simultanée des deux valeurs de $UV$, « l'un et l'autre signes » ;
\item l'exclusion du premier cas de l'alternative, que la chaîne d'inégalités n'obtient que pour $x$, $y$ de même signe et $|y|>3|x|$ ;
\item le cas $x^2=y^2=1$, qui satisfait l'équation pour tout $n$ et que la conclusion « $n=1$ » oublie.
\end{enumerate}
Pour en faire une démonstration, il faudrait au moins traiter les facteurs premiers de $n$, remplacer l'argument des deux signes par une divisibilité réellement acquise, trouver une borne valable pour tout rapport $|y|/|x|$, et énoncer l'exception $x^2=y^2=z^2=1$. Cette lecture ne voit pas comment combler les points 2 et 3 avec les moyens de la page. Les deux énoncés, corrigés de l'exception triviale, sont vrais ; ils ont été démontrés beaucoup plus tard, par d'autres voies (Darmon et Merel pour le premier, Wiles pour le second).

\section*{Résidus quadratiques : notes de travail}

\pagerange{704}{730}
Autour du manuscrit C, de petits feuillets et des billets montés sur des feuilles de garde, pour la plupart de sa main rapide de travail, portent des recherches sur les résidus quadratiques\footnote{Feuillets 346 et 347 (vues 704 et 706), puis 353 à 364 (vues 718 à 730) de la série à l'encre ; les dos blancs (vues 705, 707, 711, 721, 725, 727, 731) et le feuillet blanc sans numéro des vues 732--733 ne portent rien. Les feuillets 350 à 352 (vues 712 à 716), sur la chaleur, sont reliés parmi eux et lus avec les pièces sur la chaleur. Le billet de la vue 729 porte en outre un fragment au crayon, d'une autre main, plus grande, sur des déterminants, coupé aux deux bords ; il n'est pas lu ici.}. On note $\left(\frac{a}{p}\right)=1$ ou $-1$ selon que $a$ est ou n'est pas un carré modulo le nombre premier impair $p$ (symbole de Legendre)\footnote{Les feuillets écrivent « $a$ est résidu (non résidu) de $p$ », parfois « $a$ R $a'$ » pour « $a$ est résidu de $a'$ » (vue 730), et la congruence avec un signe à trois traits ; on écrit $\left(\frac{a}{p}\right)=\pm1$, $\equiv$ et $\pmod p$. Ces notes n'emploient pas le symbole de Legendre ; le mémoire imprimé de Legendre relié au début du volume, si.}. Presque tout tourne autour d'une question : un nombre donné est-il non-résidu d'un nombre premier \emph{plus petit que lui} ? Comme le dit un billet (vue 722), c'est un lemme dont « l'auteur » a besoin pour démontrer son « théorème fondamental ».

\subsection*{Le nombre 2, résidu ou non-résidu, selon la forme de $p$ modulo 8}

\pagerange{704}{704}
« Voici encore une autre démonstration de la propriété du nombre 2 d'être résidu ou non résidu des nombres premiers suivant leurs formes par rapport à 8 »\footnote{Vue 704, feuillet 346, complet, écrit sur ses deux tiers inférieurs et fini par un trait. Le « non » ajouté au-dessus de la première ligne de l'alinéa remplace un « non » biffé à la ligne suivante ; le 3 de « $4K+3$ » a la longue queue qui le fait lire 9, et l'arithmétique établit 3.}. La propriété est le second complément de la loi de réciprocité quadratique : $\left(\frac{2}{p}\right)=1$ si $p\equiv\pm1\pmod 8$, $-1$ si $p\equiv\pm3\pmod 8$. L'argument part de l'égalité $-1\equiv 2\cdot\frac{p-1}{2}\pmod p$ et de la connaissance de $\left(\frac{-1}{p}\right)$ et des facteurs de $\frac{p-1}{2}$.

Si $p=8K'+3$, alors $-1\equiv 2(4K'+1)$ et $-1$ est non-résidu. Le nombre $4K'+1$ a un nombre pair de facteurs premiers de la forme $4h-1$ ; ceux de la forme $4h+1$ sont résidus de $p$, et ceux de la forme $4h-1$ le sont pris avec le signe $-$ ; le produit $4K'+1$ est donc résidu, et $2$ non-résidu. Si $p=8K'+7$, $-1\equiv 2(4K'+3)$, le nombre $4K'+3$ a un nombre impair de facteurs $4h-1$, c'est $-(4K'+3)$ qui est résidu, et $1\equiv -2(4K'+3)$ montre que $2$ est résidu. Enfin, si $p=4n+1$ où $n$ contient une puissance impaire de $2$, on pose $N=2n$, qui contient une puissance paire de $2$ : $-1$ est résidu, les facteurs premiers impairs de $N$ aussi, donc $N$ l'est, et $-1\equiv 2N$ donne $2$ résidu.

C'est juste, avec une raison que la page ne donne pas : pour un facteur premier $q$ de $\frac{p-1}{2}$, on a $p\equiv 1\pmod q$, donc $\left(\frac{p}{q}\right)=1$, et la loi de réciprocité donne $\left(\frac{q}{p}\right)=1$ si $q\equiv 1\pmod 4$ et $\left(\frac{-q}{p}\right)=1$ si $q\equiv 3\pmod 4$. La démonstration suppose donc la réciprocité pour deux nombres premiers impairs. Elle couvre $p\equiv 3$ et $7\pmod 8$, et les $p\equiv 1\pmod 8$ pour lesquels la puissance de $2$ qui divise exactement $p-1$ a un exposant impair ; elle ne dit rien quand cet exposant est pair, soit pour $p\equiv 5\pmod 8$ et pour $p=2^{2m}n+1$, $n$ impair, $m\ge2$. C'est ce que dit la dernière ligne : « cette méthode n'apprend rien pour le cas $p=2^{2m}n+1$ »\footnote{La première phrase dit au contraire que la méthode ne distingue pas les $8K+5$ des $8K+1$ « dans lesquels $K$ est divisible par une puissance paire de 2 » ; or $p=8K+1$ échappe à la méthode quand $K$ contient une puissance \emph{impaire} de $2$ (par exemple $p=17$). La phrase d'ouverture contredit la ligne finale et le calcul ; la lecture de « paire » est à revoir sur l'image.}. Vérification : le second complément a été contrôlé pour tous les $p<5000$, et la ligne de partage de la méthode aussi (non couverts en dessous de 300 : $5, 13, 17, 29, 37, 53, 61, \dots$).

\subsection*{Les nombres premiers $4n+1$ résidus les uns des autres}

\pagerange{706}{706}
Sous le titre « Observations sur les nombres premiers de la forme $4n+1$ considérés comme résidus d'autres nombres premiers de même forme », la page part de deux nombres premiers $P=4N+1$ et $p=4n+1$ avec $P$ résidu de $p$ : alors $-P$ l'est aussi, et $X^2+P=Mp$\footnote{Vue 706, feuillet 347 ; en haut à gauche, un « 1 » suivi d'un trait oblique, peut-être un numéro de feuillet de sa main. La même lettre $p$ désigne sur la page le nombre premier $4n+1$ et l'un des quatre carrés de $M$ ; on les distingue ici en appelant $\pi$, $\rho$, $\sigma$, $\tau$ ces derniers. Le feuillet est écrit jusqu'au bord inférieur et s'arrête sur un « 2 » souligné ; la suite n'est pas dans le volume.}. Elle écrit $P=A^2+B^2$, $p=a^2+b^2$ (tout nombre premier $4n+1$ est somme de deux carrés) et $M=\pi^2+\rho^2+\sigma^2+\tau^2$ (tout entier est somme de quatre carrés, théorème de Lagrange), et développe exactement
\[ X^2+A^2+B^2=(a^2+b^2)(\pi^2+\rho^2+\sigma^2+\tau^2)=(a\pi\pm b\rho)^2+(b\pi\mp a\rho)^2+(a\sigma\pm b\tau)^2+(a\tau\mp b\sigma)^2. \]
Elle essaie alors d'égaler les termes un à un ($a\pi=b\rho$, $b\pi+a\rho=A$, …), trouve que $A$ serait multiple de $p$, et affirme que ces suppositions « sont les seules par lesquelles on puisse égaler terme à terme les deux membres ». L'inférence n'est pas valable : une égalité de deux sommes de carrés n'entraîne aucune égalité terme à terme, les décompositions n'étant pas uniques. La note qui suit (« Je ne vois pas d'ailleurs que l'on puisse affirmer que le produit de deux nombres qui ne seroient, l'un et l'autre, décomposables en 4 quarrés que d'une seule façon, ne le soit que de 2 … ») semble toucher cette difficulté\footnote{Une note marginale en quatre lignes serrées, appelée par un signe en face de la dernière ligne, se lit « c'est à dire par ces manières que en produisent aucune », « produisent » avec doute.}. Ce que la page cherche est vrai : pour deux nombres premiers $P$ et $p$ de la forme $4n+1$, $\left(\frac{P}{p}\right)=\left(\frac{p}{P}\right)$ ; c'est le cas $p\equiv q\equiv 1\pmod 4$ de la loi de réciprocité quadratique. La page n'y parvient pas et s'interrompt.

\subsection*{En lisant « l'auteur » : tout nombre non carré est non-résidu d'un nombre premier}

\pagerange{718}{728}
Ces billets commentent un livre qu'ils ne nomment qu'une fois : « En examinant la manière dont Gauss démontre (N\textsuperscript{o} 125) que tout nombre premier $p=4n+1$ pris négativement est non résidu d'un nombre premier moindre que $p$ … » (vue 720)\footnote{Vues 718 à 728, feuillets 353 à 363, billets montés deux par feuille de garde, parfois numérotés de bas en haut (361 au-dessus de 360, 363 au-dessus de 362). Seule la vue 720 nomme Gauss. Les autres renvois, « P. 125 », « N\textsuperscript{o} 139, § 146 » (vue 722), « §126 » (vue 724), « P. 116 » (vues 724 et 728), et « l'auteur » des vues 718, 720 (bas), 722 et 728, sont rapportés aux articles des \emph{Disquisitiones arithmeticae} par une inférence de la transcription, que rien sur ces feuillets ne confirme au-delà de la vue 720 ; « P. » pourrait être une page ou une abréviation d'article.}. Leur méthode est toujours la même : pour montrer qu'un nombre $p$ est non-résidu d'un nombre premier $q<p$, on construit un nombre $p\pm c\,a^2$ plus petit que $p$, dont tous les facteurs premiers ne peuvent être dans les classes où $c$ est résidu ; un facteur premier $q$ pris dans l'autre classe donne $p\equiv\mp c\,a^2\pmod q$, donc $\left(\frac{p}{q}\right)=\left(\frac{\mp c}{q}\right)=-1$.

\emph{La proposition générale.} « Tout nombre, excepté les quarrés positivement pris, est non résidu de quelque nombre premier » (vue 722)\footnote{Vue 722, billet 358, qui s'ouvre sur « P. 125 L'auteur se borne à démontrer le cas des nombres premiers de la forme $4n+1$ parce qu'il n'emploie que celui-là dans sa démonstration de son théorème (V. N\textsuperscript{o} 139. § 146) » ; « fondamental » est barré après « théorème », « P. », « V. » et « § » sont lus avec doute. Le billet finit sur « Retournez ».}. C'est vrai : un entier qui n'est pas un carré n'est pas résidu de tous les nombres premiers qui ne le divisent pas ; on le démontre aujourd'hui à l'aide du symbole de Jacobi et du théorème chinois\footnote{Le symbole de Jacobi (1837) est postérieur à la lecture de Gauss que suppose le billet ; la vérification de cette lecture a porté sur les entiers de $-300$ à $300$ et les nombres premiers inférieurs à $2000$.}. Elle réduit la question : admettant le cas des nombres $4n+1$, restent $4n+2$ et $4n+3$, les nombres $4n$ étant non-résidus avec $n$. Pour $4n+2$ : $4n+2+(2k+1)^2$ est de la forme $4k+3$, a donc un facteur premier $q\equiv 3\pmod 4$, et $4n+2\equiv-(2k+1)^2$ avec $\left(\frac{-1}{q}\right)=-1$ donne $4n+2$ non-résidu de $q$. C'est juste. « Je n'ai pas vu comment démontrer $-(4n+2)$ non résidu ». Au dos (vue 723), $\pm(4n+3)$ sont traités par $4n+3+4k^2$ et $4n+3-4k^2$, tous deux de la forme $4k+3$\footnote{Vue 723, dos des billets 357 et 358. Le feuillet écrit « Soit $4n+3\equiv4k^2$ » avant de considérer $4n+3-4k^2$ ; le calcul demande seulement de prendre $4k^2<4n+3$.} : d'où la « proposition remarquable » que tout nombre $4k+3$, pris avec le signe $+$ ou $-$, est non-résidu d'un nombre premier (de la forme $4k+3$). C'est juste.

\emph{Les nombres $4K+3$, avec la condition « moindre ».} Le feuillet 356 (vue 720) se propose de démontrer qu'un nombre quelconque $p=4K+3$, pris avec l'un ou l'autre signe, est non-résidu de quelque nombre premier moindre que lui. Pour $-p$ : avec $2a<\sqrt p$, le nombre $p-4a^2$ est positif, $<p$, de la forme $4K+3$, donc divisible par un premier $q\equiv3\pmod 4$, et $p\equiv 4a^2$ donne $-p$ non-résidu de $q$. Pour $+p$ : si $p=8n+3$, avec $2a<\sqrt{p/2}$, $p-8a^2\equiv 3\pmod 8$ a un facteur premier $q\equiv 3$ ou $5\pmod 8$, dont $2$ est non-résidu, et $p\equiv 2(2a)^2$ conclut\footnote{Le signe entre « immédiatement » et « que $\sqrt{p/2}$ » est noyé dans un pâté et n'est pas lu ; le calcul demande « plus petit », ce que confirme le dernier alinéa (« il suffit de prendre $2a<\sqrt{p/2}$ ») : inférence de cette lecture.} ; si $p=8K+7$, avec $a$ impair $<\sqrt{p/2}$, $p-2a^2\equiv 5\pmod 8$ a de même un facteur $q\equiv 3$ ou $5\pmod 8$, et $p\equiv 2a^2$ conclut. Tout est juste, à deux réserves près que la page ne fait pas : il faut $a\ge1$, ce qui exclut $p=3$ (aucun nombre premier plus petit ne convient) ; et, si $p$ n'est pas premier, que $q$ ne divise pas $p$. La remarque finale, que la condition « immédiatement plus petit » est superflue et qu'il suffit de $2a<\sqrt p$, $2a<\sqrt{p/2}$, est juste. Au pied du feuillet, d'une plume plus épaisse : « Le but de l'auteur étant de passer au théorème fondamental, il s'est borné à démontrer que les nombres $4k+3$ tant positifs que négatifs sont non résidus. Il faut examiner si on ne pourroit pas faire usage de la même pro[position] pour en … », où la ligne s'arrête\footnote{« pro. » est lu avec doute ; « théorème fondamental » est le nom que Gauss donne à la loi de réciprocité (contexte, non tiré des feuillets).}.

\emph{Une extension fausse.} Le feuillet 355 (même vue) étend l'énoncé qu'il prête à Gauss : la même démonstration vaudrait pour $p$ non premier, « de sorte que l'on peut affirmer que tout nombre de la forme $4n+1$ pris négativement est non résidu de quelque nombre premier moindre que lui et de la forme $4n+3$ » ; avec le feuillet 356, « tout nombre impair $2N+1$ pris avec le signe $-$ est non résidu d'un nombre premier de la forme $4K+3$ et $<2N+1$ »\footnote{Le second « $4n+1$ » de ce feuillet est lu avec doute (chiffres récrits et empâtés) ; « $2N+1$ » est écrit au-dessus de « impair ».}. C'est faux : $-5$, $-17$ et $-21$ sont résidus de tous les nombres premiers de la forme $4k+3$ plus petits qu'eux et qui ne les divisent pas ($-17\equiv 1\pmod 3$, $\equiv 4=2^2\pmod 7$, $\equiv 5=4^2\pmod{11}$) ; par le calcul de cette lecture, ce sont, avec $3$, les seules exceptions en dessous de $200\,000$. Ce que la page prête à Gauss, pour $p$ premier de la forme $4n+1$ et sans restreindre la forme de $q$, est vrai pour $p>5$ (vérifié jusqu'à $5000$) ; c'est la restriction à $q\equiv 3\pmod 4$ qui est fausse.

\emph{Les nombres $12n+7$ et $12n+5$.} Au billet 354 (vues 718 et 719) : « Ici, l'auteur ne dit rien des nombres de la forme $4k+3$ »\footnote{Vue 718. Au-dessus, le feuillet 353 ne porte de ce côté que la fin d'une phrase, « est non résidu d'un nombre premier de la même forme $4k+3$ », dont le début n'est pas dans le volume.}. Pour $p=12n+7$, avec $a$ le nombre pair immédiatement supérieur à $\sqrt{p/3}$, le nombre $3a^2-p$ est de la forme $12K+5$ ; il a un facteur premier $q\equiv 5$ ou $7\pmod{12}$ (les produits de nombres $\equiv 1$ ou $11$ restent $\equiv1$ ou $11$), et pour ces $q$ on a $\left(\frac{3}{q}\right)=-1$, avec $\left(\frac{-3}{q}\right)=-1$ si $q\equiv5$ et $+1$ si $q\equiv7\pmod{12}$ : $p\equiv 3a^2$ est non-résidu de $q$. Au dos, même chose pour $p=12n+5$ avec $p-3a^2$, $a$ pair immédiatement inférieur à $\sqrt{p/3}$. Les caractères de $3$ et $-3$ sont justes. Mais l'affirmation « $3aa-p$ sera $<$ que $p$ » est fausse pour $p=19$ : $a=4$ et $3a^2-p=29$ ; c'est la seule exception par le calcul de cette lecture, et $19$ est bien non-résidu de $7$. Au dos, il faut $a\ge 2$, ce qui exclut $p=5$.

\emph{Une page barrée.} Le revers du feuillet 353 (vue 719), écrit tête-bêche et barré d'une grande croix, attaque les nombres $120k+1$ : avec $10n+2>\sqrt{\tfrac52(120k+1)}$, le nombre $2(10n+2)^2-5(120k+1)$ serait de la forme $40k+11$, dont un facteur $q$ aurait $2$ non-résidu et $5$ résidu, d'où $120k+1$ non-résidu de $q$ ; suit une table des produits de classes $\equiv 11\pmod{40}$\footnote{Vue 719. Les parenthèses « (ou 11) », « (ou 19) » sont lues avec doute ; « $4k+11=q$ » est écrit pour $40k+11$. La table : $29\cdot39$, $9\cdot19$, $21\cdot31$, $33\cdot27$, $7\cdot13$, $37\cdot23$, $3\cdot17$, tous $\equiv 11\pmod{40}$, ce qui est exact.}. Le calcul est faux dès le départ : $2(10n+2)^2-5(120k+1)\equiv 3\pmod{40}$, et non $11$ ; pour un nombre $\equiv 3\pmod{40}$, le symbole de Jacobi $\left(\frac{10}{\cdot}\right)$ vaut $+1$, et rien ne garantit plus un facteur $q$ avec $\left(\frac{10}{q}\right)=-1$, qui est ce dont l'argument a besoin. Dans la table, deux mentions sont fausses : pour $q\equiv 27$ et $q\equiv 13\pmod{40}$, $5$ est non-résidu, non résidu ; dans ces deux paires, c'est l'autre classe ($33$, $7$) qui a $\left(\frac{10}{q}\right)=-1$.

\emph{Les nombres $3k-1$ et $8n+1$.} Le billet 357 (vue 722) montre que tout nombre $p=3k-1$ est non-résidu d'un nombre premier par $p+3a^2$, de la forme $3k-1$, dont un facteur premier $q\equiv 2\pmod 3$ a $\left(\frac{-3}{q}\right)=-1$ : juste, avec sa précaution sur la parité de $a$. Il ajoute qu'on doit pouvoir trouver ainsi $q<p$, « mais je ne sais pas le démontrer ». Sous cette forme, c'est faux : pour $p=5$, $11$, $14$, $26$ et $59$ (seules exceptions en dessous de $5000$ par le calcul de cette lecture), aucun nombre premier impair $q\equiv 2\pmod 3$ plus petit que $p$ n'a $p$ pour non-résidu. Sans restreindre la forme de $q$, l'énoncé est immédiat : $q=3$ convient pour tout $p=3k-1>3$. Au dos, les nombres à la fois $3k-1$ et $8n+1$ sont les $24n+17$, ce qui est juste\footnote{Vue 723. Le 2 de « $24n+17$ » est sous une tache et lu avec doute ; le calcul le confirme.}. Sous le titre « §126 » (vue 724), un nombre premier $p=8n+1$ est de la forme $24k+1$ ou $24k+17$ ; pour $24k+17$, qui est $12k+5$, la page reprend l'argument de $p-3a^2$ avec $a$ pair $<\sqrt{p/3}$ (renvoyant à « P. 116 » pour $3$ non-résidu des $12k+5$ et $12k+7$), et note justement que seuls les produits $12k+1$ par $12k+5$ et $12k+7$ par $12k+11$ sont de la forme $12k+5$ : c'est juste. Pour $24k+1$, elle part de $p=x^2+y^2$ et observe que $x\equiv\pm1$, $y\equiv\pm2\pmod 5$ donnerait $5\mid x^2+y^2$, puis s'arrête\footnote{Vue 724, feuillet 359, écriture plus posée. Le feuillet écrit $25n$ là où le carré de $5n$ donne $25n^2$ ; la congruence finale $\equiv 0\pmod 5$ est juste.}.

\emph{Formes et induction.} Le feuillet 361 (vue 726) énonce que si $5$ est non-résidu des nombres premiers d'une forme $F$, un produit de nombres premiers dont $5$ est résidu ne peut être de la forme $F$ ; de même pour $-3$ et les formes $3n\pm1$. C'est juste si $F$ est une réunion de classes modulo $5$ (les $5n\pm2$) : $\left(\frac{5}{q}\right)=\left(\frac{q}{5}\right)$, et un produit de nombres $\equiv\pm1\pmod 5$ reste $\equiv\pm1$ ; de même $\left(\frac{-3}{q}\right)=\left(\frac{q}{3}\right)$. Le feuillet 360 tente une induction : si tous les nombres premiers jusqu'à $P$ dont $a$ est résidu sont de la forme $ma+b^2$, prendre $e$ pair, non divisible par $a$, avec $e^2=a+fP$ et $f<P$ impair ; les facteurs de $f$ ont $a$ pour résidu, sont plus petits que $P$, donc de la forme $ma+b^2$, qui est stable par produit, d'où $P$ de même forme. L'énoncé visé est faux en général : $3$ est résidu de $11$ ($5^2\equiv 3$), et $11\equiv 2\pmod 3$ n'est pas de la forme $3m+b^2$ ; il est vrai si $a$ est un nombre premier de la forme $4k+1$, et c'est alors la loi de réciprocité. La faille de l'induction est l'existence d'un $e$ convenable : pour $a=3$, $P=11$, les solutions de $e^2\equiv 3\pmod{11}$ inférieures à $11$ sont $5$ et $6$, et la seule paire, $6$, est divisible par $3$ : aucun $e$ ne remplit les conditions\footnote{$e=6$ donne $f=3$, qui est $a$ lui-même, et la division par $f$ modulo $a$ n'a plus de sens ; un choix $e>P$ donne $f>P$, hors de l'hypothèse de récurrence.}.

\emph{Racine primitive.} Le billet 363 (vue 728), annulé de longs traits verticaux mais lisible, démontre que pour un nombre premier $p=4k+3$, si $n$ est non-résidu, $-n$ est résidu : avec une racine primitive $r$ et $n\equiv r^{2f+1}$, $-n\equiv r^{2g+1}$ donnerait $r^{2(g-f)}\equiv -1$, soit $-1$ résidu. C'est juste. « L'induction a mené l'auteur à établir pour les nombres $12n+11$ $3$ résidu et $-3$ non résidu » : juste. « Je ne vois pas que l'on puisse rien conclure pour ces nombres de la formule $\frac{4(a^3-1)}{a-1}=(2a+1)^2+3$ »\footnote{Le feuillet écrit $(2x+1)^2$ ; le dénominateur $a-1$ est coupé par un trait d'annulation et lu avec doute ; le reste du billet, un calcul sur les nombres $4k+1$, s'arrête sur une congruence dont le second membre ne se lit pas.} : l'identité est exacte, et de fait elle ne peut rien donner, puisqu'un nombre premier divisant $a^2+a+1$ a $-3$ pour résidu, ce qui n'est jamais le cas d'un $12n+11$. Le billet 362 porte six lignes d'algèbre sur $A=ma+b^2$, $a=d^2-nA$, qui donnent correctement $(1+mn)A=md^2+b^2$ et $(1+mn)a=d^2-nb^2$.

\subsection*{Racines primitives, et les nombres premiers de Fermat}

\pagerange{729}{730}
« Pour tout nombre premier $12n+5$ dans lequel $3n+1$ est premier, $3$ est racine primitive »\footnote{Vue 729, prise en largeur : dos du billet 363 (et vraisemblablement de 362). Les exemples sont en colonne, sous « tels sont » souligné ; $4\cdot19+1=77$ et $4\cdot23+1=93$ sont barrés, à juste titre ($77=7\cdot11$, $93=3\cdot31$). Trois lignes de calcul barrées suivent (« Soit à présent $3n+1=p^n$ … »), en partie illisibles.}. C'est vrai, et l'argument de la page l'établit : $p-1=4(3n+1)$ ; $3$ est non-résidu des nombres $12n+5$, donc son ordre ne divise pas $2(3n+1)$ ; s'il n'était pas $p-1$, il serait $4$, et $3^4\equiv 1$ donnerait $3^2\equiv-1$, soit $p\mid 10$, impossible pour $p>5$\footnote{La page écrit « on aura $r^{3n+1}\equiv 3$ » ; l'exposant peut aussi être $3(3n+1)$, ce qui ne change rien à $3^4\equiv1$. Elle passe de $3^4\equiv1$ à $3^2+1\equiv0$ sans écarter $3^2\equiv 1$, qui demanderait $p\mid 8$.}. Exemples de la page : $29$ et $53$ ; le calcul de cette lecture confirme $3$ racine primitive pour les $97$ nombres premiers de cette espèce inférieurs à $20\,000$.

Sous le titre « Nombres premiers de la forme $2^{2^i}+1$ » (vue 730)\footnote{Vue 730, feuillet 364, écriture posée. Le « T. » qui ouvre deux paragraphes n'est pas expliqué sur la feuille ; il introduit la réciprocité « si $a$ R $a'$, $a'$ R $a$ », vraie pour deux nombres premiers dont l'un au moins est de la forme $4k+1$, ce qui est le cas des nombres premiers de Fermat $>3$.}, la page part de la réciprocité : pour $a'>a$, $a'$ est résidu de $a$ si $a'=ma+d^2$ avec $d<a$. Elle affirme que, pour tout nombre premier $F_k=2^{2^k}+1$, les nombres premiers de Fermat plus petits sont non-résidus, donc racines primitives, parce que $2^{2^k}+1$ ne peut se mettre sous la forme $mF_j+d^2$ qu'en faisant de $2^{x}+1$ un carré, ce qui n'arrive que pour $x=3$. Entre les lignes, d'une plume plus fine : « cela est faux parce qu'il suffit que $2^{2^{i+\delta}-2^i m}+1\equiv x^2$ ». La correction est juste, et la conclusion est en partie fausse. Comme $2^{2^k}\equiv 1\pmod{F_j}$ pour $j<k$, on a $F_k\equiv 2\pmod{F_j}$ et $\left(\frac{F_j}{F_k}\right)=\left(\frac{2}{F_j}\right)$ : ce symbole vaut $-1$ pour $F_0=3$ et $F_1=5$, mais $+1$ pour $F_2=17$ et au-delà (tous $\equiv 1\pmod 8$). Ainsi $3$ et $5$ sont non-résidus de tout nombre premier de Fermat plus grand, et, $F_k-1$ étant une puissance de $2$, tout non-résidu est racine primitive : $3$ et $5$ sont racines primitives de tout nombre premier de Fermat plus grand ; en revanche $17$ est résidu de $257$ et de $65537$, et $257$ de $65537$\footnote{Que $3$ soit racine primitive de tout nombre premier de Fermat $F_k$, $k\ge1$, est à la base du test de Pépin (1877) : $F_k$ est premier si et seulement si $3^{(F_k-1)/2}\equiv-1\pmod{F_k}$. La feuille ne le cite pas et le précède sans doute.}. Le second paragraphe applique le même raisonnement (exiger un carré exact) à $a=F_i$ et $b=2^n-1$, et conclut que $3=2^2-1$ et $7=2^3-1$ sont non-résidus de $F_i$, « c'est ce que j'avois déjà trouvé d'une autre manière ». Le raisonnement a le défaut que la page vient de relever, mais la conclusion est vraie pour $i\ge1$ : $\left(\frac{3}{F_i}\right)=\left(\frac{2}{3}\right)=-1$ et, $F_i\equiv 3$ ou $5\pmod 7$, $\left(\frac{7}{F_i}\right)=-1$ ; $7$ est donc aussi racine primitive de tout nombre premier de Fermat $\ge5$ (vérifié pour $5$, $17$, $257$, $65537$).

\section*{Le mémoire imprimé de Legendre sur l'équation $4X = Y^2 \pm nZ^2$}

\pagerange{115}{134}
Au milieu des brouillons sur la courbure sont collées, sur des feuillets de montage, les pages imprimées d'un mémoire de Legendre : « Mémoire sur la détermination des fonctions $Y$ et $Z$ qui satisfont à l'équation $4(x^n-1)=(x-1)(Y^2\pm nZ^2)$, $n$ étant un nombre premier $4i\mp1$ », « Lu à l'Académie, le 11 octobre 1830 »\footnote{Vues 115 à 134, feuillets 56 à 65 de la série à l'encre, chaque page imprimée collée sur un feuillet de montage ; au pied de la première page, « T. XI. » et la signature « 11 ». Pagination imprimée 81 à 88 (vues 115 à 123), puis 9 à 19 (vues 124 à 134), le texte continuant sans lacune ; d'après les signatures (« 11 », puis « 2 », « 3 »), la transcription y voit la fin d'un cahier du volume et un tirage à part paginé de nouveau : c'est une inférence tirée des marques, non une lecture. La vue 122 est une seconde prise de vue de la page 87. Un filet imprimé ferme la page 19 et, semble-t-il, le mémoire. Les feuillets ne disent pas pourquoi cet imprimé est relié ici.}. Pour $n$ premier impair, $X=\frac{x^n-1}{x-1}$ et $m=\frac{n-1}{2}$, il s'agit de l'identité, due à Gauss\footnote{\emph{Disquisitiones arithmeticae}, art. 357 (contexte, non tiré des feuillets ; l'imprimé renvoie aux articles 511 et 512 de la \emph{Théorie des nombres} de Legendre, troisième édition).},
\[ 4X=Y^2-(-1)^{m}\,n\,Z^2, \]
avec $Y$ de degré $m$ et $Z$ de degré $m-1$ à coefficients entiers : $Y^2-nZ^2$ pour $n=4i+1$, $Y^2+nZ^2$ pour $n=4i-1$, ce que le titre condense en « $Y^2\pm nZ^2$, $n=4i\mp1$ ». Le point de départ de Legendre est que ses règles de la \emph{Théorie des nombres} supposaient les coefficients de $Y$ inférieurs à $\frac12 n$, ce qui cesse d'avoir lieu pour $n$ grand ; il en donne ici une méthode « sûre et exacte ». Il rappelle au passage que $Y=2x^m+x^{m-1}+\frac{3\mp n}{4}x^{m-2}+\dots$ n'a aucun coefficient nul : c'est exact, car $Y^2\equiv 4X=4(x-1)^{n-1}\pmod n$ donne $Y\equiv\pm2(x-1)^m\pmod n$, dont aucun coefficient n'est divisible par $n$.

\emph{La méthode.} Posons $n^*=(-1)^m n$ et $Y+Z\sqrt{n^*}=2x^m+A_1x^{m-1}+A_2x^{m-2}+\dots$ Avec $r$ racine primitive $n$-ième de l'unité et $\alpha$ parcourant les résidus quadratiques de $n$, on a $\frac12(Y+Z\sqrt{n^*})=\prod_\alpha(x-r^\alpha)$\footnote{Vues 116 à 118 (pp. 82 à 84). La note infrapaginale de la page 83 rappelle que les $\alpha$ sont aussi les carrés $1,4,9,\dots,m^2$ réduits modulo $n$. Le signe de $\sqrt n$ « peut être pris à volonté ».}. Les sommes de puissances $S_k=\sum_\alpha r^{k\alpha}$ valent $-\frac12-\frac12\sqrt{n^*}\left(\frac{k}{n}\right)$ (le signe de la racine étant fixé par $S_1$), et les identités de Newton
\[ -k\cdot\tfrac12A_k=S_k+S_{k-1}\cdot\tfrac12A_1+\dots+S_1\cdot\tfrac12A_{k-1} \]
donnent les $A_k=a_k+b_k\sqrt{n^*}$, d'où $Y=\sum a_kx^{m-k}$ et $Z=\sum b_kx^{m-k}$. La méthode est juste ; pour $n=4i-1$, il suffit, dit l'imprimé, de mettre $-n$ à la place de $n$ dans $S_1$\footnote{Vue 119 (p. 85). Dans « les valeurs des coefficients $A$, $A_2$, $A_3$ », l'indice du premier manque. Dans l'exemple I, l'imprimé écrit encore $P=-\frac12-\frac12\sqrt n$ pour $n=31$, là où la règle qu'il vient de donner demande $\sqrt{-n}$ ; les coefficients qui suivent sont bien en $\sqrt{-n}$.}.

\emph{Les exemples.} Pour $n=31$, de la forme $4i-1$, les symboles $\left(\frac{k}{31}\right)$ de l'imprimé sont exacts\footnote{Vues 120 et 121 (pp. 86 et 87). En particulier $\left(\frac{7}{31}\right)=-\left(\frac{31}{7}\right)=-\left(\frac37\right)=\left(\frac73\right)=1$, et en effet $10^2\equiv7\pmod{31}$.}, et les coefficients $A_1=1+\sqrt{-n}$, $A_2=-7+\sqrt{-n}$, …, $A_{15}=-2$ donnent
\[ Y=2x^{15}+x^{14}-7x^{13}-11x^{12}+2x^{11}+8x^{10}-3x^9-5x^8+5x^7+3x^6-8x^5-2x^4+11x^3+7x^2-x-2, \]
\[ Z=x^{14}+x^{13}-x^{12}-2x^{11}+x^9-x^8-x^7+x^6-2x^4-x^3+x^2+x, \]
qui satisfont $4X=Y^2+31Z^2$. Pour $n=37$, l'exemple II donne $Y$ de degré $18$ et $Z$ de degré $17$ avec $4X=Y^2-37Z^2$\footnote{Vues 121 et 123 (pp. 87 et 88).}. La table de la page 9 ajoute $n=41$ et $n=61$ (vue 124) ; pour $41$ des coefficients dépassent $\frac12n$, pour $61$ plusieurs dépassent $n$ ($63$, $72$, $89$, $95$). Pour $n=2521$, tous les $\left(\frac{k}{n}\right)$, $k\le 7$, valent $1$ (en effet $2521\equiv1$ modulo $8$, $3$, $5$ et $7$, et la réciprocité conclut), et $A_7=2633212+69292\sqrt n$ (vue 125) montre la croissance rapide des coefficients. Toutes ces valeurs ont été recalculées par les identités de Newton, et l'identité $4X=Y^2\mp nZ^2$ vérifiée pour $n=31$, $37$, $41$, $61$ : l'imprimé est exact partout.

\emph{Les signes.} Les coefficients également éloignés des extrêmes sont égaux au signe près, « puisque la substitution de $x^{-1}$ à la place de $x$ » ne change pas l'équation. Les pages 11 à 13 décident des signes : pour $n=4i-1$, $Y$ est antisymétrique ($x^mY(1/x)=-Y(x)$) et $Z$ symétrique, car l'autre choix donnerait $Z(1)=0$ et $4n=Y(1)^2$, ou $Y(-1)=0$ et $4=nZ(-1)^2$ ; pour $n=4i+1$, $Y$ et $Z$ sont symétriques, le coefficient moyen de $Z$ ne pouvant être nul et $Y$ antisymétrique donnant $Y(1)=0$ et $4n=-nZ(1)^2$\footnote{Vues 126 à 128 (pp. 11 à 13). Coquilles de l'imprimé, relevées par la transcription : dans la forme II de $Z$ (p. 11), le dernier terme est précédé de $+$ là où la suite des signes demande $-$ ; à la p. 12, $b^{i-1}$ pour $b_{i-1}$, et $+b_2x^2$ entre $-x$ et $-b_3x^3$ là où le calcul donne $-b_2x^2$ ; la même p. 12 rapporte le cas $n=4i+1$ à $4X=Y^2+nZ^2$, la p. 13 à $4X=Y^2-nZ^2$, qui est la bonne équation (l'argument par $x=1$ vaut avec l'un ou l'autre signe).}. C'est juste ; la raison de fond, que l'imprimé ne donne pas, est que $\alpha\mapsto-\alpha$ conserve les résidus si $n\equiv1\pmod 4$ et les échange avec les non-résidus si $n\equiv 3\pmod 4$. Les exemples s'y accordent.

\emph{Les formules générales.} Selon $n=24\lambda+1$, $+5$, $+13$, $+17$, les caractères de $2$ et $3$ sont fixés, ce qui donne $A_1$ à $A_4$ en fonction de $\lambda$ ; le passage aux $n=24\lambda-1$, $-5$, $-13$, $-17$ se fait en changeant $\lambda$ en $-\lambda$ et $n$ en $-n$\footnote{Vues 129 à 131 (pp. 14 à 16), avec le tableau des $A_1$ à $A_4$ pour $n=4i-1$.}. Toutes ces formules ont été vérifiées sur tous les nombres premiers de chaque forme jusqu'à $24\cdot 40$. Le coefficient $1-11\lambda+3\lambda^2$, de l'ordre de $n^2/192$, dépasse $n$ pour $\lambda=13$, $n=311$ ($365$) ; les seize sous-formes $120\mu+r$ sont bien celles qu'on obtient en rejetant, dans chaque cas, la valeur de $\lambda$ modulo $5$ qui rend $n$ multiple de $5$ ; et pour $n=120\mu+61$, les valeurs de $A_5$, $A_6$, puis de $Y$ et $Z$ jusqu'à $x^{m-6}$ pour $120\mu\pm61$, sont exactes\footnote{Vues 132 à 134 (pp. 17 à 19). Une coquille : pour $n=120\mu+61$, l'imprimé écrit $\left(\frac6n\right)=1$, alors que $n\equiv5\pmod8$ donne $\left(\frac2n\right)=-1$ et $\left(\frac6n\right)=-1$ ; la conséquence qu'il en tire, $S_6=S_2$, est celle de $-1$ et elle est juste. À la vue 132, « de l'ordre $\mu$ » sans exposant lisible, là où $A_6$ contient $75\mu^3$. La subdivision en $840\nu+61$, $181$, $421$, $541$, $661$, $781$ omet à juste titre $\mu=7\nu+2$, qui donne $301=7\cdot43$.}. Le seuil $75\mu^3>n^2$ pour $n>120^3/75=23\,040$ est une estimation, juste à l'ordre près.

\section*{Extraits des Mémoires de Turin}

\pagerange{734}{749}
Le volume finit sur des extraits en français, tantôt traduits, tantôt résumés, de mémoires des \emph{Mélanges} de l'Académie de Turin, d'une main régulière et penchée, peu corrigée ; les transcriptions ne s'accordent pas sur cette main, et rien sur les feuillets ne la désigne. Des remarques entre crochets, à la première personne, sont du copiste, dont les feuillets ne donnent pas l'identité\footnote{Vues 734 à 749, feuillets 365 à 372 de la série à l'encre, grands feuillets montés sur des feuilles de garde et écrits des deux côtés ; le numéro est au recto seulement. Le lot 37 voit dans les crochets des remarques de sa main à elle, le lot 38 laisse la question ouverte ; cette lecture ne tranche pas. La vue 750, dernier feuillet, est blanche, avec un « 373 » barré.}. Les dérivées partielles y sont écrites à la manière d'Euler, entre parenthèses\footnote{Le feuillet écrit $\left(\frac{d^2z}{dt^2}\right)$, $\left(\frac{ddy}{dx^2}\right)$ ; on écrit $\frac{\partial^2z}{\partial t^2}$, $\frac{\partial^2y}{\partial x^2}$. Il écrit aussi les fonctions arbitraires avec le deux-points d'Euler, $\Gamma{:}(x+ct)$, qu'on écrit $\Gamma(x+ct)$.}.

\subsection*{Lagrange, la théorie du son, et la somme de d'Alembert}

\pagerange{734}{739}
Le premier extrait, « Mémoire de Turin [2\textsuperscript{d} ?] année 1761 Pag. 16 Théorie du son La Grange »\footnote{Vues 734 et 735, feuillet 365. « 2d ? », en exposant au-dessus de « 1761 », est lu avec doute. La lettre initiale de la courbe se lit E ou P, la deuxième K ou H selon les occurrences ; l'arc « $t-ux$ » a un u lu avec doute.}, conteste que la courbe $EKS$ de Newton (\emph{Principes}, livre II, proposition 49), qui règle l'excursion de chaque particule de l'air, puisse être un cercle : les ébranlements primitifs dépendent du corps sonore, et un cercle, courbe rentrante, mettrait toutes les particules en mouvement à la fois, ce qui détruirait la propagation du son. Il va plus loin : la fonction $\varphi\left(t-\frac{xT}{\sqrt{2hA}}\right)$ qui représente les excursions doit, à l'instant initial, s'annuler identiquement hors de la petite portion d'air ébranlée, ce qu'aucune courbe « soit géométrique soit transcendante » ne peut faire ; d'où « la nécessité … d'employer un nouveau genre de fonctions … qu'on peut très bien appeler fonctions irrégulières et discontinues ». En termes actuels, c'est juste pour les fonctions analytiques : une fonction analytique réelle nulle sur un intervalle est nulle partout ; une donnée initiale à support borné ne peut donc être analytique, et il faut admettre des fonctions seulement continues, ou moins régulières encore\footnote{Le principe du prolongement analytique n'a été formulé que plus tard (Cauchy, Weierstrass) ; l'extrait ne le connaît pas.}.

Suit le « Problème » : les mouvements d'une infinité de points mobiles, les extrémités $x=0$ et $x=a$ fixes, selon $\frac{\partial^2z}{\partial t^2}=c\,\frac{\partial^2z}{\partial x^2}$. On multiplie par $M(x)\,dx$ et l'on intègre deux fois par parties :
\[ \int_0^a\frac{\partial^2z}{\partial x^2}M\,dx=\left[\frac{\partial z}{\partial x}M-z\frac{dM}{dx}\right]_0^a+\int_0^a z\,\frac{d^2M}{dx^2}\,dx ; \]
les termes aux bornes disparaissent si $M(0)=M(a)=0$. Avec $M''=KM$ et $s(t)=\int_0^a zM\,dx$, il vient $s''=cKs$ ; $M$ nul aux deux bouts impose $K<0$, $M=A\sin\bigl(x\sqrt{-K}\bigr)$ et $\sqrt{-K}=\nu\pi/a$, $\nu$ entier ; alors
\[ s=S\cos\bigl(t\sqrt{-cK}\bigr)+\frac{R}{\sqrt{-cK}}\sin\bigl(t\sqrt{-cK}\bigr),\qquad \frac{ds}{dt}=R\cos\bigl(t\sqrt{-cK}\bigr)-S\sqrt{-cK}\,\sin\bigl(t\sqrt{-cK}\bigr), \]
$S=\int_0^a Z\,M\,dx$ et $R=\int_0^a V\,M\,dx$ étant formés sur le déplacement $Z$ et la vitesse $V$ initiaux\footnote{Vues 736 et 737, feuillet 366. Écarts de la copie, que la transcription garde : l'exponentielle $e^{t-\sqrt{cK}}$ pour $e^{-t\sqrt{cK}}$ ; $M=A\sin x\sqrt{-1}$ pour $A\sin x\sqrt{-K}$ ; les constantes annoncées « $C$ et $D$ » et écrites $A$ et $B$ ; « $a\sqrt{-K}=\frac{\nu\pi}{a}$ », le dénominateur lu $a$ avec doute, alors que le calcul donne $a\sqrt{-K}=\nu\pi$ ; enfin, à la vue 737, les deux formules finales ont un signe $-$ devant le second terme de $s$ et $\int ZM\,dx$ au lieu de $\int VM\,dx$ dans celle de $ds/dt$, en désaccord avec les formules de la vue 736, qui sont les bonnes (vérifiées ici).}. C'est exact : $s$ est, à un facteur près, le coefficient de Fourier du déplacement sur le mode $\sin(\nu\pi x/a)$, et chaque mode oscille avec la pulsation $\nu\pi\sqrt c/a$, multiple de celle du mode fondamental. La « Remarque » dit que l'infinité de points demande une infinité de valeurs de $K$ : ce sont les valeurs propres $-(\nu\pi/a)^2$\footnote{Le nom de Fourier est postérieur au mémoire extrait ; la copie ne le porte pas.}.

L'extrait suivant (vues 738 et 739) porte sur la somme $\cos x+\cos2x+\cos3x+\dots$\footnote{Vues 738 et 739, feuillet 367. La référence « (1 2d. de T. P. 69) » est lue avec doute, de même que « S. 2 T. P. 317 » et « Pag. D. ».}. Lagrange l'écrit comme demi-somme de deux progressions géométriques en $e^{\pm kx\sqrt{-1}}$ et trouve $-\frac12$ ; pour la somme arrêtée au terme $\cos mx$, il trouve, exactement,
\[ \sum_{k=1}^{m}\cos kx=\frac{\cos mx-\cos(m+1)x}{2(1-\cos x)}-\frac12=-\frac12+\frac{\sin\bigl((m+\frac12)x\bigr)}{2\sin\frac x2}\qquad(x\not\equiv0\pmod{2\pi}). \]
D'Alembert objecte que le « dernier terme » $e^{\infty\sqrt{-1}}$ n'est pas nul ; Lagrange répond qu'on substitue bien $\frac{1}{1-x}$ à $1+x+x^2+\dots$ D'Alembert prend $x=45^\circ$ : les sommes partielles valent $\frac{1}{\sqrt2}$, $\frac1{\sqrt2}$, $0$, $-1$, $-1-\frac1{\sqrt2}$, $-1-\frac1{\sqrt2}$, $-1$, $0$, puis recommencent ; elles ne prennent donc que quatre valeurs, jamais $-\frac12$ (exact). Lagrange oppose $1-1+1-\dots$, qui vaudrait $0$ ou $1$ alors que $\frac{1}{1+x}$ donne $\frac12$. La remarque du copiste entre crochets note qu'il faut ajouter $-\frac12-\frac1{\sqrt2}$, $-\frac12$, $+\frac12$, $+\frac12+\frac1{\sqrt2}$ aux quatre sommes pour obtenir $-\frac12$, et que le terme $\frac{\cos mx-\cos(m+1)x}{2(1-\cos x)}$ « n'est susceptible que de quatre valeurs » : les deux affirmations sont exactes. Une autre remarque entre crochets propose de ne voir dans $x^\infty$ que « la limite de $x$, $x^2$, $x^3$ » : cela justifie $\frac{1}{1-x}$ pour $|x|<1$, mais non la somme en cosinus, où $e^{kx\sqrt{-1}}$ n'a pas de limite.

Ce qui est vrai aujourd'hui : la série $\sum\cos kx$ diverge pour tout $x$, ses termes ne tendant pas vers $0$ ; mais la moyenne de ses sommes partielles tend vers $-\frac12$ pour $x\not\equiv 0$ (sommation de Cesàro ; à $x=45^\circ$, la moyenne sur une période de huit termes vaut exactement $-\frac12$), et $\lim_{\rho\to1^-}\sum\rho^k\cos kx=-\frac12$ (sommation d'Abel), qui est exactement le calcul de Lagrange ; au sens des distributions, $\sum_{k\ge1}\cos kx=-\frac12+\pi\sum_{j}\delta(x-2\pi j)$. La série de $1-1+1-\dots$ a de même pour somme $\frac12$ au sens d'Abel et de Cesàro\footnote{Abel (1826), Cesàro (1890), la théorie des distributions (Schwartz, vers 1950) : tous postérieurs, et non cités.}.

\subsection*{Euler et Lagrange sur les cordes vibrantes}

\pagerange{740}{742}
« 3 V. des Mémoires de Turin --- Éclaircissements sur le mouvement des cordes vibrantes, par Euler »\footnote{Vues 740 à 742, feuillets 368 et 369. Le titre s'ouvre sur « 3 » suivi d'une capitale lue V. avec doute. C'est le même Euler que celui du mémoire de 1779 sur les lames, copié plus haut dans le volume. La première phrase après la « Question » ne se construit pas sur la feuille.}. La figure initiale de la corde est arbitraire, pourvu que ses ordonnées et ses pentes soient infiniment petites ; elle peut être « tirée par un mouvement libre de la main, sans qu'aucune loi de continuité y ait lieu », et la corde vibre quand même. Une remarque entre crochets donne l'intégrale $y=\Gamma(x+ct)+\Delta(x-ct)$ de $c^2\frac{\partial^2y}{\partial x^2}=\frac{\partial^2y}{\partial t^2}$ : c'est la solution de d'Alembert, exacte. Pour une corde d'épaisseur variable, $r^2\frac{\partial^2y}{\partial x^2}=\frac{\partial^2y}{\partial t^2}$ avec $r$ fonction de $x$ ; dans le cas
\[ y=x\,\Gamma\!\left(\frac\beta x-t\right)+x\,\Delta\!\left(\frac\beta x+t\right), \]
les cordes produisent toujours des vibrations régulières. C'est exact, et le calcul de cette lecture précise le cas : cette $y$ satisfait l'équation avec $r=x^2/\beta$, c'est-à-dire, à tension constante, une densité proportionnelle à $x^{-4}$ ; en posant $\xi=\beta/x$, $y/x$ satisfait l'équation de la corde uniforme en $\xi$, de sorte que, les extrémités $x_1$, $x_2$ étant fixes, tout mouvement est périodique de période $2|\beta/x_1-\beta/x_2|$.

Puis « La Grange Pag. 259 du même volume » : la courbe initiale n'a pas besoin d'être de la forme $y=\alpha\sin\pi x+\beta\sin2\pi x+\dots$ ; il suffit qu'une telle courbe en approche d'aussi près qu'on voudra, la courbe initiale n'étant alors « qu'une espèce d'asymptote » de la courbe génératrice. L'énoncé est vrai en termes actuels pour une figure initiale continue et nulle aux deux extrémités : elle est limite uniforme de sommes finies de sinus (théorème d'approximation de Weierstrass sous sa forme trigonométrique)\footnote{Weierstrass (1885), et pour la convergence de la série de Fourier elle-même Dirichlet (1829) : postérieurs, et non cités.}. La forme que la copie donne ensuite, $y=\alpha\sin\frac{\pi x}{a}+\beta\sin\frac{2\pi x}{b}+\gamma\sin\frac{3\pi x}{c}+\dots$, doit avoir partout le même dénominateur, la longueur $a$ de la corde\footnote{La transcription lit $a$, $b$, $c$, et les garde tels.}. « Pag. 275 », le développement de $\Delta(a+iy)$ en puissances de $i$ est celui de Taylor, avec les accents de Lagrange pour les dérivées : « [ce qui est bien près de la théorie des fonctions.] »\footnote{La remarque rapproche sans doute la page de la \emph{Théorie des fonctions analytiques} de Lagrange (1797) ; c'est une inférence de la transcription.}.

\subsection*{Lagrange : équations séparées}

\pagerange{743}{745}
« Mémoire de La Grange sur l'intégration de quelques équations différentielles dont les indéterminées sont séparées, mais dont chaque membre en particulier n'est point intégrable — T. 4 »\footnote{Vues 743 à 745, feuillets 369 (verso) et 370. « T. 4 » est sans doute le tome du recueil (inférence de la transcription). Des remarques entre crochets expliquent la formule d'addition des sinus et l'élimination des radicaux.}. Pour
\[ \frac{dx}{\sqrt{1-x^2}}=\frac{dy}{\sqrt{1-y^2}}\qquad\text{(A)}, \]
les membres ne sont pas intégrables algébriquement, mais l'équation a l'intégrale algébrique $\arcsin x=\arcsin y+\arcsin a$, soit
\[ x=y\sqrt{1-a^2}+a\sqrt{1-y^2}\qquad\text{(B)}, \]
par la formule d'addition des sinus ; on la retrouve par les logarithmes imaginaires, puis par une voie purement algébrique : en multipliant en croix et en intégrant par parties, $x\sqrt{1-y^2}=y\sqrt{1-x^2}+C$, qui revient à (B) avec $C=a$ (si $x=\sin u$, $y=\sin w$, c'est $\sin(u-w)=a$). Tout est exact. La même méthode donne, pour $\frac{dx}{\sqrt{\alpha+\beta x+\gamma x^2}}=\frac{dy}{\sqrt{\alpha+\beta y+\gamma y^2}}$,
\[ \left(x+\frac{\beta}{2\gamma}\right)\sqrt{\alpha+\beta y+\gamma y^2}=\left(y+\frac{\beta}{2\gamma}\right)\sqrt{\alpha+\beta x+\gamma x^2}+C, \]
qui est bien une intégrale (sa dérivée le long des solutions est nulle, vérification faite). L'extrait annonce enfin le cas d'un radicande du quatrième degré, qui est celui du théorème d'addition d'Euler pour les intégrales elliptiques, et s'arrête : « [le reste du mémoire … est copié presqu'en entier depuis la page 77 jusqu'à la p. 91 du 2\textsuperscript{me} V. du C. I. de Cousin] »\footnote{« du » et « Cousin » sont lus avec doute ; « C. I. » vaudrait « Calcul intégral » (inférence de la transcription).}.

\subsection*{Condorcet et Lagrange : variations ; les centres fixes}

\pagerange{746}{749}
« Condorcet Pag. 4 du 4\textsuperscript{me} V. de Turin »\footnote{Vue 746, feuillet 371. Les deux 4 sont lus avec doute. La caractéristique $\delta$ est un d bouclé ; la lettre donnée pour $\psi$ est un v barré, lu avec doute.} : si $Z$, expression en $x,y,z,dx,dy,\dots$, est la différentielle exacte d'une expression $B$ d'un ordre moins élevé, alors en intégrant $\delta Z$ par parties, la partie qui reste sous le signe $\int$ est identiquement nulle, et la partie hors du signe est $\delta B$. C'est, en termes actuels, la propriété que l'expression d'Euler--Lagrange d'une dérivée totale est identiquement nulle. « Pag 168 et suivantes, Mémoire de La Grange sur la méthode des variations » : $\psi=P\delta\varphi+Q\delta x+\dots$, $\int\psi+\Pi=$ constante, et le changement de $\delta$ en $d$ annule $\Pi$ ; le copiste compare entre crochets les deux manières, et conclut qu'« il faudroit avoir un mémoire précédent sous les yeux ». On rapporte ces lignes sans les commenter davantage : l'extrait ne donne pas les hypothèses.

Au verso (vue 747), deux mémoires de Lagrange (pp. 188 à 243) sur le mouvement d'un corps attiré vers un ou plusieurs centres fixes\footnote{Vues 747 à 749, feuillets 371 (verso) et 372. « la section de la méchanique » renvoie sans doute à la \emph{Mécanique analytique} de Lagrange (inférence de la transcription).} : pour deux centres, si l'un s'éloigne à l'infini, on obtient un centre newtonien et une force de grandeur et de direction constantes, ce qui est exact\footnote{Le problème des deux centres fixes est celui d'Euler ; le cas limite, un centre newtonien plus un champ uniforme, est aujourd'hui appelé problème de Stark (d'après l'effet Stark, 1913). Noms postérieurs ou absents de la feuille.}. Le « 2\textsuperscript{ème} Mémoire Pag. 218 » cherche quand $\frac{d^2x}{2dt^2}+M=0$, $\frac{d^2y}{2dt^2}+N=0$ ($M$, $N$ fonctions de $x$, $y$) admettent une intégrale première quadratique en les vitesses, $\frac{A\,dx^2+B\,dx\,dy+C\,dy^2}{2dt^2}+Z=$ constante. Il faut
\[ \frac{\partial A}{\partial x}=0,\quad \frac{\partial C}{\partial y}=0,\quad \frac{\partial A}{\partial y}+\frac{\partial B}{\partial x}=0,\quad \frac{\partial B}{\partial y}+\frac{\partial C}{\partial x}=0, \]
d'où $A=a+by+cy^2$, $C=g+hx+cx^2$, $B=k-bx-hy-2cxy$, ce que le calcul de cette lecture vérifie\footnote{La différentielle de $A\,dx^2+B\,dx\,dy+C\,dy^2$ est recopiée incomplète : il y manque le terme $\frac{\partial A}{\partial x}dx^3$, et le terme en $\frac{\partial A}{\partial y}+\frac{\partial B}{\partial x}$ y porte $dx\,dy$ au lieu de $dx^2\,dy$ ; les quatre conditions qu'on en tire sont néanmoins toutes écrites et justes. Les justifications entre crochets de la vue 748 (une fonction dont la dérivée partielle en $x$ est nulle ne contient pas $x$ ; $\frac{\partial^2A}{\partial y^2}=\frac{\partial^2C}{\partial x^2}$ force le même coefficient $c$) sont exactes.}, puis la condition d'exactitude $\frac{\partial(2AM+BN)}{\partial y}=\frac{\partial(BM+2CN)}{\partial x}$ ; en posant $D=4AC-B^2$ et $Z=D^2V$, $V$ arbitraire,
\[ M=D\left(4V\frac{\partial C}{\partial x}+2C\frac{\partial V}{\partial x}-B\frac{\partial V}{\partial y}\right),\qquad N=D\left(4V\frac{\partial A}{\partial y}+2A\frac{\partial V}{\partial y}-B\frac{\partial V}{\partial x}\right), \]
et l'intégrale $\frac{A\,dx^2+B\,dx\,dy+C\,dy^2}{2dt^2}+D^2V=$ constante. Tout est exact, pour les équations écrites avec $2dt^2$ comme à la vue 747\footnote{À la vue 749, les équations sont récrites $\frac{d^2x}{dt^2}+M=0$, sans le 2 ; avec cette forme, l'intégrale écrite n'est pas conservée (il faudrait $dt^2$ au dénominateur) : vérifié.}. Les conditions sur $A$, $B$, $C$ sont celles qu'on appelle aujourd'hui les équations d'un tenseur de Killing du plan euclidien\footnote{Nom postérieur, absent de la feuille.}.

\subsection*{Lagrange et l'équation $x^2 - ay^2 = 1$}

\pagerange{749}{749}
Sous un trait épais : « [Dans le mémoire cité (P. 37) de la Théorie des nombres on trouve (P. 46 T. 4 de Turin) »\footnote{Vue 749, verso du feuillet 372, dernière page écrite du volume, sous le cachet de la Bibliothèque royale. « 46 » et « 4 » sont lus avec doute. « la Théorie des nombres » est un ouvrage qui cite le mémoire de Lagrange ; lequel, la feuille ne le dit pas.}. Le lemme $(x^2-ay^2)(x'^2-ay'^2)=(xx'\pm ayy')^2-a(xy'\pm yx')^2$ est exact (identité de Brahmagupta). Il s'ensuit que toute puissance d'une solution de $x^2-ay^2=1$ en est une : si $p^2-aq^2=1$, alors $(p+q\sqrt a)^m=x+y\sqrt a$ donne
\[ x=\frac{(p+q\sqrt a)^m+(p-q\sqrt a)^m}{2},\qquad y=\frac{(p+q\sqrt a)^m-(p-q\sqrt a)^m}{2\sqrt a}, \]
une infinité de solutions. « La Grange démontre ensuite P. 65 » que si $p$, $q$ sont les plus petites valeurs, toutes les solutions sont de cette forme. C'est vrai, pour $a$ entier positif non carré, hypothèse que l'extrait ne rappelle pas : les solutions en entiers positifs sont exactement les $(p+q\sqrt a)^m$, $m\ge1$, où $p+q\sqrt a$ est la plus petite\footnote{Vérifié pour $a=2$, $3$, $5$, $7$, $13$ sur toutes les solutions avec $y<200\,000$. Pour $q'$, la plus petite solution après $(p,q)$, la copie écrit $q'=\frac{(p+q\sqrt a)^2+(p-q\sqrt a)^2}{2\sqrt a}$, avec $+$ au numérateur là où la formule générale donne $-$ ; la valeur juste est $q'=2pq$. La remarque finale entre crochets (pour $m=1$ les deux membres deviennent identiques) est exacte.}.

\section*{Ce que les transcriptions laissent en suspens}

\pagerange{1}{750}
Ce qui suit n'est pas la liste des mots douteux, que les transcriptions portent
tous, mais celle des points que les rédacteurs de cette lecture ont demandé de
revoir sur les images, pour une révision humaine. Rien n'y est corrigé ici : la
lecture a suivi, à chaque fois, ce que dit sa note au passage cité.

\emph{Questions d'ensemble.}
\begin{itemize}
\item La série à l'encre (1 à 372, avec un « 55 bis » à la vue 113 et le feuillet
blanc des vues 732--733 sans numéro) contre les 372 feuillets du catalogue ; et la
foliotation de la Bibliothèque, que les transcriptions ne relèvent pas à côté de
cette série.
\item La main de la copie d'Euler et de la mise au net de la critique (vues 144 à
206 et 320 à 389), celle des extraits de Turin et de leurs remarques entre
crochets (vues 734 à 749), celle des « Idées à vérifier » (vues 136 à 140).
\item Lequel des deux brouillons de la critique fut écrit le premier (vues 210 à
216 et 218 à 316) ; lequel des deux états de la comparaison des sons (vues 486 à
503 et 600 à 613) ; si les vues 593 à 613 sont une seule seconde copie.
\item Si les vues 620 à 626 sont un brouillon du mémoire « présenté en 1815 », et
si le long mémoire est le « mémoire pour 1813 » (inférences de cette lecture).
\item L'ordre des pages des calculs sur la chaleur (vues 630 à 640) et la
pagination de la pièce sur la sphère (vues 652 à 670), que la lecture établit par
le texte et non par une lecture des chiffres.
\item « L'auteur » des notes d'arithmétique, et les renvois « N.o 125 »,
« P. 116 », « §126 » (vues 718 à 728) : la transcription y voit Gauss et les
\emph{Disquisitiones}, ce que seule la vue 720 écrit.
\item Ce que sont, imprimés, l'Exposition et le mémoire sur la courbure, que le
feuillet de garde range parmi les « mémoires imprimés ».
\end{itemize}

\emph{Points de lecture, par lot de transcription.}
\begin{itemize}
\item \emph{Lot 6.} Vue 114 : les équations (1) à (3), repassées et empâtées — le
second dénominateur de (1), lu $R$ seul ; le facteur de (2), lu $\sin(2A-a)$ avec
doute ; le premier membre de (2), écrit $\frac1R-\frac1{R'}$ là où l'ordre des
sections qu'elle donne demande $\frac1{R'}-\frac1R$ (les accents sont à revoir).
\item \emph{Lot 7.} Vue 133 : l'imprimé porte-t-il bien
$\left(\frac6n\right)=1$ pour $n=120\mu+61$ (le calcul donne $-1$) ? Vues 136 à
140 : l'en-tête du lot ne se prononce pas sur la main, la note de la vue 136 les
dit « de la main de l'auteur ».
\item \emph{Lot 9.} Vue 163 : la note lue « b » dans « si sonus fundamentalis
fuerit b, sequens futurus erit cis » ; le triton sous cis est g.
\item \emph{Lot 10.} Vue 182 : les barres d'octave des deux derniers noms de
notes, qui en donnent une de plus que le calcul. Vue 199 : l'exposant surchargé lu
« $2\omega$ », que le calcul et les citations françaises des vues 210 et 224
donnent $\omega$.
\item \emph{Lot 11.} Vues 201--202 : « quinto », là où le sens et la citation
française de la vue 228 demandent le cas IV. Vue 218 : la page citée, lue
« p. 151 » ou « 131 » (le § 47 commence à la page (151) de la copie).
\item \emph{Lot 13.} Vue 242 : « P. 898 §. 83 » et « le Scholie §. 81 », là où la
mise au net (vue 332) écrit « P. 298 § 83 » et « § 21 ».
\item \emph{Lot 14.} Vue 272 : le renvoi à Chladni, « pag 4 §§ 71 », que la mise au
net (vue 354) écrit « P. 94 § 71 ».
\item \emph{Lot 16.} Vues 301 et 304 : le signe des fractions données pour
$\tan\frac12\omega$ et $\tan\omega$, contraire à celui que demandent les
coefficients ; vue 304 : l'étendue du radical, rendu
$\sqrt2\,(e^{2\omega}+e^{-2\omega})$, que le calcul veut sur tout le produit.
\item \emph{Lot 18.} Vues 352 et 354 : le double signe du facteur
$e^{\lambda\omega}\mp e^{\mu\omega}$, et l'égalité écrite
$e^{\lambda\omega}=\pm e^{-\lambda\omega}$ là où le raisonnement attend
$e^{\mu\omega}$ ; vue 354 : le « 94 » de « P. 94 », peut-être « g4 ».
\item \emph{Lot 19.} Vues 379--380 : un béquet portant « (29) », vu par
transparence à travers la vue 380 et non transcrit.
\item \emph{Lot 20.} Vue 381 : le même signe des tangentes qu'aux vues 301 et 304.
Vue 383 : le radical, et les deux numérateurs lus $e^{\omega}-e^{-\omega}$ (le
brouillon, vue 306, donne $e^{\omega}+e^{-\omega}$ pour le cosinus, ce qui est
juste). Vue 389 : « quinze », qui ressemble à « quize ». Vue 391 : le chiffre des
légendes « fig 1 », qui se lit 1 ou 7. Vue 392 : le signe lu « \&c », et le signe
de l'étape où le feuillet écrit $+1\pm e^{\omega}$ pour $-1\mp e^{\omega}$.
\item \emph{Lot 21.} Vues 409 et 413 : l'argument du dernier sinus de
l'intégrale (I), lu « $\pi\,may$ » puis « $\pi.mn.y$ ».
\item \emph{Lot 22.} Vue 429 : la comparaison entre la partie comprise entre le
second et le troisième nœud et la précédente, inversée par rapport au calcul et à
sa propre conclusion.
\item \emph{Lot 24.} Vue 461 : « dans le cas d'un nombre impair », là où la plaque
demande un nombre pair de nœuds de la lame.
\item \emph{Lot 25.} Vue 483 : les facteurs $q$ et $p$ qui manquent à deux
coefficients de la variation. Vue 487 : l'intégrale (I) dite périodique « dans le
sens des $y$ », contre « des $x$ » aux vues 409 et 413. Vue 488 : « $n=4$ »
au-dessus de la ligne, contre « $n=1$ » à la vue 601.
\item \emph{Lot 26.} Vue 510 : « 2 lignes … 4 pareilles lignes » pour la plaque
$3:5$, où le calcul demande 2 et 3. Vue 518 : « $16.25+96$ », où le calcul
demande $36$.
\item \emph{Lot 27.} Vue 540 : « 1 once 2 gros de cire », que les vues 541 et 542
contredisent (« 1 gros 2 grains »).
\item \emph{Lot 28.} Vue 549 : les renvois « N.o 19 » et « N.o 14 », que la
rédaction des vues 676 à 678 écrit 119 et 104.
\item \emph{Lot 30.} Vue 585 : le chiffre d'octave de fa$^{\sharp}$3, qui
ressemble à un 9. Vue 589 : « deux nœuds », attendu « quatre », et l'octave de
« si 3 ». Vue 596 : $\delta x$ au lieu de $\delta z$ dans la seconde condition.
\item \emph{Lot 31.} Vue 601 : « $n=1$ » (voir la vue 488).
\item \emph{Lot 32.} Vue 636 : $\cos^2\theta$ devant la dérivée seconde en
$\sin\theta$, que la dérivation du premier membre donne $\cos\theta$.
\item \emph{Lot 33.} Vue 650 : le coefficient lu « 2 » de
$\cos^2\theta\,(\partial_\theta z)^2$ dans le facteur $K^3$, que le calcul donne
égal à 1.
\item \emph{Lot 35.} Vue 684 : les dénominateurs lus « $2^2$ », « $2.3$ »,
« $2.5$ », où la règle qu'elle énonce donne $2\cdot2$, $2\cdot4$, $2\cdot6$. Vue
695 : « $A'\cos\varphi-A'\sin\varphi=0$ », où le brouillon (vue 690) porte $B'$.
\item \emph{Lot 36.} Vue 704 : « puissance paire de 2 » dans la phrase d'ouverture,
que la dernière ligne et le calcul veulent « impaire ». Vue 708 : l'exposant de
$z''$, lu $n$ là où la vue 710 et le calcul donnent $2n$. Vue 710 : l'exposant lu
« 8n » (le calcul demande $2n$), et $2^{n-1}$ là où la colonne de droite porte
$2^{2n-1}$. Vue 717 : blanche selon l'en-tête du lot et sa note de numérotation,
couverte d'autres calculs selon la note de la vue 716. Vue 719 : les lectures
« $10n+2$ » et « $4k+11$ » (pour $40k+11$). Vue 720 : le second « $4n+1$ » du
feuillet 355.
\item \emph{Lot 37.} Vue 737 : « $a\sqrt{-K}=\frac{\nu\pi}{a}$ », le dénominateur
lu $a$ avec doute (le calcul donne $a\sqrt{-K}=\nu\pi$).
\end{itemize}

\end{document}
